\documentclass[11pt]{article}
\usepackage[a4paper,margin=27mm,headheight=15pt]{geometry}
\usepackage[T1]{fontenc}
\usepackage[utf8]{inputenc}
\usepackage{lmodern,microtype,bm,graphicx}
\usepackage{amsmath,amssymb,amsthm,mathtools,mathrsfs}
\usepackage{booktabs,array,longtable,tabularx,enumitem}
\usepackage{xcolor,fancyhdr,xurl,hyperref}
\definecolor{paperink}{HTML}{185B58}
\hypersetup{colorlinks=true,linkcolor=paperink,citecolor=paperink,urlcolor=paperink,
  pdftitle={Euler Congruences: Completion, Factorial Strengthening and Sharp Dyadic Geometry},
  pdfauthor={Komorebi P.},pdfsubject={Generalized Euler congruences, factorial precision, sharp dyadic periods, formal coverage and remaining tail simplicity}}
\newcommand{\N}{\mathbb N}
\newcommand{\Z}{\mathbb Z}
\newcommand{\Q}{\mathbb Q}
\newcommand{\F}{\mathbb F}
\newcommand{\Zp}{\mathbb Z_p}
\newcommand{\Qp}{\mathbb Q_p}
\DeclareMathOperator{\coker}{coker}
\DeclareMathOperator{\im}{im}
\DeclareMathOperator{\rank}{rank}
\newcommand{\lean}[1]{{\footnotesize\nolinkurl{#1}}}
\newtheorem{theorem}{Theorem}[section]
\newtheorem{proposition}[theorem]{Proposition}
\newtheorem{lemma}[theorem]{Lemma}
\newtheorem{corollary}[theorem]{Corollary}
\theoremstyle{definition}
\newtheorem{definition}[theorem]{Definition}
\newtheorem{example}[theorem]{Example}
\theoremstyle{remark}
\newtheorem{remark}[theorem]{Remark}
\numberwithin{equation}{section}
\setlength{\parskip}{3pt}
\setlength{\emergencystretch}{2em}
\setlist{itemsep=2pt,topsep=5pt}
\allowdisplaybreaks[1]
\pagestyle{fancy}
\fancyhf{}
\fancyhead[L]{\small\scshape Euler congruences / Chapter IX}
\fancyhead[R]{\small\scshape Final strengthened edition}
\fancyfoot[L]{\small Komorebi P.}
\fancyfoot[R]{\thepage}

\title{\textbf{Euler Congruences}\\[5pt]
  \Large Completion, Factorial Strengthening,\\ and Sharp Dyadic Geometry\\[7pt]
  \large Projector-Resolved Seed Calculus IX}
\author{Komorebi P.}
\date{}

\usepackage[nameinlink,noabbrev]{cleveref}

\providecommand{\tightlist}{\setlength{\itemsep}{0pt}\setlength{\parskip}{0pt}}
\providecommand{\passthrough}[1]{#1}
\newcommand{\mcode}[1]{\ifmmode\text{\nolinkurl{#1}}\else\nolinkurl{#1}\fi}
\definecolor{ink}{HTML}{185B58}
\newcommand{\ZZ}{\mathbb Z}\newcommand{\QQ}{\mathbb Q}\newcommand{\FF}{\mathbb F}\newcommand{\NN}{\mathbb N}
\newcommand{\Wp}{\mathcal W_p}\newcommand{\Hp}{\mathcal H(\Z_p)}
\newcommand{\ordp}{\operatorname{ord}_p}\newcommand{\ord}{\operatorname{ord}}
\newcommand{\coef}[2]{[#1^{#2}]}\newcommand{\target}[1]{{\footnotesize\nolinkurl{#1}}}
\newcommand{\pageheading}[2]{{\large\bfseries\color{ink}#1}\par{\small #2}\par\medskip}
\newcommand{\sech}{\operatorname{sech}}\newcommand{\fall}[2]{(#1)_{#2}}
\newcommand{\EE}{\mathcal E}\newcommand{\PP}{\mathcal P}\newcommand{\code}[1]{{\small\nolinkurl{#1}}}
\newcommand{\Hcal}{\mathcal H}\newcommand{\Ecal}{\mathcal E}\newcommand{\Ccal}{\mathcal C}\newcommand{\Pcal}{\mathcal P}\newcommand{\Scal}{\mathcal S}\newcommand{\Ocal}{\mathcal O}\newcommand{\Rcal}{\mathcal R}\newcommand{\Vcal}{\mathcal V}\newcommand{\Wcal}{\mathcal W}
\newcommand{\End}{\operatorname{End}}\newcommand{\Ker}{\operatorname{Ker}}\newcommand{\Res}{\operatorname{Res}}\newcommand{\vp}{v_p}\newcommand{\vTwo}{v_2}\newcommand{\sTwo}{s_2}\newcommand{\DeltaH}{\Delta_h}\newcommand{\abs}[1]{\left\lvert#1\right\rvert}\newcommand{\norm}[1]{\left\lVert#1\right\rVert}\newcommand{\one}{\bm 1}\newcommand{\starh}{\mathbin{\star}}\newcommand{\e}{\mathrm e}
\newtheorem{histtheorem}[theorem]{Historical theorem}
\crefname{histtheorem}{historical theorem}{historical theorems}
\newtheorem{histproposition}[theorem]{Historical proposition}
\crefname{histproposition}{historical proposition}{historical propositions}
\newtheorem{histlemma}[theorem]{Historical lemma}
\crefname{histlemma}{historical lemma}{historical lemmas}
\newtheorem{histcorollary}[theorem]{Historical corollary}
\crefname{histcorollary}{historical corollary}{historical corollaries}
\newtheorem{histdefinition}[theorem]{Historical definition}
\crefname{histdefinition}{historical definition}{historical definitions}
\newtheorem{histexample}[theorem]{Historical example}
\crefname{histexample}{historical example}{historical examples}
\newtheorem{histremark}[theorem]{Historical remark}
\crefname{histremark}{historical remark}{historical remarks}
\newtheorem{histstandardfact}[theorem]{Historical standard fact}
\crefname{histstandardfact}{historical standardfact}{historical standardfacts}


% Operation Komorebi: shared publication profile.
\usepackage{komorebi-publication}
\renewcommand{\KomorebiTitle}{Euler Congruences}
\renewcommand{\KomorebiSubtitle}{Completion, factorial strengthening, and sharp dyadic geometry}
\renewcommand{\KomorebiSeries}{Chapter IX}
\renewcommand{\KomorebiKind}{Mathematical text}
\renewcommand{\KomorebiEdition}{}
\renewcommand{\KomorebiShortTitle}{Chapter IX}
\renewcommand{\KomorebiTitleSize}{38}
\renewcommand{\KomorebiTitleLeading}{43}
\renewcommand{\KomorebiProfile}{classical}
\renewcommand{\KomorebiOrientation}{portrait}
\KomorebiSetup
\begin{document}
\KomorebiCover

\relax
\begin{abstract}
We prove Nishibuchi's generalized Euler congruence conjecture in its full stated parameter range and develop sharper arithmetic consequences. For every prime-power packet width, a charged Hurwitz differential operator transfers integral difference estimates through the reciprocal generating series. In the dyadic case this gives mixed finite-difference bounds at every row, without a starting threshold. For the four addresses of width four, a divided-power argument adds the factorial gain $v_2(k!)$, determines every repeated-step valuation in addresses zero and one, and gives a stronger address-two estimate with exact values at all power-of-two orders. The corresponding first differences yield least periods and bijective isometric continuations onto explicit dyadic residue classes. We also determine exact Taylor radii for three addresses. The manuscript retains the complete earlier arithmetic, lifting, deformation, and analytic-regularity development, with detailed mathematical appendices. A separate appendix states the precise remaining simplicity problem for ordinary Euler tails. Formal coverage is explicit: the external conjecture and several four-address results have inherited Lean verification; the new factorial, mixed-step, all-rank, and radius arguments are written proofs with recorded audits and exact computational checks, and await dedicated Lean formalization.
\end{abstract}
\section*{Principal results}
\addcontentsline{toc}{section}{Principal results}
The original external problem is Nishibuchi's Conjecture 3.9 for generalized Euler numbers. Its full parameter range is resolved. If $p$ is prime, $m$ is positive with $m\mid p-1$ or $m=2$, $0\le j<mp$, and $q=\operatorname{lcm}(2,p-1)$, then the original Euler coefficients are $p$-adically integral and
\[
 E_{mpn+qp^r}^{(mp,j)}\equiv E_{mpn}^{(mp,j)}\pmod{p^r}
 \qquad (r\ge1,\ n\ge r).
\]
The chapter states the generating function and the conversion between ordinary degree and packet row before using this formula. The normalized coefficient is $C_{N,j}(n)$ and the original coefficient is $E_{Nn}^{(N,j)}=j!C_{N,j}(n)$. The constant coefficient of the original series is $j!$.

The strongest common-step conclusion in width four is particularly compact. For $n\ge0$ and $k,H\ge1$,
\begin{equation*}
 \sum_{i=0}^k(-1)^{k-i}\binom ki E_{4(n+iH)}^{(4,j)}
 \equiv0\pmod{2^{\,v_2(k!)+kv_2(H)+A_j(k)}},
\end{equation*}
where
\[
 A_0(k)=A_1(k)=k,\qquad
 A_2(k)=\max\{k+4,\,2k+\lfloor k/2\rfloor+1\},\qquad
 A_3(k)=k+2.
\]
There is no starting-row threshold. The exponent is exact in addresses zero and one for every order and positive step. It is exact in address two at every power-of-two order. The proof chapters also give mixed-step forms, the cancellation criterion in addresses zero and one, and exact first-displacement laws.

\begin{center}
\small
\begin{tabularx}{\linewidth}{@{}>{\raggedright\arraybackslash}p{.29\linewidth}>{\raggedright\arraybackslash}X>{\raggedright\arraybackslash}p{.21\linewidth}@{}}
\toprule
Result & Scope & Evidence\\
\midrule
Nishibuchi Conjecture 3.9 & Every parameter in the original statement; explicit $n\ge r$ & Inherited Lean proof\\
Width-four all-row common-step bound, first differences, least periods and isometry & Every address, row and positive displacement & Inherited Lean proof\\
General prime-power inverse transfer and mixed steps & Every prime-power width; signed steps for odd primes & Written proof and audit\\
Factorial gains and exact higher-order families & Four width-four addresses, with stated sharpness scopes & Written proof and audit\\
Taylor radii and broader dyadic first differences & Exact radii in addresses $0,1,2$; stated all-rank cases & Written proof and audit\\
Ordinary Euler tail simplicity & Partial families and an explicit remaining discriminant problem & Open in full generality\\
\bottomrule
\end{tabularx}
\end{center}

The arithmetic development leads to polynomial lifting, deformation and analytic regularity. The appendices give the detailed continuation arguments and the recovered foundational treatment. Where an earlier formulation is retained for comparison, its hypotheses and proof status remain explicit. The formal companion supplies the exact source map and distinguishes source review from compiler execution. The visual studies supply finite examples and diagrams derived from the theorems.

Recorded proof reviews are internal reviews. Formal proof coverage and the scope of each mathematical result are stated separately.
\clearpage



\section{Scope and conventions}
\label{sec:scope}
The central objects are sequences and polynomial equations equipped with
arithmetic coordinates. Residue projectors select addresses in a sequence,
finite differences measure compatibility between nearby addresses, and
valuation bounds control the passage from finite precision to a $p$-adic
limit. These constructions provide a common language for the Euler packet
results and the later lifting results. Each theorem nevertheless retains its
own coefficient ring, topology and nonsingularity hypotheses.

The arithmetic framework is developed together with its completed
extensions. The stated external Nishibuchi conjecture is proved in its full
specified parameter range. The dyadic refinements and polynomial deformation theorems give stronger
conclusions under the hypotheses stated with each result. The native conjectures retain their original
identifiers, 13.2 and 13.3, when discussed as research questions, although the
section numbering of this consolidated manuscript is new.

Throughout, $\N=\{0,1,2,\ldots\}$. When $p$ is prime, $\Zp$ denotes the
$p$-adic integers and $\Qp$ their fraction field. We normalize $v_p(p)=1$.
For a rational or $p$-adic quantity, a displayed lower valuation bound is
understood with the mathematical convention $v_p(0)=+\infty$.
The corresponding formal predicates are zero-safe: the quantity is zero or
its valuation meets the stated bound. An exact finite valuation assertion
requires a nonzero quantity, and this is stated where used. Congruence in
$\Zp$ means ideal divisibility. Thus
\[
 x\equiv y\pmod{p^r}\quad\Longleftrightarrow\quad
 x-y\in p^r\Zp.
\]
For vectors, ideal membership and congruence are coordinatewise.

Write $\Delta_u f(n)=f(n+u)-f(n)$ and use $\Delta_u^k$ for its $k$-fold
iterate. Euler coefficients indexed by ordinary degree and by packet row are
distinguished explicitly. A packet displacement by $u$ rows is a degree
displacement by $Nu$ in a packet of width $N$.

The main text presents theorem statements and their arguments; accompanying
source maps identify the Lean declarations and archived files that support
them. The formal companion identifies the scope of each recorded build. The
91-check polynomial record covers the polynomial development; separate
records support the arithmetic results.
The inherited formal development is retained without adding Lean declarations. The strengthening chapters add written proofs and clearly mark their formalization boundary; no historical-priority claim is made. The analytic inverse-function results
used in Section~\ref{sec:regularity} come from Mathlib.

\clearpage
{\small\setlength{\parskip}{0pt}\tableofcontents}
\clearpage
\section{Arithmetic packets and coefficient transport}
\label{sec:foundations}

We work with sequences, formal power series, and finite recurrence states.
Residue selectors isolate coefficients in a chosen congruence class, while
triangular transport expresses the same object in new coordinates. A valuation
filtration determines how many coordinates are needed at a given precision.
Together, these operations provide a common framework for the Euler
congruences and the finite lifting constructions developed later.

The original arithmetic seed notation
\[
\mathcal S=(A,a_0,\mathcal O,\mathcal R,\{P_{N,j}\},w,\tau_p)
\]
records a coefficient ring or module, a seed, allowed operators, relations,
residue selectors, a depth weight, and a topology. This notation describes the
ingredients of the construction; a category of seed presentations and a
notion of equivalence still require separate definitions. We therefore state
each theorem for a concrete object. Section~\ref{sec:open} discusses the
additional structure needed for the original native conjectures.

The Lean declarations cited in this section have namespace prefix
\lean{KomorebiPaperIX}. Their exact modules and source hashes are recorded in
the accompanying statement map. All belong to the preserved baseline
formalization; this section brings their statements and proofs together.

\subsection{Residue selectors and exact reconstruction}

Let $R$ be a commutative ring, $a:\N\to R$, $N\geq1$, and $0\leq j<N$.
Define
\[
(P_{N,j}a)(n)=
\begin{cases}
a(n),&n\equiv j\pmod N,\\
0,&\text{otherwise}.
\end{cases}
\]
This definition keeps the original index set. The compressed subsequence
$a(Nq+j)$, used below for recurrence packets, is a different representation of
the selected coefficients.

\begin{theorem}[Residue reconstruction]
\label{thm:residue-reconstruction}
The selectors satisfy
\[
P_{N,j}^2=P_{N,j},\qquad P_{N,i}P_{N,j}=0\quad(i\ne j),
\qquad \sum_{j=0}^{N-1}P_{N,j}=I.
\]
If $N\mid M$ and $M>0$, then
\[
P_{N,j}a=\sum_{\substack{0\leq r<M\\r\equiv j\pmod N}}P_{M,r}a.
\]
The same identities hold coefficientwise for formal power series.
\end{theorem}

\begin{proof}
At each index $n$, exactly one standard residue label equals $n\bmod N$.
Consequently a repeated selector changes nothing, two different selectors have
zero composite, and the sum retains $a(n)$ exactly once. If $N\mid M$, reducing
$n\bmod M$ modulo $N$ gives $n\bmod N$. This proves the refinement identity
at every coefficient. Equality of all coefficients proves the formal-series
versions.
\end{proof}

\begingroup\small\raggedright
\noindent\emph{Formal correspondence.}
\lean{ResiduePackets.packet_idempotent},
\lean{ResiduePackets.packet_orthogonal},
\lean{ResiduePackets.packet_reconstruction}, and
\lean{ResiduePackets.packet_refinement}; the series variants have the
prefix \lean{series_packet} in the same module.
The algebraic statements need no convergence hypothesis.
\par\endgroup

\subsection{Formal exponential coordinates and Hurwitz multiplication}

Over a characteristic-zero field $K$, put
\[
\mathcal E(a)=\sum_{n\geq0}\frac{a_n}{n!}t^n,
\qquad \mathcal D(A)_n=n![t^n]A.
\]
These inverse maps between sequences and $K[[t]]$ provide \emph{formal}
exponential coordinates. No analytic convergence is needed for the coefficient
identities below.

\begin{theorem}[Constructed unit-triangular coordinates]
\label{thm:egf-coordinates}
Let $h:\N\to K$ satisfy $h_0=1$, write $H=\mathcal E(h)$ and
$H^{-1}=\mathcal E(k)$, and define
\[
U_Ha=\mathcal D(H\mathcal E(a)).
\]
Then $U_H$ is an equivalence, with inverse
$U_H^{-1}a=\mathcal D(H^{-1}\mathcal E(a))$. Its coefficient formulas are
\[
(U_Ha)_n=\sum_{r=0}^{n}\binom nr h_r a_{n-r},\qquad
(U_H^{-1}a)_n=\sum_{r=0}^{n}\binom nr k_r a_{n-r}.
\]
The resulting triangular matrix has diagonal entries $1$. The transported
selectors
\[
\widehat P_{N,j}=U_HP_{N,j}U_H^{-1}
\]
are idempotent and pairwise orthogonal, and their sum is the identity.
\end{theorem}

\begin{proof}
The factorials are nonzero in $K$, so encoding and decoding cancel
coefficientwise. The constant coefficient of $H$ is $1$; hence its formal
power-series inverse exists. Multiplication by $H$ and $H^{-1}$ gives the two
inverse laws. The ordinary product coefficient formula, multiplied by $n!$,
gives the displayed binomial sums. In row $n$, the coefficient of $a_n$ is
$h_0=1$ and coefficients of $a_q$ with $q>n$ vanish. Conjugating the identities
of Theorem~\ref{thm:residue-reconstruction} by the constructed equivalence proves
the transported projector identities.
\end{proof}

\begingroup\small\raggedright
\noindent\emph{Formal correspondence.}
\lean{EGF.decode_encode}, \lean{EGF.encode_decode},
\lean{EGF.egf_product}, \lean{EGF.coordinateEquiv},
\lean{EGF.coordinate_forward}, \lean{EGF.coordinate_inverse},
and \lean{EGF.transported_reconstruction}, together with the idempotence and
orthogonality declarations in \lean{EGFCoordinates}.
\par\endgroup

In characteristic $p$, factorial division is unavailable in general, but the
coefficient operation remains meaningful over every commutative ring:
\[
(f\star g)_n=\sum_{r=0}^{n}\binom nr f_r g_{n-r},\qquad
(D^d f)_n=f_{n+d},\qquad \mathbf1_n=\begin{cases}1,&n=0,\\0,&n>0.\end{cases}
\]
Here $\star$ is Hurwitz multiplication. The symbol $D^d$ denotes a coefficient
shift, which agrees with formal differentiation after characteristic-zero
encoding.

\begin{theorem}[Hurwitz product rules]
\label{thm:hurwitz-product}
For $d\geq0$,
\[
D^d(f\star g)=\sum_{r=0}^{d}\binom dr(D^r f)\star(D^{d-r}g).
\]
If $R$ has prime characteristic $p$, then
\[
D^p(f\star g)=f\star D^p g+D^p f\star g.
\]
Over a field, if $g_0=1$, a sequence $b$ is constructed recursively with
$b\star g=\mathbf1$. If also $D^p g=Dg$ in characteristic $p$, then
$D^p b=Db$.
\end{theorem}

\begin{proof}
Pascal's identity proves the first product rule, and induction proves the
higher rule. The intermediate binomial coefficients $\binom pr$ vanish in
characteristic $p$. For the inverse, solve row $n$ using the coefficient of
$b_n$, which is $g_0=1$:
\[
b_n=\mathbf1_n-\sum_{r<n}\binom nr b_r g_{n-r}.
\]
Subtract the first and $p$th product rules for $b\star g=\mathbf1$.
The assumed identity for $g$ leaves
$(D^p b-Db)\star g=0$. A second triangular induction shows that convolution
with a sequence having constant coefficient $1$ has trivial kernel.
\end{proof}

\begingroup\small\raggedright
\noindent\emph{Formal correspondence.}
\lean{Hurwitz.shift_convolution}, \lean{Hurwitz.prime_shift_convolution},
\lean{Hurwitz.convolution_left_eq_zero},
\lean{HurwitzInverse.inverse_convolution}, and
\lean{HurwitzInverse.inverse_prime_difference}.
Working directly with coefficient sequences avoids the undefined expression
$1/n!$ in $\F_p$ when $p\mid n!$.
\par\endgroup

\subsection{Triangular inversion and a finite precision filtration}

\begin{theorem}[Forward substitution and denominator clearing]
\label{thm:triangular-foundation}
Let $T_{nq}$ be coefficients over a field with $T_{nn}\ne0$ for every $n$.
For any sequence $b$, define
\[
a_n=\frac{b_n-\sum_{q<n}T_{nq}a_q}{T_{nn}}.
\]
This constructs the unique solution of
$\sum_{q=0}^{n}T_{nq}a_q=b_n$ for all $n$. If $T_{nq}$ and $b_n$ are integers
and the solution is taken over $\Q$, then
\[
\left(\prod_{r=0}^{n}T_{rr}\right)a_n\in\Z.
\]
\end{theorem}

\begin{proof}
The recursion uses only smaller indices, and substitution verifies each row.
If a second solution agrees through row $n-1$, subtracting row $n$ and
cancelling $T_{nn}$ proves equality at $n$. For integrality, use induction and
promote every earlier clearing product to $\prod_{r<n}T_{rr}$. Multiplication
by the additional factor $T_{nn}$ cancels the only new division.
\end{proof}

\begingroup\small\raggedright
\noindent\emph{Formal correspondence.}
\lean{Triangular.solve_row}, \lean{Triangular.solve_unique}, and
\lean{Triangular.solve_denominator_integral}. The clearing product bounds the
denominator; the unscaled solution can still be nonintegral.
\par\endgroup

To isolate the filtration argument, let $p$ be any element of a commutative
ring and let $w:\N\to\N\cup\{\infty\}$. Set $p^{[r]}=p^r$ for finite $r$
and $p^{[\infty]}=0$. The latter notation means the coefficient vanishes.

\begin{theorem}[Filtered finite reconstruction]
\label{thm:filtered-foundation}
For every $n,L\in\N$ and sequences $c,S$ in that ring,
\[
\begin{split}
&\sum_{r=0}^{n}\binom nr c_rp^{[w(r)]}S_{n-r}\\
&\hspace{8mm}\equiv
\sum_{\substack{0\leq r\leq n\\w(r)<L}}
\binom nr c_rp^{[w(r)]}S_{n-r}\pmod{p^L}.
\end{split}
\]
Here congruence means divisibility of the difference by $p^L$ in the
coefficient ring. In $\Zp$, it means membership in $p^L\Zp$.
\end{theorem}

\begin{proof}
The difference of the two sums consists exactly of terms with $w(r)\geq L$.
Every such term is divisible by $p^L$, including the zero terms with infinite
weight. Their finite sum has the same divisibility.
\end{proof}

\begingroup\small\raggedright
\noindent\emph{Formal correspondence.}
\lean{Filtered.filtered_triangular_reconstruction} and
\lean{Filtered.padic_filtered_reconstruction}.
For weights $2r$, retaining $r\leq M$ gives the bound $p^{2M+2}$ through
\lean{Filtered.even_weight_tail}. A sharper kernel weight can improve this
bound, as in the classical dyadic reconstruction later in the chapter. Each
coefficient row is finite, so no hypothesis $w(r)\to\infty$ is needed here.
\par\endgroup

\subsection{Finite Mahler coordinates and weighted differences}

For $f:\N\to R$, define the step-$h$ forward difference and its iterates by
\[
\Delta_h f(n)=f(n+h)-f(n),\qquad \Delta_h^0 f=f,
\]
and put $c_r(f)=\Delta_1^r f(0)$. Thus the coefficients are determined by the
differences of $f$.

\begin{theorem}[Finite Newton--Mahler identities]
\label{thm:mahler-foundation}
For every $n,s,h\in\N$,
\[
\begin{aligned}
f(n)&=\sum_{r=0}^{n}c_r(f)\binom nr,\\
c_r(\Delta_1^s f)&=c_{s+r}(f),\\
\Delta_1^s f(n)&=\sum_{j=0}^{n}c_{s+j}(f)\binom nj,\\
\Delta_h^s f(n)&=\sum_{j=0}^{s}(-1)^{s-j}\binom sj f(n+jh).
\end{aligned}
\]
The finite interpolant
$I_Nf(n)=\sum_{r=0}^{N}c_r(f)\binom nr$ agrees with $f$ for $n\leq N$.
It has the same order-$s$ unit-step difference at $n$ whenever $n+s\leq N$.
For arbitrary coefficients $c_0,\ldots,c_N$, the corresponding binomial
interpolant has exactly those coefficients and zero coefficients above $N$.
\end{theorem}

\begin{proof}
Write $E_h f(n)=f(n+h)$ for the shift operator. The unit shift is
$I+\Delta_1$. Expanding its $n$th power and evaluating
at zero gives reconstruction. Composing difference operators gives the
coefficient shift and applying reconstruction to $\Delta_1^s f$ gives the
third identity. The binomial expansion of $(E_h-I)^s$ gives the explicit
stencil. Binomial coefficients vanish when their lower index exceeds $n$,
which proves the interpolation claim. Equality of every point of the stencil
proves equality of differences. Repeated finite differences of $\binom nr$
at zero give the exact coefficient extraction.
\end{proof}

\begingroup\small\raggedright
\noindent\emph{Formal correspondence.}
\lean{Mahler.finite_mahler_reconstruction},
\lean{Mahler.coefficient_difference},
\lean{Mahler.difference_from_coefficients},
\lean{Mahler.difference_explicit},
\lean{Mahler.finite_sample_difference}, and
\lean{Mahler.coefficient_interpolant}.
Finite samples determine the stencils contained in the sample interval;
a universal tail estimate requires additional information.
\par\endgroup

For $x\in\Zp$, extend the valuation to zero by setting
\[
\operatorname{ord}_p(x)=
\begin{cases}v_p(x),&x\ne0,\\\infty,&x=0.\end{cases}
\]
Thus $L\leq\operatorname{ord}_p(x)$ is equivalent to $x\in p^L\Zp$.

\begin{theorem}[Weighted difference bound]
\label{thm:weighted-mahler-foundation}
Let $f:\N\to\Zp$ and let $w:\N\to\N\cup\{\infty\}$ satisfy
$w(s+j)\leq\operatorname{ord}_p(c_{s+j}(f))$ for $0\leq j\leq n$.
Then
\[
\operatorname{ord}_p(\Delta_1^s f(n))\geq
\min_{0\leq j\leq n}
\left\{w(s+j)+v_p\binom nj\right\}.
\]
\end{theorem}

\begin{proof}
Apply the nonarchimedean inequality to the finite coefficient formula in
Theorem~\ref{thm:mahler-foundation}. For $0\leq j\leq n$, the integer
$\binom nj$ is positive. The order of its product with $c_{s+j}(f)$ is the
sum of the two orders, including the case where that coefficient is zero.
\end{proof}

\begingroup\small\raggedright
\noindent\emph{Formal correspondence.}
\lean{PadicBounds.nat_le_order_iff},
\lean{PadicBounds.weighted_difference_bound}, and
\lean{PadicBounds.weighted_difference_divisibility}.
An expansion on all of $\Zp$ requires a convergence theorem in addition to
these finite algebraic identities and bounds.
\par\endgroup

\subsection{Unit modes and the dyadic distinction}

\begin{theorem}[Differences of unit modes]
\label{thm:unit-modes-foundation}
Over every commutative ring,
\[
\Delta_h^s(u^n)=u^n(u^h-1)^s.
\]
For any prime $p$, units $u_i\in\Zp^\times$, coefficients $a_i\in\Zp$, and
a finite index set $I$, the sequence $f(n)=\sum_{i\in I}a_i u_i^n$ satisfies
\[
\Delta_{p-1}^s f(n)\in p^s\Zp
\qquad(n,s\in\N).
\]
\end{theorem}

\begin{proof}
Induction gives the geometric identity. Reduction of a unit to $\F_p^\times$
gives $u_i^{p-1}=1$ modulo $p$. Taking $s$ powers supplies $p^s$ in each
summand. The difference operator is linear over finite sums.
\end{proof}

\begingroup\small\raggedright
\noindent\emph{Formal correspondence.}
\lean{UnitModes.geometric_difference},
\lean{UnitModes.unit_fermat_divisibility}, and
\lean{UnitModes.finite_unit_modes_divisibility}.
The result includes $p=2$, where the displayed step is $1$. Its scope here
is finite sums of modes.
\par\endgroup

\begin{theorem}[Finite unit coordinates]
\label{thm:finite-unit-foundation}
Let $p$ be prime, $L\geq0$, and
$a\in(\Z/p^{L+1}\Z)^\times$. Define
\[
\omega_L(a)=a^{p^L},\qquad \langle a\rangle_L=a\omega_L(a)^{-1}.
\]
Then
\[
a=\omega_L(a)\langle a\rangle_L,\qquad
\omega_L(a)^{p-1}=1,\qquad \langle a\rangle_L\equiv1\pmod p.
\]
The two factors are uniquely determined by these properties and are compatible
with reduction to lower precision.
For $a\in(\Z/2^{L+2}\Z)^\times$, there is instead a unique decomposition
\[
a=\epsilon(a)\langle a\rangle_4,
\qquad \epsilon(a)\in\{1,-1\},\qquad
\langle a\rangle_4\equiv1\pmod4.
\]
\end{theorem}

\begin{proof}
The unit-group exponent $p^L(p-1)$ gives the root identity for $\omega_L$.
Its residue is that of $a$, so the remaining factor has residue $1$.
If two $(p-1)$st roots have equal residue, raising them to $p^L$ makes them
equal modulo $p^{L+1}$, while each root is fixed by this exponent. This proves
uniqueness and then compatibility under reduction. In the dyadic construction,
an odd residue is exactly $1$ or $-1$ modulo $4$. Choose that sign and multiply
$a$ by it to obtain the principal factor. The two signs remain distinct
modulo $4$, which proves uniqueness.
\end{proof}

\begingroup\small\raggedright
\noindent\emph{Formal correspondence.}
\lean{FiniteUnits.finite_unit_decomposition},
\lean{FiniteUnits.decomposition_unique},
\lean{UnitRefinement.omega_compatible},
\lean{UnitRefinement.principal_compatible}, and
\lean{DyadicUnits.dyadic_unit_decomposition} with
\lean{DyadicUnits.decomposition_unique}.
The finite $p=2$ version of the first construction has trivial root-of-unity
factor. The separate modulo-$4$ construction retains the sign that reduction
modulo $2$ cannot detect.
\par\endgroup

\subsection{Recurrence packets and finite arithmetic synchronization}

\begin{theorem}[Companion jumps and packet recurrences]
\label{thm:companion-foundation}
Let $C$ be a $d\times d$ matrix over a commutative ring and
$x_n=C^n x_0$. For $N,j,k\in\N$, put $b_k=x_{Nk+j}$. Then
\[
b_k=C^j(C^N)^kx_0,\qquad b_{k+1}=C^Nb_k.
\]
If $\chi_{C^N}(X)=X^d+\sum_{r<d}q_rX^r$, then
\[
b_{k+d}=-\sum_{r=0}^{d-1}q_r b_{k+r}.
\]
Every coordinate and every linear scalar observable satisfies the same
recurrence. No invertibility or diagonalization hypothesis is required.
\end{theorem}

\begin{proof}
The jump identity follows from the power law for one matrix. Apply the
Cayley--Hamilton identity $\chi_{C^N}(C^N)=0$ to the state $b_k$.
The leading coefficient is $1$, so the final state can be isolated. A linear
observable preserves the resulting finite sum.
\end{proof}

\begingroup\small\raggedright
\noindent\emph{Formal correspondence.}
\lean{Companion.companion_jump},
\lean{Companion.packet_charpoly_recurrence}, and
\lean{Companion.packet_observable_recurrence}.
The source-sequence versions derive the packet recurrence from the relation
$x_{n+1}=Cx_n$.
For the Fibonacci matrix $C=\left(\begin{smallmatrix}1&1\\1&0\end{smallmatrix}\right)$,
its jump trace is the Lucas number $L_N$ and its determinant is $(-1)^N$.
Thus
\[
F_{N(k+2)+j}=L_NF_{N(k+1)+j}-(-1)^NF_{Nk+j}.
\]
This specialization is
\lean{FibonacciPackets.fibonacci_packet_recurrence}; it is derived from the
Fibonacci state.
\par\endgroup

\begin{theorem}[Dyadic--odd packet factorization]
\label{thm:crt-foundation}
Let $a\geq0$ and let $M$ be a positive odd integer. The Chinese remainder map
is a ring isomorphism
\[
\Z/(2^aM)\Z\simeq\Z/2^a\Z\times\Z/M\Z.
\]
For $0\leq r<2^aM$ and any coefficient sequence $f$,
\[
P_{2^aM,r}f
=P_{2^a,r\bmod 2^a}\,P_{M,r\bmod M}f.
\]
The same identity holds for formal-series coefficient selectors.
\end{theorem}

\begin{proof}
The two moduli are coprime. An index has the prescribed combined residue
exactly when it has both component residues. Applying this equivalence to
the indicator defining each selector proves the formula.
\end{proof}

\begingroup\small\raggedright
\noindent\emph{Formal correspondence.}
\lean{CRTPackets.addressEquiv}, \lean{CRTPackets.residue_equivalence},
\lean{CRTPackets.packet_factorization}, and
\lean{CRTPackets.series_packet_factorization}.
This theorem synchronizes finite arithmetic addresses. The passage from a
family of $p$-adic objects to a global integer object is a separate question,
discussed in Section~\ref{subsec:native133}.
\par\endgroup

\subsection{Polyphase notation retained from the recovered manuscript}
\label{subsec:polyphase-exposition}

The recovered Annals text develops a broader polyphase vocabulary. We explain
it through elementary identities with explicit algebraic derivations. These
identities have not been separately formalized in Lean for this consolidation.
For a sequence $a$ and $M\geq1$, let $a_r^{[M]}(q)=a_{Mq+r}$ for $0\leq r<M$.
If $A(z)=\sum_n a_nz^n$ and $A_r(w)=\sum_q a_{Mq+r}w^q$ are ordinary generating
functions, then unique quotient and remainder give
\[
A(z)=\sum_{r=0}^{M-1}z^rA_r(z^M),\qquad
\bigl(a_r^{[M]}\bigr)_s^{[N]}=a_{r+Ms}^{[MN]}.
\]
For ordinary Cauchy multiplication, the phase-$r$ product has a term without a
carry and a term with one carry:
\[
C_r(w)=\sum_{u+v=r}A_u(w)B_v(w)
+w\sum_{u+v=r+M}A_u(w)B_v(w),\qquad 0\leq u,v<M.
\]
Adding the two phase addresses gives the formula. It describes ordinary
Cauchy multiplication; Hurwitz convolution has its own binomial weights.

For $M\geq1$ and $s,r,k\in\N$, define the affine Mahler kernel by
\[
\Theta_{M,s}(r,k)=[z^r](1+z)^s\bigl((1+z)^M-1\bigr)^k
=\sum_{j=0}^{k}(-1)^{k-j}\binom kj\binom{Mj+s}{r}.
\]
It vanishes unless $k\leq r\leq Mk+s$. For a finite binomial expansion
$f(n)=\sum_{r=0}^{D}c_r\binom nr$, the $k$th Newton coefficient of
$g(q)=f(Mq+s)$ is therefore
\[
c_k(g)=\sum_{r=k}^{\min(D,Mk+s)}c_r\Theta_{M,s}(r,k).
\]
Indeed, insert the finite stencil for $\Delta_1^k g(0)$ and interchange its
two finite sums. Coefficient extraction of the binomial expansion of
$((1+z)^M-1)^k$ gives the kernel. In the binary zero-address case,
\[
\Theta_{2,0}(k+j,k)=2^{k-j}\binom kj\qquad(0\leq j\leq k),
\]
since $(2z+z^2)^k=z^k(2+z)^k$.
Exact valuations for the general kernel remain to be determined. The
recovered infinite analytic extensions and stronger historical identities
lie beyond these finite derivations; their status is discussed in
Section~\ref{subsec:historical-scope}.

\section{Congruential Euler packets and their arithmetic}
\label{sec:arithmetic}

Two related constructions run through this section.  The congruential
Euler family is indexed by a modulus $N$ and an address $j$.  Its
exceptional dyadic member has $N=4$ and four addresses.  The classical
signed Euler numbers arise from $N=2$, $j=0$.  Keeping these indices
visible distinguishes the four-packet period theorems from the classical
edge truncations.

\subsection{Original coefficients, normalized packets, and precision}

\begin{definition}[Congruential Euler coefficients]\label{def:arith-euler}
For $N\ge1$ and $0\le j<N$, define rational coefficients by the formal identity
\begin{equation}\label{eq:arith-euler-egf}
 \sum_{d\ge0}E_d^{(N,j)}\frac{t^d}{d!}
 =\left(\sum_{n\ge0}\frac{t^{Nn}}{(Nn+j)!}\right)^{-1}.
\end{equation}
The reciprocal exists because its denominator has constant coefficient
$1/j!\ne0$.  Set
\[
 U_{N,j}(n)=E_{Nn}^{(N,j)},\qquad C_{N,j}(n)=\frac{U_{N,j}(n)}{j!}.
\]
Thus $C_{N,j}(0)=1$ and $U_{N,j}(0)=j!$.  All coefficients outside the
multiples of $N$ vanish.
\end{definition}

The normalized source and its inverse are
\begin{equation}\label{eq:arith-kernel}
 K_{N,j}(t)=j!\sum_{n\ge0}\frac{t^{Nn}}{(Nn+j)!},\qquad
 \mathcal P_{N,j}(t)=\sum_{n\ge0}C_{N,j}(n)\frac{t^{Nn}}{(Nn)!}.
\end{equation}
Their product is one.  Clearing factorials in its coefficient of $t^{Nn}$
gives the triangular system
\begin{equation}\label{eq:arith-triangular}
 \sum_{q=0}^{n}\binom{Nn+j}{Nq}C_{N,j}(q)=\delta_{n0}.
\end{equation}
In particular, for $n>0$,
\[
 C_{N,j}(n)=-\binom{Nn+j}{Nn}^{-1}
       \sum_{q=0}^{n-1}\binom{Nn+j}{Nq}C_{N,j}(q).
\]
The diagonal binomial coefficient is positive, so this recurrence determines
the sequence uniquely.  The Lean construction begins with this recurrence
and derives the reciprocal identity from it.  The resulting sequence is
then identified with the degree-indexed coefficients of
\eqref{eq:arith-euler-egf}, so both descriptions refer to the same family.

The relevant declarations are
\lean{KomorebiPaperIX.Congruential.triangular_inverse},
\lean{KomorebiPaperIX.Congruential.packet_is_actual_inverse}, and
\lean{KomorebiPaperIX.NishibuchiCompletion.generalizedEuler_at_packet}.
The cleared identity with the original positive-address exponential section
is \lean{KomorebiPaperIX.NishibuchiCompletion.original_generating_identity}.

\begin{proposition}[A denominator bound for every address]
\label{prop:arith-denominator}
For $N\ge1$, $0\le j<N$, and $n\in\N$, set
\[
 D_{N,j}(n)=\prod_{r=1}^{n}\binom{Nr+j}{Nr},\qquad D_{N,j}(0)=1.
\]
Then $D_{N,j}(n)C_{N,j}(n)\in\Z$.
\end{proposition}

\begin{proof}
Proceed by induction in the recurrence \eqref{eq:arith-triangular}.
For every $q<n$, the product $D_{N,j}(q)$ divides $D_{N,j}(n-1)$.
Multiplying the recurrence by $D_{N,j}(n-1)$ therefore makes every
previous term integral.  The diagonal factor completes the product
$D_{N,j}(n)$, proving the claim.  At $n=0$, both factors are one.
This gives a uniform denominator bound without requiring integrality of
the normalized coefficients themselves.  The Lean statements are
\lean{KomorebiPaperIX.Congruential.clearing_product} and
\lean{KomorebiPaperIX.Congruential.universal_denominator_clearing}.
\end{proof}

\begin{example}[The Bernoulli boundary address]\label{ex:arith-bernoulli}
At $N=2$, the address $j=1$ gives a second classical family.  Define the
Bernoulli numbers by
\[
 B(z)=\frac{z}{e^z-1}=\sum_{r\ge0}B_r\frac{z^r}{r!},
 \qquad B_0=1,\quad B_1=-\frac12.
\]
All quotients here are regular formal series.  The identity
\[
 2B(z)-B(2z)=\frac{z}{\sinh z}
            =\sum_{n\ge0}C_{2,1}(n)\frac{z^{2n}}{(2n)!}
\]
follows by putting the two fractions on a common denominator and using
\eqref{eq:arith-euler-egf}.  Comparing even coefficients yields
\[
 C_{2,1}(n)=(2-2^{2n})B_{2n}\qquad(n\in\N).
\]
This classical derivation restores the boundary example from the original
manuscript.  No dedicated Lean bridge for this Bernoulli specialization
was identified in the preserved sources.
\end{example}

\begin{definition}[Differences and valuation conventions]\label{def:arith-precision}
For a sequence $f$ and $h,k,n\in\N$, put
\[
 \Delta_h f(n)=f(n+h)-f(n),\qquad
 \Delta_h^0 f(n)=f(n),\qquad
 \Delta_h^{k+1}f(n)=\Delta_h^k f(n+h)-\Delta_h^k f(n).
\]
We use the extended valuation $v_p(0)=+\infty$ in mathematical statements.
Thus $v_p(x)\ge L$, including negative integers $L$, means
$x=0$ or the ordinary valuation of the nonzero rational $x$ is at least $L$.
For integral coefficients and $L\ge0$, it is equivalent to
$x\in p^L\Z_p$.  In expressions involving a step multiplier we write
$\nu_p(u)=v_p(u)$ for $u>0$ and set $\nu_p(0)=0$.  This finite convention
makes the displayed exponents meaningful at step zero, where all
positive-order differences vanish.
\end{definition}

The valuation convention at zero is encoded by
\lean{KomorebiPaperIX.GlobalEulerBoundary.precision}.  For $p$-adic
congruences, divisibility is interpreted in $\Z_p$; statements over other
coefficient rings specify the ring explicitly.  Field divisibility in
$\Q_p$ alone does not encode precision.

\subsection{The standard address window and its common reduction}

The addresses below the prime have unit diagonal coefficients.  In this
window the normalized packets are integral, and their reductions agree
at every row.

\begin{proposition}[Lucas reduction and address independence]
\label{prop:arith-lucas-window}
Let $p$ be prime, $m>0$, and $N=mp$.  For $0\le j<p$ and $n,q\in\N$,
\[
 \binom{mpn+j}{mpq}\equiv\binom{mn}{mq}\pmod p.
\]
Every $C_{mp,j}(n)$ belongs to $\Z_p$.  Its reduction is the unique
sequence $\overline C_m:\N\to\F_p$ satisfying
\[
 \sum_{q=0}^{n}\binom{mn}{mq}\overline C_m(q)=\delta_{n0}
 \quad\text{in }\F_p.
\]
Consequently, for any two addresses $0\le j,k<p$,
\[
 C_{mp,j}(n)\equiv C_{mp,k}(n)\pmod p\qquad(n\in\N).
\]
\end{proposition}

\begin{proof}
Lucas reduction separates the lowest base-$p$ digits, which are $j$ and
zero in the upper and lower arguments.  Their binomial contribution is
one, leaving the displayed reduced row.  On the diagonal the reduced
coefficient is one.  The recurrence \eqref{eq:arith-triangular} therefore
constructs an integral packet by successive division by units in $\Z_p$.
After reduction, every address satisfies the same triangular system over
$\F_p$.  Its diagonal entries are one, so its solution is unique.  This
proves both the common reduction and the congruence between addresses.
The row reduction is
\lean{KomorebiPaperIX.StandardWindow.lucas_row_reduction}.
The integral identification and the two reduction statements are
\lean{KomorebiPaperIX.WindowIndependence.standard_window_integral},
\lean{KomorebiPaperIX.WindowIndependence.coeffPadic_coe},
\lean{KomorebiPaperIX.WindowIndependence.residue_independence}, and
\lean{KomorebiPaperIX.WindowIndependence.two_window_addresses_congruent}.
\end{proof}

\subsection{The full Nishibuchi statement and its starting bound}

\begin{theorem}[Nishibuchi congruence]\label{thm:arith-nishibuchi}
Let $p$ be prime, let $m>0$, and suppose that $m\mid p-1$ or $m=2$.
Put $N=mp$ and $q=\operatorname{lcm}(2,p-1)$.  For every $0\le j<N$,
all $U_{N,j}(n)$ belong to $\Z_p$.  For $r\ge1$ and $n\ge r$,
\begin{equation}\label{eq:arith-nishibuchi}
 E_{Nn+qp^r}^{(N,j)}\equiv E_{Nn}^{(N,j)}\pmod{p^r}.
\end{equation}
Thus the positive sufficient starting index $n_0=r$ works uniformly.
If $m\mid p-1$, the stronger sufficient condition $r\le mn$ is valid.
\end{theorem}

This theorem covers the full parameter range and displacement of Conjecture
3.9 in the fixed version of Nishibuchi's paper used for the statement
comparison.  The corresponding Lean theorem is
\lean{KomorebiPaperIX.NishibuchiCompletion.conjecture_3_9}.
The stronger divisor-branch threshold is
\lean{KomorebiPaperIX.NishibuchiCompletion.divisor_period}.
The starting index is sufficient uniformly, although it need not be the
earliest one for a given tuple of parameters.

We first prove the difference estimates on the divisor branch, then derive
the period from a finite Newton expansion.

\begin{theorem}[All-address differences on the divisor branch]
\label{thm:arith-divisor-base}
Let $p$ be prime, $m>0$, $m\mid p-1$, $N=mp$, and $0\le j<N$.
Write $h=(p-1)/m$.  For every $u,k,n\in\N$,
\begin{align}
 U_{N,j}(n)&\in\Z_p,\label{eq:arith-divisor-integral}\\
 \Delta_{uh}^{k}U_{N,j}(n)&\in p^{\min(k,mn)}\Z_p,\label{eq:arith-divisor-bound}\\
 v_p\bigl(\Delta_{uh}^{k}C_{N,j}(n)\bigr)
   &\ge\min(k,mn)-v_p(j!).\label{eq:arith-divisor-normalized}
\end{align}
The factorial correction is part of the normalized statement, including
addresses $j\ge p$.
\end{theorem}

\begin{proof}
Work in the Hurwitz integral ring of formal series
\[
 \mathcal H_p=\left\{F\in\Q_p[[t]]:
            d![t^d]F\in\Z_p\text{ for every }d\right\}.
\]
For $D=d/dt$, let $B_p=(D^{p^2}-D^p)/p$.  Call $F$ flat if
$B_p^sF\in\mathcal H_p$ for every $s\ge0$.
We will prove that $j!\mathcal P_{N,j}$ is flat at every address.  The main
difficulty is to include the addresses at which the usual unit-pivot
argument fails.

Set
\[
 H_c(t)=\sum_{n\ge0}\frac{t^{Nn+c}}{(Nn+c)!},\qquad
 A_c(t)=\sum_{n\ge0}\frac{t^{Nn}}{(Nn+c)!},\qquad
 R_c(t)=\frac{A_{c-1}(t)}{A_c(t)}\quad(1\le c<N).
\]
The denominator of this last quotient has nonzero constant coefficient.
The section identities give
\[
 R_cH_c=tH_{c-1},\qquad R_c(0)=c,\qquad
 j!\mathcal P_{N,j}=\mathcal P_{N,0}\prod_{c=1}^{j}R_c.
\]
Cancellation of consecutive source sections proves the product identity
and ties the ratios directly to the family under study.

At the singular addresses $c=ap$, $1\le a<m$, use the weighted ring
\[
 \mathcal W_p=\left\{\sum_{d\ge0}a_dT^d:
          p^{\lfloor d/p\rfloor}a_d\in\Z_p\right\}.
\]
The key estimate is $R_{ap}(\log(1+T))\in\mathcal W_p$.  Its arithmetic
input comes from the cyclic group ring $\Z[\Z/N\Z]$.
For its canonical element $x$ and $I=(p,x^{p-1}-1)$, the relation
$(x^{p-1})^p=1$ implies $I^p\subseteq(p)I$.  Induction yields
$I^q\subseteq(p^{v_p(q!)})$.  Grouping the falling factorial
$(x)_d=\prod_{i=0}^{d-1}(x-i)$ into consecutive blocks of length $p$
then gives an integral residue-coordinate factor
$p^{v_p(\lfloor d/p\rfloor!)}$.  The identity
\[
 v_p(d!)=\lfloor d/p\rfloor+v_p(\lfloor d/p\rfloor!)
\]
accounts for precisely the weight in $\mathcal W_p$ after dividing by $d!$.
Comparison of the section and group-ring recurrences identifies these
coordinates with the logarithmically substituted section.

To pass from the section to the singular ratio, factor out its initial zero:
for $c=ap$, the section has the form
\[
 H_c(\log(1+T))=\frac{T^c}{c!}\,V(T),\qquad V(0)=1,
 \qquad V\in\mathcal W_p.
\]
The weight of $c!$ compensates for the coefficient shift by $c$.
The weighted ring is closed under inversion of series with constant
coefficient one and under $L_T=\log(1+T)(1+T)d/dT$.  With
$\ell(T)=\log(1+T)/T$, interpreted by coefficient shifting, differentiation
of the displayed factorization gives
\[
 R_c(\log(1+T))=c(1+T)\ell(T)+(L_TV)V^{-1}.
\]
Since $p\ell(T)\in\mathcal W_p$ and $p\mid c$, both terms are weighted.
This is the singular estimate
\lean{KomorebiPaperIX.SingularRatioAmice.paper_singular_amice_estimate}.

To deduce flatness from the weighted estimate, use
$\Theta=(1+T)d/dT$.  Both $\Theta$ and
$(\Theta^{p^2}-\Theta^p)/p$ preserve $\mathcal W_p$.
For the first operator, its coefficient formula is
$(d+1)a_{d+1}+da_d$; whenever the floor weight rises, the factor $d+1$
provides the missing $p$.  For the second operator, the proof uses its
full Leibniz correction and the scalar congruence
$p^2\mid b^{p^2}-b^p$.  Weighted monomials factor as
\[
 a_dT^d=(p^{\lfloor d/p\rfloor}a_d)
          (T^p/p)^{\lfloor d/p\rfloor}T^{d\bmod p}.
\]
The operator preserves these monomials and hence finite weighted
polynomials.  Each fixed output coefficient depends on only finitely many
input coefficients, which extends the result to the formal series.
Substitution $T=e^t-1$ intertwines $D$ with $\Theta$ and maps
$\mathcal W_p$ into $\mathcal H_p$, because
\[
 d![t^d]F(e^t-1)=[T^0]\Theta^dF.
\]
It follows that the singular ratios are flat.  This transfer is proved in
\lean{KomorebiPaperIX.WeightedDworkTransfer.singular_ratio_flat}.

At a nonsingular address $2\le c<N$ with $p\nmid c$, differentiation of the section
identities gives
\[
 tR_c'=R_c(1+R_{c-1}-R_c).
\]
Writing $q=R_{c-1}-(c-1)$ and $y=R_c-c$, this becomes
\[
 (tD+c)y=cq+yq-y^2.
\]
On degrees divisible by $p$, the inverse of $tD+c$ divides the degree-$d$
coefficient by the unit $d+c$.  Starting with $y_0=0$, form
\[
 y_{s+1}=(tD+c)^{-1}(cq+y_sq-y_s^2).
\]
Flatness and support are preserved. Indeed, put $E_c=tD+c$ and $C_p=pD^{p^2}-D^p$.
The commutator is $[B_p,E_c]=C_p$, and iteration gives
\[
 B_p^rE_cy=E_cB_p^ry+rC_pB_p^{r-1}y.
\]
For the inverse equation $E_cy=f$, induction on $r$ makes the
correction term Hurwitz-integral. Both operators preserve support in
degrees divisible by $p$. On that support the degree multiplier $d+c$
is a unit, so integrality of $E_cB_p^ry$ implies integrality of
$B_p^ry$. Thus the inverse preserves full flatness, as well as support.
 Since the difference of two right
sides factors as $(z-y)(q-z-y)$ and its second factor has zero constant
term, agreement with the solution gains one coefficient at each iteration.
Consequently $[t^d]y_s=[t^d]y$ for $d\le s$.
Every coefficient of $B_p^r y$ can therefore be read from a sufficiently
long finite iterate.  The argument is entirely formal.
This is the content of
\lean{KomorebiPaperIX.RiccatiFlatPropagation.riccati_solution_flat}.

The standard window supplies the initial ratio $R_1$.  The weighted
estimate handles $p\mid c$, and Riccati propagation handles $p\nmid c$.
Flatness is closed under products, so the exact factorization establishes
flatness of $j!\mathcal P_{N,j}$.  Extracting its supported Hurwitz coefficients
and applying the proved saturation and step-multiple lemmas yields
\eqref{eq:arith-divisor-integral}--\eqref{eq:arith-divisor-bound}.

Here is the coefficient extraction explicitly. If $a_d=d![t^d]F$,
then
\[
 d![t^d]B_p^kF=p^{-k}\Delta_{p(p-1)}^ka_{d+pk}.
\]
For $F=j!\mathcal P_{N,j}$ and $k\le mn$, take $d=mpn-pk$.
Support in multiples of $mp$ turns this into
$p^{-k}\Delta_h^kU_{N,j}(n)$, so flatness gives divisibility by $p^k$.
For $k>mn$, expand $\Delta_h^{k-mn}$ applied to
$\Delta_h^{mn}U_{N,j}(n)$; every shifted row is at least $n$, so
each summand is divisible by $p^{mn}$. Finally
$\Delta_{uh}=((I+\Delta_h)^u-I)$, whose $k$th power for $u>0$
is an integral combination of powers of $\Delta_h$ of order at least $k$.
The case $u=0$ gives zero for $k>0$, and order zero is the original
integral coefficient. This proves the bound for all $u$ and $k$.
Finally $\Delta_{uh}^kU_{N,j}=j!\Delta_{uh}^kC_{N,j}$ gives
\eqref{eq:arith-divisor-normalized}, with the zero case retained.
The resulting statements are
\lean{KomorebiPaperIX.AllResidueEuler.all_residue_flat},
\lean{KomorebiPaperIX.AllResidueEuler.global_unnormalized_bound}, and
\lean{KomorebiPaperIX.AllResidueEuler.global_normalized_precision}.
\end{proof}

\begin{example}[Why the factorial loss is necessary]
\label{ex:arith-normalization-loss}
Take $p=3$, $m=2$, $N=6$, and $j=3$.  The recurrence gives
\[
 C_{6,3}(1)=-\frac1{84},\qquad
 C_{6,3}(2)=\frac{703}{5460},\qquad
 C_{6,3}(2)-C_{6,3}(1)=\frac{64}{455}.
\]
The last difference has $3$-adic valuation zero.  Here $h=1$, so at
$n=1$ the order-one bound without factorial loss would require valuation
at least one.  The corrected normalized bound is sharp, since
\[
 v_3\bigl(C_{6,3}(2)-C_{6,3}(1)\bigr)=0=1-v_3(3!).
\]
Multiplication by $3!$ gives
\[
 U_{6,3}(2)-U_{6,3}(1)=\frac{384}{455},\qquad
 v_3\bigl(U_{6,3}(2)-U_{6,3}(1)\bigr)=1,
\]
consistent with the unnormalized theorem.  Thus the factorial correction
is already necessary at a single step in the all-address parameter range.
These exact values and the refutation of a universal bound without loss
are formalized in
\lean{KomorebiPaperIX.Normalization.normalized_first},
\lean{KomorebiPaperIX.Normalization.normalized_second},
\lean{KomorebiPaperIX.Normalization.normalized_difference},
\lean{KomorebiPaperIX.Normalization.difference_valuation},
\lean{KomorebiPaperIX.Normalization.factorial_loss_is_sharp},
\lean{KomorebiPaperIX.Normalization.unnormalized_difference},
\lean{KomorebiPaperIX.Normalization.unnormalized_difference_valuation}, and
\lean{KomorebiPaperIX.Normalization.universal_no_loss_refuted}.
\end{example}

\begin{corollary}[A common period in the standard window]
\label{cor:arith-window-period}
Let $p$ be prime, $m>0$, $m\mid p-1$, and $0\le j<p$.  Put
$h=(p-1)/m$.  For every $n\ge1$ and $u\in\N$,
\[
 \overline C_m(n+uh)=\overline C_m(n),\qquad
 C_{mp,j}(n+uh)\equiv C_{mp,j}(n)\pmod p.
\]
Thus $h$ is a sufficient period from row one, common to all addresses
in the standard window.  It need not be the least period.
\end{corollary}

\begin{proof}
Since $j<p$, the factorial $j!$ is a unit in $\Z_p$.  Also $mn\ge1$.
The order-one case of Theorem~\ref{thm:arith-divisor-base} therefore gives
valuation at least one for the normalized difference at every step $uh$.
This proves the congruence, and Proposition~\ref{prop:arith-lucas-window}
identifies its reduction with the equality for $\overline C_m$.
The preserved Lean development also proves this period directly from the
Hurwitz inverse in characteristic $p$:
\lean{KomorebiPaperIX.WindowPeriod.uniform_standard_window_period},
\lean{KomorebiPaperIX.WindowPeriod.multiple_period}, and
\lean{KomorebiPaperIX.WindowPeriod.standard_window_period_divisibility}.
\end{proof}

\begin{lemma}[Prime-power period from finite differences]
\label{lem:arith-newton-period}
Let $R$ be a commutative ring, $p$ a prime, $f:\N\to R$, $h,n\in\N$,
and $r\ge1$.  If
\[
 p^{\min(k,r)}\mid\Delta_h^k f(n)\quad\text{for every }k\in\N,
\]
then $p^r\mid f(n+p^{r-1}h)-f(n)$.
\end{lemma}

\begin{proof}
The finite Newton expansion gives
\[
 f(n+p^{r-1}h)-f(n)
 =\sum_{k=1}^{p^{r-1}}\binom{p^{r-1}}k\Delta_h^k f(n).
\]
Terms with $k\ge r$ have the required factor from the difference.
For $1\le k<r$, the identity
\[
 k\binom{p^{r-1}}k=p^{r-1}\binom{p^{r-1}-1}{k-1}
\]
and $v_p(k)<k$ imply
$p^r\mid\binom{p^{r-1}}k p^k$ in $\Z$.
Mapping that integer divisibility to $R$ proves the claim term by term.
The argument includes $r=1$.
The declaration is
\lean{KomorebiPaperIX.PeriodFromDifferences.prime_power_period}.
\end{proof}

\begin{proof}[Proof of Theorem~\ref{thm:arith-nishibuchi}]
On the divisor branch, $m\mid q$ and the packet step $q/m$ is an integer
multiple of $h=(p-1)/m$.  If $r\le mn$, Theorem~\ref{thm:arith-divisor-base}
supplies the hypotheses of Lemma~\ref{lem:arith-newton-period}.
Its packet displacement is $p^{r-1}q/m$, which converts to the original
degree as follows:
\[
 mp\left(n+p^{r-1}\frac qm\right)=mpn+qp^r.
\]
The choice $n\ge r$ implies $mn\ge r$ because $m\ge1$.

For an odd prime, $2\mid p-1$, so the option $m=2$ introduces no further
case.  The sole remaining branch is $p=2,m=2,N=4$.  The all-row theorem
proved next gives $2^k\mid\Delta_1^kC_{4,j}(n)$ at every row and address.
After multiplying by $j!$, Lemma~\ref{lem:arith-newton-period} gives packet
period $2^{r-1}$ and degree displacement $4\cdot2^{r-1}=2^{r+1}=qp^r$.
It also gives integrality.  Hence the complete statement follows.
The formal degree conversion is
\lean{KomorebiPaperIX.NishibuchiCompletion.original_degree_shift};
the all-parameter degree theorem is
\lean{KomorebiPaperIX.NishibuchiCompletion.full_period_degree}.
\end{proof}

\subsection{All-row dyadic differences and displacement amplification}

For the remainder of the four-packet discussion abbreviate
$C_j=C_{4,j}$ and $U_j=U_{4,j}$.  Removing the original row cap requires
an unshifted coefficient operator.

\begin{theorem}[All-row dyadic higher differences]\label{thm:arith-dyadic-differences}
For $0\le j<4$ and all $u,k,n\in\N$,
\begin{align}
 C_j(n)&\in\Z_2,\\
 \Delta_u^kC_j(n)&\in 2^{k(1+\nu_2(u))}\Z_2,\label{eq:arith-allrow-c}\\
 \Delta_u^kU_j(n)&\in
 2^{k(1+\nu_2(u))+v_2(j!)}\Z_2.\label{eq:arith-allrow-u}
\end{align}
There is no restriction $k\le n$ and no minimum with the row $n$.
\end{theorem}

\begin{proof}
The Hurwitz coefficients of the normalized source $K_{4,j}$ at degrees
$4n$ are
\[
 g_j(n)=\binom{4n+j}{4n}^{-1}.
\]
Their explicit expressions are
\[
 \begin{aligned}
 g_0(n)&=1,&
 g_1(n)&=\frac1{4n+1},\\
 g_2(n)&=\frac1{(4n+1)(2n+1)},&
 g_3(n)&=\frac3{(4n+1)(2n+1)(4n+3)}.
 \end{aligned}
\]
Every denominator is odd.  An affine sequence with even slope satisfies
$2^k\mid\Delta_1^k f(n)$ for every $k,n$.  The exact product identity
\[
 \Delta_1^s(fg)(n)=\sum_{i=0}^s\binom si
          \Delta_1^if(n)\Delta_1^{s-i}g(n+i)
\]
proves closure of this property under products.  Applying the formula
with the factors interchanged to $fg=1$ isolates $f(n+s)\Delta_1^sg(n)$.
Multiplication by $f(n+s)^{-1}$ and induction then prove closure under
pointwise unit inverses.
Thus all four $g_j$ have the property.  The explicit formulas and difference
bounds are formalized in
\lean{KomorebiPaperIX.DyadicFourKernel.source_coe} and
\lean{KomorebiPaperIX.DyadicFourKernel.source_kummer}.

For an integer charge $c$, put
\[
 A_c=\frac{D^4-c}{2}.
\]
A series is flat of charge $c$ when applying any power $A_c^s$ to it gives
an element of $\mathcal H_2$.  The required product formula is
\begin{equation}\label{eq:arith-charged-product}
 A_{c+d}(FG)=(A_cF)G+F(A_dG)
       +2F'G'''+3F''G''+2F'''G'.
\end{equation}
All middle coefficients are integral.  Derivatives preserve
$\mathcal H_2$ and commute with $A_c$.  Induction on the number of
iterations in \eqref{eq:arith-charged-product} proves that multiplication
adds charges.  If $FG=1$, the same formula yields
\[
 A_{-c}G=-G^2A_cF-G(2F'G'''+3F''G''+2F'''G').
\]
Induction, starting from the independently established integrality of
the reciprocal $G$, proves that inversion negates the charge.
Finally $A_{c+2}=A_c-1$, so increasing a charge by two preserves flatness.
These are
\lean{KomorebiPaperIX.ChargedDworkClosure.B_mul},
\lean{KomorebiPaperIX.ChargedDworkClosure.flat_inverse}, and
\lean{KomorebiPaperIX.ChargedDworkClosure.flat_charge_shift}.

Packing $g_j$ into degrees divisible by four gives $K_{4,j}$.
Its source property is precisely charge-one flatness.  Inversion gives
charge $-1$, and the charge shift restores charge one for the
reciprocal $\mathcal P_{4,j}$.  If $\operatorname{dec}(F)(d)=d![t^d]F$, the exact
coefficient identity is
\[
 \operatorname{dec}(A_1^sF)(d)
       =2^{-s}\Delta_4^s\operatorname{dec}(F)(d).
\]
At $d=4n$, packing turns the right side into
$2^{-s}\Delta_1^sC_j(n)$.  This extraction requires no negative degree
and therefore holds at every $n$, including zero.  It proves the unit-step
bound \lean{KomorebiPaperIX.DyadicFourAllRows.base_step_bound}.

To incorporate $v_2(u)$, expand a displacement using finite Newton
coefficients.  For any prime $p$, $u>0$, and $1\le a\le u$,
\[
 p^{v_p(u)+1}\mid\binom ua p^a,
\]
as follows from $a\binom ua=u\binom{u-1}{a-1}$ and $v_p(a)<a$.
Repeated displacement, retaining the remaining base-difference order
at each induction step, gives a factor $p^{k(1+v_p(u))}$ from $k$
displacements.  This induction is formalized in
\lean{KomorebiPaperIX.StepAmplification.repeated_precision}.
Applying it with $p=2$ proves \eqref{eq:arith-allrow-c}; multiplication
by $j!$ proves \eqref{eq:arith-allrow-u}.  Zero step and zero order follow
from their definitions.  The final statements are
\lean{KomorebiPaperIX.DyadicFourFinalStrengthening.normalized_step_bound}
and
\lean{KomorebiPaperIX.DyadicFourFinalStrengthening.unnormalized_step_bound}.
\end{proof}

\begin{corollary}[Step-sensitive divisor-branch bound]\label{cor:arith-amplified}
Under the hypotheses of Theorem~\ref{thm:arith-divisor-base}, for all
$u,k,n\in\N$,
\[
 \Delta_{u(p-1)/m}^{k}U_{mp,j}(n)
       \in p^{\min(k(1+\nu_p(u)),mn)}\Z_p.
\]
For $C_{mp,j}$ the corresponding valuation lower bound is
$\min(k(1+\nu_p(u)),mn)-v_p(j!)$.
\end{corollary}

\begin{proof}
Apply the same repeated-displacement theorem with the fixed cap $mn$ to
\eqref{eq:arith-divisor-bound}.  The argument preserves this cap and
improves the contribution from the order of the difference.
Division by the fixed factorial gives the rational valuation statement,
with the zero case included.
These are
\lean{KomorebiPaperIX.AllResidueStepAmplification.unnormalized_precision}
and
\lean{KomorebiPaperIX.AllResidueStepAmplification.normalized_precision}.
\end{proof}

\begin{corollary}[Dyadic periods from row zero]\label{cor:arith-allrow-period}
For $0\le j<4$, $r\ge1$, and $n\ge0$,
\[
 U_j(n+2^{r-1})\equiv U_j(n)\pmod{2^r}.
\]
For $j=2,3$ and $r\ge2$, $2^{r-2}$ is also a sufficient packet period.
The respective original degree displacements are $2^{r+1}$ and $2^r$.
\end{corollary}

\begin{proof}
Use Lemma~\ref{lem:arith-newton-period} and the unit-step case of
Theorem~\ref{thm:arith-dyadic-differences}.  At $j=2,3$ the factorial
contributes an additional power of two.  The original degree versions are
\lean{KomorebiPaperIX.DyadicFourFinalStrengthening.all_rows_period_degree}
and
\lean{KomorebiPaperIX.DyadicFourFinalStrengthening.all_rows_half_period_degree}.
The next theorem sharpens these sufficient periods to a complete
classification.
\end{proof}

\subsection{Exact displacement valuations and least packet periods}

\begin{lemma}[Finite Newton dominance]\label{lem:arith-dominance}
Let $f:\N\to\Q$ and write $a_k=\Delta_1^kf(0)$.  Suppose that $c\in\N$,
$a_1\ne0$, $v_2(a_1)=c$, and
\[
 v_2(a_k)\ge c+1+\lfloor\log_2k\rfloor\qquad(k\ge2).
\]
Then for every $n\ge0$ and $T>0$,
\[
 f(n+T)-f(n)\ne0,\qquad
 v_2\bigl(f(n+T)-f(n)\bigr)=v_2(T)+c.
\]
\end{lemma}

\begin{proof}
Finite Newton reconstruction at an arbitrary row gives
\[
 \Delta_1^kf(n)=\sum_{i=0}^{n}\binom ni a_{k+i}.
\]
The monotonicity of $\lfloor\log_2k\rfloor$ therefore preserves the
higher-coefficient bound after translation.  In addition,
$\Delta_1f(n)-a_1\in2^{c+1}\Z_2$, so the first difference is nonzero
of exact valuation $c$.

Now expand the displacement:
\[
 f(n+T)-f(n)=T\Delta_1f(n)
          +\sum_{k=2}^{T}\binom Tk\Delta_1^kf(n).
\]
For $1\le k\le T$, the binomial denominator identity implies
\[
 v_2\binom Tk\ge v_2(T)-v_2(k)
            \ge v_2(T)-\lfloor\log_2k\rfloor.
\]
Every term of the remaining sum consequently has valuation at least
$v_2(T)+c+1$.  Its sum cannot cancel the leading term, which has exact
valuation $v_2(T)+c$.  This proves nonvanishing and the claimed valuation
using only finite sums.
The general theorem is
\lean{KomorebiPaperIX.NewtonDominance.exact_displacement}.
\end{proof}

\begin{theorem}[Exact four-packet displacement law]\label{thm:arith-exact}
Set
\[
 (c_0,c_1,c_2,c_3)=(1,1,4,2),\qquad
 (d_0,d_1,d_2,d_3)=(1,1,5,3).
\]
For $0\le j<4$, $n\ge0$, and every positive integer $T$,
\begin{align*}
 C_j(n+T)-C_j(n)&\ne0,&
 v_2\bigl(C_j(n+T)-C_j(n)\bigr)&=v_2(T)+c_j,\\
 U_j(n+T)-U_j(n)&\ne0,&
 v_2\bigl(U_j(n+T)-U_j(n)\bigr)&=v_2(T)+d_j.
\end{align*}
\end{theorem}

\begin{proof}
The first Newton coefficients, derived from
\eqref{eq:arith-triangular}, are
\[
 a_{0,1}=-2,\quad a_{1,1}=-\frac65,\quad
 a_{2,1}=-\frac{16}{15},\quad a_{3,1}=-\frac{36}{35}.
\]
Their exact valuations are $c_j$.  Theorem~\ref{thm:arith-dyadic-differences}
gives $v_2(a_{j,k})\ge k$ for all orders.  To obtain the stronger
logarithmic envelope required by Lemma~\ref{lem:arith-dominance}, only
a finite initial segment needs improvement.  The recurrence proves the
following exact valuations:
\begin{center}
\begin{tabular}{cll}\toprule
Address $j$ & Further orders $k$ & $v_2(a_{j,k})$ in that order\\\midrule
$0$ & $2$ & $3$\\
$1$ & $2$ & $3$\\
$2$ & $2,3,4,5,6$ & $6,9,13,16,22$\\
$3$ & $2,3,4$ & $8,8,15$\\\bottomrule
\end{tabular}
\end{center}
For all subsequent $k$, the elementary inequality
$c_j+1+\lfloor\log_2k\rfloor\le k$ completes the envelope.
The entries follow by exact rational computation from the defining
recurrence, and their valuations are proved in Lean.  Together with the
inequality for all subsequent orders, they establish the bound at every
order.
The uniform envelope is
\lean{KomorebiPaperIX.DyadicFourPrefix.tail_precision}.

Apply Lemma~\ref{lem:arith-dominance} to $C_j$.  Multiplication by the
nonzero factorial $j!$ adds $v_2(j!)$ and preserves nonzeroness.
Hence $d_j=c_j+v_2(j!)$, giving the four valuation gains for the original
coefficients.
The declarations are
\lean{KomorebiPaperIX.DyadicFourLeastPeriods.normalized_exact_displacement}
and \lean{KomorebiPaperIX.DyadicFourLeastPeriods.exact_displacement}.
\end{proof}

\begin{corollary}[All positive periods and their least element]
\label{cor:arith-least-period}
For $r\in\N$ and $0\le j<4$, a positive integer $T$ is a packet period
of $U_j$ modulo $2^r$ if and only if
\[
 2^{\max(0,r-d_j)}\mid T.
\]
Here a packet period means $U_j(n+T)\equiv U_j(n)\pmod{2^r}$ for every
$n\ge0$.  Consequently the least positive packet periods are
\[
 \begin{array}{c|cccc}
 j&0&1&2&3\\\hline
 \text{least period}&2^{\max(0,r-1)}&2^{\max(0,r-1)}&
                 2^{\max(0,r-5)}&2^{\max(0,r-3)}.
 \end{array}
\]
\end{corollary}

\begin{proof}
By the exact displacement law, the congruence at any row is equivalent to
$v_2(T)+d_j\ge r$.  This is exactly the displayed integer
divisibility criterion.  Its least positive solution is the indicated
power of two, which also covers $r=0$ and $r\le d_j$.
The all-period and least-element statements are
\lean{KomorebiPaperIX.DyadicFourLeastPeriods.period_iff_dvd} and
\lean{KomorebiPaperIX.DyadicFourLeastPeriods.least_packet_period}.
\end{proof}

For example, $U_2(n)\equiv2\pmod{32}$ and $U_3(n)\equiv6\pmod8$
for every row.  At precision $2^{10}$ the least packet periods are
$512,512,32,128$.  A packet displacement $T$ has original degree
displacement $4T$, as proved by
\lean{KomorebiPaperIX.DyadicFourLeastPeriods.original_degree_exact_displacement}.
Minimality here concerns packet shifts.  Shifts of the full degree-indexed
sequence, including its zero coefficients off the packet support, require
a separate analysis.  The new strengthening chapters determine all repeated-step higher differences in addresses zero and one and a power-of-two-order family in address two. The full exact classification of arbitrary higher-order differences
is also a separate question from the first-displacement law.

\subsection{Classical odd-length coordinates and their congruence ladder}

The classical signed Euler numbers have generating function
$\operatorname{sech}t=\sum_{n\ge0}E_{2n}t^{2n}/(2n)!$.
They are the $N=2$, $j=0$ member of
\eqref{eq:arith-euler-egf}.  The coordinate theorem reconstructs
them from finite alternating sums of odd powers.

\begin{definition}[Odd-length source and inverse coefficients]
\label{def:arith-odd-source}
For a positive odd integer $p$, define
\[
 S_d(p)=\sum_{a=0}^{p-1}(-1)^a(2a+1)^d.
\]
Let $\gamma_m$ be the integral Hurwitz coefficients of $\operatorname{sech}^2$:
\[
 \operatorname{sech}^2t=\sum_{m\ge0}\gamma_m\frac{t^{2m}}{(2m)!},
 \qquad
 \gamma_m=\sum_{b=0}^{2m}\binom{2m}{b}E_bE_{2m-b},
\]
where $E_b=0$ for odd $b$.  Then $\gamma_0=1$.  Define
\[
 T^{\mathrm{odd}}_{p,M}(n)=\sum_{m=0}^{\min(M,n)}
       \binom{2n}{2m}\gamma_m p^{2m}S_{2(n-m)}(p).
\]
Here $p$ need only be odd; primality is needed only to interpret the
result as a prime-adic congruence ladder.
\end{definition}

The source and inverse coefficients are
\lean{KomorebiPaperIX.OddPower.source} and
\lean{KomorebiPaperIX.EulerCoefficients.inverseEven}; the truncation is
\lean{KomorebiPaperIX.OddEuler.truncation}.

The same inverse coefficients can be written in the signed zigzag notation
of the original manuscript.  Define $Z_r$ by
\[
 \sec t+\tan t=\sum_{r\ge0}Z_r\frac{t^r}{r!}.
\]
Then
\begin{equation}\label{eq:arith-signed-zigzag}
 \gamma_m=(-1)^m Z_{2m+1}\qquad(m\in\N).
\end{equation}
Indeed, the odd coefficients come from $\tan t$, whose derivative is
$\sec^2t$.  Changing the sign of the squared variable in the even series
for $\operatorname{sech}^2t$ gives $\sec^2t$, so its degree-$2m$
Hurwitz coefficient is $(-1)^m\gamma_m$.  Differentiation identifies
this coefficient with $Z_{2m+1}$, proving the formula.  The Lean bridge is
\lean{KomorebiPaperIX.Zigzag.inverse_is_signed_zigzag}; it uses the
formal generating-function definition of the zigzag numbers.

\begin{theorem}[Classical odd-length reconstruction]
\label{thm:arith-odd-reconstruction}
For every positive odd integer $p$ and every $n\in\N$,
\[
 E_{2n}=\sum_{m=0}^{n}\binom{2n}{2m}\gamma_m p^{2m}S_{2(n-m)}(p).
\]
Equivalently, $E_{2n}=\sum_{q=0}^n V_{nq}S_{2q}(p)$ for the lower
triangular matrix
\[
 V_{nq}=\binom{2n}{2q}\gamma_{n-q}p^{2(n-q)}\quad(q\le n),
 \qquad V_{nq}=0\quad(q>n).
\]
It has $V_{nn}=1$ and $p^{2(n-q)}\mid V_{nq}$ when $q\le n$.
\end{theorem}

\begin{proof}
The exponential generating series of the source is the finite sum
\[
 \mathcal S_p(t)=\sum_{a=0}^{p-1}(-1)^a e^{(2a+1)t}.
\]
Multiplication by $\cosh t$ telescopes.  Since $p$ is odd,
\[
 \mathcal S_p(t)=\frac{1+e^{2pt}}{2}\operatorname{sech}t.
\]
Taking its even part gives
\[
 \sum_{n\ge0}S_{2n}(p)\frac{t^{2n}}{(2n)!}
       =\cosh^2(pt)\operatorname{sech}t.
\]
Multiplication by the reciprocal $\operatorname{sech}^2(pt)$
yields $\operatorname{sech}t$.  The EGF product rule now gives the
reconstruction coefficient by coefficient.  Replacing $m$ by $n-q$
produces the displayed lower triangular row.  Its diagonal and weights
follow from $\gamma_0=1$ and the explicit factor $p^{2(n-q)}$.
The source-series identity is
\lean{KomorebiPaperIX.OddPower.odd_source_even_egf}; the reconstruction
and row statement are
\lean{KomorebiPaperIX.OddEuler.odd_euler_reconstruction} and
\lean{KomorebiPaperIX.OddEuler.row_action}.
\end{proof}

\begin{corollary}[Odd-prime truncation ladder]\label{cor:arith-odd-ladder}
For every positive odd integer $p$ and all $M,n\in\N$,
\[
 p^{2M+2}\mid E_{2n}-T^{\mathrm{odd}}_{p,M}(n)\quad\text{in }\Z.
\]
If $M\ge n$, the truncation equals $E_{2n}$ exactly.  In particular,
for an odd prime $p$ this is a congruence modulo $p^{2M+2}$ in $\Z_p$.
\end{corollary}

\begin{proof}
Every omitted summand has $m\ge M+1$, hence contains the integer factor
$p^{2m}$ divisible by $p^{2M+2}$.  All its other factors are integral.
The finite sum of these omitted terms proves the assertion, with the
empty-tail case giving exact equality.  These are
\lean{KomorebiPaperIX.OddEuler.odd_euler_supercongruence} and
\lean{KomorebiPaperIX.OddEuler.full_depth_exact}.
\end{proof}

The reconstruction and its truncation bound already hold over the integers.
For the dyadic construction we use a different source.  Its original bound
of $2^{2M+2}$ follows from integral reciprocal coefficients, and the next
subsection sharpens that bound to the optimal uniform exponent $4M+3$.

\subsection{The optimal classical dyadic edge tower}

Return now to the classical signed Euler numbers
\[
 \operatorname{sech}t=\frac1{\cosh t}
       =\sum_{n\ge0}E_{2n}\frac{t^{2n}}{(2n)!}.
\]
Thus $E_{2n}=E_{2n}^{(2,0)}$ in the convention
\eqref{eq:arith-euler-egf}.  In particular $E_0=1$, $E_2=-1$, and $E_4=5$.
The following construction concerns these $N=2$ coefficients, not $C_{4,0}$.

\begin{definition}[Dyadic truncation]\label{def:arith-truncation}
Define $\beta_m$ by the formal unit-series identity
\[
 \sum_{m\ge0}\beta_m z^{2m}=\left(\frac{z}{\sinh z}\right)^2,
 \qquad \beta_0=1,
\]
where the quotient means the inverse of the unit series
$(\sinh z/z)^2$.  Define
\[
 A(n)=\frac{3^{2n+2}-1}{4(2n+1)(2n+2)},\qquad
 S_m(n)=\beta_m2^{2m}\frac{(2n)!}{(2n-2m)!}A(n-m)
       \quad(0\le m\le n),
\]
and
\begin{equation}\label{eq:arith-tm}
 T_M(n)=\sum_{m=0}^{\min(M,n)}S_m(n)\qquad(M,n\in\N).
\end{equation}
\end{definition}

These definitions correspond to
\lean{KomorebiPaperIX.Dyadic.beta},
\lean{KomorebiPaperIX.Dyadic.sourceEven}, and
\lean{KomorebiPaperIX.DyadicReconstruction.truncation}.
The formal summand uses a falling factorial, and
\lean{KomorebiPaperIX.DyadicReconstruction.falling_factorial_fraction}
proves its equality with the quotient displayed above.

\begin{lemma}[Exact reconstruction and source integrality]
\label{lem:arith-reconstruction}
For every $n\in\N$,
\[
 E_{2n}=\sum_{m=0}^{n}S_m(n).
\]
Every $\beta_m$ is $2$-adically integral and every $A(n)$ is a nonzero
$2$-adic unit.  If $M\ge n$, then $T_M(n)=E_{2n}$.
\end{lemma}

\begin{proof}
The ordinary kernel coefficients of $(\sinh z/z)^2$, in $z^2$, are
\[
 \frac{2^{2m+1}}{(2m+2)!}.
\]
Legendre's strict factorial bound
$v_2((2m+2)!)<2m+2$ makes them integral.  Their constant coefficient is
one.  The coefficient recurrence for its reciprocal therefore
proves integrality of every $\beta_m$.
For the source factor, lifting the exponent at the even power $2n+2$
gives
\[
 v_2(3^{2n+2}-1)=2+v_2(2n+2).
\]
Its denominator has the same valuation, since $2n+1$ is odd.  Thus
$A(n)$ is a nonzero unit.

Its exponential generating function is
\[
 \mathcal A(t)=\sum_{n\ge0}A(n)\frac{t^{2n}}{(2n)!}
             =\frac{\cosh(3t)-\cosh t}{4t^2}.
\]
The quotients by powers of $t$ are well defined because the required
initial coefficients vanish.  The hyperbolic identity
\[
 \mathcal A(t)\cosh t=\left(\frac{\sinh(2t)}{2t}\right)^2
\]
then gives the exact formal-series product
\[
 \left(\sum_{m\ge0}\beta_m(2t)^{2m}\right)\mathcal A(t)
       =\operatorname{sech}t.
\]
Taking the factorial-scaled coefficient of $t^{2n}$ produces the stated
finite sum.  Full depth leaves no omitted terms.
The relevant declarations are
\lean{KomorebiPaperIX.Dyadic.beta_integral},
\lean{KomorebiPaperIX.Dyadic.source_valuation},
\lean{KomorebiPaperIX.DyadicSeries.dyadic_reconstruction_series}, and
\lean{KomorebiPaperIX.DyadicReconstruction.dyadic_reconstruction}.
\end{proof}

\begin{lemma}[Exact weighted inverse coefficient]\label{lem:arith-kappa}
Let $\kappa_m=\beta_m2^{2m}(2m)!$.  For every $m\ge1$,
\[
 \kappa_m=-2^{4m}(2m-1)B_{2m}\ne0,\qquad
 v_2(\kappa_m)=4m-1,
\]
where $B_{2m}$ is the usual even Bernoulli number.
\end{lemma}

\begin{proof}
Let $Q(z)=z\coth z$, the regular formal series with constant coefficient
one.  The elementary differential identity
\[
 \left(\frac z{\sinh z}\right)^2=Q(z)-zQ'(z)
\]
and the Bernoulli expansion
\[
 Q(z)=1+\sum_{m\ge1}\frac{2^{2m}B_{2m}}{(2m)!}z^{2m}
\]
give
$\beta_m=-2^{2m}(2m-1)B_{2m}/(2m)!$.
Multiplication by $2^{2m}(2m)!$ proves the identity for $\kappa_m$.
The Bernoulli denominator theorem gives $v_2(B_{2m})=-1$ for $m>0$: its denominator contains each prime $q$ with $(q-1)\mid 2m$ exactly once, including $2$\cite{DLMF24}.
Since $2m-1$ is odd, the valuation is exactly $4m-1$.
This also proves nonvanishing.  In Lean, the hyperbolic identity is proved
after clearing removable zeros.  The library Bernoulli theorem then
determines the valuation of $\kappa_m$.
The declarations are
\lean{KomorebiPaperIX.DyadicSharpTower.kappa_bernoulli} and
\lean{KomorebiPaperIX.DyadicSharpTower.kappa_valuation}.
\end{proof}

\begin{theorem}[Optimal uniform dyadic tower]\label{thm:arith-sharp-tower}
For all $M,n\in\N$,
\begin{equation}\label{eq:arith-sharp-tail}
 E_{2n}-T_M(n)\in2^{4M+3}\Z_2.
\end{equation}
At the first omitted row $n=M+1$, the error is nonzero and
\[
 v_2\bigl(E_{2M+2}-T_M(M+1)\bigr)=4M+3.
\]
Thus the exponent $4M+3$ cannot be increased uniformly in the Euler row.
\end{theorem}

\begin{proof}
Rewrite the summand exactly as
\[
 S_m(n)=\kappa_m\binom{2n}{2m}A(n-m).
\]
For $m>M$, its first factor is divisible by $2^{4m-1}$, the binomial
coefficient is integral, and $A(n-m)$ is a unit.  Hence every omitted
summand is divisible by $2^{4M+3}$.  Summing the finite tail proves
\eqref{eq:arith-sharp-tail}.  If $M\ge n$, the tail is empty and the
statement includes the resulting zero error.

At $n=M+1$ there is exactly one omitted summand.  Here
$\binom{2n}{2n}=1$ and $A(0)=1$, so the error is $\kappa_{M+1}$.
Lemma~\ref{lem:arith-kappa} gives its exact valuation $4M+3$ and
nonvanishing.  The statements are
\lean{KomorebiPaperIX.DyadicSharpTower.dyadic_supercongruence_sharp},
\lean{KomorebiPaperIX.DyadicSharpTower.first_omitted_error}, and
\lean{KomorebiPaperIX.DyadicSharpTower.first_omitted_error_valuation}.
\end{proof}

\begin{proposition}[Row-dependent summand weights]\label{prop:arith-row-weight}
For $1\le m\le n$, the summand is nonzero and
\[
 v_2(S_m(n))=4m-1+v_2\binom{2n}{2m}.
\]
More generally, if $L\in\N$ satisfies
\[
 L\le4m-1+v_2\binom{2n}{2m}\qquad\text{for every }M<m\le n,
\]
then $E_{2n}-T_M(n)\in2^L\Z_2$.
\end{proposition}

\begin{proof}
Use the same factorization of $S_m(n)$, the nonzero binomial coefficient
for $m\le n$, the exact valuation of $\kappa_m$, and the unit property
of $A(n-m)$.  The tail conclusion follows by addition in $2^L\Z_2$.
These are \lean{KomorebiPaperIX.DyadicSharpTower.summand_valuation} and
\lean{KomorebiPaperIX.DyadicSharpTower.dyadic_supercongruence_row_weight}.
\end{proof}

The minimum summand weight bounds the tail valuation from below, but
cancellation may raise that valuation. Appendix~\ref{app:root:scope} gives a later exact factorization and finite-polynomial reduction for each fixed depth. A universal explicit root and simplicity classification for that factor remains open.  Uniform optimality follows
from the equality already established at $n=M+1$.

\subsection{Isometric continuation of the four dyadic packets}

The exact displacement law also determines the images of the packets
after continuous extension to the dyadic integers.

\begin{theorem}[Isometric packet bijections]\label{thm:arith-isometry}
For $0\le j<4$, define on natural numbers
\[
 f_j(n)=\frac{U_j(n)-j!}{2^{d_j}},\qquad (d_0,d_1,d_2,d_3)=(1,1,5,3).
\]
There is a unique continuous map $\Phi_j:\Z_2\to\Z_2$ whose value at
every $n\in\N$ equals the embedded rational $f_j(n)$.
It is an isometric bijection and $\Phi_j(0)=0$:
\[
 |\Phi_j(x)-\Phi_j(y)|_2=|x-y|_2\qquad(x,y\in\Z_2).
\]
In particular, every element of $\Z_2$ has exactly one preimage.
\end{theorem}

\begin{proof}
For $n>0$, the displacement law at row zero gives
$v_2(U_j(n)-j!)=v_2(n)+d_j$, so $f_j(n)\in\Z_2$.
The value at zero is zero.  Applying the same law to the smaller of two
unequal natural numbers shows
\[
 |f_j(n)-f_j(m)|_2=|n-m|_2.
\]
The natural numbers are dense in $\Z_2$, and the target is complete.
The map therefore extends continuously.  The distance identity extends
on the dense product of natural-number points, making the extension
an isometry.  Density also proves uniqueness among all continuous maps
with these prescribed values.

To prove surjectivity, use compactness.  For an isometry $g$ of a compact
metric space into itself, the sequence of iterates $g^n(x)$ has a
convergent subsequence.  Two sufficiently close iterates, with exponents
$a<b$, yield $d(x,g^{b-a}(x))$ equally small after removing $a$ isometric
iterations.  Since $b-a>0$, the latter point is in the image of $g$.
Thus the image is dense.  It is also compact and hence closed, so it is
the whole space.  Apply this to $\Phi_j$ on $\Z_2$.

The general ingredients are
\lean{KomorebiPaperIX.DensePadicIsometry.existsUnique_extensionEquiv} and
\lean{KomorebiPaperIX.CompactSelfIsometry.surjective}.
The packet specialization is
\lean{KomorebiPaperIX.DyadicFourIsometry.existsUnique_isometry};
uniqueness among continuous maps is
\lean{KomorebiPaperIX.DyadicFourIsometry.unique_continuous}.
\end{proof}

\begin{corollary}[Exact image of the original coefficients]
\label{cor:arith-affine-image}
The continuous extension of the original sequence is
\[
 \widetilde U_j(x)=j!+2^{d_j}\Phi_j(x).
\]
It is injective, its image is exactly $j!+2^{d_j}\Z_2$, and
\[
 |\widetilde U_j(x)-\widetilde U_j(y)|_2
     =2^{-d_j}|x-y|_2.
\]
Every point of that affine image has exactly one preimage in $\Z_2$.
\end{corollary}

\begin{proof}
The formula agrees with $U_j$ on every natural number.  Continuity and
the distance equality follow from the corresponding properties of
$\Phi_j$.  Its surjectivity gives precisely the displayed affine image,
and multiplication by the nonzero scalar $2^{d_j}$ preserves injectivity.
The declarations are
\lean{KomorebiPaperIX.DyadicFourIsometry.originalExtension_nat},
\lean{KomorebiPaperIX.DyadicFourIsometry.originalExtension_dist}, and
\lean{KomorebiPaperIX.DyadicFourIsometry.originalExtension_range}.
\end{proof}

The four rational packet sequences thus extend to bijections onto their
exact dyadic images.  The proof uses density, completeness, and the
displacement law. The strengthening chapters supply the further factorial and Taylor-coefficient argument for analyticity, including exact radii in addresses zero, one and two.  The analytic polynomial solution maps constructed
later rest on separate hypotheses and a separate proof.

\clearpage
% Integrated strengthening, based on the frozen proof revision of 2026-10-06.
% No preamble. All valuations have v_p(0)=+infinity.
\section{Threshold-free reciprocal transfer at prime-power section orders}
\label{sec:prime-power-transfer-final}

The completion of Nishibuchi's conjecture fixes the requested eventual
congruence. The stronger question is whether the sufficient starting row
can be removed, and how much precision the difference order and the step
actually provide. This section answers the first question at every
prime-power section order. Subsequent sections give a factorial improvement
and exact results at section order four.

Throughout these sections, the normalized source and its reciprocal are
\begin{equation}
 G_{N,j}(t)=j!\sum_{n\geq0}\frac{t^{Nn}}{(Nn+j)!},\qquad
 F_{N,j}(t)=G_{N,j}(t)^{-1}
   =\sum_{n\geq0}C_{N,j}(n)\frac{t^{Nn}}{(Nn)!},
 \label{eq:final-normalization}
\end{equation}
where $N\geq1$ and $0\leq j<N$. In particular, $C_{N,j}(0)=1$.
The coefficients in the original convention satisfy
\begin{equation}
 E_{Nn}^{(N,j)}=j!C_{N,j}(n),\qquad
 E_d^{(N,j)}=0\quad(N\nmid d).
 \label{eq:final-original-normalization}
\end{equation}
Thus every valuation of a nonzero difference gains $v_p(j!)$ on returning
to the original coefficients. This factor must be retained when stating
optimal periods or the final congruence.

Write $S f(n)=f(n+1)$ and $\Delta_H=S^H-1$. A mixed difference
$\Delta_{H_1}\cdots\Delta_{H_k}$ acts on the row variable $n$, rather than
on the full Euler subscript. We use $v_p(0)=+\infty$. Orders and starting
rows are nonnegative integers; steps are positive integers unless stated
otherwise. An empty mixed difference is the original value, and a zero
ordinary step makes every nonempty product containing it zero.

\subsection{The Hurwitz lattice and translated differential flatness}

\begin{definition}[Hurwitz integral series]
For a prime $p$, let
\[
 \mathcal H_p=\left\{\sum_{d\geq0}f_d\frac{t^d}{d!}:
                    f_d\in\mathbb Z_p\right\}
       \subset\mathbb Q_p[[t]].
\]
For $N=p^a$ and $\lambda\in\mathbb Z_p$, put
\[
 D=\frac{d}{dt},\qquad L_\lambda=\frac{D^N-\lambda}{p}.
\]
A series $A$ is $\lambda$-flat if $L_\lambda^k A\in\mathcal H_p$ for every
$k\geq0$.
\end{definition}

Products preserve $\mathcal H_p$ because the coefficient at factorial
degree $d$ is
\[
 (AB)_d=\sum_{i=0}^d\binom di a_i b_{d-i}.
\]
The derivative $D$ acts by the coefficient shift $(DA)_d=a_{d+1}$.
Consequently every nonnegative power of $D$ preserves the lattice.
If $A\in\mathcal H_p$ has unit constant term, its reciprocal belongs to
the same lattice: with $B=A^{-1}$,
\begin{equation}
 b_0=a_0^{-1},\qquad
 b_d=-a_0^{-1}\sum_{i=1}^d\binom di a_i b_{d-i}.
 \label{eq:hurwitz-reciprocal-recurrence}
\end{equation}
If $A$ is supported at degrees divisible by $N$, this recurrence also
proves that $B$ has the same support.

\begin{lemma}[Translated product formula]
\label{lem:translated-hurwitz-product}
For $N=p^a$, $a\geq1$, and $\lambda,\mu\in\mathbb Z_p$,
\begin{equation}
 L_{\lambda+\mu}(AB)
  =(L_\lambda A)B+A(L_\mu B)
       +\sum_{i=1}^{N-1}c_i(D^iA)(D^{N-i}B),
 \qquad c_i=\frac1p\binom Ni\in\mathbb Z.
 \label{eq:translated-product}
\end{equation}
Flatness types add under multiplication.
\end{lemma}

\begin{proof}
The identity $(1+X)^{p^a}=1+X^{p^a}$ over $\mathbb F_p$ shows that every
interior binomial coefficient $\binom Ni$ is divisible by $p$.
Subtract $(\lambda+\mu)AB$ from the ordinary $N$th derivative product
formula and divide by $p$ to obtain \eqref{eq:translated-product}.

For completeness, the iterated formula is an integral multinomial
formula, rather than a product rule with unspecified remainders. On the
tensor of the two factors use the mutually commuting operators
\[
 L_\lambda\otimes1,\quad 1\otimes L_\mu,\quad
 c_iD^i\otimes D^{N-i}\quad(1\leq i<N).
\]
Applying multiplication after the $k$th power of their sum gives exactly
$L_{\lambda+\mu}^k(AB)$. Expanding it yields terms indexed by
$a,b,e_1,\ldots,e_{N-1}\geq0$, with $a+b+\sum_i e_i=k$:
\begin{equation}
 \frac{k!}{a!b!\prod_i e_i!}\prod_i c_i^{e_i}
 (D^\alpha L_\lambda^a A)(D^\beta L_\mu^b B),
 \quad
 \alpha=\sum_i i e_i,\quad
 \beta=\sum_i(N-i)e_i.
 \label{eq:translated-multinomial}
\end{equation}
Every scalar is integral. If $A$ and $B$ have the stated flatness types,
both factors in every term belong to $\mathcal H_p$, as does their product.
\end{proof}

\begin{lemma}[Reciprocal transfer]
\label{lem:translated-hurwitz-inverse}
If $A$ is $\lambda$-flat and has unit constant term, then $A^{-1}$ is
$(-\lambda)$-flat.
\end{lemma}

\begin{proof}
Set $B=A^{-1}$. Its order-zero integrality was established in
\eqref{eq:hurwitz-reciprocal-recurrence}. Assume by induction that
$L_{-\lambda}^bB\in\mathcal H_p$ for $b<k$. Since $AB=1$,
\[
 L_0^k(AB)=\left(\frac{D^N}{p}\right)^k1=0\qquad(k>0).
\]
Use \eqref{eq:translated-multinomial} with $\mu=-\lambda$.
Exactly one term has inverse order $b=k$: it has $a=0$ and all $e_i=0$,
and is $A L_{-\lambda}^k B$. In every other term $b<k$, so its second
factor is integral by induction and derivative-shift integrality. Its
first factor is integral by the flatness of $A$, and its scalar is
integral. Hence $A L_{-\lambda}^k B\in\mathcal H_p$. Multiplying by the
already integral $B$ proves $L_{-\lambda}^k B\in\mathcal H_p$.
This induction takes place in the entire Hurwitz lattice. It places no
restriction on the coefficient degree or the starting row.
\end{proof}

\subsection{The actual Euler source}

The packed factorial coefficient of $G_{N,j}$ is
\begin{equation}
 g_{N,j}(n)=\frac{j!(Nn)!}{(Nn+j)!}
   =\prod_{i=1}^j\left(1+\frac Ni n\right)^{-1}.
 \label{eq:actual-prime-power-kernel}
\end{equation}
The empty product is one. If $N=p^a$ and $1\leq i<N$, then
$v_p(i)<a$, so $N/i\in p\mathbb Z_p$. Every affine denominator in
\eqref{eq:actual-prime-power-kernel} is therefore a unit on $\mathbb Z_p$.

For $s\in p\mathbb Z_p$, an exact finite-difference identity is
\begin{equation}
 \Delta_1^k\frac1{1+sn}
  =\frac{(-s)^k k!}{\prod_{h=0}^k(1+s(n+h))}.
 \label{eq:affine-reciprocal-exact}
\end{equation}
It follows by induction, subtracting adjacent rational functions; each
successive subtraction adds the factor $-(k+1)s$ in the numerator.
Its denominator is a unit, so it belongs to $p^k\mathbb Z_p$.
The discrete Leibniz formula
\begin{equation}
 \Delta_1^k(fg)(n)
 =\sum_{r=0}^k\binom kr
       (\Delta_1^r f)(n)(\Delta_1^{k-r}g)(n+r)
 \label{eq:discrete-leibniz-final}
\end{equation}
then proves $\Delta_1^k g_{N,j}(n)\in p^k\mathbb Z_p$ for every $k,n$.
The coefficient at degree $Nn$ of $L_1^kG_{N,j}$ is precisely
$p^{-k}\Delta_1^k g_{N,j}(n)$, and the other residue classes vanish.
Thus the actual source is $1$-flat.

\begin{theorem}[All-row prime-power transfer]
\label{thm:all-prime-power-transfer}
Let $p$ be prime, $a\geq1$, $N=p^a$, and $0\leq j<N$.
The coefficients $C_{N,j}(n)$ belong to $\mathbb Z_p$ for all $n\geq0$.
Set $A_{N,j}(n)=(-1)^n C_{N,j}(n)$ and
$T_H=S^H-(-1)^H$. For $k\geq1$ and positive $H_1,\ldots,H_k$,
\begin{align}
 v_p(\Delta_{H_1}\cdots\Delta_{H_k}A_{N,j}(n))
  &\geq k+\sum_{\nu=1}^k v_p(H_\nu),
 \label{eq:prime-power-signed-bound}\\
 v_p(T_{H_1}\cdots T_{H_k}E_{Nn}^{(N,j)})
  &\geq v_p(j!)+k+\sum_{\nu=1}^k v_p(H_\nu).
 \label{eq:prime-power-original-twist}
\end{align}
At $p=2$, the ordinary differences of $C_{N,j}$ and of the original
coefficients satisfy the same bounds. These dyadic statements also hold
for $N=1$.
\end{theorem}

\begin{proof}
Reciprocal transfer gives $L_{-1}^kF_{N,j}\in\mathcal H_p$.
Extraction of the packed coefficient gives
$(S+1)^k C_{N,j}(n)\in p^k\mathbb Z_p$. Moreover,
\[
 \Delta_1^k A_{N,j}(n)=(-1)^{n+k}(S+1)^k C_{N,j}(n),
\]
so $A_{N,j}$ has the ordinary $p^k$ difference property.
At $p=2$, $L_1=L_{-1}-1$; its integral binomial expansion on
$F_{N,j}$ gives $\Delta_1^k C_{N,j}(n)\in2^k\mathbb Z_2$ directly.
For $N=1$ the interior sum in the product formula is empty, so the same
argument applies without change; equivalently, $F_{1,0}(t)=e^{-t}$.

To retain the valuation of a step, let $B$ be any sequence satisfying
$\Delta_1^m B\in p^m\mathbb Z_p$ at every order and row. Use the finite
operator identity
\[
 \Delta_H=\sum_{i=1}^H\binom Hi\Delta_1^i.
\]
The relation $i\binom Hi=H\binom{H-1}{i-1}$ implies
\[
 v_p\left(p^i\binom Hi\right)
 \geq i+v_p(H)-v_p(i)\geq1+v_p(H).
\]
Multiplying these identities for the $k$ steps, every term of total
order $M$ has its scalar multiplied by $p^M$, and is divisible by
$p^{k+\sum_\nu v_p(H_\nu)}$. This proves the mixed-step bound.
Finally,
\[
 \Delta_H A_{N,j}(n)=(-1)^{n+H}T_H C_{N,j}(n).
\]
Iteration contributes only a global sign. Multiplying by $j!$ proves
\eqref{eq:prime-power-original-twist}.
\end{proof}

\begin{corollary}[Threshold-free dyadic congruence]
\label{cor:all-dyadic-final}
For $N=2^a$, $a\geq0$, all permitted $j$, and $n\geq0$,
\begin{equation}
 v_2(\Delta_{H_1}\cdots\Delta_{H_k}E_{Nn}^{(N,j)})
  \geq v_2(j!)+k+\sum_{\nu=1}^k v_2(H_\nu).
 \label{eq:all-dyadic-final-bound}
\end{equation}
In particular, $\Delta_H^k$ gains $k(1+v_2(H))$ bits before normalization.
For $r>v_2(j!)$, the row displacement
$H=2^{\max(0,r-v_2(j!)-1)}$ is a sufficient period modulo $2^r$ from
row zero. If $r\leq v_2(j!)$, all original coefficients vanish modulo
$2^r$.
\end{corollary}

\begin{remark}[Necessary boundaries]
At an odd prime the ordinary step-one difference cannot replace the
sign-twisted operator in the general theorem. For $j=0$,
$C_{N,0}(0)=1$ and $C_{N,0}(1)=-1$, so their ordinary difference is
$-2$, a unit at every odd prime. The proof also does not extend
unchanged to non-prime-power $N$, whose interior binomial coefficients
need not be divisible by the chosen prime. Neither observation is a
counterexample to Nishibuchi's original eventual congruence.
\end{remark}

\section{Factorial precision at section order four}
\label{sec:factorial-four-final}

Write $C_j=C_{4,j}$ and $F_j=F_{4,j}$. The all-row theorem removes the
threshold, but at $N=4$ it leaves a further factorial gain unused.
The exact algebra below proves that gain for all four addresses and
identifies several infinite families where it is sharp.

\begin{theorem}[Four-address factorial bound]
\label{thm:four-factorial-bound}
For $j\in\{0,1,2,3\}$, $n\geq0$, $k\geq1$, and positive mixed steps,
\begin{equation}
 v_2(\Delta_{H_1}\cdots\Delta_{H_k}C_j(n))
 \geq k+v_2(k!)+\sum_{\nu=1}^k v_2(H_\nu)+\beta_j,
 \quad (\beta_0,\beta_1,\beta_2,\beta_3)=(0,0,3,1).
 \label{eq:factorial-four-lower}
\end{equation}
Address two also satisfies
\begin{equation}
 v_2(\Delta_{H_1}\cdots\Delta_{H_k}C_2(n))
 \geq 2k+v_2(k!)+\left\lfloor\frac{k}{2}\right\rfloor
             +\sum_{\nu=1}^k v_2(H_\nu).
 \label{eq:factorial-four-j2-filtration}
\end{equation}
For the original coefficients, the offsets in
\eqref{eq:factorial-four-lower} are $(0,0,4,2)$.
No positive-order offset is asserted for the empty difference.
\end{theorem}

The proof of the first two addresses uses one classical arithmetic
input: Frobenius's exact congruence for the ordinary Euler numbers in
the $\operatorname{sech}$ convention, as stated and derived by
Gessel~\cite{gessel}, Section~9. The address-two and address-three
bounds below are finite Bernoulli calculations and do not require that
input. A divided-power argument is necessary: ordinary integral moments
alone do not justify division by $k!$.

\subsection{Euler moments and an integral divided-power algebra}

Let $\mathcal L$ be the ordinary Euler evaluation functional on a
variable $\mathcal E$, defined by
\[
 \mathcal L(\mathcal E^d)=\operatorname{Euler}_d,
 \qquad \sum_{d\geq0}\operatorname{Euler}_d\frac{t^d}{d!}
                  =\operatorname{sech}t.
\]
The functional is linear; it is not a multiplicative evaluation.
Independent variables mean that moments of polynomials in separate
variables factor. Odd moments vanish. Frobenius's congruence is
\begin{equation}
 v_2\left(\sum_{r=0}^m(-1)^{m-r}\binom mr
                 \operatorname{Euler}_{2n+2r}\right)
       =m+v_2(m!)\qquad(n,m\geq0).
 \label{eq:classical-frobenius-exact}
\end{equation}
For $m=0$, the assertion is that every even Euler number is odd.
With $P=(\mathcal E^2-1)/2$ and
$\gamma_m(P)=P^m/m!$, the case $n=0$ says
\begin{equation}
 \mathcal L(\gamma_m(P))\in\mathbb Z_2^{\times}
                 \qquad(m\geq0).
 \label{eq:euler-pd-odd-moments}
\end{equation}
Introduce an independent Euler variable $\mathcal F$ and
$Q=(\mathcal F^2-1)/2$.

Put $K=\mathbb Q_2(i)$, $R=\mathbb Z_2[i]$,
$a=(1+i)/2$, and $b=(1-i)/2$. The exact identity
\begin{equation}
 F_0(t)=\frac{2}{\cosh t+\cos t}
       =\operatorname{sech}(at)\operatorname{sech}(bt)
 \label{eq:four-euler-factorization}
\end{equation}
gives, for $U=a\mathcal E+b\mathcal F$,
$\mathcal L(U^{4n})=C_0(n)$. Although $a$ and $b$ have negative
valuation, their denominators disappear from the following exact
polynomial identities:
\begin{align}
 U^2&=\mathcal E\mathcal F+i(P-Q),\nonumber\\
 H_0:=\frac{U^4-1}{2}
   &=P+Q-\frac{P^2}{2}-\frac{Q^2}{2}
                 +3PQ+i\mathcal E\mathcal F(P-Q).
 \label{eq:four-euler-H}
\end{align}

Let $\Gamma_R(P,Q)$ be the $R$-subalgebra generated by all
$\gamma_m(P)$ and $\gamma_r(Q)$. Adjoin $\mathcal E,\mathcal F$ with
\[
 \mathcal E^2=1+2P,\qquad \mathcal F^2=1+2Q,
\]
and call the resulting algebra $\mathcal A_0$. Its elements are finite
expressions; no completion is being presumed. It is generated as a
$\Gamma_R(P,Q)$-module by $1,\mathcal E,\mathcal F,\mathcal E\mathcal F$.
The functional maps this module to $R$. For an even basis term this is
\eqref{eq:euler-pd-odd-moments} and independence. For a basis term with
an odd Euler factor the moment is zero, since $P$ and $Q$ are even
polynomials in their respective variables.

Every displayed component of $H_0$ has all its divided powers in
$\mathcal A_0$. Besides the generators, the needed identities are
\begin{align}
 \gamma_m(P^2/2)
     &=\frac{(2m)!}{2^m m!}\gamma_{2m}(P),\nonumber\\
 \gamma_m(3PQ)&=3^m m!\gamma_m(P)\gamma_m(Q),\nonumber\\
 \gamma_m(P-Q)&=\sum_{r=0}^m(-1)^r
                   \gamma_{m-r}(P)\gamma_r(Q).
 \label{eq:components-divided-powers}
\end{align}
The first scalar counts partitions into pairs and is an integer.
For the last term of $H_0$, multiply the third identity by
$(i\mathcal E\mathcal F)^m$ and reduce the Euler powers using the
integral quadratic relations. Finally,
\[
 \gamma_k(X+Y)=\sum_{r=0}^k\gamma_r(X)\gamma_{k-r}(Y)
\]
proves $\gamma_k(H_0)\in\mathcal A_0$. Applying the functional and
putting $W=U^4-1=2H_0$ gives
\begin{equation}
 \mathcal L(W^m)\in2^m m!R\qquad(m\geq0).
 \label{eq:euler-W-factorial}
\end{equation}
These moments are rational. The identity
$R\cap\mathbb Q_2=\mathbb Z_2$, obtained from the linear independence
of $1,i$ over $\mathbb Q_2$, descends every such divisibility statement
to $\mathbb Z_2$. No field property of $R/2R$ is needed.

For every starting row,
\begin{equation}
 \Delta_1^k C_0(n)
   =\mathcal L(U^{4n}W^k)
   =\sum_{s=0}^n\binom ns\mathcal L(W^{k+s}).
 \label{eq:base-row-finite-expansion}
\end{equation}
Each summand lies in $2^k k!\mathbb Z_2$, proving the factorial
bound for address zero from row zero onward.

\subsection{A finite step-expansion lemma}

The conversion from unit steps to arbitrary mixed steps is used several
times, so it is useful to make it explicit. Write $\Delta=\Delta_1$.
For $H>0$,
\begin{equation}
 \Delta_H=\sum_{m=1}^H\binom Hm\Delta^m.
 \label{eq:finite-step-expansion}
\end{equation}
A mixed term with local orders $m_1,\ldots,m_k\geq1$ and
$M=\sum_\nu m_\nu$ therefore has valuation at least
\[
 M+v_2(M!)+\beta
    +\sum_\nu\bigl(v_2(H_\nu)-v_2(m_\nu)\bigr)
\]
if the unit-order estimate has offset $\beta$ for every positive order.
The inequalities
\[
 m-1\geq v_2(m),\qquad M\geq k,\qquad v_2(M!)\geq v_2(k!)
\]
make this at least $k+v_2(k!)+\beta+\sum_\nu v_2(H_\nu)$.
Every expansion is finite. This proves the desired mixed-step transfer
without an exchange of limits.

\subsection{An elementary Bernoulli denominator bound}

Let $B_m$ denote the Bernoulli numbers with $B_1=-1/2$.
The identity obtained by substituting $u=e^t-1$ into
$\log(1+u)/u$ is
\begin{equation}
 B_m=\sum_{r=0}^m(-1)^r\frac{r!}{r+1}
                                  \left\{\begin{matrix}m\\r\end{matrix}\right\},
 \label{eq:bernoulli-stirling-denominator}
\end{equation}
where the braces are Stirling numbers of the second kind. Each
$2r!/(r+1)$ is dyadically integral. Indeed, if $s=v_2(r+1)>0$, then
$r\geq2^s-1$ and
\[
 v_2(r!)\geq v_2((2^s-1)!)=2^s-1-s\geq s-1.
\]
The last inequality is $2^s\geq2s$; the case $s=0$ is immediate.
Consequently
\begin{equation}
 2B_m\in\mathbb Z_2\qquad(m\geq0).
 \label{eq:elementary-bernoulli-bound}
\end{equation}
Thus the transfer factors below require no appeal to the
von Staudt--Clausen theorem.

\subsection{Address one and the scaled-binomial transfer}

The address-one source is
$G_{4,1}(t)=(\sinh t+\sin t)/(2t)$. The identities $a+b=1$,
$a-b=i$, and $ab=1/2$ give
\begin{equation}
 D_1(t):=G_{4,1}(t)F_0(t)
 =\frac12\left(\frac{\tanh(at)}{at}
                         +\frac{\tanh(bt)}{bt}\right).
 \label{eq:address-one-tangent-factor}
\end{equation}
For example, its numerator follows from
\[
 \sinh t+\sin t
 =2b\sinh(at)\cosh(bt)+2a\cosh(at)\sinh(bt),
\]
while $\cosh t+\cos t=2\cosh(at)\cosh(bt)$.
Only degrees divisible by four survive. Its actual Hurwitz packets are
\begin{equation}
 D_1(t)=\sum_{i\geq0}d_i\frac{t^{4i}}{(4i)!},\qquad
 d_i=\left(-\frac14\right)^i
 \frac{2^{4i+2}(2^{4i+2}-1)B_{4i+2}}{(4i+2)(4i+1)}.
 \label{eq:address-one-packets}
\end{equation}
The constant is one. Since $v_2(4i+2)=1$, the elementary Bernoulli bound
shows $v_2(d_i)\geq2i$. Hurwitz inversion therefore gives
\[
 D_1(t)^{-1}=\sum_{i\geq0}h_i\frac{t^{4i}}{(4i)!},
 \qquad h_0=1,\quad v_2(h_i)\geq2i.
\]
This follows by induction from the inverse recurrence: for packet index
$i$, every product $d_rh_{i-r}$ has valuation at least $2i$, and its
binomial multiplier is integral.
Because $F_1=F_0D_1^{-1}$,
\begin{equation}
 C_1(n)-C_0(n)
     =\sum_{i\geq1}h_i\binom{4n}{4i}C_0(n-i).
 \label{eq:address-one-exact-transfer}
\end{equation}
At a natural $n$ this is a finite sum, with $i\leq n$.

Here is the polynomial estimate that supplies one additional bit.
For $i\geq1$, put $P_i(x)=4^i\binom{4x}{4i}$. Its constant is zero,
and its degree-$d$ coefficient is
\begin{equation}
 [x^d]P_i(x)=\frac{4^i4^d s(4i,d)}{(4i)!},
 \label{eq:scaled-binomial-coefficients}
\end{equation}
where $s(m,d)$ is a signed Stirling number of the first kind.
In each product contributing to $s(4i,d)$, at most $d-1$ of the
$2i-1$ even integers in $1,\ldots,4i-1$ can be omitted. Thus
\[
 v_2(s(4i,d))\geq\max(0,2i-d).
\]
Writing $s_2(i)$ for the number of ones in the binary expansion of $i$,
and using $v_2((4i)!)=4i-s_2(i)$, gives
\begin{equation}
 v_2([x^d]P_i)
 \geq2d-2i+s_2(i)+\max(0,2i-d)
 \geq d+s_2(i)\geq d+1\quad(d\geq1).
 \label{eq:scaled-binomial-degree-bound}
\end{equation}
For $k\geq1$, $\Delta_1^k x^d$ vanishes when $d<k$ and is an
integer polynomial with values divisible by $k!$. To see the latter,
expand $x^d=\sum_r\left\{\begin{smallmatrix}d\\r\end{smallmatrix}\right\}
r!\binom xr$ and apply
$\Delta_1^k\binom xr=\binom{x}{r-k}$.
It follows that
\begin{equation}
 \Delta_1^kP_i(x)\in2^{k+1}k!\mathbb Z_2
                      \qquad(k\geq0,\ x\in\mathbb Z_2).
 \label{eq:scaled-binomial-difference-bound}
\end{equation}
For $k=0$ the bound follows directly from the coefficient estimate.

To apply discrete products in \eqref{eq:address-one-exact-transfer}
without a moving truncation, first extend $C_0$ to $\mathbb Z_2$ by its
uniformly convergent Newton series. The all-row bound
$v_2(\Delta_1^m C_0(0))\geq m$ proves this convergence. Its factorial
bounds extend to every dyadic starting point by continuity and
approximation with natural rows. The sum in
\eqref{eq:address-one-exact-transfer} is then uniformly convergent on
$\mathbb Z_2$: both the binomial polynomial and $C_0(x-i)$ are
integral-valued and $v_2(h_i)\geq2i$.

For the pointwise product $P_i(x)C_0(x-i)$, the discrete Leibniz formula
has scalar divisibility
\[
 \binom kr(2^{r+1}r!)(2^{k-r}(k-r)!)=2^{k+1}k!.
\]
Multiplication by $h_i/4^i\in\mathbb Z_2$ and uniform summation now give
\begin{equation}
 v_2(\Delta_1^k(C_1-C_0)(n))
               \geq k+v_2(k!)+1\qquad(k,n\geq0).
 \label{eq:address-one-extra-bit}
\end{equation}
The same extra bit holds for positive mixed steps by
\eqref{eq:finite-step-expansion}. Combining this with the address-zero
bound proves Theorem~\ref{thm:four-factorial-bound} for address one.

\subsection{Central factorials and integral Bernoulli divided powers}

\begin{lemma}[Bernoulli divided-power moments]
\label{lem:bernoulli-central-pd}
Let $Z$ be the ordinary Bernoulli variable,
$\mathcal L(Z^r)=B_r$, with exponential transform $t/(e^t-1)$.
For $P=Z(Z+1)$ and $\mu_m=\mathcal L(P^m/m!)$,
\begin{equation}
 \mu_m\in\mathbb Z_2,\qquad
 v_2(\mu_m)\geq s_2(m)-1\quad(m\geq1).
 \label{eq:bernoulli-central-moment-bound}
\end{equation}
Moreover, $\mu_m$ is odd exactly when $m=0$ or $m$ is a positive power
of two.
\end{lemma}

\begin{proof}
Define the central-factorial basis polynomials
\[
 P_r(X)=\prod_{a=0}^{r-1}(X-a(a+1)).
\]
At $X=Z(Z+1)$, they have the exact factorization
\begin{equation}
 P_r(Z(Z+1))
   =\prod_{a=0}^{r-1}(Z-a)(Z+a+1)
   =(2r)!\binom{Z+r}{2r}.
 \label{eq:central-factorial-substitution}
\end{equation}
The Bernoulli transform gives
$\mathcal L((1+T)^Z)=\log(1+T)/T$, whence
$\mathcal L(\binom Zj)=(-1)^j/(j+1)$. Finite Vandermonde expansion
therefore gives
\begin{align}
 \mathcal L\left(\binom{Z+r}{2r}\right)
 &=(-1)^r\sum_{q=0}^r(-1)^q\binom rq\frac1{r+q+1}\nonumber\\
 &=(-1)^r\frac{(r!)^2}{(2r+1)!}.
 \label{eq:central-factorial-evaluation}
\end{align}
The last equality is the rational identity obtained by expanding
$\int_0^1t^r(1-t)^r\,dt$. This use of an ordinary finite integral makes
no assertion about a positive or bounded $p$-adic measure. In
particular,
\[
 \mathcal L(P_r(P))=(-1)^r\frac{(r!)^2}{2r+1}.
\]

Expand $X^m=\sum_{r=0}^mT(m,r)P_r(X)$. Since
$XP_r(X)=P_{r+1}(X)+r(r+1)P_r(X)$, these integer coefficients satisfy
\begin{equation}
 T(m+1,r)=T(m,r-1)+r(r+1)T(m,r),
 \quad T(0,0)=1.
 \label{eq:central-factorial-recurrence}
\end{equation}
Every $r(r+1)$ is even. Induction, including the boundary indices,
proves $2^{m-r}\mid T(m,r)$ for $0\leq r\leq m$.
For $m>0$ the $r=0$ term vanishes. Each remaining summand of
$\mathcal L(P^m/m!)$ has valuation at least
\begin{equation}
 m-r+2v_2(r!)-v_2(m!)
   =r+s_2(m)-2s_2(r)\geq s_2(m)-1\geq0.
 \label{eq:central-factorial-term-bound}
\end{equation}
Here $r\geq2s_2(r)-1$ follows from
$r\geq2^{s_2(r)}-1$ and $2^u\geq2u$ for $u\geq1$.
This proves integrality and the stronger bound.

If $m$ is not a power of two, $s_2(m)\geq2$ makes $\mu_m$ even.
If $m$ is a positive power of two, the only central indices which can
contribute an odd residue are $r=1$ and $r=3$, since
$r-2s_2(r)=-1$ only at these two positive integers. The $r=1$ term is
odd: $T(m,1)=2^{m-1}$ and $v_2(m!)=m-1$. For $m\geq4$ a power of
two, the $r=3$ term is even. More explicitly,
\[
 \frac{T(m,3)}{2^{m-3}}=h_{m-3}(1,3,6),
 \qquad
 \sum_{d\geq0}h_d(1,3,6)t^d
       =\frac1{(1-t)(1-3t)(1-6t)}.
\]
Reduction modulo two gives $(1-t)^{-2}$, whose degree $m-3$
coefficient has parity $m-2=0$. At $m=1,2$ the $r=3$ term is absent,
and the $r=1$ term is again odd. Finally $\mu_0=1$.
\end{proof}

\subsection{Address two and its stronger filtration}

The variable $\mathcal B=2Z+1$ has transform
\[
 e^t\frac{2t}{e^{2t}-1}=\frac{t}{\sinh t}.
\]
Its odd moments vanish, and
$P=(\mathcal B^2-1)/4=Z(Z+1)$ has the divided-power moments of
Lemma~\ref{lem:bernoulli-central-pd}. Introduce an independent copy
$\mathcal D$ and $Q=(\mathcal D^2-1)/4$. The exact factorization
\begin{equation}
 F_2(t)=\frac{t^2}{\cosh t-\cos t}
       =\frac{at}{\sinh(at)}\frac{bt}{\sinh(bt)}
 \label{eq:address-two-factorization}
\end{equation}
shows that $U=a\mathcal B+b\mathcal D$ satisfies
$\mathcal L(U^{4n})=C_2(n)$. Direct expansion gives
\begin{equation}
 H_2:=\frac{U^4-1}{4}
     =P+Q-P^2-Q^2+6PQ+i\mathcal B\mathcal D(P-Q).
 \label{eq:address-two-H}
\end{equation}

Let $\mathcal A_2$ be the divided-power algebra over $R$ generated by
$P,Q$, adjoining $\mathcal B,\mathcal D$ with
$\mathcal B^2=1+4P$, $\mathcal D^2=1+4Q$.
The functional is integral on its module basis
$1,\mathcal B,\mathcal D,\mathcal B\mathcal D$ by the Bernoulli lemma,
independence, and vanishing of odd moments. Every component in
\eqref{eq:address-two-H} has integral divided powers, by the same
finite identities as before; $P^2$ is even easier than $P^2/2$.
Consequently
\begin{equation}
 \mathcal L((U^4-1)^m)\in4^m m!\mathbb Z_2,
 \qquad
 v_2(\Delta_1^kC_2(n))\geq2k+v_2(k!).
 \label{eq:address-two-basic-factorial}
\end{equation}
The starting-row conclusion follows from the finite base expansion
\eqref{eq:base-row-finite-expansion}, now with $W=4H_2$.

The low Newton coefficients of the actual reciprocal are
\begin{equation}
 C_2(0)=1,\quad C_2(1)=-\frac1{15},\quad C_2(2)=\frac{13}{45},
 \qquad a_1=-\frac{16}{15},\quad a_2=\frac{64}{45},
 \label{eq:address-two-initial-newton}
\end{equation}
where $a_m=\Delta_1^m C_2(0)$. Thus $v_2(a_1)=4$ and $v_2(a_2)=6$.
For $m\geq3$, \eqref{eq:address-two-basic-factorial} bounds $a_m$ by
$2m+v_2(m!)$. The finite shifted Newton formula
\begin{equation}
 \Delta_1^kC_2(n)=\sum_{m=k}^{n+k}a_m\binom n{m-k}
 \label{eq:shifted-newton-finite}
\end{equation}
therefore gives
$v_2(\Delta_1^kC_2(n))\geq k+v_2(k!)+3$ for every $k\geq1$.
For $k\geq3$ this is $2k\geq k+3$; for $k=1,2$ it uses the exact
initial coefficients and the higher tail bound. The finite step
expansion proves the offset three in
\eqref{eq:factorial-four-lower}.

There is a still stronger filtration. Put
\[
 X=P+Q,\qquad \tau=i\mathcal B\mathcal D,\qquad \epsilon=1+\tau.
\]
Then the exact identities are
\begin{equation}
 H_2=\epsilon X+2G,\quad
 G=-\gamma_2(P)-\gamma_2(Q)+3PQ-\tau Q,\quad
 \epsilon^2=2V,\quad V=\tau-2X-8PQ.
 \label{eq:address-two-filtration-algebra}
\end{equation}
Every divided power of $G$ is integral component by component.
The finite addition identity gives
\begin{equation}
 \gamma_k(H_2)=\sum_{r=0}^k2^r\gamma_r(G)
                  \epsilon^{k-r}\gamma_{k-r}(X)
          \in2^{\lfloor k/2\rfloor}\mathcal A_2,
 \label{eq:address-two-filtration-pd}
\end{equation}
because the $r$th term belongs to
$2^{r+\lfloor(k-r)/2\rfloor}\mathcal A_2$.
Applying the functional, rational descent, and the base-index expansion
proves the unit-step version of
\eqref{eq:factorial-four-j2-filtration}.
For a mixed term of total order $M$, its lower bound exceeds the target
at order $k$ by at least
\[
 2(M-k)+\left\lfloor\frac M2\right\rfloor
          -\left\lfloor\frac k2\right\rfloor
       +v_2(M!)-v_2(k!)-\sum_\nu v_2(m_\nu)\geq0.
\]
Indeed $\sum_\nu v_2(m_\nu)\leq M-k$. This proves the stronger mixed
bound. The earlier offset-three bound is retained, since it is stronger
at the lowest orders.

\subsection{Address three and cancellation of the apparent pole}

The address-three source is
$G_{4,3}(t)=3(\sinh t-\sin t)/t^3$. Addition identities give
\[
 \sinh t-\sin t
 =2a\sinh(at)\cosh(bt)+2b\cosh(at)\sinh(bt),
 \qquad \cosh t-\cos t=2\sinh(at)\sinh(bt).
\]
Consequently
\begin{equation}
 D_3(t):=G_{4,3}(t)F_2(t)
 =\frac32\left(\frac{\coth(at)}{at}
                            +\frac{\coth(bt)}{bt}\right).
 \label{eq:address-three-cotangent-factor}
\end{equation}
This expression is a power series: its two polar terms cancel because
$a^{-2}+b^{-2}=0$. Its constant is one, only degrees divisible by four
survive, and the Bernoulli expansion gives
\begin{equation}
 D_3(t)=\sum_{i\geq0}d_i\frac{t^{4i}}{(4i)!},\qquad
 d_i=3\left(-\frac14\right)^i
          \frac{2^{4i+2}B_{4i+2}}{(4i+2)(4i+1)}.
 \label{eq:address-three-packets}
\end{equation}
Again $v_2(d_i)\geq2i$, and the inverse coefficients $h_i$ have the same
bound. Since $F_3=F_2D_3^{-1}$,
\begin{equation}
 C_3(n)-C_2(n)=\sum_{i\geq1}h_i\binom{4n}{4i}C_2(n-i).
 \label{eq:address-three-exact-transfer}
\end{equation}
The scaled-binomial argument
\eqref{eq:scaled-binomial-difference-bound} applies unchanged.
The basic bound for $C_2$ includes order zero; its continuous extension
is available from the threshold-free theorem. Thus the shifted sums
converge uniformly on $\mathbb Z_2$, and the pointwise product rule
proves
\begin{equation}
 v_2(\Delta_1^k(C_3-C_2)(n))\geq k+v_2(k!)+1
                                      \qquad(k,n\geq0).
 \label{eq:address-three-extra-bit}
\end{equation}
At every positive order $C_2$ has three extra bits, so this gives the
one-bit offset for $C_3$. Finite step expansion finishes the proof of
Theorem~\ref{thm:four-factorial-bound}.

\section{Sharp higher-difference families}
\label{sec:sharp-higher-final}

The lower bounds do not mean that every mixed stencil has exactly the
stated valuation. Addresses zero and one have a complete classification
of when the leading factorial precision is attained. Address two has
exact mixed precision at every positive power-of-two order.

\begin{theorem}[Exact factorial precision at addresses zero and one]
\label{thm:zero-one-exact-mixed}
Let $j\in\{0,1\}$, $k\geq1$, $n\geq0$, and $H_1,\ldots,H_k>0$.
Put
\[
 L=k+v_2(k!)+\sum_{\nu=1}^k v_2(H_\nu),\qquad
 e=\#\{\nu:2\mid H_\nu\}.
\]
The mixed difference has valuation exactly $L$ if $k$ is odd or $e$ is
even. If $k$ is even and $e$ is odd, it has valuation at least $L+1$.
In particular, for every homogeneous positive step,
\begin{equation}
 v_2(\Delta_H^kC_j(n))=k+v_2(k!)+k v_2(H).
 \label{eq:zero-one-homogeneous-exact}
\end{equation}
\end{theorem}

\begin{proof}
Return to the Euler variables of address zero and put
$X=P+Q$, $\tau=i\mathcal E\mathcal F$, $\epsilon=1+\tau$.
Direct rearrangement of \eqref{eq:four-euler-H} gives
\begin{equation}
 H_0=\gamma_2(X)+\epsilon X+2G,\qquad
 G=-\gamma_2(P)-\gamma_2(Q)+PQ-\tau Q,
 \qquad \epsilon^2=2(\tau-P-Q-2PQ).
 \label{eq:zero-parity-decomposition}
\end{equation}
The divided powers of $G$ are integral by the same component formulas
used above. Thus modulo $2\mathcal A_0$, $\gamma_k(H_0)$ equals the
divided power of $\gamma_2(X)+\epsilon X$. The identity
\[
 \gamma_r(\gamma_2(X))
   =\frac{(2r)!}{2^r r!}\gamma_{2r}(X)
\]
has an odd integer scalar. Every term involving $\epsilon^j$ for
$j\geq2$ vanishes modulo two. The surviving zero and first powers give
\begin{equation}
 \gamma_k(H_0)\equiv
       \gamma_{2k}(X)+\epsilon\gamma_{2k-1}(X)
                                  \pmod{2\mathcal A_0}.
 \label{eq:zero-parity-leading-pd}
\end{equation}
The multiplication coefficient in the second term is $2k-1$, also odd.
By the odd moments \eqref{eq:euler-pd-odd-moments} and independence,
\[
 \mathcal L(\gamma_m(P+Q))\equiv m+1\pmod2.
\]
The $\tau$-weighted moment vanishes because it contains an odd power of
each Euler variable. Therefore the first term of
\eqref{eq:zero-parity-leading-pd} evaluates to one modulo two and the
second to zero. Hence $\mathcal L(\gamma_k(H_0))$ is odd at every order,
including order zero. This proves
\[
 v_2(\Delta_1^k C_0(0))=k+v_2(k!).
\]
In the base-row expansion \eqref{eq:base-row-finite-expansion}, every
$s>0$ term has at least one extra bit after division by $2^k k!$.
Thus the same exact valuation holds at every starting row.
The extra-bit transfer \eqref{eq:address-one-extra-bit} proves it for
address one as well.

It remains to track possible ties in the mixed-step expansion.
For local orders $m_\nu\geq1$ and total $M$, the excess over $L$ is
at least
\begin{equation}
 \left(M-k-\sum_\nu v_2(m_\nu)\right)
                       +\left(v_2(M!)-v_2(k!)\right).
 \label{eq:zero-one-mixed-excess}
\end{equation}
The first bracket is nonnegative, and equality requires every
$m_\nu$ to be one or two. Two or more occurrences of two make the
second bracket positive, since the intervening factorial interval then
contains an even number. Exactly one occurrence of two can tie only
when $k$ is even; its binomial coefficient can have the minimally
allowed valuation only if its associated step is even.

The all-one multi-index contributes one after normalization by $2^L$.
When $k$ is even, each even step supplies exactly one further tied term,
again with normalized residue one. For example,
$\binom H2=H(H-1)/2$ attains $v_2(H)-1$ exactly when $H$ is even;
all remaining factors and the exact normalized order-$k+1$ difference
are odd. Thus the total normalized parity is one for odd $k$, and
$1+e$ for even $k$. This proves the classification.
In the homogeneous case, $e$ is zero or $k$, so it is even whenever
$k$ is even, proving \eqref{eq:zero-one-homogeneous-exact}.
\end{proof}

\begin{remark}[A mixed cancellation]
The exceptional class can gain more than one bit. The exact values
\[
 \Delta_1\Delta_2C_0(0)=-33728,\qquad
 \Delta_1\Delta_2C_1(0)=-\frac{150208}{585}
\]
both have valuation six, whereas their baseline $L$ is four.
The theorem classifies attainment of the baseline; it does not give a
complete deeper valuation in the cancellation class.
\end{remark}

\begin{theorem}[Exact address-two power orders]
\label{thm:two-exact-power-orders}
For every starting row $n\geq0$ and all positive mixed steps,
\begin{equation}
 v_2(\Delta_{H_1}\cdots\Delta_{H_k}C_2(n))
       =V_2(k)+\sum_{\nu=1}^k v_2(H_\nu)
 \label{eq:two-power-orders-exact}
\end{equation}
whenever $k$ is a positive power of two, where
\[
 V_2(1)=4,\qquad V_2(2)=6,\qquad
 V_2(k)=2k+v_2(k!)+\frac{k}{2}\quad(k\geq4).
\]
For the original coefficients, the right side gains one bit.
\end{theorem}

\begin{proof}
For $k=2^r\geq4$, put $q=k/2$, which is even. Normalize the exact
expansion \eqref{eq:address-two-filtration-pd} by $2^q$ before reducing
modulo two. Its $u$th auxiliary term has excess
$u+\lfloor(k-u)/2\rfloor-q=\lfloor u/2\rfloor$.
Only auxiliary indices zero and one can therefore survive:
\begin{equation}
 \frac{\gamma_k(H_2)}{2^q}
 \equiv V^q\gamma_k(X)
       +\epsilon V^{q-1}G\gamma_{k-1}(X)
                                \pmod{2\mathcal A_2}.
 \label{eq:two-power-leading-normalized}
\end{equation}
Here $V\equiv\tau\pmod2$ and $\tau^2\equiv1\pmod2$.
Since $q$ is even and $\epsilon\tau\equiv\epsilon\pmod2$, this reduces
to $\gamma_k(X)+\epsilon G\gamma_{k-1}(X)$.
Moreover, $G\equiv\gamma_2(X)+\tau Q\pmod2$.
After evaluation, every term with a single $\tau$ vanishes, so the
second contribution reduces to
\[
 \mathcal L\bigl((\gamma_2(X)+Q)\gamma_{k-1}(X)\bigr).
\]
The $\gamma_2(X)$ product has scalar $\binom{k+1}{2}$, which is even
because $4\mid k$. The other contribution expands as
\begin{equation}
 \mathcal L(Q\gamma_{k-1}(P+Q))
    =\sum_{a=0}^{k-1}(k-a)\mu_a\mu_{k-a}.
 \label{eq:two-power-Q-evaluation}
\end{equation}
An odd summand would require $k-a$ to be odd and $\mu_{k-a}$ to be
odd. By Lemma~\ref{lem:bernoulli-central-pd}, this requires $k-a=1$.
Then $a=k-1$ is not a power of two and $\mu_a$ is even. Hence the
second contribution is zero modulo two.

For the first contribution, independence gives
$\mathcal L(\gamma_k(P+Q))=\sum_{a=0}^k\mu_a\mu_{k-a}$.
The parity generating series of these convolution moments is
\[
 (1+t+t^2+t^4+t^8+\cdots)^2
        =1+t^2+t^4+t^8+\cdots\quad\text{in }\mathbb F_2[[t]].
\]
Its coefficient is one at every $k=2^r\geq4$. Thus the normalized
moment in \eqref{eq:two-power-leading-normalized} is odd and
\[
 v_2(\Delta_1^kC_2(0))=2k+v_2(k!)+k/2.
\]
Inserting $U^{4n}=(1+4H_2)^n$ introduces only corrections with strictly
higher valuation, by the filtered bound at orders $k+s$ for $s>0$.
The equality consequently holds at every starting row.

For mixed steps the all-one local multi-index gives the claimed leading
term. A higher term of total order $M>k$ has excess at least
\[
 2(M-k)+\left\lfloor\frac M2\right\rfloor-\frac k2
       +v_2(M!)-v_2(k!)-\sum_\nu v_2(m_\nu)
       \geq M-k\geq1.
\]
Hence the leading term cannot cancel.

At order one, $a_1$ in \eqref{eq:address-two-initial-newton} has
valuation four and the Newton tail starts at valuation six; its shifted
value remains exact. In the step expansion a local order two can lose
at most one bit through its binomial coefficient, so every correction
is still above $4+v_2(H_1)$. At order two, $a_2$ has valuation six and
the next coefficient has lower bound eight from the stronger
filtration. A total-order-three correction loses at most one bit,
so it stays above the leading valuation six, while later orders gain
still more. The same argument proves exactness for both mixed steps.
\end{proof}

\begin{remark}[What the power orders do not decide]
The exact orders in Theorem~\ref{thm:two-exact-power-orders} form an
infinite family. They do not supply a complete valuation formula at
non-power-of-two orders in address two, nor an exact higher-difference
formula at address three. Numerical patterns at other orders remain
bounded observations unless proved separately.
\end{remark}

\section{Exact first differences, least periods, and dyadic images}
\label{sec:sharp-periods-final}

The factorial estimates are stronger than needed for the optimal
first-difference law. A general Newton criterion gives the exact law
from the threshold-free bound and finitely many rational coefficients.
This route also yields all addresses at ranks two and eight, and an
infinite family at every larger dyadic rank.

\begin{lemma}[A first-difference criterion]
\label{lem:newton-exact-isometry-criterion}
Let $C:\mathbb Z_{\geq0}\to\mathbb Z_2$ have Newton coefficients
$a_k=\Delta_1^kC(0)$. Suppose for an integer $b\geq0$ that
\begin{equation}
 v_2(a_1)=b,\qquad
 v_2(a_k)\geq b+1+\lfloor\log_2k\rfloor\quad(k\geq2).
 \label{eq:newton-isometry-criterion}
\end{equation}
Then for every $n\geq0$ and $H>0$,
\begin{equation}
 v_2(C(n+H)-C(n))=b+v_2(H).
 \label{eq:newton-isometry-conclusion}
\end{equation}
The Newton interpolation extends continuously to $\mathbb Z_2$ and
satisfies the same identity for any two distinct dyadic arguments.
\end{lemma}

\begin{proof}
At a natural argument the exact Newton expansion is finite.
Vandermonde gives
\[
 \binom{n+H}{k}-\binom nk
   =\sum_{i=1}^k\binom Hi\binom n{k-i}.
\]
For $1\leq i\leq H$,
$v_2(\binom Hi)\geq v_2(H)-v_2(i)
 \geq v_2(H)-\lfloor\log_2k\rfloor$; terms with $i>H$ vanish.
Thus each Newton summand with $k\geq2$ has valuation at least
$b+v_2(H)+1$. The $k=1$ summand is $Ha_1$, of exact valuation
$b+v_2(H)$, proving \eqref{eq:newton-isometry-conclusion}.

The coefficient hypothesis implies $a_k\to0$ dyadically, so the
Newton series converges uniformly on $\mathbb Z_2$, where
$\binom xk$ is integral-valued. For distinct $x,y\in\mathbb Z_2$,
approximate the pair by natural integer pairs with their difference
valuation eventually equal to $v_2(x-y)$. The corresponding function
differences have the fixed exact valuation $b+v_2(x-y)$.
A limit of elements with that fixed valuation has the same valuation:
the normalized residues are units, and $\mathbb Z_2^\times$ is closed.
\end{proof}

\begin{theorem}[Exact four-address first differences]
\label{thm:rank-four-exact-periods}
For every $n\geq0$, $H>0$, and $j\in\{0,1,2,3\}$,
\begin{align}
 v_2(C_j(n+H)-C_j(n))&=v_2(H)+b_j,
      &(b_0,b_1,b_2,b_3)&=(1,1,4,2),\label{eq:rank-four-normalized-isometry}\\
 v_2(E_{4(n+H)}^{(4,j)}-E_{4n}^{(4,j)})&=v_2(H)+B_j,
      &(B_0,B_1,B_2,B_3)&=(1,1,5,3).
 \label{eq:rank-four-original-isometry}
\end{align}
Consequently, the least positive row period modulo $2^r$ is exactly
\begin{equation}
 P_j(r)=2^{\max(0,r-B_j)}\qquad(r\geq1).
 \label{eq:rank-four-least-period-final}
\end{equation}
The corresponding displacement in the original Euler subscript is
$4P_j(r)$.
\end{theorem}

\begin{proof}
The actual reciprocal recurrence is
\begin{equation}
 C_{N,j}(n)=-\sum_{q=1}^n\binom{Nn}{Nq}
           g_{N,j}(q)C_{N,j}(n-q),\qquad C_{N,j}(0)=1.
 \label{eq:exact-reciprocal-row-recurrence}
\end{equation}
For $N=4$, the four kernels are
\[
 1,\qquad\frac1{4q+1},\qquad
 \frac1{(4q+1)(2q+1)},\qquad
 \frac3{(4q+1)(2q+1)(4q+3)}.
\]
The threshold-free theorem gives $v_2(a_k)\geq k$ in every address.
The following exact values cover all orders where that bound does not
already imply \eqref{eq:newton-isometry-criterion}:
\begin{center}
\begin{tabular}{c r l r}
$j$&$k$&$a_k$&$v_2(a_k)$\\\hline
0&1&$-2$&1\\
0&2&$72$&3\\
1&1&$-6/5$&1\\
1&2&$184/45$&3\\
2&1&$-16/15$&4\\
2&2&$64/45$&6\\
2&3&$-14848/1365$&9\\
2&4&$71426048/69615$&13\\
2&5&$-25509691392/85085$&16\\
2&6&$149113198870528/765765$&22\\
3&1&$-36/35$&2\\
3&2&$256/231$&8\\
3&3&$-65792/35035$&8\\
3&4&$161644544/4849845$&15
\end{tabular}
\end{center}
They are finite rational evaluations of
\eqref{eq:exact-reciprocal-row-recurrence}, followed by
$a_k=\sum_{h=0}^k(-1)^{k-h}\binom kh C_j(h)$.
For addresses zero and one, $k\geq3$ is covered by the generic bound;
for address two, $k\geq7$ is covered; for address three, $k\geq5$ is
covered. In each range the elementary inequality is
$k-\lfloor\log_2k\rfloor\geq b_j+1$.
Lemma~\ref{lem:newton-exact-isometry-criterion} proves the normalized
identity. Multiplication by $j!$ proves the original identity.

A row displacement $H$ preserves the original coefficient modulo
$2^r$ at even one row exactly when $v_2(H)+B_j\geq r$.
The least positive such integer is
$2^{\max(0,r-B_j)}$, and the same displacement works at every row.
This proves both optimality and the absence of a starting threshold.
\end{proof}

\begin{corollary}[Exact image and residue enumeration]
\label{cor:four-exact-images}
Let $\widetilde C_j$ denote the dyadic Newton interpolation.
The map
\[
 T_j(x)=\frac{\widetilde C_j(x)-1}{2^{b_j}}
\]
is a bijective isometry $\mathbb Z_2\to\mathbb Z_2$, and
\begin{equation}
 \widetilde C_j(\mathbb Z_2)=1+2^{b_j}\mathbb Z_2,
 \qquad j!\widetilde C_j(\mathbb Z_2)=j!+2^{B_j}\mathbb Z_2.
 \label{eq:four-interpolation-images-final}
\end{equation}
For $r>B_j$, one least period modulo $2^r$ visits every residue in
$j!+2^{B_j}\mathbb Z/2^r\mathbb Z$ exactly once. For $r\leq B_j$,
there is one residue and the least row period is one.
\end{corollary}

\begin{proof}
The exact difference identity gives
$v_2(T_j(x)-T_j(y))=v_2(x-y)$ and $T_j(0)=0$.
For every $s\geq1$, the map induces an injection of
$\mathbb Z_2/2^s\mathbb Z_2$ into itself. It is a bijection because
both sets have $2^s$ elements. These quotient bijections are compatible
with reduction. Inverting them gives a compatible preimage system for
every dyadic integer, proving surjectivity. Multiplication by $j!$
changes the image ideal by exactly $v_2(j!)$.
The finite quotient bijection is also precisely the residue enumeration
in a least row period.
\end{proof}

\subsection{An infinite exact family at all dyadic ranks}

\begin{theorem}[Lower-half exact laws]
\label{thm:lower-half-all-dyadic-isometry}
Let $N=2^a\geq4$ and $0\leq j<N/2$. Then for all $n\geq0$ and $H>0$,
\begin{equation}
 v_2(C_{N,j}(n+H)-C_{N,j}(n))=1+v_2(H).
 \label{eq:lower-half-all-dyadic-exact}
\end{equation}
Its normalized interpolation is a bijection onto $1+2\mathbb Z_2$.
For the original coefficients, the least row period modulo $2^r$ is
$2^{\max(0,r-1-v_2(j!))}$.
\end{theorem}

\begin{proof}
Every slope $N/i$ in the source has valuation at least two, since
$i\leq j<N/2$. Hence $g(1)\equiv1\pmod4$ and
$g(2)\equiv1\pmod8$. The first two Newton coefficients from the
reciprocal recurrence are
\begin{equation}
 a_1=-g(1)-1,\qquad
 a_2=\binom{2N}{N}g(1)^2-g(2)+2g(1)+1.
 \label{eq:lower-half-initial-newton}
\end{equation}
The first has valuation exactly one. Also
$\binom{2N}{N}\equiv6\pmod8$ for every $N=2^a\geq2$.
To see this, write
\[
 B_a=\binom{2^{a+1}}{2^a}
       =2\prod_{h=1}^{2^a-1}\left(1+\frac{2^a}{h}\right).
\]
The even-$h$ factors give $B_{a-1}$; the odd-$h$ factors are one
modulo four for $a\geq2$. Since $B_{a-1}$ has valuation one,
$B_a\equiv B_{a-1}\pmod8$, and $B_1=6$.
Therefore $a_2\equiv6-1+2+1=0\pmod8$.
The global Mahler bound covers all $k\geq3$, completing the criterion
with $b=1$. The period and image deductions are the same as above.
The proof also covers $N=2,j=0$, whose kernel is identically one.
\end{proof}

\subsection{Complete exact laws at ranks two and eight}

\begin{proposition}[Small-rank complete isometry table]
\label{prop:two-four-eight-isometry-table}
Every address at $N\in\{2,4,8\}$ satisfies
\[
 v_2(C_{N,j}(n+H)-C_{N,j}(n))=b_{N,j}+v_2(H).
\]
The following entries are ordered by $j=0,\ldots,N-1$:
\begin{center}
\begin{tabular}{c l l}
$N$&normalized shifts $b_{N,j}$&original shifts $B_{N,j}=b_{N,j}+v_2(j!)$\\\hline
2&$1,2$&$1,2$\\
4&$1,1,4,2$&$1,1,5,3$\\
8&$1,1,1,1,4,3,2,2$&$1,1,2,2,7,6,6,6$
\end{tabular}
\end{center}
For each entry, the least original row period is
$2^{\max(0,r-B_{N,j})}$, its Euler subscript displacement is $N$ times
that period, and the exact-image and residue-enumeration conclusions
hold with the corresponding shifts.
\end{proposition}

\begin{proof}
Rank four was proved above. At ranks two and eight the finite
criterion is verified by the actual rational recurrence
\eqref{eq:exact-reciprocal-row-recurrence}. The following valuation
lists record every initial Newton coefficient needed in addition to
the generic tail $v_2(a_k)\geq k$:
\begin{center}
\begin{tabular}{c c c l}
$N$&$j$&$b_{N,j}$&$v_2(a_1),v_2(a_2),\ldots$ through the finite cutoff\\\hline
2&0&1&$1,3$\\
2&1&2&$2,5,9,11$\\
8&0&1&$1,3$\\
8&1&1&$1,3$\\
8&2&1&$1,3$\\
8&3&1&$1,3$\\
8&4&4&$4,6,9,13,16,20$\\
8&5&3&$3,6,10,13,17$\\
8&6&2&$2,6,8,13$\\
8&7&2&$2,11,8,15$
\end{tabular}
\end{center}
The first entries can also be read directly from
$a_1=-1-\binom{N+j}{j}^{-1}$. At rank eight they are respectively
\[
 -2,\ -\frac{10}{9},\ -\frac{46}{45},\ -\frac{166}{165},\
 -\frac{496}{495},\ -\frac{1288}{1287},\
 -\frac{3004}{3003},\ -\frac{6436}{6435}.
\]
All later entries in the table are finite rational arithmetic, with
full numerator and denominator records retained in the accompanying
exact certificate \texttt{all\_rank\_extensions\_verified.json}.
Every entry with $k\geq2$ meets
$b_{N,j}+1+\lfloor\log_2k\rfloor$. Beyond the displayed cutoff, the
inequality $k\geq b_{N,j}+1+\lfloor\log_2k\rfloor$ holds at every
order. Thus the finite table and the proved infinite tail establish
the full hypothesis of Lemma~\ref{lem:newton-exact-isometry-criterion}.
This is a finite coefficient proof followed by a universal tail
argument, rather than extrapolation from sampled rows or steps.
\end{proof}

\begin{remark}[Failure of a universal upper-half exact law]
The all-rank lower bound remains valid in every address. The exact
first-difference law does not automatically extend to all upper-half
addresses. Exact rational recurrence calculations give
\begin{center}
\begin{tabular}{c c c c}
$(N,j)$&$v_2(C(1)-C(0))$&$v_2(C(2)-C(0))$&naive $b+v_2(2)$\\\hline
$(16,9)$&5&7&6\\
$(32,18)$&6&6&7\\
$(32,19)$&5&7&6
\end{tabular}
\end{center}
For example, at $(N,j)=(16,9)$ the first difference is
$-2042976/2042975$, whereas the two-row difference is
$-2215603929728/2215923048625$. These exact counterexamples refute
the additional universal isometry proposal. They preserve the
threshold-free lower bound and the original conjecture.
\end{remark}

\section{Analytic interpolation and exact Taylor radii}
\label{sec:analytic-radii-final}

The threshold-free Mahler theorem gives continuous interpolation.
Analytic interpolation needs a stronger estimate, and the factorial
bounds provide it directly. The ordinary Taylor radius is determined
by the coefficient growth, including infinite exact subsequences.

\begin{theorem}[Taylor coefficient bounds and radii]
\label{thm:four-taylor-radii}
For $j=0,1,2,3$, the dyadic interpolation has an ordinary power series
\[
 \widetilde C_j(x)=1+\sum_{\ell\geq1}d_{j,\ell}x^\ell,
 \qquad v_2(d_{j,\ell})\geq\ell+\beta_j.
\]
Its Taylor series converges for $|x|_2<2$.
At addresses zero and one,
\begin{equation}
 v_2(d_{j,\ell})=\ell\quad(\ell\geq1),
 \qquad R_0=R_1=2.
 \label{eq:zero-one-exact-taylor-radius}
\end{equation}
At address two,
\begin{equation}
 v_2(d_{2,\ell})\geq
       \max\left(2\ell+\left\lfloor\frac\ell2\right\rfloor,
                          \ell+3\right),
 \qquad R_2=2^{5/2}.
 \label{eq:two-exact-taylor-radius}
\end{equation}
The radius at address three is at least two; its exact radius is not
asserted here. The stated exact radii at addresses zero, one, and two
hold at every center in $\mathbb Z_2$.
\end{theorem}

\begin{proof}
Put $a_m=\Delta_1^mC_j(0)$ and expand the falling factorial in the
Newton series. The ordinary coefficients are
\begin{equation}
 d_{j,\ell}=\sum_{m\geq\ell}\frac{a_m}{m!}s(m,\ell).
 \label{eq:newton-taylor-stirling}
\end{equation}
Since $s(m,\ell)$ is integral and
$v_2(a_m/m!)\geq m+\beta_j$, the coefficient sum converges and has
valuation at least $\ell+\beta_j$. The sum exchange is justified on
$\mathbb Z_2$ by these uniform geometric coefficient bounds.
The ordinary power series agrees there with the Newton interpolation
and converges whenever $|x|_2<2$.

For $j=0,1$, the term $m=\ell$ in
\eqref{eq:newton-taylor-stirling} has exact valuation $\ell$, since
$s(\ell,\ell)=1$ and the Newton valuation is exact. Every later term
has valuation at least $\ell+1$. Thus
$v_2(d_{j,\ell})=\ell$, and the coefficient radius is exactly two.
At a center $x_0\in\mathbb Z_2$, the coefficient of $y^\ell$ in the
translated series is
\begin{equation}
 \sum_{m\geq\ell}d_{j,m}\binom m\ell x_0^{m-\ell}.
 \label{eq:taylor-translation-coefficient}
\end{equation}
Its initial term again has exact valuation $\ell$, and all later
terms have at least one more bit. The exact radius is therefore two
at every integral center.

For address two, the stronger factorial filtration gives the lower
bound in \eqref{eq:two-exact-taylor-radius}. At each
$\ell=2^r\geq4$, Theorem~\ref{thm:two-exact-power-orders} makes the
$m=\ell$ term in \eqref{eq:newton-taylor-stirling} have valuation
$2\ell+\ell/2=5\ell/2$. Every $m>\ell$ term has strictly greater
lower bound $2m+\lfloor m/2\rfloor$. Hence
\[
 v_2(d_{2,\ell})=5\ell/2\qquad(\ell=2^r\geq4).
\]
The universal lower bound and this unbounded exact subsequence prove
that the coefficient radius is exactly $2^{5/2}$. Translation by an
integral center preserves the same unique leading coefficients at
these degrees in \eqref{eq:taylor-translation-coefficient}, proving
the exact radius at every such center.
\end{proof}

\begin{remark}[Boundary interpretation]
For addresses zero and one, at an argument of absolute value two every
positive-degree Taylor term has valuation zero, so that ordinary Taylor
series diverges. For address two, over a complete extension admitting
arguments of absolute value $2^{5/2}$, the power-of-two-degree terms
have valuation zero and fail to tend to zero. These statements concern
the Taylor expansions just constructed; they do not prohibit a
different analytic continuation. On $\mathbb Q_2$-valued arguments,
open-disc notation must be read with its discrete value group. The
coefficient radius is the precise invariant in the theorem.
\end{remark}

\section{The final congruence in the original Euler convention}
\label{sec:final-original-congruence}

The preceding bounds can be stated as one uniform congruence without
suppressing the address normalization. This is the strengthened formula
for the four-address family, valid from the first row.

\begin{corollary}[Final common-step congruence]
\label{cor:final-euler-congruence}
Let $j\in\{0,1,2,3\}$, $n\geq0$, and $k,H\geq1$. Define
\[
 A_0(k)=A_1(k)=k,\qquad
 A_2(k)=\max\left(k+4,\ 2k+\left\lfloor\frac k2\right\rfloor+1\right),
 \qquad A_3(k)=k+2.
\]
Then
\begin{equation}
 \boxed{\displaystyle
 \sum_{i=0}^k(-1)^{k-i}\binom ki E_{4(n+iH)}^{(4,j)}
       \equiv0\pmod{2^{\,v_2(k!)+k v_2(H)+A_j(k)}}.}
 \label{eq:boxed-final-euler-congruence}
\end{equation}
For $j=0,1$, this displayed exponent is the exact valuation for every
$n,k,H$ in the stated range. For $j=2$, the exact positive
power-of-two-order valuations are given by
Theorem~\ref{thm:two-exact-power-orders}, with one added bit.
\end{corollary}

\begin{proof}
The left side is $j!\Delta_H^kC_j(n)$.
At addresses zero and one, $v_2(j!)=0$ and the exact homogeneous
formula gives the exponent. At address two, $v_2(2!)=1$; the maximum
of \eqref{eq:factorial-four-lower} and
\eqref{eq:factorial-four-j2-filtration} yields
$\max(k+4,2k+\lfloor k/2\rfloor+1)$ after removing the common
$v_2(k!)+k v_2(H)$. At address three, the one-bit normalized offset
and $v_2(3!)=1$ give $k+2$.
\end{proof}

For mixed steps the same corollary holds with $k v_2(H)$ replaced by
$\sum_\nu v_2(H_\nu)$ and the binomial stencil replaced by
$j!\Delta_{H_1}\cdots\Delta_{H_k}C_j(n)$.
At addresses zero and one, exactness is then governed by the parity
classification in Theorem~\ref{thm:zero-one-exact-mixed}.
At first order these bounds are exact in every address, with original
valuation offsets $(1,1,5,3)$; the least periods are those in
\eqref{eq:rank-four-least-period-final}.

\subsection{Proof coverage and remaining questions}

The results in these sections are mathematical theorems with detailed
proofs and preserved audit evidence. Their formal status must be read
statement by statement. The recovered Paper~IX release already contains
Lean results for the rank-four threshold-free homogeneous bound, the
sharp first-difference laws, the least periods, and the dyadic isometry
and image conclusions. Those inherited results should not be described
as missing merely because the earlier short completion supplement did
not include their later modules.

The all-rank dyadic and general prime-power transfer proof presented
here, the factorial improvements, the mixed leading-parity
classification, the Bernoulli power-order sharpness, the all-rank
lower-half exact laws, the complete rank-eight extension, and the exact
Taylor-radius arguments have not received new Lean formalization in
this strengthening. A mathematical audit and an exact computation are
neither a Lean compilation nor external peer review. The separate Lean
guide records the recovered formal declarations and any fresh release
checks without extending their coverage to these new claims.

The completed external statement is Nishibuchi's Conjecture~3.9 in
its prescribed parameter family and coefficient normalization.
These stronger congruences enlarge its dyadic branch and remove its
starting threshold there. They do not settle every Euler problem.
In particular, the following questions remain separate:
\begin{itemize}
 \item a complete exact higher-difference law at all non-power-of-two
       orders in address two;
 \item complete higher-difference valuations and the exact ordinary
       Taylor radius in address three;
 \item a classification of the upper-half dyadic addresses having an
       exact first-difference metric law at arbitrary larger ranks;
 \item the distinct universal simplicity or nonvanishing question
       for the later Euler-tail residuals, and the original native
       Paper~IX conjectures beyond the completed external statement.
\end{itemize}
The exact counterexamples above rule out the naive universal upper-half
isometry, and mixed cancellations rule out unconditional baseline
sharpness for arbitrary stencils. No global historical-priority or
novelty assertion is made for these strengthened statements.

\clearpage
\section{Finite obstructions, compatible lifts, and deformation}
\label{sec:lifting}

The arithmetic results above concern particular families. We now study the
underlying correction problem, from finite polynomial lifts to continuous
parameter charts and quantitative deformation. Throughout, we distinguish
three objects: a residual obstruction, the set of correcting digits, and a
root in a completion. A vanishing obstruction at one precision gives a lift
to the next. Continuing indefinitely requires additional hypotheses.

Unless a different coefficient ring is specified, a polynomial system is a
finite list of integer polynomials. All divisibility conditions on vectors are
coordinatewise. For a positive integer $q$, write $S_q=\Z/q\Z$ and use an
overline for reduction modulo $q$. For a prime $p$, write
$\F_p=S_p$, normalize $|p|_p=p^{-1}$, and put
\[
 B_L(x)=x+p^L\Zp^d.
\]
An integer center is included in this notation by its natural image in
$\Zp^d$. All dimensions below are finite and may be zero. A statement over
$S_q$ allows composite $q$ and $q=1$ when explicitly stated. A cardinality
formula in terms of vector-space rank uses prime $p$.

\subsection{The polynomial correction equation}

For $F=(F_1,\ldots,F_e)\in\Z[X_1,\ldots,X_d]^e$, define its evaluated
Jacobian by the formal partial derivatives:
\[
 DF(x)_{ij}=(\partial_jF_i)(x),\qquad
 J_x=DF(x)\bmod q:S_q^d\longrightarrow S_q^e.
\]
The obstruction therefore uses the Jacobian of the given polynomials. The
following algebraic Taylor identity relates this linear map to polynomial
evaluation.

\begin{lemma}[Polynomial Taylor remainder]\label{lem:polynomial-taylor}
Let $R$ be a commutative ring, $P\in R[X_1,\ldots,X_d]$, $a\in R^d$, and
$c\in R$. There is a polynomial $Q\in R[U_1,\ldots,U_d]$, with $Q(0)=0$,
such that for every $u\in R^d$,
\begin{equation}\label{eq:polynomial-taylor}
 P(a+cu)=P(a)+c\sum_{j=1}^d(\partial_jP)(a)u_j+c^2Q(u).
\end{equation}
In particular, when $R=\Z$ and $a,u,c$ are integral, subtracting the first
two terms on the right from $P(a+cu)$ gives a multiple of $c^2$.
\end{lemma}

\begin{proof}
Apply polynomial induction using constants, addition, and multiplication by a
variable. Constants have zero remainder, and remainders add. For the
multiplication step, suppose the assertion is known for $P$ and write
$D_P(U)=\sum_j(\partial_jP)(a)U_j$. Multiplication of its expansion by
$a_i+cU_i$ gives constant term $P(a)a_i$, linear term
$c(a_iD_P(U)+P(a)U_i)$, and remainder
\[
 c^2\bigl(Q(U)(a_i+cU_i)+D_P(U)U_i\bigr).
\]
The product rule identifies the linear term with the derivative of
$PX_i$. The displayed remainder vanishes at $U=0$. This constructs the
polynomial remainder before evaluation at $u$, even over rings where
division by $c$ would be invalid.
\end{proof}

The general polynomial identity is formalized by
\lean{Komorebi.PolynomialTaylorRing.exists_eval_scaled_remainder}.
Its integer divisibility form is
\lean{Komorebi.PolynomialTaylor.taylor_dvd}.

\begin{theorem}[Finite polynomial obstruction]\label{thm:finite-polynomial-obstruction}
Let $q\ge1$, $L\ge1$, $x\in\Z^d$, and suppose
$F(x)\in q^L\Z^e$. Define the exact integer residual $r_x$ by
$F(x)=q^Lr_x$, and put
\[
 \mathcal O_x=S_q^e/\im J_x,\qquad o_x=[\overline{r_x}].
\]
Then $\mathcal O_x$ is finite and
\begin{equation}\label{eq:finite-obstruction}
 o_x=0\quad\Longleftrightarrow\quad
 \exists y\in\Z^d:\quad
 y\equiv x\pmod{q^L},\qquad F(y)\equiv0\pmod{q^{L+1}}.
\end{equation}
No surjectivity or nonsingularity assumption is required.
\end{theorem}

\begin{proof}
Every point in the lower residue class has the form $y=x+q^Lt$ with
$t\in\Z^d$. By Lemma~\ref{lem:polynomial-taylor} and $2L\ge L+1$,
\[
 F(x+q^Lt)\equiv q^L\bigl(r_x+DF(x)t\bigr)\pmod{q^{L+1}}.
\]
The integer $q^L$ is nonzero. Cancelling it in the corresponding divisibility
condition proves
\begin{equation}\label{eq:digit-correction}
 F(x+q^Lt)\equiv0\pmod{q^{L+1}}
 \quad\Longleftrightarrow\quad
 \overline{r_x}+J_x\overline t=0.
\end{equation}
Integer vectors reduce onto $S_q^d$. The equation on the right is solvable
exactly when $\overline{r_x}$ belongs to $\im J_x$, since that image is
closed under negation. This is precisely $o_x=0$. Finiteness follows because
$\mathcal O_x$ is a quotient of the finite module $S_q^e$.
\end{proof}

The formal theorem
\lean{Komorebi.PolynomialObstruction.finite_polynomial_obstruction}
presents the image using the integer-domain map
$\Z^d\xrightarrow{DF(x)}\Z^e\to S_q^e$. Reduction onto $S_q^d$ identifies
that image with $\im J_x$. The class-preserving comparison is explicit in
\lean{Komorebi.PolynomialLiftCounts.obstructionEquiv}; the change of
presentation does not change the distinguished residual class.

\begin{proposition}[Affine depth zero]\label{prop:affine-obstruction}
Let $F(x)=Ax-b$ for an integer linear map $A:\Z^d\to\Z^e$. The criterion
of Theorem~\ref{thm:finite-polynomial-obstruction} holds for every $L\ge0$,
with $J=A\bmod q$.
\end{proposition}

\begin{proof}
The affine identity
$F(x+q^Lt)=F(x)+q^LAt$ has no quadratic remainder. The proof of
\eqref{eq:digit-correction} therefore needs no inequality involving $2L$.
The residual and the finite quotient are obtained as above.
\end{proof}

The affine statement is formalized by
\lean{Komorebi.NativeAffine.finite_affine_obstruction}.
Positive depth is essential in the nonlinear theorem: for $F(X)=X^2-1$,
$q=2$, $L=0$, and $x=0$, the Jacobian is zero and the residual class is
nonzero, yet $y=1$ is a root modulo two in the same residue modulo one.

\begin{proposition}[Independence of the integer representative]
\label{prop:representative-independence}
Under the hypotheses of Theorem~\ref{thm:finite-polynomial-obstruction},
let $x'=x+q^Ls$ and write $F(x')=q^Lr_{x'}$. Then
\[
 J_{x'}=J_x,\qquad
 \overline{r_{x'}}=\overline{r_x}+J_x\overline s.
\]
Consequently the two residual classes coincide in the common quotient
$S_q^e/\im J_x$.
\end{proposition}

\begin{proof}
Polynomial evaluation respects coordinatewise congruence. Apply this first
to every formal partial derivative: $x'\equiv x\pmod q$ gives
$J_{x'}=J_x$. Next apply \eqref{eq:polynomial-taylor} to $F$ at scale $q^L$.
Subtract the two exact residual identities and cancel $q^L$ to obtain the
displayed relation modulo $q$. Its additional term belongs to the common
Jacobian image, so it vanishes in the quotient.
\end{proof}

This comparison is formalized by
\lean{Komorebi.PolynomialRepresentative.representative_independence}.
This proves independence of the integer representative within a fixed
polynomial presentation. It makes no identification of the raw cokernels
of different equation lists. For example, adjoining a zero equation leaves
the roots unchanged while adding an $S_q$ summand to the raw cokernel.

\subsection{One correction shared by two systems}

The next obstruction concerns compatibility between individually attainable
corrections. Its algebraic form works over an arbitrary commutative ring.
Let $U,V,W$ be $R$-modules and let $J:U\to V$, $K:U\to W$ be linear.
Define
\begin{equation}\label{eq:common-quotient}
 C(J,K)=U/(\ker J+\ker K).
\end{equation}
An attainable value $Ju$ has label $[u]\in C(J,K)$; this is independent of
the choice of its preimage because $\ker J$ is killed in the quotient.
Similarly, $Kv$ has label $[v]$. Denote these maps by
$\lambda_J:\im J\to C(J,K)$ and $\lambda_K:\im K\to C(J,K)$.

\begin{theorem}[Common correction and fiber product]\label{thm:common-correction}
For every $u,v\in U$,
\begin{equation}\label{eq:common-correction}
 [u]=[v]\text{ in }C(J,K)
 \quad\Longleftrightarrow\quad
 \exists w\in U:\ Jw=Ju,\quad Kw=Kv.
\end{equation}
Both label maps are surjective, and the image of the joint map is exactly
the fiber product:
\begin{equation}\label{eq:joint-fiber-product}
 \im\bigl(U\to\im J\times\im K,\ u\mapsto(Ju,Ku)\bigr)
 =\im J\mathbin{\times}_{C(J,K)}\im K.
\end{equation}
If either $V$ or $W$ is finite, then $C(J,K)$ is finite.
\end{theorem}

\begin{proof}
If $[u]=[v]$, write $u-v=s+t$ with $s\in\ker J$ and $t\in\ker K$.
Then $w=u-s=v+t$ satisfies both required equalities. Conversely, a common
$w$ gives
\[
 u-v=(u-w)+(w-v)\in\ker J+\ker K.
\]
Every class $[u]$ is the label of $Ju$ and of $Ku$, proving surjectivity.
For attainable $a=Ju$ and $b=Kv$, the first equivalence says exactly that
equal labels are equivalent to a simultaneous preimage. This proves
\eqref{eq:joint-fiber-product}. Finally, a finite target has finite image,
and the associated surjective label map makes $C(J,K)$ finite.
\end{proof}

These statements are
\lean{Komorebi.SharedCorrection.quotient_mk_eq_iff_common} and
\lean{Komorebi.SharedCorrection.joint_range_eq_fiber}, together with the
finite-target corollaries in the same module. The quotient also has the
expected universal property: every linear map from $U$ annihilating both
kernels factors uniquely through $C(J,K)$. Equivalently, compatible linear
observations on $\im J$ and $\im K$ factor uniquely through their label
maps. This is formalized by
\lean{Komorebi.SharedCorrection.common_paired_universal}.

\begin{corollary}[Simultaneous polynomial lifting]\label{cor:shared-polynomial}
Let $F\in\Z[X_1,\ldots,X_d]^e$ and
$G\in\Z[X_1,\ldots,X_d]^f$ have a common lower root $x$ modulo $q^L$,
where $q\ge1$ and $L\ge1$. Use the integer-domain reduced Jacobians
$J:\Z^d\to S_q^e$ and $K:\Z^d\to S_q^f$. Suppose $u,v\in\Z^d$ are
separately valid corrections for $F$ and $G$, respectively. Then
\[
 [u]=[v]\text{ in }\Z^d/(\ker J+\ker K)
\]
if and only if a single $y\equiv x\pmod{q^L}$ solves both systems modulo
$q^{L+1}$. Each separate label is independent of the chosen valid
correction. The common quotient is finite.
\end{corollary}

\begin{proof}
For each system, \eqref{eq:digit-correction} translates a valid next root
into a linear correction equation. Once $u$ corrects $F$, another digit
$w$ corrects $F$ precisely when $Jw=Ju$. The analogous condition for $G$
is $Kw=Kv$. Theorem~\ref{thm:common-correction} now gives the criterion.
Two valid corrections of one block differ by a vector in its Jacobian
kernel, so they have equal labels. Finiteness follows from either finite
residue target. An unconditional formulation quantifies over separately
valid $u,v$ with equal labels, thereby including their existence in the
criterion rather than assuming it.
\end{proof}

See \lean{Komorebi.PolynomialShared.joint_lift_iff_class_eq} and
\lean{Komorebi.PolynomialShared.finite_polynomial_pair_compatibility}.

For affine integer systems $F(X)=AX-b$ and $G(X)=BX-c$, where
$A:\Z^d\to\Z^e$ and $B:\Z^d\to\Z^f$ are linear,
Corollary~\ref{cor:shared-polynomial} holds for $q\ge1$ and every
$L\ge0$, including depth zero. Here $x$ is a common root modulo $q^L$,
and $J$ and $K$ are the reductions of $A$ and $B$ with domain $\Z^d$.
Given separately valid corrections $u,v$, a simultaneous next lift exists
if and only if their classes in $\Z^d/(\ker J+\ker K)$ agree. The quotient
is finite, and each label is independent of the chosen valid correction.
Indeed, Proposition~\ref{prop:affine-obstruction} supplies the correction
equation without a positive-depth restriction, and the proof of
Theorem~\ref{thm:common-correction} then applies unchanged. The formal
statement is \lean{Komorebi.AffinePair.joint_lift_iff_class_eq}.
Quantifying over separately valid corrections with equal labels also
includes their existence in the criterion, as in
\lean{Komorebi.AffinePair.finite_affine_pair_compatibility}.
Separate liftability alone is insufficient. At $x=0$ modulo two,
$F(X)=X$ and $G(X)=X-2$ both lift to modulus four, but their correcting
digits are respectively $0$ and $1$ modulo two. No single residue solves
both equations modulo four.

\subsection{Correction torsors and exact residue counts}

Integer representatives and residue classes at the next level must also be
distinguished. For $q\ge1$ and $L\ge0$, the map
\[
 S_q^d\longrightarrow
 \{y\in(\Z/q^{L+1}\Z)^d:y\equiv x\pmod{q^L}\},
 \qquad t\longmapsto x+q^L\widetilde t
\]
is a bijection, where $\widetilde t$ is any integer representative.
Independence of the representative and injectivity both follow by cancelling
$q^L$ in integer divisibility; surjectivity follows by expressing the lower
congruence as a multiple of $q^L$. This is the equivalence
\lean{Komorebi.ResidueLiftDigits.digitEquivResidue}.

At positive polynomial depth, restrict it to the correction fiber
\[
 \mathcal T_x=\{t\in S_q^d:\overline{r_x}+J_xt=0\}
\]
and the lift set
\[
 \mathcal L_{x,L}=
 \{y\in(\Z/q^{L+1}\Z)^d:y\equiv x\pmod{q^L},\ F(y)=0\}.
\]
The polynomial predicate in the second display is evaluation in the residue
ring. Evaluation on canonical integer representatives gives the same
predicate, because integer polynomial evaluation respects congruence.

\begin{theorem}[Torsor and lift cardinalities]\label{thm:lift-torsor-count}
Under the hypotheses of Theorem~\ref{thm:finite-polynomial-obstruction},
there is an equivalence $\mathcal T_x\simeq\mathcal L_{x,L}$.
Every nonempty fiber is a torsor for $\ker J_x$ under addition of
correction digits. If $q=p$ is prime and
$\rho_x=\rank_{\F_p}J_x$, then
\begin{equation}\label{eq:lift-count}
 \#\mathcal L_{x,L}=
 \begin{cases}p^{d-\rho_x},&o_x=0,\\0,&o_x\ne0,\end{cases}
 \qquad
 \#\mathcal O_x=p^{e-\rho_x}.
\end{equation}
\end{theorem}

\begin{proof}
Equation~\eqref{eq:digit-correction} identifies exactly which points of the
residue fiber satisfy the equations, proving the restricted equivalence.
If $u,t\in\mathcal T_x$, subtraction gives $J_x(t-u)=0$. Conversely,
adding any element of $\ker J_x$ preserves the correction equation. The
unique kernel vector carrying $u$ to $t$ is $t-u$, so the action is free
and transitive. Choosing one valid digit identifies the fiber with the
kernel, without making that choice part of its translation action.

For prime $p$, rank-nullity gives
$\dim\ker J_x=d-\rho_x$ and
$\dim(\F_p^e/\im J_x)=e-\rho_x$. The cardinality of an
$n$-dimensional $\F_p$-space is $p^n$. Combine these identities with the
obstruction criterion and the residue equivalence.
\end{proof}

The cardinality formula is formalized by
\lean{Komorebi.PolynomialLiftCounts.finite_polynomial_lift_count}.
The explicit torsor and chosen-point equivalence are developed in
\lean{Komorebi.FiniteFiberCounts.translate_exists_unique} and
\lean{Komorebi.PolynomialLiftCounts.liftEquivKernel}.
For composite $q$, the torsor and kernel cardinality remain valid, while the
vector-space rank formula is unavailable. Multiplication by two on
$\Z/4\Z$ already has a kernel of size two.

\subsection{Smooth compatible towers and exact roots}

Surjectivity removes every residual obstruction and persists throughout a
positive-depth residue ball. This supplies the additional hypothesis needed
to iterate the finite theorem.

\begin{theorem}[Compatible polynomial tower]\label{thm:smooth-tower}
Let $q\ge1$, $L\ge1$, $F\in\Z[X_1,\ldots,X_d]^e$, and
$x\in\Z^d$ satisfy $F(x)\in q^L\Z^e$. Suppose $J_x:S_q^d\to S_q^e$
is surjective. There is a single sequence $z_n\in\Z^d$ such that
\begin{align*}
 z_0&=x,& F(z_n)&\in q^{L+n}\Z^e,\\
 z_{n+1}-z_n&\in q^{L+n}\Z^d,& J_{z_n}&=J_x.
\end{align*}
For $a\le b$, it also satisfies $z_b-z_a\in q^{L+a}\Z^d$.
\end{theorem}

\begin{proof}
At level $N=L+n$, divide the residual by $q^N$. Surjectivity supplies
$t\in S_q^d$ with $J_{z_n}t=-\overline{F(z_n)/q^N}$. Choose an integer
representative and apply \eqref{eq:digit-correction} to obtain $z_{n+1}$.
The new point is congruent to $z_n$ modulo $q$. Evaluation of each formal
partial derivative therefore gives $J_{z_{n+1}}=J_{z_n}$.

Construct the sequence recursively, retaining the root condition and the
fixed Jacobian at each step. Surjectivity at every later stage then follows
from the single initial assumption. The cross-level congruence follows by
induction, adding consecutive differences whose divisibilities all imply divisibility
by $q^{L+a}$.
\end{proof}

The full tower statement is
\lean{Komorebi.PolynomialIntegerTower.exists_compatible_tower_strong}.
The proof constructs a compatible path using classical choice. It does not
claim an optimized procedure for selecting its digits.

\begin{theorem}[Completion and smooth branching]\label{thm:smooth-padic-root}
In Theorem~\ref{thm:smooth-tower}, let $q=p$ be prime. The tower can be
chosen with a limit $a\in\Zp^d$ such that
\[
 F(a)=0,\qquad z_n\longrightarrow a,\qquad
 a-z_n\in p^{L+n}\Zp^d\quad(n\ge0).
\]
At each stage, the number of distinct next-level residue roots over $z_n$
is exactly $p^{d-e}$, and the finite obstruction module has one element.
If $d=e$, there is exactly one root of $F$ in $B_L(x)$.
\end{theorem}

\begin{proof}
Cross-level compatibility bounds the supremum distance between later terms
by $p^{-(L+n)}$. Completeness gives a limit, and closed norm bounds retain
the congruence to every $z_n$. View the coefficients of the original integer
polynomials in $\Zp$. Their evaluations satisfy
both $F_i(z_n)\to F_i(a)$, by continuity, and $F_i(z_n)\to0$, by their
growing divisibility. Uniqueness of limits gives exact vanishing.

Surjectivity gives $\rank J_x=e\le d$. Theorem~\ref{thm:lift-torsor-count}
then gives $p^{d-e}$ next residue roots and obstruction cardinality one at
every stage.

For square uniqueness, let $a,b\in B_L(x)$ be arbitrary exact roots.
Their canonical integer truncations are genuine roots in the respective
finite residue rings, since evaluation commutes with reduction. They agree
at depth $L$. If they agree at depth $N\ge L$, both next residues belong
to the same lift set, which has cardinality $p^{d-d}=1$. Induction gives
agreement at every depth. The residue maps separate $p$-adic integers, so
$a=b$. Existence was already supplied by the tower.
\end{proof}

The endpoint construction is
\lean{Komorebi.PolynomialPadicLift.exists_padic_root_with_tower}, the
branching statement is
\lean{Komorebi.PolynomialSmoothCounts.exists_regular_tower}, and uniqueness
for all roots in the ball is
\lean{Komorebi.PolynomialPadicUnique.exists_unique_padic_root}.

\begin{remark}\label{rem:singular-tower-boundary}
The one-step singular criterion and the smooth infinite theorem have
different hypotheses. The polynomial $X^2-5$ has odd roots modulo two and
four but has no root modulo eight. Thus a previous zero obstruction does
not force the next one to vanish. Also, rectangular surjectivity alone
does not imply uniqueness: $F(X,Y)=X$ has the distinct dyadic roots
$(0,0)$ and $(0,2)$ in the same residue ball modulo two. These examples show
why the hypotheses of the tower theorem and its uniqueness clause are needed.
\end{remark}

\subsection{Coefficient stability and continuous root charts}

\begin{proposition}[Smooth coefficient stability]\label{prop:smooth-coefficient-stability}
Let $p$ be prime, $M\ge1$, $F,H\in\Z[X_1,\ldots,X_d]^e$, and
$c\in p^M\Z$. Suppose $DF(x)\bmod p$ is surjective for an integer
$x\in\Z^d$. If $a\in\Zp^d$ satisfies $F(a)=0$ and
$a\equiv x\pmod p$, then there is a root $b$ of $G=F+cH$ with
$b-a\in p^M\Zp^d$.

If $d=e$ and $M\ge L\ge1$, every pair of roots $a,b\in B_L(x)$ of
$F$ and $G$, respectively, satisfies $b-a\in p^M\Zp^d$. The perturbed
root is unique in that ball whenever it exists.
\end{proposition}

\begin{proof}
Since $p\mid c$, differentiation gives
$DG(x)\equiv DF(x)\pmod p$. Truncate $a$ to an integer vector $y$ at
depth $M$. Then $F(y)\equiv0\pmod{p^M}$, hence also
$G(y)\equiv0\pmod{p^M}$. Congruence of $y$ and $x$ modulo $p$ preserves
the reduced Jacobian. The smooth root theorem gives
$b\equiv y\equiv a\pmod{p^M}$.

In the square case this constructed root lies in $B_L(x)$ because $M\ge L$.
Square uniqueness for $G$ identifies it with every competing root in that
ball. The construction therefore gives the precision assertion for an
arbitrary root there.
\end{proof}

See \lean{Komorebi.PolynomialRootStability.exists_nearby_perturbed_root}
and \lean{Komorebi.PolynomialRootStability.perturbed_roots_congruent}.
The scalar $c$ may be negative or zero. The weighted theorems below allow
arbitrary integral $p$-adic coefficients and give an explicit first-order
displacement.

\begin{theorem}[Continuous charts with a selected Jacobian minor]
\label{thm:continuous-parameter-chart}
Let $p$ be prime, $L\ge1$,
$F\in\Z[X_1,\ldots,X_{e+k}]^e$, and $x\in\Z^{e+k}$ satisfy
$F(x)\in p^L\Z^e$. Suppose the full reduced Jacobian is surjective.
There is a permutation $\sigma$ of the coordinates such that, writing
$\sigma x=(y_0,t_0)$ with $e$ solved coordinates and $k$ parameter
coordinates, the parameter ball
$T=t_0+p^L\Zp^k$ is homeomorphic to the entire root locus
\[
 \mathcal R=\{z\in B_L(x):F(z)=0\}.
\]
The inverse homeomorphism is the actual free-coordinate projection
$z\mapsto\operatorname{free}(\sigma z)$. If $h:T\to\mathcal R$ is the
forward map, then for every $N\ge0$,
\begin{equation}\label{eq:chart-congruence}
 h(t)\equiv h(u)\pmod{p^N}
 \quad\Longleftrightarrow\quad t\equiv u\pmod{p^N}.
\end{equation}
Equivalently, $\sigma h(t)=(\phi(t),t)$ for a continuous
$\phi:T\to y_0+p^L\Zp^e$ whose graph is the entire reordered root locus.
\end{theorem}

\begin{proof}
First consider a square integer polynomial map $P$ whose reduced Jacobian
at $x$ is surjective. Its Jacobian is invertible over $\F_p$. A digit
comparison using Taylor expansion shows that equal outputs modulo
$p^{N+1}$ force equal inputs modulo $p^{N+1}$ whenever the inputs already
agree modulo $p^N$ and belong to $B_1(x)$. Induction and integer truncation
give, for arbitrary $a,b\in B_1(x)$,
\[
 P(a)\equiv P(b)\pmod{p^N}
 \quad\Longleftrightarrow\quad a\equiv b\pmod{p^N}.
\]

More generally, for a rectangular map with surjective reduced Jacobian,
\begin{equation}\label{eq:polynomial-ball-image}
 P(B_L(x))=P(x)+p^L\Zp^e.
\end{equation}
One inclusion follows from polynomial congruence. For integer targets in
the right-hand ball, subtract the target from the constant terms and use
the smooth root theorem. For arbitrary $p$-adic targets, approximate by
integer truncations. The input ball is compact, so its continuous image is
closed. It contains all these approximants and hence their limit. This
argument does not assume that independently chosen preimages form a
compatible sequence.

For square $P$, the resulting bijection between the two compact residue
balls is continuous, and its target is Hausdorff. It is therefore a
homeomorphism, with the same congruence preservation for its inverse.

Now suppose a chosen set of $e$ columns of $DF(x)\bmod p$ is independent,
and put those coordinates first. Apply the square result to the augmented
polynomial map
\[
 \Psi(Y,T)=(F(Y,T),T).
\]
Its reduced Jacobian has upper blocks $(D_YF,D_TF)$ and lower blocks
$(0,I_k)$, so it is invertible. For $t\in T$, the target $(0,t)$ belongs
to $\Psi(x)+p^L\Zp^{e+k}$ because $F(x)\equiv0\pmod{p^L}$.
The inverse of $\Psi$ at that target gives the unique root with parameters
$t$. The inverse congruence property gives one direction of
\eqref{eq:chart-congruence}; projection onto the parameter coordinates
gives the other.

Finally, surjectivity of the original rectangular Jacobian guarantees a
basis among its columns. Extend the corresponding injection of coordinate
indices to a permutation $\sigma$. Renaming the variables by
$\sigma^{-1}$ transports polynomial evaluation and every formal partial
derivative. The fixed-coordinate construction therefore applies to the
original system after undoing the permutation.
\end{proof}

The main theorem is
\lean{Komorebi.PolynomialGeneralChart.exists_parameter_homeomorph}.
The square local inverse and surjectivity for arbitrary targets are
\lean{Komorebi.PolynomialLocalInverse.localHomeomorph} and
\lean{Komorebi.PolynomialLocalSurjective.exists_padic_preimage}.
The center $x$ is only an approximate root: the chart need not satisfy
$\phi(t_0)=y_0$ unless $F(x)=0$ exactly. The assertion selects suitable
coordinates; it does not permit an arbitrary prescribed solved block.

\subsection{Prescribed parameters and their obstruction}

Suppose now that the split into solved variables and parameter variables
is fixed in advance. Let
\[
 F\in\Z[Y_1,\ldots,Y_m,T_1,\ldots,T_k]^e,\qquad x=(y_0,t_0),
\]
and assume $F(x)=q^Lr_x$ with $q\ge1$ and $L\ge1$. Write the reduced
Jacobian as
\[
 DF(x)\bmod q=(J_x\ K_x),\qquad
 J_x:S_q^m\to S_q^e,\quad K_x:S_q^k\to S_q^e.
\]
For a prescribed parameter digit $v\in S_q^k$, let $\mathcal L_{x,L,v}$
be the set of next residue roots over $x$ whose parameter coordinates are
$t_0+q^L\widetilde v$ modulo $q^{L+1}$.

\begin{theorem}[Prescribed-parameter obstruction and count]
\label{thm:prescribed-obstruction}
Define
\[
 o_x(v)=[\overline{r_x}+K_xv]\in S_q^e/\im J_x.
\]
Then $\mathcal L_{x,L,v}$ is nonempty if and only if $o_x(v)=0$.
When nonempty, it is a torsor for $\ker J_x$. If $q=p$ is prime and
$\rho_x=\rank_{\F_p}J_x$, then
\begin{equation}\label{eq:prescribed-count}
 \#\mathcal L_{x,L,v}
 =\begin{cases}p^{m-\rho_x},&o_x(v)=0,\\0,&o_x(v)\ne0,\end{cases}
 \qquad \#(\F_p^e/\im J_x)=p^{e-\rho_x}.
\end{equation}
The criterion holds uniformly for every prescribed $v$ without a
surjectivity hypothesis.
\end{theorem}

\begin{proof}
A candidate for the next lift has the form
$x+q^L(\widetilde u,\widetilde v)$, with solved digit $u\in S_q^m$.
Its Taylor correction equation is
\[
 \overline{r_x}+J_xu+K_xv=0.
\]
Solvability is exactly vanishing of the displayed cokernel class. Restrict
the residue-digit equivalence to the prescribed parameter coordinates.
This identifies the lift fiber with the solved-digit fiber.
Translation by $\ker J_x$ is free and transitive on any nonempty such
fiber. The prime-field count follows from rank-nullity as before.
\end{proof}

The formal statements start from the lower root and construct its residual:
\lean{Komorebi.PolynomialPrescribedObstruction.finite_prescribed_obstruction}
and \lean{Komorebi.PolynomialPrescribedObstruction.finite_prescribed_lift_count}.
For $F(Y,T)=T$ at the dyadic origin, the full Jacobian $(0\ 1)$ is
surjective but $J_x=0$. Prescribing $v=1$ forces $T=2\pmod4$, so no root
exists; prescribing $v=0$ gives two choices of $Y$ modulo four over the
initial zero residue. This example explains why the criterion depends on
the solved block, rather than merely on the full Jacobian.

\begin{theorem}[Every compatible integer parameter path]
\label{thm:prescribed-integer-tower}
In the preceding notation, assume $J_x$ is surjective. For every supplied
integer path $t_n\in\Z^k$ satisfying
\[
 t_0=\operatorname{free}(x),\qquad
 t_{n+1}-t_n\in q^{L+n}\Z^k,
\]
there is a compatible polynomial tower $z_n\in\Z^{m+k}$ with
\[
 z_0=x,\quad \operatorname{free}(z_n)=t_n,\quad
 F(z_n)\in q^{L+n}\Z^e,\quad
 z_{n+1}-z_n\in q^{L+n}\Z^{m+k}.
\]
The free-coordinate equality is exact in integers. Both reduced Jacobian
blocks remain fixed, and $z_b-z_a\in q^{L+a}\Z^{m+k}$ whenever $a\le b$.
\end{theorem}

\begin{proof}
At level $N=L+n$, write the supplied parameter increment exactly as
$t_{n+1}-t_n=q^Nw$ with $w\in\Z^k$. Surjectivity of the solved block
gives a digit $u$ satisfying
$J_{z_n}u=-\overline{F(z_n)/q^N}-K_{z_n}\overline w$.
Choose an integer representative $\widetilde u$ and set
$z_{n+1}=z_n+q^N(\widetilde u,w)$. Its parameter coordinates equal the
supplied target exactly. Taylor expansion proves the next root condition,
while congruence modulo $q$ preserves the full reduced Jacobian.
Construct the sequence recursively, preserving the exact parameter
coordinates and the fixed solved Jacobian. The earlier cross-level
argument proves all the stated integer congruences.
\end{proof}

This is
\lean{Komorebi.PolynomialPrescribedIntegerTower.exists_prescribed_tower_strong}.
In particular, no assumption about solvability at each later level has been
added to the initial surjectivity condition.

\begin{theorem}[Arbitrary $p$-adic parameter endpoints]
\label{thm:prescribed-padic-tower}
Let $q=p$ be prime in Theorem~\ref{thm:prescribed-integer-tower}.
For every $\tau\in t_0+p^L\Zp^k$, there are an integer tower $z_n$ and
$b\in\Zp^{m+k}$ such that
\[
 F(b)=0,\qquad\operatorname{free}(b)=\tau,\qquad
 z_n\longrightarrow b,\qquad b-z_n\in p^{L+n}\Zp^{m+k}.
\]
The initial value is $z_0=x$. At every stage and for every proposed
parameter digit $v\in\F_p^k$, the corresponding next residue fiber has
exactly $p^{m-e}$ elements. If $m=e$, then for each $\tau$ there is exactly
one such root in $B_L(x)$, independently of all choices of integer
approximants and correcting digits.
\end{theorem}

\begin{proof}
Set the first integer parameter approximation equal to the prescribed
$t_0$. For $n\ge1$, use canonical representatives of
$\tau\bmod p^{L+n}$. The initial congruence and truncation precision give
\[
 \tau-t_n\in p^{L+n}\Zp^k.
\]
Subtracting consecutive relations gives compatibility in $\Zp$.
Faithfulness of integer divisibility under $\Z\hookrightarrow\Zp$
gives compatibility in $\Z$, and the same norm bounds give $t_n\to\tau$.
Theorem~\ref{thm:prescribed-integer-tower} and completion now produce $b$
with exact polynomial vanishing and all retained precisions. Each free
coordinate has both the corresponding coordinate of $b$ and that of
$\tau$ as its limit, proving exact parameter equality.

The solved Jacobian remains surjective, so every shifted obstruction
$[\overline r+K_{z_n}v]$ vanishes. Formula~\eqref{eq:prescribed-count}
therefore gives $p^{m-e}$ choices for every parameter digit at every stage.
When $m=e$, the fixed-coordinate chart in the proof of
Theorem~\ref{thm:continuous-parameter-chart} contains every root in the
initial ball with the specified parameters. Its uniqueness proves the
endpoint assertion. This is uniqueness of the actual endpoint, not of the
integer representatives used to approach it.
\end{proof}

The endpoint theorems are
\lean{Komorebi.PolynomialPrescribedPadicTower.exists_prescribed_padic_root_with_tower}
and \lean{Komorebi.PolynomialPrescribedPadicTower.exists_unique_prescribed_padic_root}.
The separate supplied-path theorem
\lean{Komorebi.PolynomialPrescribedPadicTower.exists_root_with_prescribed_path}
accepts any compatible integer parameter path with limit $\tau$.
For $m>e$, endpoint uniqueness is false in general: the system
$F(Y_1,Y_2,T)=Y_1$ has two distinct roots $(0,0,0)$ and $(0,2,0)$ in the
same dyadic ball with identical parameter value.

\subsection{Weighted lifting for a nonunit Jacobian}

The remaining theorems in this section allow arbitrary coefficients and
centers in $R=\Zp$. They concern square solved Jacobians invertible over
$\Qp$, even when their reductions modulo $p$ are singular. The exponent
controlling an integral inverse is retained explicitly.

\begin{theorem}[Weighted nonlinear lifting]\label{thm:weighted-lifting}
Let $p$ be prime, $d,s,r\in\N$, $F\in R[X_1,\ldots,X_d]^d$, and
$a\in R^d$. Set $J=DF(a)$. Suppose
\begin{equation}\label{eq:weighted-hypotheses}
 J\in\operatorname{GL}_d(\Qp),\qquad
 p^sJ^{-1}\in M_d(R),\qquad F(a)\in p^rR^d,\qquad r>2s.
\end{equation}
Then there is exactly one root $b$ of $F$ in $a+p^{s+1}R^d$, and it
satisfies
\begin{equation}\label{eq:weighted-displacement}
 b-a\in p^{r-s}R^d.
\end{equation}
Uniqueness holds throughout the full ball of depth $s+1$.
\end{theorem}

\begin{proof}
Let $A=p^sJ^{-1}\in M_d(R)$. Inclusion into $\Qp$ and its injectivity
give the two integral identities
\[
 AJ=JA=p^sI_d.
\]
These identities are also an equivalent integral certificate for the
first two hypotheses of \eqref{eq:weighted-hypotheses}. In particular,
$A$ is injective: $Aw=0$ implies $p^sw=JAw=0$, hence $w=0$.

Precondition the polynomial system by $A$ and set $c=p^{s+1}$.
Its derivative at $a$ is $AJ=p^sI_d$. Lemma~\ref{lem:polynomial-taylor}
gives a polynomial vector $Q$ with $Q(0)=0$ and
\[
 (AF)(a+cu)=(AF)(a)+p^{2s+1}u+p^{2s+2}Q(u).
\]
Since $A$ is integral, $(AF)(a)\in p^rR^d$. The strict integer inequality
$r>2s$ permits $k=r-(2s+1)\ge0$ and a vector $C\in p^kR^d$ such that
$(AF)(a)=p^{2s+1}C$. The normalized root equation is therefore
\begin{equation}\label{eq:weighted-fixed-point}
 u=-C-pQ(u).
\end{equation}

Integral polynomial evaluation is $1$-Lipschitz for the supremum
$p$-adic metric: a common divisibility of input differences is preserved
under addition and multiplication by integral elements. Consequently the
right side of \eqref{eq:weighted-fixed-point} has contraction constant
$|p|_p<1$ on the complete space $R^d$. It has a unique fixed point $u$.

To obtain the sharper displacement, begin with $u\in R^d$. If
$u\in p^nR^d$, then $Q(0)=0$ and polynomial congruence give
$Q(u)\in p^nR^d$. For $n<k$, both $C$ and $pQ(u)$ are divisible by
$p^{n+1}$. Equation~\eqref{eq:weighted-fixed-point} therefore improves
the divisibility of $u$ by one digit. Induction gives $u\in p^kR^d$.
Thus $b=a+p^{s+1}u$ satisfies \eqref{eq:weighted-displacement}, because
$s+1+k=r-s$. Its normalized preconditioned value vanishes, and injectivity
of $A$ gives $F(b)=0$.

Every competing root $z\in a+p^{s+1}R^d$ can be written
$z=a+p^{s+1}w$. Its root equation implies the same fixed-point
equation after cancellation of the nonzero scalar $p^{2s+1}$.
Contraction uniqueness gives $w=u$, proving uniqueness on the entire
stated ball.
\end{proof}

The existence and uniqueness statements are
\lean{Komorebi.PolynomialWeightedLift.exists_weighted_root} and
\lean{Komorebi.PolynomialWeightedLift.exists_unique_weighted_root}.
The inverse hypothesis in Lean specifies an integral matrix whose image
is exactly $p^sJ^{-1}$; it does not substitute a potentially larger
determinant valuation for $s$.

\begin{example}[Strict threshold and displacement boundary]
\label{ex:weighted-boundary}
For $p=2$, $F(X)=X^2+3$, and $a=1$, one has $F(a)=4$ and $DF(a)=2$.
Thus $s=1$ and $r=2$ satisfy every displayed hypothesis except the strict
inequality. No root exists in $\Z_2$, since none exists modulo eight.
Hence $r\ge2s$ cannot replace $r>2s$ uniformly. Also,
$F(X)=p^sX+p^r$ has root $-p^{r-s}$, whose displacement is not divisible
by $p^{r-s+1}$. The guaranteed displacement cannot uniformly gain a digit.
\end{example}

\subsection{First-order coefficient and parameter motion}

\begin{theorem}[Weighted coefficient deformation]\label{thm:coefficient-motion}
Let $F,H\in R[X_1,\ldots,X_d]^d$, $F(a)=0$, and $J=DF(a)$. Suppose
$J\in\operatorname{GL}_d(\Qp)$, $p^sJ^{-1}\in M_d(R)$, and $M>2s$.
For $G=F+p^MH$, there is exactly one root $b\in a+p^{s+1}R^d$.
It satisfies $b-a\in p^{M-s}R^d$ and, in $\Qp^d$,
\begin{equation}\label{eq:coefficient-first-order}
 b-a=-p^MJ^{-1}H(a)+\varepsilon,
 \qquad\varepsilon\in p^{2M-3s}R^d.
\end{equation}
Only the original Jacobian $DF(a)$ is subject to an inverse hypothesis.
\end{theorem}

\begin{proof}
Retain the original certificate $A=p^sJ^{-1}$ and put $c=p^{s+1}$.
Taylor expansion and a polynomial difference identity yield polynomial
vectors $Q_F,Q_H$, both zero at the origin, such that
\[
 (AF)(a+cu)=p^{2s+1}u+p^{2s+2}Q_F(u),\qquad
 (AH)(a+cu)=AH(a)+cQ_H(u).
\]
Consequently
\[
 (AG)(a+cu)=p^{2s+1}\bigl(C+u+pQ(u)\bigr),
\]
where
\[
 C=p^{M-(2s+1)}AH(a),\qquad
 Q=Q_F+p^{M-s-1}Q_H,\qquad Q(0)=0.
\]
All exponents are nonnegative. The same contraction and divisibility
argument as in Theorem~\ref{thm:weighted-lifting} gives the root,
the displacement $p^{M-s}$, and full-ball uniqueness. This constructs the
root without introducing an inverse hypothesis for $DG(a)$.

For the error estimate, let $\delta=b-a$ and $D=M-s$. Taylor's theorem
for the preconditioned polynomial gives
\[
 E_F=(AF)(b)-(AF)(a)-p^s\delta\in p^{2D}R^d.
\]
Polynomial congruence gives
$E_H=(AH)(b)-(AH)(a)\in p^DR^d$.
The equations $F(a)=0$ and $G(b)=0$ imply
\[
 p^s\delta+p^MAH(a)=-E_F-p^ME_H\in p^{2D}R^d,
\]
because $M+D\ge2D$. Cancel $p^s$ in the integral domain $R$ to obtain
\[
 \delta+p^{M-s}AH(a)\in p^{2M-3s}R^d.
\]
This integral vector is $\varepsilon$. After inclusion into $\Qp$,
$p^{M-s}A=p^MJ^{-1}$ gives \eqref{eq:coefficient-first-order}.
\end{proof}

The complete theorem is
\lean{Komorebi.PolynomialRootMotion.exists_unique_root_motion}.
The error depth exceeds the displacement depth by $M-2s\ge1$. It is a
guaranteed divisibility, not an assertion of exact valuation for every
system.

\begin{theorem}[Joint coefficient and parameter deformation]
\label{thm:joint-deformation}\leavevmode\par
Let $F,H\in R[Y_1,\ldots,Y_e,T_1,\ldots,T_k]^e$ and
$a=(y_0,t_0)\in R^{e+k}$ satisfy $F(a)=0$. Define the actual blocks
\[
 J=D_YF(a)\in M_e(R),\qquad K=D_TF(a)\in M_{e\times k}(R).
\]
Suppose $J\in\operatorname{GL}_e(\Qp)$,
$p^sJ^{-1}\in M_e(R)$, and $M>2s$. For every $v\in R^k$, set
\[
 \tau=t_0+p^Mv,\qquad G=F+p^MH.
\]
There is exactly one $y\in y_0+p^{s+1}R^e$ with $G(y,\tau)=0$.
It satisfies
\begin{equation}\label{eq:joint-first-order}
 \begin{aligned}
  y-y_0&\in p^{M-s}R^e,\\
  y-y_0&=-p^MJ^{-1}\bigl(H(a)+Kv\bigr)+\varepsilon,
  \qquad\varepsilon\in p^{2M-3s}R^e.
 \end{aligned}
\end{equation}
The equality involving $J^{-1}$ is in $\Qp^e$, with an integral error
vector. The parameter equality is exact.
\end{theorem}

\begin{proof}
Use the augmented polynomial systems
\[
 \mathcal F(Y,T)=(F(Y,T),T-t_0),\qquad
 \mathcal H(Y,T)=(H(Y,T),-v).
\]
Then $\mathcal F(a)=0$, and its derivative is
\[
 B=D\mathcal F(a)=\begin{pmatrix}J&K\\0&I_k\end{pmatrix}.
\]
For $A=p^sJ^{-1}\in M_e(R)$ and $c=p^s$, define
\[
 C=\begin{pmatrix}A&-AK\\0&cI_k\end{pmatrix}.
\]
The identities $AJ=JA=cI_e$ give $CB=BC=cI_{e+k}$. In the second
product, the upper right block is $-JAK+cK=0$; in the first, it is
$AK-AK=0$. Thus augmentation preserves the original inverse bound $s$.

Theorem~\ref{thm:coefficient-motion} gives a unique root $b$ of
$\mathcal F+p^M\mathcal H$ in $a+p^{s+1}R^{e+k}$, with displacement
depth $M-s$. Its last equations say exactly
$b_T=t_0+p^Mv=\tau$. Its first equations are $G(b)=0$.
The integral first-order expression uses
\[
 C\binom{H(a)}{-v}
 =\binom{A(H(a)+Kv)}{-p^sv}.
\]
Projection onto the solved coordinates gives
\[
 b_Y-y_0+p^{M-s}A(H(a)+Kv)\in p^{2M-3s}R^e.
\]
The same inclusion identity as before converts this to
\eqref{eq:joint-first-order}.

Finally, if $z\in y_0+p^{s+1}R^e$ is any competing solved root with
parameters $\tau$, then $(z,\tau)$ belongs to the full augmented ball.
Indeed, $M>2s$ implies $M\ge s+1$. Full-point uniqueness therefore gives
$(z,\tau)=b$. This comparison proves uniqueness on exactly the solved
ball stated in the theorem.
\end{proof}

The exact solved-coordinate statement is
\lean{Komorebi.PolynomialJointDeformation.exists_joint_deformation},
with its unique-existence interface
\lean{Komorebi.PolynomialJointDeformation.exists_unique_joint_deformation}.
The triangular inverse identities hold over arbitrary commutative rings,
so they apply in particular over the $p$-adic integers.

\begin{example}[First-order sign and error precision]\label{ex:motion-boundary}
For $p=2$, $s=0$, $M=2$, $F(Y,T)=Y+T$, $H=1$, $a=0$, and $v=1$,
the prescribed parameter is $\tau=4$ and the solved root is $y=-8$.
Here $J=K=1$, and the forcing term is $H(a)+Kv=2$. Omitting the parameter
term would predict $-4$, and reversing its sign would predict $0$.
The corresponding errors $-4$ and $-8$ are both outside
$2^{2M}\Z_2=16\Z_2$. Thus either change would violate the stated error
bound.

The error exponent cannot uniformly be increased. For
$F(Y,T)=Y+Y^2+T$, $H=0$, $a=0$, $p=2$, $s=0$, $M=1$, and $v=1$,
there is an even root $y$ of $y+y^2+2=0$. Its equation forces
$y\equiv2\pmod4$ and $y+2=-y^2$, so
\[
 4\mid y+2,\qquad 8\nmid y+2.
\]
Thus the error depth $2M-3s=2$ is attained even for parameter-only motion.
The same quadratic with a constant coefficient perturbation gives the
coefficient-only boundary. These examples refute a uniform extra error
digit; they do not assert optimality for every individual choice of
$p,s,M$.
\end{example}

The combined results establish finite singular obstructions, smooth
compatible towers, full continuous root charts after a suitable coordinate
choice, and quantitative deformation for a rationally invertible solved
Jacobian. Rationally rank-deficient persistence remains a separate problem.
The analytic solution-map theorem in the next section supplies the
regularity conclusion for these same derivative blocks.

\section{Analytic regularity of the local solution map}
\label{sec:regularity}
The deformation estimates control how far an exact root moves under a small
perturbation. We now show that its local dependence on the parameters is
analytic, and compute the derivative from the polynomial Jacobian. The
analytic argument takes place over a complete normed field. For integral
$p$-adic data, we first include the coefficients and base point in $\Qp$;
$\Zp$ itself is not a vector space over $\Qp$.

Let $\mathbb K$ be a complete nontrivially normed field. Let $e,k\in\N$, and
write the variables as $(Y,T)\in\mathbb K^e\times\mathbb K^k$. For
\[
 F\in\mathbb K[Y_1,\ldots,Y_e,T_1,\ldots,T_k]^e,
 \qquad a=(y_0,t_0),
\]
define the matrices of evaluated formal partial derivatives
\begin{equation}
 J=D_YF(a)\in M_e(\mathbb K),\qquad
 K=D_TF(a)\in M_{e\times k}(\mathbb K).
 \label{eq:regularity-blocks}
\end{equation}
Rows index equations and columns index variables. We first establish that
these formal derivatives are the derivatives of the evaluation map.

\begin{lemma}[Derivative of polynomial evaluation]
\label{lem:actual-polynomial-derivative}
For $P\in\mathbb K[X_1,\ldots,X_n]$, polynomial evaluation is analytic at
every point. At $a\in\mathbb K^n$ its strict Fr\'echet derivative is
\[
 v\longmapsto\sum_{j=1}^n (\partial_jP)(a)v_j.
\]
For a finite system, its strict derivative is the continuous linear map
represented by the matrix of evaluated formal partial derivatives.
\end{lemma}
\begin{proof}
For the derivative statement, use polynomial induction on constants,
addition, and multiplication by a variable. Constants have zero derivative.
The sum rule agrees with the additivity of formal partial differentiation.
For $PX_j$, the product rule gives
\[
 d(PX_j)_a(v)=a_j\,dP_a(v)+P(a)v_j,
\]
which agrees with the evaluated formal product rule for every coordinate.
The finite product rule for maps to coordinate spaces then assembles the
system derivative. Each matrix defines a continuous linear map because the
coordinate spaces are finite-dimensional over the complete field.
Polynomial analyticity follows from the corresponding polynomial induction
in Mathlib and its finite-coordinate assembly.
\end{proof}
The derivative bridge and analytic statements are
\lean{Komorebi.PolynomialAnalyticEvaluation.hasStrictFDerivAt_evalSystem},
\lean{Komorebi.PolynomialAnalyticEvaluation.analyticAt_evalSystem}, and
\lean{Komorebi.PolynomialAnalyticEvaluation.analyticOnNhd_evalSystem}.

\begin{lemma}[The augmented derivative]
\label{lem:augmented-analytic-inverse}
If $J\in\operatorname{GL}_e(\mathbb K)$, the polynomial map
\[
 \Phi(Y,T)=(F(Y,T),T)
\]
has invertible strict derivative at $a$, with
\begin{equation}
 D\Phi(a)=\begin{pmatrix}J&K\\0&I_k\end{pmatrix},\qquad
 D\Phi(a)^{-1}=\begin{pmatrix}J^{-1}&-J^{-1}K\\0&I_k\end{pmatrix}.
 \label{eq:analytic-block-inverse}
\end{equation}
The local inverse supplied by this derivative is analytic near $\Phi(a)$.
\end{lemma}
\begin{proof}
The lower coordinates of $\Phi$ are the parameter variables. Differentiating
them gives the blocks $0$ and $I_k$. Lemma~\ref{lem:actual-polynomial-derivative}
identifies the remaining blocks with~\eqref{eq:regularity-blocks}.
Multiplication of the two matrices in~\eqref{eq:analytic-block-inverse},
in both orders, gives the identity. The resulting linear equivalence is
continuous, as is its inverse. Apply the strict inverse-function theorem to
$\Phi$. The analytic inverse theorem applies to the same open partial
homeomorphism, because $\Phi$ is analytic and its derivative is that
continuous linear equivalence. It supplies a convergent local power-series
representation of this inverse.
\end{proof}
The augmented Jacobian is identified in
\lean{Komorebi.PolynomialJointSystem.jacobian_augmented}.
The continuous inverse block and its restriction are
\lean{Komorebi.PolynomialAnalyticBlocks.upperEquiv} and
\lean{Komorebi.PolynomialAnalyticBlocks.inverse_solved_zero}.
Analyticity of the same constructed inverse is
\lean{Komorebi.PolynomialAnalyticInverse.analyticAt_inverse}.

\begin{theorem}[Analytic local root graph]
\label{thm:analytic-local-graph}
Suppose $F(a)=0$ and $J\in\operatorname{GL}_e(\mathbb K)$. There are open
neighborhoods $t_0\in U\subseteq\mathbb K^k$ and
$y_0\in V\subseteq\mathbb K^e$, and a map
$\psi:\mathbb K^k\to\mathbb K^e$, such that
\begin{gather}
 \psi(t_0)=y_0,\qquad \psi(U)\subseteq V,\qquad
 \psi\text{ is analytic at every point of }U,\label{eq:analytic-psi}\\
 F(y,t)=0\quad\Longleftrightarrow\quad y=\psi(t)
 \qquad(t\in U,\ y\in V).\label{eq:analytic-root-graph}
\end{gather}
At the base parameter, the strict Fr\'echet derivative is
\begin{equation}
 \boxed{D\psi(t_0)=-J^{-1}K.}
 \label{eq:implicit-derivative}
\end{equation}
In particular, $\psi$ is differentiable on $U$.
\end{theorem}
\begin{proof}
Let $g$ be the analytic local inverse of $\Phi$ supplied by
Lemma~\ref{lem:augmented-analytic-inverse}. Since $F(a)=0$,
$\Phi(a)=(0,t_0)$ and $g(0,t_0)=a$. Define
\[
 \psi(t)=\pi_Yg(0,t),
\]
where $\pi_Y$ projects to the solved coordinates. This gives the base value
and analyticity at $t_0$. The chain rule and
\eqref{eq:analytic-block-inverse} give
\[
 D\psi(t_0)v
 =\pi_Y\begin{pmatrix}J^{-1}&-J^{-1}K\\0&I_k\end{pmatrix}
       \binom{0}{v}
 =-J^{-1}Kv.
\]

To obtain the root equivalence on a product neighborhood, we choose the
neighborhoods as follows. The source of the open partial homeomorphism defined by
$\Phi$ is open around $a$; its inverse is $g$. Choose open sets $V$ and $U_0$
containing $y_0$ and $t_0$ such that
$V\times U_0$ lies in that source. Restrict $U_0$ further by the preimage of
an open analytic inverse neighborhood under $t\mapsto(0,t)$, the open locus
where this same $\psi$ is analytic, and the interior of $\psi^{-1}(V)$.
All of these sets contain $t_0$. For the last set, use the continuity of
$\psi$ at $t_0$ and $\psi(t_0)=y_0\in V$. Their intersection is an open
neighborhood $U$ with $\psi(U)\subseteq V$.

If $t\in U$, $y\in V$ and $F(y,t)=0$, then
$\Phi(y,t)=(0,t)$. The left inverse identity forces
$g(0,t)=(y,t)$, hence $y=\psi(t)$. Conversely, the right inverse identity
$\Phi(g(0,t))=(0,t)$ says both that its solved equations vanish and that its
parameter coordinates are exactly $t$. Thus
$g(0,t)=(\psi(t),t)$ and $F(\psi(t),t)=0$.
Analyticity at every point of $U$ implies differentiability there.
\end{proof}
This is
\lean{Komorebi.PolynomialAnalyticImplicit.exists_analytic_implicit}.
Differentiability is also included explicitly in the statement contract
\lean{Komorebi.PolynomialRegularityContracts.actual_padic_analytic_root_graph}.

\begin{theorem}[Regularity of a supplied root selection]
\label{thm:selection-regularity}
Under the hypotheses of Theorem~\ref{thm:analytic-local-graph}, let
$h:\mathbb K^k\to\mathbb K^e$ be continuous at $t_0$, with
$h(t_0)=y_0$. Suppose $F(h(t),t)=0$ for every $t$ in some neighborhood of
$t_0$. Then $h$ and $\psi$ agree in a neighborhood of $t_0$. Consequently
$h$ is analytic there as a germ and has strict derivative
$Dh(t_0)=-J^{-1}K$.
\end{theorem}
\begin{proof}
Continuity at $t_0$ places $h(t)$ inside $V$ for all sufficiently close
parameters. Intersect this neighborhood with $U$ and the neighborhood on
which the root equation holds. Equation~\eqref{eq:analytic-root-graph} gives
$h(t)=\psi(t)$ throughout the intersection. Analyticity and strict
differentiability at a point are unchanged by agreement on a neighborhood.
\end{proof}
The formal result is
\lean{Komorebi.PolynomialAnalyticSelection.continuous_selection_regularity}.
Its formulation derives $F(a)=0$ from the root equation near the base
parameter and the specified base value.

\begin{corollary}[Regularity throughout an open parameter set]
\label{cor:selection-open-regularity}
Let $U\subseteq\mathbb K^k$ be open, and let $h$ be continuous on $U$.
Suppose that for every $t\in U$,
\[
 F(h(t),t)=0,
 \qquad J_t=D_YF(h(t),t)\in\operatorname{GL}_e(\mathbb K).
\]
Then $h$ is analytic at every point of $U$ and its strict derivative is
\[
 Dh(t)=-J_t^{-1}D_TF(h(t),t)\qquad(t\in U).
\]
\end{corollary}
\begin{proof}
At each $t\in U$, openness turns continuity on $U$ into continuity at $t$
and makes the root equation hold on a neighborhood of $t$. Apply
Theorem~\ref{thm:selection-regularity} at $(h(t),t)$.
\end{proof}
This is
\lean{Komorebi.PolynomialAnalyticSelection.continuous_selection_analyticOnNhd}.
The pointwise invertibility assumption is part of this corollary.

\begin{corollary}[Integral coefficients and centers]
\label{cor:integral-analytic-graph}
Let $p$ be any prime. Suppose $F\in\Zp[Y,T]^e$, $a\in\Zp^{e+k}$,
and $F(a)=0$. Let $\iota:\Zp\hookrightarrow\Qp$ be the inclusion.
If $\iota(D_YF(a))\in\operatorname{GL}_e(\Qp)$, the polynomial system
obtained by including the original coefficients in $\Qp$ has the analytic
local root graph of Theorem~\ref{thm:analytic-local-graph} at $\iota(a)$.
The strict derivative at the base is
\[
 -\iota(D_YF(a))^{-1}\iota(D_TF(a)).
\]
The open neighborhoods and the output map are in the ambient spaces over
$\Qp$.
\end{corollary}
\begin{proof}
Coefficient inclusion commutes with evaluation and formal partial
differentiation. Consequently both Jacobian blocks after inclusion are the
entrywise inclusions of the original blocks, and the original exact root
remains an exact root. Apply Theorem~\ref{thm:analytic-local-graph} to the system with included
coefficients.
\end{proof}
This is
\lean{Komorebi.PolynomialRegularityIntegral.exists_analytic_integral_implicit}.
It requires no invertibility modulo $p$. In particular, it includes solved
Jacobians singular modulo $p$ but invertible over the fraction field.

\begin{example}[The derivative sign and a nonlinear control]
For $F(Y,T)=Y+T$, the exact branch is $\psi(T)=-T$ and $J=K=1$.
The derivative is $-1$. Uniqueness of the derivative rules out $+1$ over
$\Qp$, including $p=2$. The nonlinear polynomial $F(Y,T)=Y+Y^2+T$
has $F(0,0)=0$ and $J=K=1$ there. Applying the formal theorem to this polynomial constructs an analytic
local root graph with derivative $-1$. Thus the example checks the existence
conclusion as well as the derivative formula.
The declarations and the deliberately rejected positive-sign statement are
in the regularity control and statement-contract files.
\end{example}

The results cover arbitrary finite $e$ and $k$, including zero. The local
inverse machinery explicitly handles a zero-dimensional domain. The analytic theorem does not assign a uniform numerical radius to its
neighborhoods. The earlier
weighted lifting and joint deformation theorems retain their quantitative
radii and error bounds, while the present statements establish regularity
of the local solution germ. A direct Lean identification with each
historical subtype-valued chart is an additional interface task. General
persistence at a fraction-field singular solved block is a separate problem.

\section{Open questions and further directions}
\label{sec:open}

The arithmetic results above resolve Nishibuchi's conjecture throughout its
stated parameter range. The local analyticity and differentiability theorem
is also complete under its stated invertibility hypothesis. What remains is
a set of extensions and sharper classifications, together with the broader
questions raised by the original native seed-calculus program.

\subsection{Singular persistence beyond an invertible solved block}
\label{subsec:singular-open}

The weighted lifting theorem allows the solved Jacobian to be noninvertible
modulo $p$ provided that it is invertible over $\Qp$ with a controlled inverse.
This includes a substantial class of lifts often called singular modulo $p$.
The next problem is to understand Jacobians that remain rank-deficient over
$\Qp$, where the inverse construction is no longer available.

For a fixed polynomial system $F\in\Z[X_1,\ldots,X_d]^e$, a precision $L\geq1$,
and a root $\bar x$ modulo $p^L$, let
\[
\mathcal L_{F,L}(\bar x)
=\{x'\in(\Z/p^{L+1}\Z)^d:
F(x')=0,\ x'\equiv\bar x\pmod{p^L}\}.
\]
The finite obstruction theorem decides whether this set is empty from the
residual digit and reduced Jacobian. When it is nonempty, the theorem also
determines its cardinality and affine kernel structure. The conclusion
concerns one level: vanishing of the obstruction permits the next step,
but does not guarantee that subsequent steps are possible.

For example, $X^2-5$ has the root $1$ modulo $2$, and both $1$ and $3$ lift it
to roots modulo $4$. Every odd square is $1$ modulo $8$, so neither branch
continues to a root modulo $8$. Thus even a successful first correction can
lead to branches that stop at the following level.

We seek structural hypotheses and higher obstruction data that control an
entire singular branch, including its behavior under coefficient and parameter
deformation. One route is to propagate suitable conditions from level to
level. Another is to prove existence at all levels in a form that yields
compatible branches. The finite obstruction calculus supplies the tests for
individual levels, but a persistence theorem needs a further argument linking
them. Any general formulation must allow multiple branches, later
obstructions, and changes in the rank of the Jacobian.

\subsection{Native 13.2: presentation invariance}
\label{subsec:native132}

The original conjecture asks for a canonical finite obstruction packet whose
vanishing is necessary and sufficient for a one-level lift in a finitely
presented arithmetic seed system. For a fixed finite polynomial presentation,
the chapter supplies the precise obstruction
\[
o_{F,\bar x}=[r]\in
\coker\!\left(J_{F,\bar x}:\F_p^d\longrightarrow\F_p^e\right),
\]
where $r$ is the residual digit. We have proved independence from the chosen
integer representative. Invariance under a change of equations or presentation
requires more.

The raw cokernel already reveals the difficulty.
The presentations $F=(X)$ and $G=(X,0)$ describe the same zero locus over every
coefficient ring. At zero their Jacobians are, respectively,
\[
\F_p\xrightarrow{\,1\,}\F_p,
\qquad
\F_p\xrightarrow{\,(1,0)\,}\F_p^2.
\]
Their raw cokernels are $0$ and $\F_p$. Both distinguished obstruction classes
vanish, but the ambient obstruction groups are not isomorphic. The example
rules out the raw cokernel as a presentation-independent object while leaving
open the possibility of an intrinsic obstruction construction.

A precise form of Native 13.2 therefore needs four ingredients:
\begin{enumerate}
\item a category of seed objects, including how the operators, relations,
projectors, weights, and topology enter an object;
\item specified changes of presentation and a proof that they preserve the
actual lift problem;
\item an obstruction object that respects dependencies among the relations,
together with comparison maps or a universal property;
\item a precise meaning of ``finite'' and ``canonical'', and a proof that
vanishing is equivalent to the intended one-step continuation.
\end{enumerate}
For commutative algebra presentations, the relation module offers a natural
way to account for dependencies among listed equations. To extend this
approach to the full seed vocabulary, we would need an encoding that preserves
the lift problem, together with a proof of that preservation. This is one of
the links still needed between the concrete obstruction theorem and Native 13.2.

\subsection{Native 13.3: finite compatibility, completion, and descent}
\label{subsec:native133}

The original proposed form is
\[
X\simeq X_2^{\mathrm{wt}}\times_C\prod_{p>2}'X_p.
\]
It leaves $X$, the distinguished subsets defining the restricted product,
the maps to $C$, the ambient category, and the universal meaning of
``smallest'' unspecified. In particular, one must decide whether $X$ is a
global object or a completed object.

The linear compatibility theorem provides one concrete part of this program.
If two local observations
$J:U\to V$ and $K:U\to W$ are linear maps from the \emph{same} correction
module, define
\[
C_{J,K}=U/(\ker J+\ker K).
\]
The attainable observations $\im J$ and $\im K$ have canonical labels in this
quotient. Equality of labels is equivalent to simultaneous realization by a
single correction. The joint image is exactly the fiber product over
$C_{J,K}$, with the proved linear universal property.
The declarations are
\begingroup\small\raggedright
\lean{Komorebi.SharedCorrection.labels_eq_iff_common},
\lean{Komorebi.SharedCorrection.joint_range_eq_fiber}, and
\lean{Komorebi.SharedCorrection.common_universal}.
\par\endgroup
To apply this result arithmetically, one must construct the common correction
module and the observation maps, then show that the resulting compatibility
condition agrees with the intended arithmetic one.

\paragraph{Why the global/completed distinction is necessary.}
Consider constant sequences with integer initial value. Their global parameter
space is $\Z$, while their local parameter spaces are $\Z_p$. Under the
canonical local maps, the image of the global objects is the diagonal
\[
D=\{(n)_p:n\in\Z\}\subseteq\prod_p\Z_p.
\]
This diagonal is dense: a basic open set specifies finitely many congruences,
and the Chinese remainder theorem provides an integer satisfying them. It is
proper: the tuple with $2$-adic coordinate zero and all odd-adic coordinates
one is not the image of an integer. Indeed an integer zero in $\Z_2$ must be
zero, contradicting its $3$-adic coordinate.

Consequently this canonical global diagonal cannot be an equality locus
\[
\{(a,b):f(a)=g(b)\}
\subseteq \Z_2\times\prod_{p>2}\Z_p
\]
for continuous maps to a Hausdorff space $C$. Such an equality locus is closed,
as the inverse image of the closed diagonal in $C\times C$, whereas $D$ is
proper and dense. This topological argument, retained from the native progress
note, concerns the stated topological interpretation. It has not been
formalized in Lean for this manuscript.

Finite CRT reconstruction leads to the familiar completed object
$\widehat\Z=\prod_p\Z_p$, but does not identify which local tuples come from a
single global integer. A complete form of Native 13.3 must therefore specify
finite-level compatibility, transition maps and completion, and any additional
global descent condition. If recurrence, spectral, or Frobenius information
is to supply that condition, it must be expressed through defined maps or
constraints on the local data, with a precise account of how they enter $C$.

The accompanying \nolinkurl{native/Native_Conjectures_Progress.md} records
further results in restricted settings: an intrinsic finite
$\operatorname{Ext}^1$ obstruction for finitely presented commutative integer
algebras, finite CRT reconstruction for polynomial solution sets, and a
characterization of rectangular finite relations by equality of common
labels. The obstruction result is a stated specialization of a classical
external theorem; the CRT and rectangularity results have prose proofs in
the note. These are preserved mathematical progress, with the restrictions
and source records given there, rather than Lean completions of the general
Native 13.2 and 13.3 statements.

\subsection{Arithmetic refinements}
\label{subsec:arithmetic-open}

Write $U_{N,j}(n)=j!C_{N,j}(n)$ and $U_j=U_{4,j}$, and use
$v_p(0)=\infty$ when discussing valuation bounds. The following classification
questions remain open in this program.

\begin{enumerate}
\item \textbf{The remaining higher-difference classification.}
The new strengthening chapters determine every repeated-step valuation in addresses $j=0,1$, give the exact criterion for equality in the common mixed-step lower bound there, and determine address-two valuations at all power-of-two orders. What remains is a complete exact formula in address two at the other orders, a full classification in address three, and the amount of additional cancellation in mixed-step cases where the leading term vanishes. Numerical patterns through a finite range do not settle these questions.

\item \textbf{Upper-half addresses at larger dyadic rank.}
The exact first-difference laws now include $N=2,4,8$ and every lower-half address for $N=2^a\ge4$. Larger upper-half addresses need a finer classification. Explicit counterexamples show that extrapolating one low-order Newton coefficient to a universal metric law is false. A least period is measured in packet rows; its Euler-degree displacement is $NP$.

\item \textbf{Exact affine Mahler-kernel valuations.}
For the kernel of Section~\ref{subsec:polyphase-exposition}, determine exact
values of
\[
v_p\bigl(\Theta_{M,s}(r,k)\bigr),\qquad
\Theta_{M,s}(r,k)=[z^r](1+z)^s((1+z)^M-1)^k,
\]
beyond the explicit special families. The minimum valuation of the summands
gives a lower bound. Equality requires a separate argument excluding
cancellation.

\item \textbf{Classical tail root classification and simplicity.}
The later tail program gives an exact factorization, finite preparation for each fixed depth, and several infinite regular-wall and simplicity families, as summarized in Appendix~\ref{app:root:scope}. A universal explicit root classification and universal simplicity of the remaining analytic factor are still open. In the exceptional address-one quadratic reduction, the residual task is a nonvanishing discriminant condition. Empty tails remain exactly zero. This problem is distinct from the solved Nishibuchi congruence conjecture.
\end{enumerate}

These refinements build on the proved estimates and exact cases, including
the sharp uniform classical dyadic bound and the resolution of Nishibuchi's
conjecture.

\subsection{Historical material and interface work}
\label{subsec:historical-scope}

The original Paper IX, with corrected proof-status labels, survives as a
28-page PDF. Its recovered formalization supports the framework, Euler
coordinates, Mahler differences,
unit coordinates, recurrence packets, and CRT results retained in this
manuscript. The exact original TeX bytes, original figure-generation files, and complete
historical experiment package were not recovered with that PDF.
The reconstructed Annals text is reproduced in Appendix~\ref{app:recovered},
with its separate provenance and local qualifications. Their older
figure and test counts remain historical records; those computations have
not been reproduced for this consolidation.

A reconstructed textual copy of the Annals manuscript was also recovered.
It includes finite-field mode and subgroup inversion, Fourier invariance,
Carmichael and matrix transfer, Amice pullbacks, depleted analytic Euler
branches, and extensive polyphase formulas. The sharp classical dyadic kernel
and truncation claims from that text have since been connected to checked
Lean proofs. The finite polyphase formulas used for explanation here have
short explicit algebraic derivations. The remaining claims in that larger
historical text have not all been formalized.
The source archive retains the text with its provenance, and the statement
map distinguishes these historical claims from the certified claims in the
present manuscript.

Remaining recovery work includes the precise older odd-prime edge identities,
their stated ranges, and the binary digit-sum form of the checked weights.
For each item, the work is to recover the formula with its stated range,
check the proof, and identify any corresponding Lean declaration. Where the
historical source is still missing, the record remains incomplete.

One further task is to connect the older chart constructions to the new
regularity theorem. That theorem applies to polynomials with arbitrary
coefficients in the ambient field $\Qp$ and identifies the analytic
germ of every continuous ambient-field root selection under its hypotheses.
Several older rectangular charts were defined using integer coefficients,
integer centers, or subtype-valued maps. Comparing each such historical
interface explicitly with the new ambient formulation remains to be done.
These comparisons would connect the existing constructions to the proved
analytic implicit-function theorem. General singular persistence remains the
separate problem described in Section~\ref{subsec:singular-open}.

\section{Formal evidence and reproduction}
\label{sec:verification}

The verification record supports the manuscript in several ways. The Lean source states the quantified theorem and its hypotheses. A
compiler report records acceptance of that source by a specified Lean and
Mathlib environment. An axiom audit identifies the axioms used by the resulting
named declarations. Statement contracts and mathematical controls check
selected interfaces and boundary cases. Finally, source correspondence reviews
compare the formal objects with the intended mathematical statement. Each check has a specific role: numerical examples test particular cases,
whereas a universal claim requires a proof, and correspondence with an
external conjecture requires a comparison of the statements.

This edition revises the exposition and checks the evidence index against
the preserved sources. The Lean sources are unchanged. Their compiler
acceptance is supported by the earlier successful reports, whose source and
log hashes are checked by the accompanying index. The most recent compilation
covers local analytic regularity and its polynomial dependencies; earlier
reports cover the arithmetic and native obstruction developments.

The proofreading pass also restores several results omitted from the first
consolidation and records the comparison with the original paper. Its
coverage matrix and semantic review appear in
\texttt{consolidation/editorial/}. The Bernoulli boundary is accompanied by
a classical derivation, without a claim of a dedicated Lean certificate.
The reviews distinguish this addition from restored results already present
in the checked sources.

\subsection{Verification frontiers}

Table~\ref{tab:verification-frontiers} gives the recorded scope of each major
verification frontier. An \emph{inventoried module} is a source file in that
frontier's formalization list. Older lists sometimes include files of examples
or positive controls in this group. A \emph{check} is an entry in the report's
run list, including expected rejection of invalid claims. An \emph{axiom
closure} is the recorded axiom dependency set of an explicitly named
declaration; definitions and abbreviations contribute to this count as well
as theorems. None of these columns is a count of distinct mathematical results
across the chapter.

\begin{table}[htbp]
\centering
\small
\begin{tabular}{@{}lrrr@{}}
\toprule
Preserved frontier & Modules & Checks & Axiom closures\\
\midrule
Original completion baseline & 79 & $79+37+3$ & 1000\\
Dyadic strengthening and isometry & 22 & 38 & 246\\
Standalone finite affine obstruction & 2 & 4 & 17\\
Shared correction quotient & 6 & 13 & 60\\
Polynomial one-step lifting & 9 & 18 & 91\\
Exact polynomial lift counts & 9 & 17 & 106\\
Smooth infinite towers & 14 & 22 & 127\\
Coefficient deformation & 17 & 26 & 159\\
Continuous parameter charts & 25 & 40 & 244\\
Weighted lifting & 30 & 50 & 304\\
First-order root motion & 35 & 58 & 326\\
Prescribed parameter obstruction & 39 & 65 & 375\\
Prescribed parameter towers & 45 & 74 & 408\\
Joint coefficient and parameter motion & 49 & 82 & 460\\
Local analytic regularity & 55 & 91 & 515\\
\bottomrule
\end{tabular}
\caption{Recorded verification scopes. Rows overlap extensively through
inherited sources and must not be added. The baseline separates 79 module
checks, 37 control checks, and three additional recorded runs; its nine
dependency revision comparisons are recorded separately.}
\label{tab:verification-frontiers}
\end{table}

The baseline inventory contains 836 theorem commands, 162 definitions and two
abbreviations, matching its 1000 audited names. The latest polynomial inventory
contains 322 theorem commands, 76 definitions and seven abbreviations in its
55 formalization modules. Its eight positive-control support modules and
13 statement-contract modules bring the full audit to 415 theorem commands,
92 definitions and eight abbreviations, hence 515 closures. The 91 checks also
include 12 negative controls. Six new regularity modules contribute 30 theorem
commands and eight definitions; all 49 inherited formalization modules are
unchanged in that release.

The arithmetic strengthening inventory exposes a small historical metadata
error. Its summary fields say 179 theorems and 12 definitions, while its
module entries list 215 theorem commands, ten lemma commands,
20 definitions and one abbreviation. These 246 named entries exactly match
the preserved axiom audit. The consolidation uses the entries and audit, and
records the stale summary fields explicitly. The original inventory remains
unchanged as historical evidence. This is a correction to the reporting
metadata, not a newly discovered proof or a modification of a checked theorem.

\subsection{Claim-to-source correspondence}

The source references attached to the mathematical sections identify the
principal declarations for each result. Their complete names, exact source
hashes, roles and verification frontiers are collected in
\texttt{consolidation/audit/DECLARATIONS.csv}. The accompanying
\nolinkurl{consolidation/audit/EVIDENCE_INDEX.json} records the corresponding inventories,
reports and retained logs. In particular, the following interfaces distinguish
the principal completed results from generic auxiliary lemmas:
\begin{itemize}
\item The external conjecture is assembled by
\lean{NishibuchiCompletion.conjecture_3_9} in the original completion
baseline. Its coefficient identification, integrality, parameter split and
explicit threshold are part of that proof chain.
\item Exact $N=4$ packet periods are supplied by
\lean{DyadicFourLeastPeriods.least_packet_period}; the isometric extension
is supplied by \lean{DyadicFourIsometry.existsUnique_isometry}. The
normalization of the original coefficients is an explicit part of these
interfaces.
\item The classical sharp edge tower is supplied by
\lean{DyadicSharpTower.dyadic_supercongruence_sharp}, together with its
first omitted error valuation. Its optimality concerns the stated uniform
truncation bound.
\item Prescribed parameter lifting uses the actual residual and actual
partial Jacobians in
\lean{PolynomialPrescribedObstruction.prescribed_lift_iff_obstruction_zero}.
The quotient criterion and exact lift counts are not assumptions of the
later tower results.
\item Joint deformation is supplied by
\lean{PolynomialJointDeformation.exists_unique_joint_deformation}.
Local regularity is supplied by
\lean{PolynomialAnalyticImplicit.exists_analytic_implicit}, with the
actual polynomial derivative and the derivative formula $-J^{-1}K$.
\item \lean{PolynomialAnalyticSelection.continuous_selection_regularity}
transfers the analytic germ and strict derivative to any continuous local
root selection satisfying the stated root and base-point conditions.
\end{itemize}
Here arithmetic declaration names have namespace prefix
\texttt{KomorebiPaperIX.}; polynomial and native declarations have prefix
\texttt{Komorebi.}. The section maps and CSV retain the full names.

The declaration index covers every explicitly named theorem, lemma,
definition and abbreviation in the 14 release-frontier inventories, together
with the earlier standalone affine report. Byte-identical copies of a source
file are collapsed, while all their locations remain recorded. If the same
qualified name occurs in different source versions, the versions remain
separate and are marked. Anonymous examples, compiler-generated declarations,
failed development attempts, and the contents of earlier nested historical
snapshots are not counted as named results in this index. Negative-control
source files are checked for integrity but are not recorded as proved
mathematical assertions. Consequently, the number of CSV rows must not be
quoted as a global unique theorem count.

\subsection{External statement and library dependencies}

The original conjecture correspondence was reviewed against Nishibuchi's
arXiv version \texttt{2605.10677v2}, dated 1 June 2026. The retained source
record identifies the official version, the full title, the PDF checksum and
the extracted-text checksum. The detailed review compares Definition~2.1,
Notation~3.2 and Conjecture~3.9 with the actual reciprocal EGF coefficients,
the factor $j!$, the residue range, the degree displacement, and the threshold
in the formal capstone. Nishibuchi's conjecture is the conclusion to be
proved; it is not an axiom imported into Lean.

The full external paper is not redistributed in this source release.
Consolidation preserves the earlier source correspondence audit and its
version record. The earlier consolidation recorded a blocked request to the official arXiv endpoints.
For the preserved prior edition, the official unversioned PDF was inspected;
its internal version is v2, dated 1 June 2026. Definition 2.1 and Remark 2.3
(printed page 3), Notation 3.2 (page 5), and Conjecture 3.9 (page 8)
confirm the parameter, index, degree-displacement and normalization match.
Some direct versioned-endpoint requests were blocked in that earlier check. The later strengthening source review reached the primary PDF and HTML through the official record and confirmed the same v2 statement. These are dated source-access records, not mathematical hypotheses.
This source check is distinct from the inherited compiler acceptance.
No broad literature or priority claim is made.

The mathematical environment is Lean \texttt{4.34.0-rc2}, with Mathlib revision
\nolinkurl{aeb72f1e9981b6d35b426b16dca7759be3e7ba06}.
The allowed axiom dependencies are
\texttt{propext}, \texttt{Classical.choice} and \texttt{Quot.sound}.
The reports do not describe a rebuild of Mathlib from source. The local
regularity audit separately records the applicable Mathlib analytic inverse
and strict derivative interfaces and their hypotheses. In particular, the
field-valued analysis is performed over a complete nontrivially normed field;
its integral specialization embeds the data into $\Qp$. It does not use
$\Zp$ as a vector space over $\Qp$.

\subsection{Reproduction and limits of the consolidation check}

The archive contains the mathematical sources, verifier scripts, controls,
contracts, prior logs and exact inventories. The pinned Lean and Mathlib kit
is supplied separately. From the consolidated archive root, the latest
polynomial gate can be reproduced with a new output directory:
\begin{verbatim}
python3 replay_paper9_regularity.py \
  --kit /absolute/path/to/pinned-kit \
  --output /absolute/path/to/new-polynomial-run
\end{verbatim}
The preserved arithmetic, native and shared-correction chain has a separate
entry point:
\begin{verbatim}
python3 replay_shared_correction.py --all \
  --kit /absolute/path/to/pinned-kit \
  --output /absolute/path/to/new-arithmetic-run
\end{verbatim}
These commands reproduce the recorded checks; they were not rerun during
manuscript consolidation or proofreading. The second command rebuilds
the preserved baseline before the arithmetic continuation. Each output path
must be new. The reports retain the original full commands, runtime
fingerprints and output hashes. The executable-path compatibility adapter
used in the recorded environment restores runtime path discovery; it does
not alter Lean's mathematical kernel or the source statements.

For the narrower evidence reconciliation carried out for this edition, use
\begin{verbatim}
python3 consolidation/tools/evidence_index.py \
  --source . --output /absolute/path/to/new-evidence-index
\end{verbatim}
This check compared 644 source-file occurrences with their recorded hashes
and 713 retained log occurrences with their recorded hashes. Four additional
logs from the older standalone affine report were fingerprinted, but that
legacy report does not contain prior log hashes; those four are not promoted
to hash comparisons. Every named declaration in the selected inventories
matches the appropriate recorded standard-axiom dependency set. These are
integrity and provenance checks, not another Lean compilation.

Expected negative-control failures are retained as evidence that specific
invalid strengthenings are rejected. Examples include an unjustified extra
precision digit, omitted parameter terms, a wrong derivative sign, and
uniqueness without the appropriate hypotheses. Development failures and
interrupted runs remain historical records; only the successful designated
release reports support the verification table. Only accepted declarations in the designated successful reports support
claims of formal verification.

Finally, the proof of the original stated conjecture does not resolve every
extension mentioned in this manuscript. The open questions concerning
rationally rank-deficient singular points, presentation invariance, global
or completed descent, and further exact arithmetic have no certificate
supplied by these counts. Section~\ref{sec:open} states the remaining implications explicitly.

\subsection{Inherited checks for the prior final edition}
The preserved prior final-edition review reconstructed the arguments in its 54 written
proof blocks under their stated hypotheses. The two expanded paragraphs
in the all-address proof make the inverse-flatness commutator and the
coefficient-extraction saturation explicit. The recovered historical text
received five local repairs: integral observation vectors for matrix
congruences, positive precision for matrix amplification, integral Mahler
coefficients for the sparsification consequence, the corrected modulo-seven
example, and the even-subsequence range for source changes. Its remaining
historical and external-source evidence limits are retained in
Appendix~\ref{app:recovered}; preservation of its 82 proof bodies is not
a blanket formal certificate.

An independent Python gate using arbitrary-precision integers and
\texttt{fractions.Fraction} passed 116,492 deterministic assertions in nine
finite families: packet construction, foundational identities, divisor-branch
congruences, four-packet sharpness, classical truncations, polynomial digit
fibers, shared affine quotients, finite towers and charts, and weighted
deformation. Nineteen invalid strengthenings were rejected by concrete
controls. The exact ranges, coefficient domains, scripts and output hashes
are in the accompanying \nolinkurl{verification/computation} manifests.
These finite tests do not establish a universal theorem.

The unchanged early-row script separately reproduced 312 finite period
cases and its odd-prime row-zero boundary control. Other research-script
replays, unchanged baseline regression gates and recovered-text controls
are recorded separately in the same audit directory; their totals are
not added to the independent gate. Lean and Lake were unavailable in this
execution environment. No fresh compiler run, infinite $p$-adic computation,
analytic-continuation computation or hardware benchmark is asserted.


\subsection{Checks for the strengthened Euler edition}
The new proof routes are recorded in the threshold-free strengthening checkpoint, whose manifest pins 52 files. Its six principal proof audits are internal computational-assistant reviews and identify the exact proof revisions they checked. The manuscript delta review compares the new chapters and current scope statements with those proofs. It does not replace the earlier Lean kernel or constitute external peer review.

The dedicated arithmetic experiment computed the four reciprocal rows through index 512 modulo $2^{4096}$, with independent exact-rational inversion through index 16. It checked 525,312 first differences, 1,141,424 repeated-step factorial bounds, 139,000 mixed-step bounds, and 2,048 Newton coefficients. The supplemental experiments cover additional mixed-step and sharpness cases and retain their separate tested ranges. Their outputs and source hashes are preserved under \nolinkurl{preserved/threshold_free/computation}. All these are bounded implementation checks; the written arguments establish the universal claims.

The visual guide has its own exact computation contract and figure manifest. Its modular recurrence is compared with independent rational inversion, and each theorem-derived plot is marked as such. The fresh static Lean review reconciles the unchanged formal sources and inherited compiler logs, with counts stated in the Lean guide. No Lean or Lake executable was available for this release, and no new compiler acceptance is asserted for the new proof routes.

The release includes the editable three-document sources, exact figure generator, the strengthened proof checkpoint, unchanged earlier source packages, and manifests for the final outputs. PDF compilation, figure arithmetic, manifest validation, source reconciliation, and written proof auditing are separate checks with separate scopes.


\clearpage
\appendix
\section{Exact classical Euler tails and the remaining simplicity problem}
\label{app:root:scope}

This appendix records the related classical Euler tail program at its
stationary-separation checkpoint of 5 October 2026. Its purpose is to state
the strongest proved reduction and its exact remaining obligation. The
universal simplicity problem below is a separate question from
Nishibuchi's Conjecture~3.9, which the congruence results of this chapter
resolve. It is also separate from the presentation and gluing questions
called Native~13.2 and Native~13.3 in the original Paper~IX.

The tail results have mathematical proofs, supplemented in the indicated
families by complete coefficient certificates and proved bounds for every
omitted coefficient. Bounded sampled values support implementation checks;
they do not supply the universal quantifiers. The latest continuation did
not run Lean. Its mathematical claims must therefore be read with the
formal-coverage distinctions in the Lean guide.

\subsection{The ordinary Euler packet identity}
\label{app:root:identity}

To distinguish ordinary signed Euler numbers from the generalized numbers
of the main text, write
\[
 \operatorname{sech}t=\sum_{n\ge0}\mathcal E_n\frac{t^n}{n!},
 \qquad \frac{t}{e^t-1}=\sum_{n\ge0}B_n\frac{t^n}{n!},
 \qquad B_1=-\frac12.
\]
Throughout this appendix, $v_2(2)=1$ and $v_2(0)=\infty$.
Let $s_2(n)$ be the number of ones in the binary expansion of the
nonnegative integer $n$. Define
\[
 A(N)=\frac{3^{2N+2}-1}{4(2N+1)(2N+2)},\qquad
 \left(\frac z{\sinh z}\right)^2=\sum_{m\ge0}\beta_m z^{2m},
 \qquad \kappa_m=2^{2m}(2m)!\beta_m.
\]
The initial values are $\beta_0=\kappa_0=1$. For $m\ge1$, differentiation
of $z\coth z$ gives
\begin{equation}
 \beta_m=-\frac{2^{2m}(2m-1)B_{2m}}{(2m)!},\qquad
 \kappa_m=-2^{4m}(2m-1)B_{2m}.
 \label{eq:root:beta}
\end{equation}
The classical dyadic Bernoulli denominator law
$v_2(B_{2m})=-1$ and Legendre's formula yield
\begin{equation}
 v_2(\beta_m)=s_2(m)-1,\qquad
 v_2(\kappa_m)=4m-1.\label{eq:root:weights}
\end{equation}
The source identity
\[
 \sum_{N\ge0}A(N)\frac{t^{2N}}{(2N)!}
 =\frac{\sinh(2t)\sinh t}{2t^2}
\]
implies, after multiplication by $4t^2/\sinh^2(2t)$,
\begin{equation}
 \mathcal E_{2N}=\sum_{m=0}^N
    \kappa_m\binom{2N}{2m}A(N-m).
 \label{eq:root:packet}
\end{equation}
For a retained depth $M\ge0$, put
\[
 \mathcal T_M(N)=\mathcal E_{2N}
   -\sum_{m=0}^{\min(M,N)}\kappa_m\binom{2N}{2m}A(N-m).
\]
An empty tail, $N\le M$, is exactly zero.

\begin{proposition}[Sharp uniform classical truncation bound]
For every $M,N\ge0$,
\[
 v_2\bigl(\mathcal T_M(N)\bigr)\ge4M+3.
\]
For each fixed $M$, equality holds at $N=M+1$.
\end{proposition}
\begin{proof}
The source $A(n)$ is a dyadic unit for every nonnegative integer $n$.
Each nonempty summand has $m\ge M+1$, so
\eqref{eq:root:weights} and integrality of the binomial coefficient give
the asserted bound. At $N=M+1$, only the first omitted summand remains;
its binomial coefficient and $A(0)$ both equal one.
\end{proof}

Thus later row-dependent cancellation can improve a particular tail,
but cannot increase the uniform exponent required to hold at every row.

\subsection{A finite carry minimum and an exact all-row factorization}
\label{app:root:factorization}

Let $\mathbb Z_2\langle x\rangle$ denote the restricted power series:
their integral coefficients tend to zero dyadically. Extend $A$ by
\[
 \mathscr A(x)=\frac{9^{x+1}-1}{8(x+1)(2x+1)},
 \qquad \mathscr A(-1)=-\frac{\log9}{8}.
\]
The apparent singularity at $x=-1$ is removable. Since
$v_2(\log9)=3$, its convergent expansion is
\begin{equation}
 \mathscr A(x)=\frac{\log9/8}{1+2x}
   \sum_{h\ge0}\frac{(\log9)^h}{(h+1)!}(x+1)^h.
 \label{eq:root:source}
\end{equation}
It follows coefficientwise that
$\mathscr A\in\mathbb Z_2\langle x\rangle^\times$,
$\mathscr A\equiv1\pmod2$, and $\mathscr A(N)=A(N)$.

Set $r=M+1\ge1$ and define
\[
 P_q(x)=\prod_{i=0}^{q-1}(x-i)(2x-(2i+1)),\qquad P_0=1,
\]
\[
 H_r(x)=\sum_{q\ge0}8^q\frac{\beta_{r+q}}{\beta_r}
                    P_q(x)\mathscr A(x-q).
\]
The coefficient valuations tend to infinity. At any nonnegative integer
$x$, terms with $q>x$ vanish. Cancelling the factorials in
\eqref{eq:root:packet} therefore gives the exact identity
\begin{equation}
 \mathcal T_{r-1}(N)=\kappa_r\binom{2N}{2r}H_r(N-r)
 \qquad(N\ge r).\label{eq:root:factor}
\end{equation}
This is a finite rational identity at every original Euler row.

The scalar packet weights have exact valuations
\[
 \lambda_r(q)=3q+s_2(r+q)-s_2(r).
\]
Put
\[
 \mu_r=\min_{q\ge0}\lambda_r(q),\qquad
 J_r=\{q:\lambda_r(q)=\mu_r\},\qquad G_r=2^{-\mu_r}H_r.
\]
Since $\lambda_r(0)=0$ and
$\lambda_r(q)\ge3q+1-s_2(r)$, every minimizing address satisfies
\begin{equation}
 0\le q\le\left\lfloor\frac{s_2(r)-1}{3}\right\rfloor.
 \label{eq:root:cutoff}
\end{equation}
In particular $\mu_r\le0$, and the minimum is found by a finite carry
calculation. Reduction of the whole series gives
\begin{equation}
 \overline G_r(X)=\sum_{q\in J_r}X^{\underline q}
       \quad\text{in }\mathbb F_2[X].
 \label{eq:root:reduction}
\end{equation}
Its highest degree is $\max J_r$, with leading coefficient one.
Thus the normalization is the actual Gauss normalization.

\begin{theorem}[Exact classical tail valuation]
For every $r\ge1$ and $N\ge r$,
\begin{equation}
 v_2\bigl(\mathcal T_{r-1}(N)\bigr)
 =4r-1+v_2\binom Nr+\mu_r+v_2G_r(N-r).
 \label{eq:root:exactvaluation}
\end{equation}
There are unique factors
\[
 G_r(X)=W_r(X)U_r(X),
\]
where $W_r\in\mathbb Z_2[X]$ is monic of degree $d_r=\max J_r$ and
$U_r\in\mathbb Z_2\langle X\rangle^\times$ satisfies $U_r\equiv1\pmod2$.
Consequently $v_2G_r(x)=v_2W_r(x)$ for every $x\in\mathbb Z_2$, and
\[
 \deg W_r\le\left\lfloor\frac{s_2(r)-1}{3}\right\rfloor.
\]
\label{thm:root:finitepreparation}
\end{theorem}
\begin{proof}
Equation~\eqref{eq:root:factor}, the exact weight in
\eqref{eq:root:weights}, and the digit-sum identity
$v_2\binom{2N}{2r}=v_2\binom Nr$ prove
\eqref{eq:root:exactvaluation}. The binomial coefficients themselves
need not be equal.

For preparation, start from the monic polynomial
$g=\overline G_r$ and its zero-one lift $W_1$, with $U_1=1$.
Suppose $G_r-W_kU_k$ is coefficientwise divisible by $2^k$.
Its divided reduction is a finite polynomial because $G_r$ is restricted.
Divide it by $g$ in $\mathbb F_2[X]$:
\[
 \overline{(G_r-W_kU_k)/2^k}=gq+s,\qquad\deg s<d_r.
\]
With zero-one coefficient lifts, set
$W_{k+1}=W_k+2^ks$ and $U_{k+1}=U_k+2^kq$.
The error is now divisible by $2^{k+1}$, the polynomial remains monic of
degree $d_r$, and the unit reduces to one. The sequences converge in the
Gauss norm. Uniqueness of quotient and remainder at successive binary
digits proves uniqueness of the factors. For $d_r=0$, take $W_r=1$.
The degree bound follows from \eqref{eq:root:cutoff}.
\end{proof}

This theorem reduces every depth to a finite polynomial without discarding
the infinite analytic tail. It does not explicitly classify every
$W_r$, decide every repeated root from a fixed precision, or prove that
all its dyadic roots are simple.

\subsection{Regular walls: exact root emergence}
\label{app:root:regular}

For a wall label $j\ge1$ and $0\le q<j$, define
\[
 e_j(q)=3q+s_2(j-1)-s_2(j-q-1),\qquad
 E_j=\min_{q<j}e_j(q),
\]
\[
 \Theta(j)=3j+1+s_2(j-1),\qquad L_j=\Theta(j)-E_j.
\]
Here $E_j$ is a carry minimum, not an Euler number. A wall label is
\emph{regular} when $e_j(q)>0$ for $1\le q<j$.
For $r=2^L-j$ with $L>L_j$, put $\delta=L-L_j$.
The carry identities give
\[
 \mu_r=E_j-\delta,\qquad J_r=\{j\},\qquad
 G_r(x)=P_j(x)U(x-j)+2^\delta V(x),
\]
where $U\equiv1\pmod2$. For a regular label, also
$E_j=0$ and $V\equiv1\pmod2$.

For $0\le i<j$, let
\[
 d_i=v_2(i!)+v_2((j-1-i)!),\qquad
 k_i=1+\left\lfloor\log_2\max(i,j-1-i)\right\rfloor,
\]
\[
 B_i=i+2^{k_i}\mathbb Z_2,\qquad t_i=d_i+k_i.
\]
At $j=1$ use $d_0=k_0=t_0=0$. The balls $B_i$ partition
$\mathbb Z_2$. Indeed a residue represented twice among $0,\ldots,j-1$
splits into its two children; every other represented residue remains
one ball. On each such ball,
\[
 P_j(i+2^{k_i}y)=2^{t_i}yQ_i(y),\qquad
 Q_i\in\mathbb Z_2\langle y\rangle^\times,
 \qquad Q_i\equiv1\pmod2.
\]

\begin{theorem}[Regular-wall root law]
For every regular label $j$ and every integer $\delta\ge1$,
\[
 v_2G_r(x)=
 \begin{cases}
  \delta,&x\in B_i,\ \delta<t_i,\\
  d_i+v_2(x-\alpha_i),&x\in B_i,\ \delta\ge t_i.
 \end{cases}
\]
In the second case there is exactly one root $\alpha_i$ in $B_i$,
it is simple, and $v_2(\alpha_i-i)=\delta-d_i$.
Thus the number of dyadic roots is exactly
\[
 \#\{i:0\le i<j,\ t_i\le\delta\}.
\]
\label{thm:root:regular}
\end{theorem}
\begin{proof}
If $\delta<t_i$, the unit perturbation has smaller valuation than
the skeleton term throughout $B_i$, so the valuation is exactly
$\delta$. Otherwise, after the ball substitution and division by
$2^{t_i}$, the coefficientwise reduction is $y+c$ for a constant
$c\in\mathbb F_2$. A restricted integral series with this reduction
is a dyadic isometry: division of its difference by $y-z$ gives an
odd linear coefficient plus even higher contributions. Binary lifting
therefore gives its unique root and exact distance law. Evaluating at
$x=i$ gives the stated displacement, since
$G_r(i)=2^\delta V(i)$ has exact valuation $\delta$.
The partition proves exhaustion of the roots.
\end{proof}

The largest emergence threshold is
\[
 \max_i t_i=v_2((j-1)!)+\lceil\log_2 j\rceil.
\]
For $j=2^a$, all thresholds equal $j-1$. There are no roots before
that distance and exactly $j$ simple roots at and after it, within the
positive-distance range. The equality and pre-wall cases are also
resolved: regular labels $j\ge2$ have no roots there; $j=1$ has one
simple odd root at equality.

For example, $j=5$ has $(t_0,\ldots,t_4)=(6,3,4,3,6)$.
The exact root counts at distances $1,2,3,4,5,\ge6$ are
$0,0,2,3,3,5$. The older sufficient stages $3,5,7$ are not the
sharp emergence stages, which are $3,4,6$.
The first irregular label, $j=17$, has $V\equiv1+X\pmod2$;
the regular-wall theorem does not apply. Complete coefficient
certificates close distances $1$ through $19$, and the general separation
theorem gives all seventeen simple roots for $\delta\ge20$.
Those two kinds of evidence together resolve that family.

\subsection{What the later simplicity reductions prove}
\label{app:root:families}

For a general depth $r$, call $j_r=\max J_r$ its \emph{true terminal}.
This definition depends on the actual carry minimum; rewriting the same
integer as another wall expression does not change $j_r$.
The universal simplicity target is
\begin{equation}
 \text{for every integer }r\ge1\text{ and }x\in\mathbb Z_2,
 \qquad G_r(x)=0\ \Longrightarrow\ G_r'(x)\ne0.
 \label{eq:root:simplicitytarget}
\end{equation}
The following are completed restricted results of the later continuations.
Their family statements include a proof that the indicated label is the
true terminal throughout the displayed ranges.

\begin{center}
\begin{tabular}{lll}
\toprule
Depth $r$ & Odd multiplier range & Outer scales\\
\midrule
$B2^L-129$ & $0<B<2^{89}$ & $L\ge392$\\
$B2^L-161$ & $0<B<2^{932}$ & $L\ge487$\\
\bottomrule
\end{tabular}
\end{center}
All dyadic roots are simple in both infinite integer families.
For arbitrary larger admissible multipliers, a multiple root with true
terminal $129$ can occur only at $L=515$ and in two necessary multiplier
classes modulo $2^{90}$. With true terminal $161$, it can occur only at
$L=640$ and in one necessary multiplier class modulo $2^{935}$.
These classes are necessary conditions. The records do not assert
existence of a singular ordinary Euler index in them.

For terminal $161$, the neighboring phases are resolved uniformly:
$L=639$ has exactly one simple root and $L=641$ exactly three simple
roots, for every positive odd multiplier with that true terminal.
The middle phase is reduced, before the deeper certificate is applied,
to
\[
 B\equiv17659\pmod{2^{18}},\qquad
 x\equiv10202\pmod{2^{15}}.
\]
The $935$-bit multiplier condition composes this whole-phase reduction
with an inherited $917$-bit local parameter certificate. It is neither
a new $935$-bit reconstruction nor an all-precision interpolation of
the undepleted Euler family.

The structural argument also extends beyond lower address one.
For a certified true-terminal wall, let $a$ be its last lower minimizing
address, $k=j_r-a$ the regular complement, and $t=v_2(k)$.
Under the regular lower-address hypotheses of the source theorem,
all roots above emerged complement balls are simple. Every root is
simple at every wall distance in the band
\[
 \Theta(a)\le t<\Theta(a)+v_2(a!).
\]
Equality walls are simple even outside that band. This includes
certified terminal-$258$ walls with $a=2$, $k=256$, and $L\ge777$.
For $a=3$, $k=4096$, only $L=16385$ remains a possible critical scale.
Outside the band, the theorem gives necessary local valuation conditions;
it does not exclude the remaining unit cancellation. Irregular lower
addresses need their own prefix analysis.

\subsection{The actual stationary discriminant}
\label{app:root:stationary}

The current focused obstruction is the admitted ordinary-index family
\begin{equation}
 r=(17659+2^{18}T)2^{640}-161,
 \qquad T\in\mathbb Z_{\ge0},\qquad j_r=161.
 \label{eq:root:focusedfamily}
\end{equation}
On its surviving chart define the \emph{actual}, undepleted function
\[
 K_r(Z)=2^{-178}G_r(10202+2^{15}Z).
\]
The complete coefficient theorem gives
\begin{equation}
 K_r(Z)=a_r+2b_rZ+u_rZ^2+256Z^3R_r(Z),
 \quad a_r,b_r\in\mathbb Z_2,\quad
 u_r\in\mathbb Z_2^\times,\quad
 R_r\in\mathbb Z_2\langle Z\rangle.
 \label{eq:root:quadraticclass}
\end{equation}
Its half derivative is a dyadic isometry. It therefore has a unique
critical point $c_r\in\mathbb Z_2$. Set
$\Phi_r^{\mathrm{act}}=K_r(c_r)$.
The earlier independent-undepletion notation for this same scalar is
$\Phi_r(2^{2r-25})$.
Excluding a multiple root on the surviving chart is exactly the assertion
$\Phi_r^{\mathrm{act}}\ne0$.

\begin{proposition}[Exact quadratic compression]
There is a unique preparation
\[
 K_r(Z)=\mathcal U_r(Z)\mathcal W_r(Z),\qquad
 \mathcal W_r(Z)=Z^2+2\beta_rZ+\alpha_r,\qquad
 \mathcal U_r\equiv u_r\pmod{256}.
\]
With $\delta_r=\beta_r^2-\alpha_r$ and
$\tau_r=\mathcal U_r'(c_r)/\mathcal U_r(c_r)$,
\begin{equation}
 \delta_r=-\frac{\Phi_r^{\mathrm{act}}}{\mathcal U_r(c_r)}
  \left(1-\frac{\tau_r^2\Phi_r^{\mathrm{act}}}{4\mathcal U_r(c_r)}\right),
 \qquad v_2(\delta_r)=v_2(\Phi_r^{\mathrm{act}}).
 \label{eq:root:discriminant}
\end{equation}
The equality includes the zero case.
\end{proposition}
\begin{proof}
The same successive monic polynomial division used in
Theorem~\ref{thm:root:finitepreparation}, started at precision eight,
preserves the quadratic and unit reductions of
\eqref{eq:root:quadraticclass}. Write
$W=Z^2+2\beta_rZ+\alpha_r$ and $w=W(c_r)$.
Stationarity gives $W'(c_r)=-\tau_rw$, while
$w=\Phi_r^{\mathrm{act}}/\mathcal U_r(c_r)$.
The quadratic identity
$\beta_r^2-\alpha_r=W'(c_r)^2/4-W(c_r)$ proves the formula.
Since $\tau_r\in256\mathbb Z_2$, its bracket is a unit.
\end{proof}

This compression retains the unit-cofactor derivative. It does not
assert that $c_r$, $\alpha_r$, or $\beta_r$ is rational or algebraic
over $\mathbb Q$. The remaining arithmetic obligation is therefore
\begin{equation}
 \delta_r\ne0\quad\text{for every admitted ordinary index in
 \eqref{eq:root:focusedfamily}.}
 \label{eq:root:remaining}
\end{equation}
Even this assertion would still leave the terminal-$129$ and higher-address
obligations in the universal target \eqref{eq:root:simplicitytarget}.

\subsection{Approximation reaches every depth; separation remains open}
\label{app:root:approximation}

The latest continuation retains the exact Bernoulli undepletion factor
\[
 e_{r,q}=\frac{1-2^{2r-1}}{1-2^{2r-1}4^q},\qquad q\ge0.
\]
It is a dyadic unit, including $e_{r,0}=1$. Multiplication by this
factor preserves the all-$q$ error of the shifted Jacobi moment model.
Consequently the actual model critical values satisfy, for every $m\ge25$,
\begin{equation}
 v_2\bigl(\Phi_{r,m}^{\mathrm{act}}-\Phi_r^{\mathrm{act}}\bigr)
 \ge L_m=8m-4s_2(m)-178,
 \qquad L_m\longrightarrow\infty.
 \label{eq:root:actualapprox}
\end{equation}
The estimate includes every packet and controls motion of the critical
point. There is no ceiling at the undepletion exponent $2r-25$.
For each fixed admitted $r$ it implies the exact equivalence
\begin{equation}
 \Phi_r^{\mathrm{act}}\ne0
 \quad\Longleftrightarrow\quad
 \text{there exists }m\ge25\text{ such that }
 v_2(\Phi_{r,m}^{\mathrm{act}})<L_m.
 \label{eq:root:stricttest}
\end{equation}
Indeed a passing strict test has the same valuation as the limit; conversely
any finite limit valuation is eventually smaller than $L_m$.
The existence of a passing depth for every ordinary index is unproved.
Nonzero models at every depth are compatible with an exactly zero limit.

The degree-two preparation also reduces the old representative-height
cost to a linear cost, about $2P$ bits at precision $P$. This remains
larger than the $P$ bits needed for a strict nonzero residue certificate.
The fixed rational control $(Z-1/3)^2$ has zero critical value while
its least-residue integer discriminant representatives are nonzero at
every even $P\ge6$. Thus a nonzero integer representative alone cannot
prove a nonzero dyadic discriminant.

One sufficient new arithmetic input would be an odd-denominator
projective quadratic
\[
 A Z^2+2BZ+C\equiv A\mathcal W_r(Z)\pmod{2^P},\qquad A\text{ odd},
 \qquad 0<|B^2-AC|<2^P,
\]
at a proved precision for every admitted ordinary $r$.
Then $(B^2-AC)/A^2\equiv\delta_r\pmod{2^P}$ cannot vanish
modulo $2^P$. No depth-independent theorem supplying this inequality
has been established. The tested recurrence and exterior-contact
identities retain balancing terms and do not force a contradiction.
Finite exponential models approximate the actual target at every depth,
but their order grows with precision, their clearing polynomial has
zeros on the actual chart, and their proved construction-specific height
cost does not give the required discriminant bound.

\subsection{Status and source map}
\label{app:root:evidence}

\begin{center}
\begin{tabular}{>{\raggedright\arraybackslash}p{0.64\linewidth}>{\raggedright\arraybackslash}p{0.26\linewidth}}
\toprule
Claim & Status at this edition\\
\midrule
Nishibuchi Conjecture~3.9 in its stated parameter range
 & Resolved in the main text\\
Sharp uniform ordinary Euler tail exponent $4M+3$
 & Proved, with earlier formal coverage\\
Exact tail factorization and finite polynomial at every depth
 & Proved; later prose development\\
Regular-wall emergence and indicated integer simplicity families
 & Proved within their stated hypotheses\\
Actual quadratic compression and unbounded stationary approximation
 & Proved; no fresh Lean build\\
Nonzero actual stationary discriminant for every admitted index
 & Open\\
Universal simplicity of every $G_r$
 & Open\\
Native~13.2 presentation invariance and Native~13.3 global/completed gluing
 & Separate open general programs\\
\bottomrule
\end{tabular}
\end{center}

The tail factorization, regular-wall proofs, finite preparation, and
coefficient cutoffs are recorded in the preserved
\nolinkurl{sources/euler_v4.md} of the regular-address checkpoint.
The two integer families and multiplier conditions are proved in its
\nolinkurl{next/terminal129/PROOF.md} and
\nolinkurl{next/critical_phases/PROOF.md};
\nolinkurl{local_parameter/RECENTER_PROOF.md} supplies the inherited
local precision certificate. The higher-address extension is in
\nolinkurl{next/higher_address/regular_address_transversality.md}.
The stationary checkpoint records the exact discriminant identity in
\nolinkurl{stationary_separation/ordinary/QUADRATIC_COMPRESSION_AND_HEIGHT.md}
and actual convergence in
\nolinkurl{stationary_separation/root/ACTUAL_LIMIT_AND_SEPARATION.md}.
Its synthesis and proof reviews retain the failed
implications and the precise remaining arithmetic input.

The source archive's \nolinkurl{evidence/root_scope_inventory.md} gives
the full checkpoint paths and content hashes used for this appendix.
The 2026-10-05 complete-work archive preserves the original layouts and
historical releases. Its stored hashes and verification receipts are
evidence of integrity and recorded execution; they are not substitutes
for a proof of the open assertions above.

\clearpage
\section{Detailed continuation derivations}\label{app:continuation}
The following sections develop the detailed arguments behind polynomial
lifting, correction torsors, compatible towers and local regularity. Each
construction retains its stated hypotheses. The formal companion identifies
the corresponding declarations and verification records. The main text
distinguishes the completed results from the remaining questions.
\hypertarget{cont1:paper-ix-finite-polynomial-lifting-and-shared-corrections}{%
\subsection{Chapter IX: finite polynomial lifting and shared
corrections}\label{cont1:paper-ix-finite-polynomial-lifting-and-shared-corrections}}

The nonlinear Taylor argument supplies the correction formula used in
one-step polynomial lifting. The statements concern a fixed
finite list of integer polynomials. They hold for every positive integer
modulus, including composite moduli and modulus one. No
Jacobian-surjectivity assumption is needed.

\hypertarget{cont1:statement-and-notation}{%
\subsubsection{Statement and notation}\label{cont1:statement-and-notation}}

Let \[
F=(F_1,\ldots,F_e)\in\mathbb Z[X_1,\ldots,X_d]^e,
\quad q\ge1,\quad L\ge1,
\quad x\in\mathbb Z^d,
\quad F(x)\equiv0\pmod{q^L}.
\] Write \(R_q=\mathbb Z/q\mathbb Z\), let \(\rho_q\) denote
coordinatewise reduction, and put \[
r_x=F(x)/q^L\in\mathbb Z^e,
\qquad A_x=DF(x):\mathbb Z^d\longrightarrow\mathbb Z^e,
\qquad J_x=\rho_q\circ A_x.
\] Here \(A_x\) is computed from the actual formal partial derivatives
of the input polynomials. The Lean map \(J_x\) has domain
\(\mathbb Z^d\) and codomain \(R_q^e\), viewed as integer modules. It
factors through reduction to \(R_q^d\); the cokernel below is presented
directly using its integer-domain image.

Define the finite obstruction module and its distinguished class by \[
\mathcal O_x=R_q^e/\operatorname{im}J_x,
\qquad o_x=[\rho_q(r_x)]\in\mathcal O_x.
\] The checked criterion is \[
\boxed{
 o_x=0
 \iff
 \exists y\in\mathbb Z^d:\quad
 y\equiv x\pmod{q^L},\qquad F(y)\equiv0\pmod{q^{L+1}}.
}
\] The residual exists by the lower congruence and is unique because
\(q>0\). The quotient is finite since its ambient module \(R_q^e\) is
finite. The theorem permits zero variables or zero equations as well.

\hypertarget{cont1:the-checked-taylor-bridge}{%
\subsubsection{The checked Taylor
bridge}\label{cont1:the-checked-taylor-bridge}}

For every \(P\in\mathbb Z[X_1,\ldots,X_d]\), integer \(a\), and
\(x,t\in\mathbb Z^d\), Lean proves \[
\boxed{
 a^2\mid P(x+at)-P(x)-a\sum_{i=1}^{d}(\partial_iP)(x)t_i.
}
\] The proof uses polynomial induction: constants, addition, and
multiplication by a variable. The product rule for the actual partial
derivatives controls the multiplication step. It does not assume a
Taylor identity or a liftability condition.

Taking \(a=q^L\), the inequality \(2L\ge L+1\) gives \[
F(x+q^Lt)\equiv q^L(r_x+A_xt)\pmod{q^{L+1}}.
\] Since \(q^L\ne0\) in the integers, cancellation in divisibility
yields \[
F(x+q^Lt)\equiv0\pmod{q^{L+1}}
\iff \rho_q(r_x)+J_xt=0.
\] Every integer representative of the lower residue has the form
\(x+q^Lt\), and this correction equation has a solution exactly when
\(o_x=0\).

The positive-depth hypothesis is necessary for the general nonlinear
theorem. For \(F(X)=X^2-1\), \(q=2\), \(L=0\), \(x=0\), the Jacobian is
zero and the residual class is nonzero, yet \(y=1\) is a solution modulo
two in the same residue modulo one. Affine systems retain the previously
checked \(L\ge0\) theorem.

\hypertarget{cont1:independence-of-the-integer-representative}{%
\subsubsection{Independence of the integer
representative}\label{cont1:independence-of-the-integer-representative}}

Suppose \(x'=x+q^Ls\) represents the same lower residue, with exact
residual \(r_{x'}\). Polynomial evaluation respects coordinatewise
congruences, including when applied to each partial derivative. Hence \[
J_{x'}=J_x,
\qquad
\rho_q(r_{x'})=\rho_q(r_x)+J_xs.
\] The two cokernels therefore have the same ambient module and the same
defining submodule, and \[
[\rho_q(r_{x'})]=[\rho_q(r_x)].
\] Lean states this using one explicit common reduced map, avoiding an
implicit comparison between differently indexed quotient types. This
establishes independence of the integer representative within the fixed
polynomial presentation.

Changing the polynomial presentation is a separate question. Appending a
zero equation can change the raw cokernel by adding a copy of \(R_q\),
even though it does not change the solutions.

\hypertarget{cont1:two-polynomial-blocks-and-one-shared-lift}{%
\subsubsection{Two polynomial blocks and one shared
lift}\label{cont1:two-polynomial-blocks-and-one-shared-lift}}

Let \(G\in\mathbb Z[X_1,\ldots,X_d]^f\) be a second fixed system with
the same variables and the same lower root \(x\). Let
\(K_x=\rho_q\circ DG(x)\), and define \[
\mathcal C_x=\mathbb Z^d/(\ker J_x+\ker K_x).
\] This quotient is finite: it is a quotient of the finite image of
\(J_x\), or equivalently of the finite image of \(K_x\).

Suppose \(u,v\in\mathbb Z^d\) are separately valid corrections: \[
F(x+q^Lu)\equiv0\pmod{q^{L+1}},
\qquad
G(x+q^Lv)\equiv0\pmod{q^{L+1}}.
\] Then \[
\boxed{
\begin{gathered}
[u]=[v]\text{ in }\mathcal C_x
\iff
\exists y\equiv x\pmod{q^L}:\\
F(y)\equiv0\pmod{q^{L+1}},\quad G(y)\equiv0\pmod{q^{L+1}}.
\end{gathered}
}
\] The labels are independent of the choices of valid \(u\) and \(v\).
An unconditional version asks for separately valid corrections with
equal labels; it does not assume that each block is liftable.

Taylor transport turns a common correction \(t\) into the simultaneous
equations \(J_xt=J_xu\) and \(K_xt=K_xv\). The previously checked
common-quotient theorem gives precisely the displayed equivalence.
Separate solvability alone does not imply joint solvability; the
controls include \(F(X)=X\), \(G(X)=X-2\), starting at \(x=0\) modulo
two and lifting to modulus four.

\hypertarget{cont2:paper-ix-correction-torsors-and-exact-polynomial-lift-counts}{%
\subsection{Chapter IX: correction torsors and exact polynomial lift
counts}\label{cont2:paper-ix-correction-torsors-and-exact-polynomial-lift-counts}}

The correction-digit torsor gives exact prime-field lift counts. The counts apply to distinct next-level
residue classes satisfying the actual polynomial equations. They do not
count integer representatives, of which there can be infinitely many.

\hypertarget{cont2:the-theorem}{%
\subsubsection{The theorem}\label{cont2:the-theorem}}

Let \(p\) be prime, \(L\ge1\), and \[
F\in\mathbb Z[X_1,\ldots,X_d]^e,
\qquad x\in\mathbb Z^d,
\qquad F(x)\equiv0\pmod{p^L}.
\] Write \[
r_x=F(x)/p^L,\qquad b_x=r_x\bmod p,
\qquad A_x=DF(x)\bmod p:\mathbb F_p^d\longrightarrow\mathbb F_p^e,
\qquad \rho_x=\operatorname{rank}_{\mathbb F_p}A_x.
\] The actual lift set is \[
\mathcal L_{x,L}=
\left\{y\in(\mathbb Z/p^{L+1}\mathbb Z)^d:
 y\equiv x\pmod{p^L},\quad F(y)=0\right\},
\] where polynomial evaluation in the displayed set is modulo
\(p^{L+1}\). Define \[
\mathcal O_x=\mathbb F_p^e/\operatorname{im}A_x,
\qquad o_x=[b_x].
\] Lean proves \[
\boxed{
\#\mathcal L_{x,L}=
\begin{cases}
p^{d-\rho_x},&o_x=0,\\
0,&o_x\ne0,
\end{cases}
\qquad
\#\mathcal O_x=p^{e-\rho_x}.
}
\] These formulas permit singular Jacobians and zero variables or zero
equations. Nonsingularity is not assumed. In particular, rank alone
gives the size of a nonempty lift fiber; the residual obstruction
determines whether it is empty.

\hypertarget{cont2:why-these-are-actual-residue-class-counts}{%
\subsubsection{Why these are actual residue-class
counts}\label{cont2:why-these-are-actual-residue-class-counts}}

For every positive integer \(q\), including composite moduli and modulus
one, and every \(L\ge0\), the checked map \[
\Phi_{x,L}:(\mathbb Z/q\mathbb Z)^d\longrightarrow
\{y\bmod q^{L+1}:y\equiv x\pmod{q^L}\},
\qquad t\longmapsto x+q^L\widetilde t
\] is a bijection. Here \(\widetilde t\) is the canonical integer
representative of the digit vector. The entire residue fiber therefore
has \(q^d\) elements. Injectivity is proved by cancelling \(q^L\) in the
corresponding integer divisibility statement; surjectivity comes from
the lower-residue congruence.

The formal type \mcode{ActualLifts} is a subtype of these genuine
\texttt{ZMod\ (q\^{}(L+1))} residue vectors. Its polynomial predicate is
evaluated on canonical integer representatives and tested for
divisibility by \(q^{L+1}\). The checked polynomial congruence theorem
makes that predicate independent of the integer representative. Thus
using canonical representatives does not discard or duplicate residue
classes.

At positive depth, the preceding actual multivariable Taylor theorem
gives \[
F(x+p^L\widetilde t)\equiv
p^L(r_x+DF(x)\widetilde t)\pmod{p^{L+1}}.
\] The new bridge proves that the evaluated integer Jacobian reduces to
the actual \(\mathbb F_p\)-linear map \(A_x\). Restricting
\(\Phi_{x,L}\) therefore gives an explicit equivalence \[
\boxed{\{t\in\mathbb F_p^d:b_x+A_xt=0\}\simeq\mathcal L_{x,L}.}
\] The proof uses the actual formal partial derivatives of the original
polynomial list.

\hypertarget{cont2:the-torsor}{%
\subsubsection{The torsor}\label{cont2:the-torsor}}

Let \(\mathcal T_x=\{t:b_x+A_xt=0\}\). Its natural kernel action is \[
k\cdot t=k+t,\qquad k\in\ker A_x.
\] For every pair \(u,t\in\mathcal T_x\), exactly one kernel vector
sends \(u\) to \(t\), namely \(t-u\). Thus every nonempty correction
fiber is a torsor for \(\ker A_x\). Choosing one valid digit \(u\) gives
\[
\mathcal T_x\simeq\ker A_x,
\qquad t\longmapsto t-u.
\] Lean constructs an explicit \mcode{AddTorsor} structure for the
digit fiber, proves the unique translation property, and supplies a
chosen-point equivalence from the actual residue lift set to the kernel.
A basepoint is required for this equivalence, but the translation action
itself does not depend on choosing one.

The torsor and residue bijection remain valid over every positive
modulus at positive polynomial depth. The displayed powers involving
vector-space rank are restricted to prime moduli.

\hypertarget{cont2:connecting-the-old-obstruction-to-the-new-count}{%
\subsubsection{Connecting the old obstruction to the new
count}\label{cont2:connecting-the-old-obstruction-to-the-new-count}}

The earlier obstruction used the integer-domain map \[
J_x=\rho_p\circ DF(x):\mathbb Z^d\to\mathbb F_p^e.
\] The new count uses \(A_x:\mathbb F_p^d\to\mathbb F_p^e\). Reduction
of integer vectors onto finite vectors is surjective, so the two maps
have the same image after restricting scalars. Lean constructs a
class-preserving integer-linear equivalence \[
\mathbb F_p^e/\operatorname{im}J_x
\simeq
\mathbb F_p^e/\operatorname{im}A_x.
\] Consequently the cardinality theorem counts the original obstruction
group, preserving its distinguished residual class. It is not merely a
formula for a separately introduced comparison object.

Rank-nullity gives \(\dim\ker A_x=d-\rho_x\), and the quotient dimension
is \(e-\rho_x\). Cardinalities of finite vector spaces complete both
formulas.

\hypertarget{cont3:paper-ix-compatible-polynomial-towers-and-p-adic-roots}{%
\subsection{Chapter IX: compatible polynomial towers and p-adic
roots}\label{cont3:paper-ix-compatible-polynomial-towers-and-p-adic-roots}}

The nonsingular construction passes from one-step
polynomial lifting to an infinite compatible tower and an actual p-adic
root. Surjectivity is required only of the reduced Jacobian at the
initial root. The proof derives its persistence, constructs the tower,
proves convergence and verifies the original polynomial equations at the
limit.

\hypertarget{cont3:main-theorem}{%
\subsubsection{Main theorem}\label{cont3:main-theorem}}

Let \(p\) be prime, \(L\ge1\), and \[
F\in\mathbb Z[X_1,\ldots,X_d]^e,\qquad x\in\mathbb Z^d,
\qquad F(x)\equiv0\pmod{p^L}.
\] Assume the actual reduced Jacobian \[
A_x=DF(x)\bmod p:\mathbb F_p^d\longrightarrow\mathbb F_p^e
\] is surjective. Then Lean proves that there is a point \[
\boxed{a\in\mathbb Z_p^d,\qquad F(a)=0,\qquad
 a\equiv x\pmod{p^L}.}
\] Here the coefficients of the original integer polynomials are cast
into \(\mathbb Z_p\), and their multivariable evaluation is exactly
zero. This is an existence theorem for a fixed polynomial presentation
with the stated Jacobian hypothesis.

The stronger checked statement constructs \(z_n\in\mathbb Z^d\),
\(n\ge0\), with \[
\begin{aligned}
z_0&=x,\\
F(z_n)&\equiv0\pmod{p^{L+n}},\\
z_{n+1}&\equiv z_n\pmod{p^{L+n}},\\
DF(z_n)\bmod p&=A_x,\\
z_n&\longrightarrow a\quad\text{in }\mathbb Z_p^d,\\
a-z_n&\in p^{L+n}\mathbb Z_p^d.
\end{aligned}
\] The last inclusion retains every prescribed precision. It is a
divisibility lower bound, not an assertion that the error has exactly
that valuation. An approximation can already equal the limiting root.

\hypertarget{cont3:why-the-single-initial-hypothesis-suffices}{%
\subsubsection{Why the single initial hypothesis
suffices}\label{cont3:why-the-single-initial-hypothesis-suffices}}

The one-step Taylor theorem already supplies the exact correction
equation \[
b_n+A_{z_n}t=0,\qquad
b_n=F(z_n)/p^{L+n}\bmod p.
\] Surjectivity supplies a correcting digit for every residual \(b_n\).
The chosen correction gives an actual new integer point \(z_{n+1}\),
using the proved nonlinear Taylor remainder.

Every coordinate change is divisible by \(p^{L+n}\), hence by \(p\).
Each formal partial derivative is itself an integer polynomial, so its
value is unchanged modulo \(p\). Therefore \[
A_{z_{n+1}}=A_{z_n}=A_x.
\] The proof assumes no separate surjectivity or solvability statement
at later levels. A dependent recursion chooses each next point from the
proved one-step existence theorem. It produces one compatible sequence,
rather than unrelated witnesses at different precisions.

The integer tower construction works more generally for every positive
modulus \(q\), including composite \(q\) and \(q=1\), with a surjective
map over \(\mathbb Z/q\mathbb Z\). The p-adic completion and the
rank-based branching formulas in this block use prime \(p\).

\hypertarget{cont3:convergence-and-exact-polynomial-vanishing}{%
\subsubsection{Convergence and exact polynomial
vanishing}\label{cont3:convergence-and-exact-polynomial-vanishing}}

Compatibility implies, for \(m\ge n\), \[
p^{L+n}\mid z_m-z_n.
\] After casting to \(\mathbb Z_p\), this bounds each coordinate
distance by \(p^{-(L+n)}\), so the sequence is Cauchy. Completeness
supplies its limit. Closed norm bounds show that the limit retains the
congruence at every stage.

The generic convergence theorem is proved independently of the
polynomial application. It also proves \[
p^{L+n}\mid v_n\ \text{for all }n
\quad\Longrightarrow\quad v_n\longrightarrow0\text{ in }\mathbb Z_p.
\] Apply this to each integer sequence \(F_i(z_n)\). Continuity of
multivariable polynomial evaluation also gives \(F_i(z_n)\to F_i(a)\).
Uniqueness of limits then yields \(F_i(a)=0\).

The cast-evaluation identity is explicitly checked, so the continuous
function is the evaluation of the original polynomials. Convergence and
polynomial vanishing are conclusions, not assumptions of the main
existence theorem.

\hypertarget{cont3:exact-number-of-choices-at-each-level}{%
\subsubsection{Exact number of choices at each
level}\label{cont3:exact-number-of-choices-at-each-level}}

Surjectivity implies \[
\operatorname{rank}_{\mathbb F_p}A_x=e,\qquad e\le d.
\] Combining this with the preceding exact residue-lift count gives \[
\boxed{\#\mathcal L_{z_n,L+n}=p^{d-e}\quad\text{for every }n.}
\] These are distinct solutions modulo \(p^{L+n+1}\) lying above the
selected point modulo \(p^{L+n}\). The original finite obstruction group
at each smooth point has cardinality one.

The theorem constructs one path through these finite choices. When
\(d>e\), it does not claim that all paths have the same limit. The
generic uniqueness theorem concerns the limit of a fixed sequence only.

\hypertarget{cont3:boundary-controls}{%
\subsubsection{Boundary controls}\label{cont3:boundary-controls}}

The positive control theorems and rejected invalid proofs cover three
distinctions.

\begin{enumerate}
\def\labelenumi{\arabic{enumi}.}
\tightlist
\item
  \textbf{An initial lift does not guarantee an infinite tower without
  the Jacobian condition.} \(X^2-5\) has the odd lower root modulo two
  and a lift modulo four, but no root modulo eight. Lean proves that no
  integer tower can continue through all precisions in this example.
\item
  \textbf{Surjective rectangular Jacobians do not give uniqueness of all
  roots.} The system \(F(X,Y)=X\) has a surjective Jacobian and distinct
  dyadic roots \((0,0)\) and \((0,2)\), both in the same residue modulo
  two.
\item
  \textbf{The retained initial precision cannot be increased without
  justification.} For \(F(X)=X-2\), zero is a lower root modulo two with
  surjective derivative, but its exact root two does not retain zero
  modulo four.
\end{enumerate}

The controls are theorem checks of these explicit examples. They
supplement the universal proof audit and do not replace it.

For several polynomial blocks, the hypothesis applies to their combined
system. Surjectivity of each block separately does not imply
surjectivity of the combined Jacobian or shared liftability. The
preceding shared-correction obstruction remains relevant in that
situation.

\hypertarget{cont4:polynomial-roots-under-coefficient-perturbation}{%
\subsection{Polynomial roots under coefficient
perturbation}\label{cont4:polynomial-roots-under-coefficient-perturbation}}

The argument proves uniqueness of roots for square systems and nearby-root stability
for finite systems of integer polynomials. The polynomials and their
partial derivatives are evaluated directly. No root-uniqueness or
perturbation-stability premise is assumed.

\hypertarget{cont4:notation}{%
\subsubsection{Notation}\label{cont4:notation}}

Let \(p\) be prime, \(L\ge1\), and \[
F,H\in\mathbb Z[X_1,\ldots,X_d]^e,
\qquad x\in\mathbb Z^d.
\] Write \[
A_x=DF(x)\bmod p:\mathbb F_p^d\longrightarrow\mathbb F_p^e,
\qquad B_L(x)=x+p^L\mathbb Z_p^d.
\] For an integer \(c\), the perturbed system is the actual polynomial
list \[
G=F+cH.
\] There is no restriction on the sign of \(c\), and \(c=0\) is allowed.
Divisibility conclusions give lower bounds on precision and allow zero
differences.

\hypertarget{cont4:unique-roots-for-square-systems}{%
\subsubsection{1. Unique roots for square
systems}\label{cont4:unique-roots-for-square-systems}}

If \(d=e\), \(A_x\) is surjective, and \(F(x)\equiv0\pmod{p^L}\), then
\[
\boxed{\exists!\,a\in B_L(x):F(a)=0.}
\] The stronger comparison lemma proves that any two exact roots of
\(F\) in \(B_L(x)\) coincide, without needing a separate lower-root
hypothesis.

\textbf{Proof implemented.} Map every arbitrary p-adic root into the
actual residue rings \(\mathbb Z/p^n\mathbb Z\), and choose its
canonical integer representative. Evaluation commutes with reduction, so
these representatives are genuine finite polynomial roots. The preceding
count theorem gives exactly \(p^{d-d}=1\) next-level residue lift.
Inductively any two roots agree at every precision. The p-adic residue
maps separate points, so the roots are equal. Existence is the preceding
checked tower theorem.

This extends uniqueness from the limit of a selected sequence to
\textbf{all exact roots in the prescribed ball}.

\hypertarget{cont4:coefficient-perturbations-preserve-smooth-liftability}{%
\subsubsection{2. Coefficient perturbations preserve smooth
liftability}\label{cont4:coefficient-perturbations-preserve-smooth-liftability}}

If \(p^L\mid c\), the actual identities give \[
G(x)=F(x)+cH(x),\qquad DG(x)\equiv DF(x)\pmod p.
\] Thus if \(F(x)\equiv0\pmod{p^L}\) and \(A_x\) is surjective, then \[
\exists b\in B_L(x):G(b)=0.
\] This theorem allows rectangular systems. The evaluation, derivative
and Jacobian identities are proved from the polynomial definition
\(G_i=F_i+C(c)H_i\), rather than supplied as hypotheses. Algebraic
congruence helpers also permit composite moduli; exact p-adic root
existence uses prime \(p\).

\hypertarget{cont4:nearby-roots-for-rectangular-systems}{%
\subsubsection{3. Nearby roots for rectangular
systems}\label{cont4:nearby-roots-for-rectangular-systems}}

Let \(M\ge1\) and \(p^M\mid c\). Suppose \(A_x\) is surjective and
\(a\in\mathbb Z_p^d\) is an exact root of \(F\) with
\(a\equiv x\pmod p\). Then \[
\boxed{\exists b\in\mathbb Z_p^d:\quad G(b)=0,
\qquad b\equiv a\pmod{p^M}.}
\] No uniqueness is claimed when \(d>e\).

\textbf{Proof implemented.} Truncate \(a\) to an integer vector \(y\)
modulo \(p^M\). It is an actual lower root of \(F\), and its reduced
Jacobian agrees with \(A_x\). Apply the coefficient-perturbed existence
theorem at precision \(M\), then combine the two congruences through
\(y\).

\hypertarget{cont4:full-precision-for-the-unique-square-root}{%
\subsubsection{4. Full precision for the unique root of a square
system}\label{cont4:full-precision-for-the-unique-square-root}}

Suppose \(d=e\), \(M\ge L\ge1\), \(p^M\mid c\), and \(A_x\) is
surjective. If \[
a,b\in B_L(x),\qquad F(a)=0,\qquad G(b)=0,
\] then \[
\boxed{b-a\in p^M\mathbb Z_p^d.}
\] In particular, for \(G=F+p^MH\), no p-adic digits below depth \(M\)
are changed in the selected root.

The formalization also proves a combined existence theorem. Starting
only with \(F(x)\equiv0\pmod{p^L}\), it constructs an original root
\(a\in B_L(x)\) and proves that there is exactly one \(b\in B_L(x)\)
satisfying \(G(b)=0\) and \(b\equiv a\pmod{p^M}\). The comparison
theorem shows that the precision bound holds for every root of \(G\) in
that ball, so uniqueness is not obtained merely by adding a restrictive
displacement condition.

\textbf{Proof implemented.} Section 3 produces a nearby root \(b'\).
Since \(M\ge L\), it belongs to \(B_L(x)\). The perturbed Jacobian
remains surjective, and square uniqueness identifies \(b'\) with the
arbitrary root \(b\).

\hypertarget{cont5:paper-ix-continuous-p-adic-parameter-charts}{%
\subsection{Chapter IX: continuous p-adic parameter
charts}\label{cont5:paper-ix-continuous-p-adic-parameter-charts}}

\hypertarget{cont5:main-statement}{%
\subsubsection{Main statement}\label{cont5:main-statement}}

Let \(p\) be prime, \(L\ge1\), \(e,q\ge0\), and \[
F\in\mathbb Z[X_1,\ldots,X_{e+q}]^e,\qquad
x\in\mathbb Z^{e+q},\qquad F(x)\equiv0\pmod{p^L}.
\] Assume the actual reduced Jacobian \[
A_x=DF(x)\bmod p:\mathbb F_p^{e+q}\longrightarrow\mathbb F_p^e
\] is surjective. All vector congruences are coordinatewise.

There is a coordinate permutation \(\sigma\), with
\((\sigma z)_i=z_{\sigma(i)}\). Write \(\sigma z=(y,t)\), using \(e\)
solved coordinates and \(q\) free coordinates. Define \[
(y_0,t_0)=\sigma x,\qquad
B=x+p^L\mathbb Z_p^{e+q},\qquad
T=t_0+p^L\mathbb Z_p^q,\qquad
\mathcal R=\{z\in B:F(z)=0\}.
\] There is a homeomorphism \[
h:T\xrightarrow{\ \sim\ }\mathcal R,\qquad
h^{-1}(z)=\operatorname{free}(\sigma z).
\] Equivalently, a continuous function
\(\phi:T\to y_0+p^L\mathbb Z_p^e\) gives the entire root locus through
\(\sigma h(t)=(\phi(t),t)\). Every parameter has exactly one
corresponding root in \(B\). For every \(N\ge0\), \[
\boxed{h(t)\equiv h(u)\pmod{p^N}\iff t\equiv u\pmod{p^N}.}
\]

The coordinate choice is a conclusion. A suitable minor and a
preexisting chart are not hypotheses. The theorem covers the original
root locus after undoing the coordinate reordering. Since \(x\) is an
approximate root, it does not assert \(\phi(t_0)=y_0\) unless \(F(x)=0\)
exactly. The cases \(p=2\), \(L=1\), \(e=0\), and \(q=0\) are included.

\hypertarget{cont5:supporting-theorems}{%
\subsubsection{Supporting theorems}\label{cont5:supporting-theorems}}

For a rectangular integer polynomial system with a surjective reduced
Jacobian at \(x\), and \(L\ge1\), \[
F(x+p^L\mathbb Z_p^d)=F(x)+p^L\mathbb Z_p^e.
\] This covers arbitrary p-adic output targets. No lower-root hypothesis
is needed.

When \(d=e\), the actual restricted polynomial map is a homeomorphism.
Moreover, for arbitrary \(a,b\in x+p\mathbb Z_p^d\) and every \(N\ge0\),
\[
F(a)\equiv F(b)\pmod{p^N}\iff a\equiv b\pmod{p^N}.
\] Its inverse therefore preserves every congruence depth.

\hypertarget{cont5:proof-structure}{%
\subsubsection{Proof structure}\label{cont5:proof-structure}}

\begin{enumerate}
\def\labelenumi{\arabic{enumi}.}
\tightlist
\item
  The actual integer Taylor formula gives one additional input digit
  from one additional output digit when the reduced square Jacobian is
  surjective. Its finite domain and codomain have the same size, so it
  is injective. Induction gives every integer precision.
\item
  Canonical integer truncations of arbitrary p-adic vectors preserve
  evaluation modulo each power of \(p\). Applying the integer result to
  truncations proves the p-adic congruence equivalence.
\item
  Subtracting an integer target from the constant terms leaves the
  actual Jacobian unchanged. The previous root-lifting theorem supplies
  a preimage. The image of the compact input ball is closed and contains
  the convergent truncations of every admissible p-adic target, hence
  contains the target. No compatibility of arbitrarily chosen preimages
  is assumed.
\item
  In the square case, the resulting continuous bijection from a compact
  ball to a Hausdorff ball is a homeomorphism. Its forward map is actual
  polynomial evaluation, and its inverse solves the output equation.
\item
  For a specified surjective solved-variable block, apply this inverse
  to the actual augmented polynomial map \[
  \Psi(y,t)=(F(y,t),t).
  \] Its derivative has an identity free-coordinate block. The inverse
  at \((0,t)\) gives the root. The lower-root congruence places this
  target in the output ball. Injectivity proves uniqueness and full
  graph coverage. Projection identifies the inverse chart and reflects
  congruences.
\item
  The columns of a surjective finite-dimensional linear map contain a
  basis. Select those columns and extend their injection to a coordinate
  permutation. Polynomial renaming by \(\sigma^{-1}\) transports actual
  evaluations and partial derivatives. The fixed-coordinate chart then
  gives the main theorem on the entire original root locus.
\end{enumerate}

\hypertarget{cont5:scope}{%
\subsubsection{Scope}\label{cont5:scope}}

Coefficients are integers; roots and parameters are arbitrary p-adic
points. This supplies continuous polynomial parameter charts.
Analyticity, differentiability, arbitrary p-adic coefficient interfaces,
weighted lifting for nonunit Jacobians, and first-order root-motion
estimates are separate.

Full Jacobian surjectivity need not allow an arbitrary prescribed
parameter choice. At the dyadic origin, \(F(X,T)=X^2-T\) has a
surjective full Jacobian but zero \(X\)-derivative. Prescribing \(T=2\)
gives no root even modulo \(4\). Suitable coordinates instead solve for
\(T\).

\hypertarget{cont5:verification-and-provenance}{%
\subsubsection{Verification and
provenance}\label{cont5:verification-and-provenance}}

The eight new mathematical modules contain 75 declarations: 60 theorems
and 15 definitions. The active polynomial dependency set contains 25
mathematical modules. Counts including inherited support and theorem
contracts are recorded separately in the final report.

The exact universal statements are Lean theorems, not conclusions from
sampled numerical tests. The fresh gate rebuilds project sources with
warnings treated as errors, checks expanded theorem contracts, audits
explicit declarations, and checks rejected invalid extensions. Only
propext, Classical.choice, and Quot.sound are permitted as axiom
dependencies.

Pinned runtime: Lean 4.34.0-rc2, Mathlib revision
\mcode{aeb72f1e9981b6d35b426b16dca7759be3e7ba06}.

See \nolinkurl{inventory.json}, \nolinkurl{verification/polynomial_charts_verification.json},
and \nolinkurl{INDEPENDENT_PROOF_AUDIT.md} for actual evidence.
RECOVERY\_RECORD.md explains workspace maintenance, the intact previous
release, and reconstruction of the unpublished chart sources. Lost logs
are not represented as recovered.

\hypertarget{cont6:paper-ix-weighted-nonlinear-lifting}{%
\subsection{Chapter IX: weighted nonlinear
lifting}\label{cont6:paper-ix-weighted-nonlinear-lifting}}

\hypertarget{cont6:the-theorem}{%
\subsubsection{The theorem}\label{cont6:the-theorem}}

Let \(p\) be prime, \(R=\mathbb Z_p\), \(K=\mathbb Q_p\), and
\(d,s,r\in\mathbb N\). Let
\(F=(F_1,\ldots,F_d)\in R[X_1,\ldots,X_d]^d\), \(a\in R^d\), and let \[
J=DF(a)=\bigl(\partial_jF_i(a)\bigr)_{i,j}.
\] Assume \[
J\in\operatorname{GL}_d(K),\qquad p^sJ^{-1}\in M_d(R),\qquad
F(a)\in p^rR^d,\qquad r>2s.
\] Then \[
\boxed{\exists!\,b\in a+p^{s+1}R^d:\ F(b)=0,}
\qquad
\boxed{b-a\in p^{r-s}R^d.}
\] All vector divisibility assertions are coordinatewise. The uniqueness
predicate is the entire ball of depth \(s+1\): competitors need not be
assumed to have depth \(r-s\). The theorem permits arbitrary integral
p-adic coefficients and centers, every prime including two, and
dimension zero. It retains the supplied inverse-denominator exponent
\(s\); there is no replacement by a determinant valuation. Integrality
of \(p^sJ^{-1}\) is expressed in Lean by existence of an integral matrix
with exactly that image under the coefficient inclusion
\(R\hookrightarrow K\).

\hypertarget{cont6:proof}{%
\subsubsection{Proof}\label{cont6:proof}}

Write \(A=p^sJ^{-1}\in M_d(R)\). Then \[
AJ=JA=p^sI.
\] The formal matrix bridge proves the equivalence with the stated
fraction-field hypotheses, including nonsingularity. In particular,
multiplication by \(A\) is injective, since multiplication by \(J\)
sends \(Av=0\) to \(p^sv=0\).

Put \(G=AF\), \(t=p^{s+1}\), and \(k=r-(2s+1)\). Taylor expansion of the
actual preconditioned polynomials gives one polynomial system \(Q\) such
that \[
G(a+tu)=G(a)+p^{2s+1}u+p^{2s+2}Q(u),\qquad Q(0)=0.
\] This expansion is proved first over every commutative ring. Its
remainder is a polynomial valid for every \(u\), not a pointwise
quotient depending on \(u\). Since \(G(a)\in p^rR^d\), it can be written
\(p^{2s+1}C\) with \(C\in p^kR^d\). Thus \[
G(a+tu)=p^{2s+1}\bigl(C+u+pQ(u)\bigr).
\]

Integral polynomial evaluation is 1-Lipschitz in the supremum p-adic
metric. Consequently \(T(u)=-C-pQ(u)\) has contraction constant
\(|p|_p<1\) on the complete space \(R^d\), and has exactly one fixed
point \(u\).

For completeness, the sharp precision is additional to this fixed-point
existence. Inductively, if \(u\in p^nR^d\), then polynomial congruence
and \(Q(0)=0\) imply \(Q(u)\in p^nR^d\). For \(n<k\), the equation
\(u=-C-pQ(u)\) therefore gives \(u\in p^{n+1}R^d\). Starting at \(n=0\),
this proves \(u\in p^kR^d\). Hence \(b=a+tu\) satisfies \[
b-a\in p^{s+1+k}R^d=p^{r-s}R^d.
\] The normalization and injectivity of \(A\) give \(F(b)=0\).

Every other \(z\in a+p^{s+1}R^d\) has the form \(a+tw\). If \(F(z)=0\),
the same normalization, with cancellation of the nonzero scalar
\(p^{2s+1}\), gives \(w=T(w)\). Fixed-point uniqueness implies \(w=u\),
hence \(z=b\).

\hypertarget{cont6:boundary-and-sharpness}{%
\subsubsection{Boundary and sharpness}\label{cont6:boundary-and-sharpness}}

The strict inequality cannot in general be weakened to \(r\ge2s\): for
\(p=2\), \(F(X)=X^2+3\), and \(a=1\), we have \(F(a)=4\), \(DF(a)=2\),
and may take \(s=1,r=2\). Yet no root exists in \(\mathbb Z_2\), since
there is none modulo eight. The auxiliary Lean controls prove this
obstruction and the exact residual depth.

The displacement exponent cannot uniformly be increased by one. For
\(F(X)=p^sX+p^r\) at \(a=0\), with \(r>2s\), its root is \(b=-p^{r-s}\),
which is not divisible by \(p^{r-s+1}\). The controls prove the linear
root and the absence of that extra digit for all primes.

\hypertarget{cont7:paper-ix-weighted-first-order-root-motion}{%
\subsection{Chapter IX: weighted first-order root
motion}\label{cont7:paper-ix-weighted-first-order-root-motion}}

\hypertarget{cont7:the-theorem}{%
\subsubsection{The theorem}\label{cont7:the-theorem}}

Let \(p\) be prime, \(R=\mathbb Z_p\), \(K=\mathbb Q_p\), and
\(d,s,M\in\mathbb N\). Let \(F,H\in R[X_1,\ldots,X_d]^d\), let
\(a\in R^d\) satisfy \(F(a)=0\), and set \[
J=DF(a),\qquad G=F+p^M H.
\] Assume \[
J\in\operatorname{GL}_d(K),\qquad p^sJ^{-1}\in M_d(R),\qquad M>2s.
\] Then there is exactly one root \(b\) of \(G\) in \(a+p^{s+1}R^d\). It
satisfies \[
b-a\in p^{M-s}R^d,
\] and there is \(\varepsilon\in p^{2M-3s}R^d\) such that, in \(K^d\),
\[
\boxed{b-a=-p^M J^{-1}H(a)+\varepsilon.}
\] The inverse in this formula belongs to the original Jacobian
\(DF(a)\). The theorem constructs the perturbed root from these
hypotheses and proves uniqueness throughout the full ball of depth
\(s+1\). It does not require an additional Jacobian hypothesis on \(G\).
Coefficients and centers are arbitrary integral p-adic elements. Every
prime and finite square dimension are included, including \(p=2\) and
dimension zero.

The error depth exceeds the displacement depth by \[
(2M-3s)-(M-s)=M-2s\ge1.
\] Thus the first-order term determines strictly more digits than the
displacement bound alone. These are guaranteed divisibility bounds;
equality of valuations is not asserted for every system.

\hypertarget{cont7:root-construction-with-the-original-inverse}{%
\subsubsection{Root construction with the original
inverse}\label{cont7:root-construction-with-the-original-inverse}}

Write \(A=p^sJ^{-1}\in M_d(R)\), so \(AJ=JA=p^sI\). Let \(t=p^{s+1}\)
and \(k=M-(2s+1)\ge0\). The actual polynomial Taylor and difference
identities give polynomial systems \(Q_F,Q_H\), both vanishing at zero,
such that \[
(AF)(a+tu)=p^{2s+1}u+p^{2s+2}Q_F(u),
\] \[
(AH)(a+tu)=(AH)(a)+tQ_H(u).
\] Consequently \[
(AG)(a+tu)=p^{2s+1}\bigl(C+u+pQ(u)\bigr),
\] where \[
C=p^kAH(a),\qquad
Q=Q_F+p^{M-s-1}Q_H,\qquad Q(0)=0.
\] All exponents are nonnegative under \(M>2s\). Actual integral
polynomial evaluation is 1-Lipschitz, so \(u\mapsto-C-pQ(u)\) is a
strict contraction on the complete space \(R^d\). Its unique fixed point
inherits the depth \(k\) of \(C\), because \(Q(0)=0\). Thus \(b=a+tu\)
has displacement depth \(s+1+k=M-s\).

The certificate \(JA=p^sI\) makes multiplication by \(A\) injective, so
the normalized zero is an actual root of \(G\). Every root in \(a+tR^d\)
has the form \(a+tw\) and satisfies the same fixed-point equation. This
proves uniqueness on the entire stated ball using the original matrix
certificate.

\hypertarget{cont7:the-first-order-error}{%
\subsubsection{The first-order error}\label{cont7:the-first-order-error}}

Put \(\delta=b-a\) and \(L=M-s\). The displacement form of Taylor's
theorem, proved over any commutative ring, gives \[
E_F=(AF)(b)-(AF)(a)-p^s\delta\in p^{2L}R^d.
\] Polynomial congruence gives \[
E_H=(AH)(b)-(AH)(a)\in p^LR^d.
\] Since \(F(a)=0\) and \(G(b)=0\), \[
p^s\delta+p^MAH(a)=-E_F-p^ME_H.
\] We have \(M+L\ge2L\), so the right side lies in \(p^{2L}R^d\).
Cancelling the nonzero scalar \(p^s\) proves the integral estimate \[
\delta+p^{M-s}AH(a)\in p^{2M-3s}R^d.
\] Define this integral vector to be \(\varepsilon\). Under
\(R\hookrightarrow K\), the exact identity \(p^{M-s}A=p^MJ^{-1}\) gives
the displayed first-order formula. The Lean bridge checks matrix
evaluation and all coefficient inclusions explicitly.

\hypertarget{cont7:sharpness-and-the-strict-threshold}{%
\subsubsection{Sharpness and the strict
threshold}\label{cont7:sharpness-and-the-strict-threshold}}

For \(p=2\), take \(F(X)=X+X^2\), \(H(X)=1\), \(a=0\), \(s=0\), \(M=1\).
There is a unique root \(b\in2\mathbb Z_2\) of \(X+X^2+2\). Its equation
gives \(b+2=-b^2\). It also forces \(b\equiv2\pmod4\). Therefore \[
v_2(b+2)=2=2M-3s.
\] The accompanying Lean controls show that the error is divisible by
four and is not divisible by eight. This refutes a uniform replacement
of \(2M-3s\) by \(2M-3s+1\), even for a one-variable unit Jacobian.

The prior strict-threshold example also applies to perturbations: take
\(F(X)=X^2-1\), \(H(X)=1\), \(a=1\), \(p=2\), \(s=1\), \(M=2\). Then
\(F(a)=0\), \(DF(a)=2\), and \(F+2^2H=X^2+3\) has no root modulo eight,
hence none in \(\mathbb Z_2\). The new explanatory identification uses
the previously checked no-root polynomial. The strict inequality
\(M>2s\) cannot uniformly be replaced by \(M\ge2s\).

\hypertarget{cont8:paper-ix-prescribed-parameter-polynomial-lifting}{%
\subsection{Chapter IX: prescribed-parameter polynomial
lifting}\label{cont8:paper-ix-prescribed-parameter-polynomial-lifting}}

\hypertarget{cont8:setup-and-exact-lift-set}{%
\subsubsection{Setup and exact lift
set}\label{cont8:setup-and-exact-lift-set}}

Let \(q>0\), \(L\ge1\), and \(m,k,e\ge0\). Write
\(S=\mathbb Z/q\mathbb Z\). Let \[
F\in\mathbb Z[Y_1,\ldots,Y_m,T_1,\ldots,T_k]^e,
\qquad x=(y_0,t_0)\in\mathbb Z^{m+k},
\qquad F(x)\equiv0\pmod{q^L}.
\] The actual lower root supplies the unique integral residual \[
r=q^{-L}F(x)\in\mathbb Z^e,\qquad \bar r=r\bmod q\in S^e.
\] Let the actual reduced Jacobian blocks be \[
J=(D_YF)(x)\bmod q:S^m\to S^e,
\qquad K=(D_TF)(x)\bmod q:S^k\to S^e.
\] For a prescribed parameter digit \(v\in S^k\), choose its canonical
integral representative \(\widehat v\). Define \(\mathcal L_v\) to be
the set of actual residue vectors
\(z\in(\mathbb Z/q^{L+1}\mathbb Z)^{m+k}\) satisfying \[
z\equiv x\pmod{q^L},\qquad F(z)\equiv0\pmod{q^{L+1}},\qquad
z_T\equiv t_0+q^L\widehat v\pmod{q^{L+1}}.
\] The formal predicate uses canonical integer representatives for
evaluation and the lower congruence. The previously checked polynomial
congruence theorem and digit-to-residue equivalence establish the actual
residue-class interpretation. Changing an integral representative of a
parameter digit by a multiple of \(q\) does not change its prescribed
residue modulo \(q^{L+1}\).

\hypertarget{cont8:the-obstruction-theorem}{%
\subsubsection{The obstruction theorem}\label{cont8:the-obstruction-theorem}}

Define the finite solved-block cokernel and its parameter-dependent
class by \[
\mathcal C_x=S^e/\operatorname{im}J,
\qquad o_x(v)=[\bar r+Kv]\in\mathcal C_x.
\] Then \[
\boxed{\mathcal L_v\ne\varnothing\iff o_x(v)=0.}
\] More precisely, the checked digit-to-lift equivalence is \[
\boxed{\mathcal L_v\cong\{u\in S^m:Ju=-\bar r-Kv\}.}
\] It sends a solved digit \(u\) to the actual next residue represented
by \(x+q^L(\widehat u,\widehat v)\). The obstruction depends on the
solved block \(J\), not on the range of the full map \((J\ K)\).

Every nonempty \(\mathcal L_v\) has a checked free and transitive action
of \(\ker J\). The action is defined without choosing a base lift.
Choosing one valid solved digit then gives an explicit bijection
\(\mathcal L_v\cong\ker J\). In particular, for every positive modulus,
\[
o_x(v)=0\quad\Longrightarrow\quad\#\mathcal L_v=\#\ker J.
\] These results permit composite \(q\), the degenerate modulus \(q=1\),
and zero dimensions. Neither \(J\) nor the full reduced Jacobian is
assumed surjective.

\hypertarget{cont8:exact-prime-field-counts}{%
\subsubsection{Exact prime-field
counts}\label{cont8:exact-prime-field-counts}}

When \(q=p\) is prime, put \(\rho=\operatorname{rank}_{\mathbb F_p}J\).
Then \[
\boxed{\#\mathcal L_v=
\begin{cases}
p^{m-\rho},&o_x(v)=0,\\
0,&o_x(v)\ne0,
\end{cases}}
\qquad
\boxed{\#\mathcal C_x=p^{e-\rho}.}
\] The count is of distinct actual residue classes at level \(L+1\). The
solved dimension \(m\), rather than the total number \(m+k\) of
variables, appears in the exponent because the parameter digit is fixed.
These rank formulas are asserted only over prime fields.

The main wrapper assumes only the original lower polynomial congruence.
It constructs one residual \(r\) and supplies the criterion and count
uniformly for every \(v\), together with the exact obstruction
cardinality.

\hypertarget{cont8:proof-and-dependency-chain}{%
\subsubsection{Proof and dependency
chain}\label{cont8:proof-and-dependency-chain}}

The actual derivative calculation gives \[
(DF(x)\bmod q)(u,v)=Ju+Kv.
\] Taylor expansion at depth \(L\ge1\) makes the quadratic remainder
divisible by \(q^{L+1}\). Therefore the digit \((u,v)\) is a polynomial
lifting correction exactly when \[
\bar r+Ju+Kv=0.
\] The quotient class \([\bar r+Kv]\) vanishes exactly when this
equation admits \(u\). This uses additive inverses in the image of
\(J\), so the sign agrees with the actual Taylor equation.

The earlier unrestricted correction-digit equivalence is then restricted
to the predicate that the parameter digit equals \(v\). Taking free
coordinates commutes with forming digit residues, and digit residues are
injective in the digit. Consequently the restricted equivalence enforces
exactly the actual parameter residue condition on \(\mathcal L_v\).

Translation by \(\ker J\) acts freely and transitively on the
solved-digit fiber. Transporting this explicit action through the
residue equivalence yields the torsor on \(\mathcal L_v\). Rank-nullity
over \(\mathbb F_p\) then gives the two displayed cardinalities. No
numerical enumeration is used to establish the general theorem.

\hypertarget{cont8:why-the-solved-block-matters}{%
\subsubsection{Why the solved block
matters}\label{cont8:why-the-solved-block-matters}}

Take \(F(Y,T)=T\), \(q=2\), \(L=1\), and \(x=(0,0)\). The full Jacobian
\((0\ 1)\) is surjective, but \(J=0\). Prescribing \(v=1\) forces
\(T\equiv2\pmod4\), which cannot satisfy \(F=0\pmod4\). Prescribing
\(v=0\) gives exactly two lifts, corresponding to \(Y\equiv0,2\pmod4\).

The Lean controls check this actual polynomial example. Thus
full-Jacobian surjectivity alone does not make every prescribed
coordinate choice admissible. The automatic coordinate selection in the
earlier chart theorem is compatible with this obstruction.

\hypertarget{cont9:paper-ix-compatible-towers-with-prescribed-parameters}{%
\subsection{Chapter IX: compatible towers with prescribed
parameters}\label{cont9:paper-ix-compatible-towers-with-prescribed-parameters}}

The construction gives actual compatible polynomial towers from initial
solved-block surjectivity, for every compatible integer parameter path.
For a prime base, it constructs such a path for every p-adic parameter
in the initial ball and proves an exact polynomial endpoint with all
retained precisions. The formal companion identifies the six supporting modules and their
recorded compiler and axiom-audit results.

\hypertarget{cont9:statement-and-notation}{%
\subsubsection{Statement and notation}\label{cont9:statement-and-notation}}

Let \(q>0\), \(L\ge1\), and \[
F\in\mathbb Z[Y_1,\ldots,Y_m,T_1,\ldots,T_k]^e,
\qquad x=(y_0,t_0)\in\mathbb Z^{m+k},
\qquad F(x)\in q^L\mathbb Z^e.
\] Write \(S_q=\mathbb Z/q\mathbb Z\), and let the actual reduced
Jacobian blocks be \[
J_x=D_YF(x)\bmod q:S_q^m\longrightarrow S_q^e,
\qquad K_x=D_TF(x)\bmod q:S_q^k\longrightarrow S_q^e.
\] Assume \(J_x\) is surjective. A vector divisibility condition below
is coordinatewise.

\textbf{Integer prescribed-path theorem.} For every supplied path
\(t:\mathbb N\to\mathbb Z^k\) satisfying \[
t(0)=t_0,\qquad t(n+1)-t(n)\in q^{L+n}\mathbb Z^k,
\] there is a single tower \(z:\mathbb N\to\mathbb Z^{m+k}\) with \[
z(0)=x,\qquad \operatorname{free}(z(n))=t(n),
\] \[
F(z(n))\in q^{L+n}\mathbb Z^e,\qquad
z(n+1)-z(n)\in q^{L+n}\mathbb Z^{m+k}.
\] The parameter equality is exact in integers. For \(a\le b\), \[
z(b)-z(a)\in q^{L+a}\mathbb Z^{m+k},\qquad
DF(z(n))\bmod q=DF(x)\bmod q.
\] Consequently \(J_{z(n)}=J_x\) at every level. No later surjectivity
or per-step solvability is assumed. This theorem includes positive
composite bases and \(q=1\).

\textbf{Arbitrary p-adic target theorem.} Now let \(q=p\) be prime. For
every \[
\tau\in t_0+p^L\mathbb Z_p^k,
\] there are a compatible integer tower \(z\) as above and
\(b\in\mathbb Z_p^{m+k}\) such that \[
z(n)\longrightarrow b,\qquad F(b)=0,\qquad
\operatorname{free}(b)=\tau,
\qquad b-z(n)\in p^{L+n}\mathbb Z_p^{m+k}\quad(n\ge0).
\] A compatible integer parameter path converging to \(\tau\) is
constructed in the proof, with its initial representative fixed exactly.
A separate theorem accepts every supplied compatible integer path with
limit \(\tau\) and retains its exact coordinates in the constructed
tower.

\textbf{Exact branching.} At any stage \(n\), put \(N=L+n\). For every
parameter digit \(v\in\mathbb F_p^k\), let \[
\mathcal L_{n,v}=\left\{\bar w\in(\mathbb Z/p^{N+1}\mathbb Z)^{m+k}:
\begin{array}{l}
\bar w\equiv z(n)\pmod{p^N},\\
F(\bar w)=0,\\
\operatorname{free}(\bar w)=t(n)+p^Nv\pmod{p^{N+1}}
\end{array}\right\}.
\] Here \(p^Nv\) means any integer representative of \(v\), multiplied
by \(p^N\); the residue class does not depend on the representative.
Then \[
\#\mathcal L_{n,v}=p^{m-e}\qquad\text{for every }n\text{ and every }v.
\] Surjectivity implies \(e\le m\). This is a count of distinct actual
next-level residue classes. It does not count integer representatives or
entire infinite towers.

\textbf{Square solved-block uniqueness.} If \(m=e\), every prescribed
\(\tau\) has exactly one root in the original ball: \[
\exists!\,b\in x+p^L\mathbb Z_p^{e+k}:\quad
F(b)=0,\qquad\operatorname{free}(b)=\tau.
\] Thus all valid integer parameter approximations and all lifting
choices with the same target have the same endpoint. This does not
assert uniqueness of the integer representatives in the tower. All prime
statements include \(p=2\).

\hypertarget{cont9:proof}{%
\subsubsection{Proof}\label{cont9:proof}}

At depth \(N\ge1\), suppose \(F(a)=q^Nr\) and a desired next integer
parameter target satisfies \(t'-\operatorname{free}(a)=q^Nw\), with
\(w\in\mathbb Z^k\). Surjectivity supplies a solved digit \(u\in S_q^m\)
satisfying \[
J_a u=-\bar r-K_a\bar w.
\] Choose an integer representative \(\widetilde u\), and set \[
a'=a+q^N(\widetilde u,w).
\] The parameter quotient is the supplied integer \(w\). Therefore
\(\operatorname{free}(a')=t'\) exactly. The checked polynomial Taylor
identity gives \[
F(a')\equiv F(a)+q^NDF(a)(\widetilde u,w)
\equiv0\pmod{q^{N+1}}.
\] Its quadratic remainder is divisible by \(q^{2N}\), and
\(2N\ge N+1\). Also \(a'\equiv a\pmod q\), so evaluation of every
polynomial partial derivative gives the same reduced Jacobian. Induction
on the level, carrying the root, exact parameters, and fixed solved
Jacobian as its state, constructs the entire integer tower. The
existence proof uses ordinary classical choice; it does not supply an
optimized computational lifting algorithm.

For a p-adic target \(\tau\), use \(t(0)=t_0\) and, for \(n\ge1\), the
canonical integer representatives of \(\tau\bmod p^{L+n}\). The initial
congruence hypothesis and the truncation precision imply \[
\tau-t(n)\in p^{L+n}\mathbb Z_p^k.
\] Subtracting successive precision relations gives compatibility in
\(\mathbb Z_p\); the checked faithful integer-divisibility theorem gives
compatibility in \(\mathbb Z\). These precision relations also give
convergence to \(\tau\).

The compatible vector-tower completion theorem constructs \(b\), proves
convergence, and retains every stage precision. Continuous polynomial
evaluation and the increasing root precision give \(F(b)=0\). Each
parameter coordinate has both the corresponding coordinate of \(b\) and
\(\tau\) as its limit, so uniqueness of limits gives exact parameter
equality.

For every proposed next parameter digit, the previously checked
obstruction is \([\bar r+K_av]\in S_p^e/\operatorname{im}J_a\). The
fixed solved Jacobian is surjective, so this class vanishes for every
digit at every stage. The actual lift fiber is a torsor for
\(\ker J_a\). Rank-nullity gives its cardinality \(p^{m-e}\).

When \(m=e\), the new solved-Jacobian definition is definitionally the
existing square block used by the checked continuous parameter chart.
Every root in the initial ball with the given parameters is that chart
point. Applying this to any two roots proves equality, including
endpoints obtained by different constructions.

\hypertarget{cont9:scope-controls}{%
\subsubsection{Scope controls}\label{cont9:scope-controls}}

The inherited checked example \(F(Y,T)=T\), \(p=2\), \(x=0\), \(L=1\)
has a surjective full Jacobian but zero solved block. Prescribing
parameter digit \(1\) forces \(T=2\pmod4\), so no next root exists.
Full-Jacobian surjectivity cannot replace the hypothesis on the
prescribed solved block.

The new checked example \(F(Y_1,Y_2,T)=Y_1\) has a surjective solved
block \((1,0)\). Nevertheless \((0,0,0)\) and \((0,2,0)\) are distinct
exact roots in the same modulo-two ball with identical parameters.
Uniqueness therefore requires the square solved case. The corresponding
negative control deliberately attempts the stronger false conclusion and
must be rejected by Lean.

Depth \(L\ge1\) is essential for the Taylor step used here. The general
singular case can have a nonzero next obstruction after an earlier zero
obstruction. Neither the finite one-step criterion nor this smooth
theorem establishes unrestricted singular persistence.

\hypertarget{cont10:paper-ix-simultaneous-coefficient-and-parameter-deformation}{%
\subsection{Chapter IX: simultaneous coefficient and parameter
deformation}\label{cont10:paper-ix-simultaneous-coefficient-and-parameter-deformation}}

The construction combines coefficient perturbation with prescribed parameter
motion. The original solved Jacobian controls existence, local
uniqueness, displacement and the first-order error. The formal companion identifies the four supporting modules and their
recorded verification scope.

\hypertarget{cont10:the-theorem}{%
\subsubsection{The theorem}\label{cont10:the-theorem}}

Let \(p\) be any prime, \(R=\mathbb Z_p\), and \(Q=\mathbb Q_p\). Let \[
F,H\in R[Y_1,\ldots,Y_e,T_1,\ldots,T_k]^e,
\qquad a=(y_0,t_0)\in R^{e+k},\qquad F(a)=0.
\] Here \(e,k\ge0\) are arbitrary finite dimensions. There are as many
solved coordinates as equations. Define the actual derivative blocks at
the original root: \[
J=D_YF(a)\in M_e(R),\qquad K=D_TF(a)\in M_{e\times k}(R).
\] Suppose, for natural numbers \(s,M\), that \[
J\in\operatorname{GL}_e(Q),\qquad p^sJ^{-1}\in M_e(R),\qquad M>2s.
\] For every \(v\in R^k\), set \[
\tau=t_0+p^Mv,\qquad G=F+p^MH.
\] Then there is exactly one \(y\in y_0+p^{s+1}R^e\) with \[
G(y,\tau)=0.
\] It satisfies \[
y-y_0\in p^{M-s}R^e
\] and there is an integral vector \(\varepsilon\in p^{2M-3s}R^e\) such
that, in \(Q^e\), \[
\boxed{\displaystyle
 y-y_0=-p^MJ^{-1}\bigl(H(a)+Kv\bigr)+\varepsilon.
}
\] The same original exponent \(s\) appears throughout. No inverse or
rank assumption on the perturbed Jacobian is imposed. The statement
includes arbitrary integral p-adic coefficients and centers, the dyadic
prime, and zero dimensions.

The full-point theorem constructs \(b=(y,\tau)\) and also proves \[
b-a\in p^{M-s}R^{e+k}\cap p^{s+1}R^{e+k}.
\] Parameter equality \(b_T=\tau\) is exact. Uniqueness concerns every
root in the stated solved-coordinate ball with these parameters, not
only roots obtained by the construction. No global uniqueness is
claimed.

\hypertarget{cont10:actual-augmented-system}{%
\subsubsection{Actual augmented system}\label{cont10:actual-augmented-system}}

Define polynomial systems with \(e+k\) equations and variables by \[
\mathcal F(Y,T)=\bigl(F(Y,T),T-t_0\bigr),\qquad
\mathcal H(Y,T)=\bigl(H(Y,T),-v\bigr).
\] Both are actual polynomials over \(R\). Direct evaluation proves \[
\mathcal F(a)=0,
\qquad
(\mathcal F+p^M\mathcal H)(b)=0
\iff F(b)+p^MH(b)=0\ \text{and}\ b_T=t_0+p^Mv.
\] Actual polynomial partial differentiation gives \[
D\mathcal F(a)=B=\begin{pmatrix}J&K\\0&I_k\end{pmatrix}.
\] These identities are proved from evaluation and partial derivatives,
rather than supplied as additional hypotheses.

Choose the integral matrix \(A=p^sJ^{-1}\). The existing
inverse-certificate equivalence gives \[
AJ=JA=p^sI_e.
\] Put \(c=p^s\) and define an integral block matrix \[
C=\begin{pmatrix}A&-AK\\0&cI_k\end{pmatrix}.
\] Then \[
CB=\begin{pmatrix}AJ&AK-AK\\0&cI_k\end{pmatrix}=cI_{e+k},
\] \[
BC=\begin{pmatrix}JA&-JAK+cK\\0&cI_k\end{pmatrix}=cI_{e+k}.
\] The second equality uses \(JA=cI_e\). The generic block proof works
over any commutative ring. The p-adic specialization also explicitly
proves that \(B\) is invertible over \(Q\) and that \(p^sB^{-1}\) is
integral. There is no denominator loss from appending the parameter
equations.

\hypertarget{cont10:construction-error-and-uniqueness}{%
\subsubsection{Construction, error and
uniqueness}\label{cont10:construction-error-and-uniqueness}}

Apply the previously checked weighted coefficient-perturbation
construction to \(\mathcal F,\mathcal H,a,C\). Its hypotheses follow
from the actual augmented identities and the two-sided certificate. It
constructs \(b\), proves the full displacement bound \(p^{M-s}\), and
proves uniqueness throughout \(a+p^{s+1}R^{e+k}\). The augmented root
equivalence gives the required original equations and exact parameter
values.

The checked integral first-order estimate gives \[
b-a+p^{M-s}C\mathcal H(a)\in p^{2M-3s}R^{e+k}.
\] The decisive block action is \[
C\binom{H(a)}{-v}
=
\binom{A\bigl(H(a)+Kv\bigr)}{-p^sv}.
\] Projecting the solved coordinates therefore gives \[
y-y_0+p^{M-s}A\bigl(H(a)+Kv\bigr)
\in p^{2M-3s}R^e.
\] The existing rational-conversion theorem turns \(p^{M-s}A\) into
\(p^MJ^{-1}\) and retains an integral error vector. This proves the
displayed first-order formula with the original derivative blocks.

For any competing solved-coordinate root \(z\in y_0+p^{s+1}R^e\), the
joined point \((z,\tau)\) lies in the augmented uniqueness ball:
\(M>2s\) implies \(M\ge s+1\), so \(\tau-t_0=p^Mv\in p^{s+1}R^k\).
Full-point uniqueness gives \((z,\tau)=b\), hence \(z=y\). This is the
additional comparison needed to obtain uniqueness in exactly the ball
stated in the solved-coordinate theorem.

\hypertarget{cont10:checked-controls-and-precision-boundary}{%
\subsubsection{Checked controls and precision
boundary}\label{cont10:checked-controls-and-precision-boundary}}

For \(p=2\), \(s=0\), \(F(Y,T)=Y+T\), \(H=1\), \(a=0\), and \(v=1\),
both original blocks satisfy \(J=K=1\). At \(M=1\), the exact parameter
is \(\tau=2\), and the unique solved value is \(y=-4\). The combined
correction is \(-2(1+1)=-4\), with zero error.

The sign controls use \(M=2\): the actual augmented root is
\((y,\tau)=(-8,4)\), and the required error ideal is \(16\mathbb Z_2\).
Omitting \(Kv\) predicts \(-4\), with error \(-4\notin16\mathbb Z_2\).
Reversing its sign predicts \(0\), with error \(-8\notin16\mathbb Z_2\).
Both failures are checked. At binary depth \(M=1\), the reversed-sign
difference would lie in the allowed error ideal, so that shallower
example alone would not distinguish the sign. This control-design
limitation was identified and corrected before the release gate.

For parameter-only motion, use \[
F(Y,T)=Y+Y^2+T,\quad H=0,\quad a=0,\quad v=1,\quad p=2,\ s=0,\ M=1.
\] A checked transfer from the prior scalar quadratic root constructs an
even \(y\) with \[
y+y^2+2=0,\qquad 4\mid y+2,\qquad 8\nmid y+2.
\] Thus the error exponent \(2M-3s=2\) is attained even when only
parameters move. No uniform extra error digit can be added to the
general theorem. This example establishes that limitation; it does not
assert optimality at every choice of \(p,s,M\).

The inherited controls continue to record the necessity of the strict
weighted threshold and other earlier scope boundaries. General
rationally rank-deficient systems are outside the inverse hypothesis.

\hypertarget{cont11:paper-ix-analytic-regularity-of-the-local-parameter-solution}{%
\subsection{Chapter IX: analytic regularity of the local parameter
solution}\label{cont11:paper-ix-analytic-regularity-of-the-local-parameter-solution}}

The argument establishes both analyticity and differentiability of the
actual local polynomial solution map at a root with an invertible solved
Jacobian. The statements apply to every prime, including \(p=2\). They
also hold over any complete nontrivially normed field.

Let \(\mathbb K\) be such a field, let \(e,k\in\mathbb N\), and let \[
F\in\mathbb K[Y_1,\ldots,Y_e,T_1,\ldots,T_k]^e,
\qquad a=(y_0,t_0)\in\mathbb K^e\times\mathbb K^k.
\] Evaluation of \(F\) means coordinatewise evaluation of its actual
polynomials. Define the actual derivative blocks \[
J=D_YF(a)\in M_e(\mathbb K),
\qquad K=D_TF(a)\in M_{e\times k}(\mathbb K),
\] where each matrix entry is the corresponding formal partial
derivative evaluated at \(a\). The hypotheses are exactly \[
F(a)=0,
\qquad J\in\operatorname{GL}_e(\mathbb K).
\] Neither analyticity of a proposed solution nor a derivative formula
is assumed.

\textbf{Local analytic graph theorem.} There exist a map \[
\psi:\mathbb K^k\longrightarrow\mathbb K^e
\] and open sets \(U\subseteq\mathbb K^k\), \(V\subseteq\mathbb K^e\)
such that \[
t_0\in U,\qquad y_0\in V,\qquad \psi(t_0)=y_0,
\qquad \psi(U)\subseteq V.
\] The map is analytic at every point of \(U\), and its graph describes
all roots in the product neighborhood: \[
\forall t\in U\;\forall y\in V,
\qquad F(y,t)=0\quad\Longleftrightarrow\quad y=\psi(t).
\] At the base parameter it has the strict Fréchet derivative \[
\boxed{D\psi(t_0)=-J^{-1}K.}
\] Consequently it is differentiable on \(U\). The map is defined on the
ambient parameter space, but its root, analyticity and uniqueness
conclusions are restricted to the specified neighborhood. Uniqueness is
among solved coordinates in \(V\).

The main Lean theorem is
\mcode{Komorebi.PolynomialAnalyticImplicit.exists_analytic_implicit}.
The matrix in its \mcode{HasStrictFDerivAt} conclusion is the
continuous linear map represented by the displayed matrix.
\mcode{differentiableOn_of_analytic_graph} records the
differentiability consequence.

\textbf{Proof chain.} The scalar polynomial derivative bridge is proved
by induction on actual multivariate polynomials: constants, sums, and
multiplication by a variable. For a scalar polynomial \(P\), its strict
derivative is \[
v\longmapsto\sum_j \frac{\partial P}{\partial x_j}(a)v_j.
\] Assembling the coordinates gives the strict derivative of the
evaluated system as its actual Jacobian matrix. Polynomial analyticity
is supplied by the pinned Mathlib theorem for multivariate polynomial
evaluation, followed by finite coordinate assembly.

Consider the actual augmented polynomial map \[
\Phi(y,t)=(F(y,t),t).
\] Its derivative and inverse derivative at \(a\) are \[
D\Phi(a)=
\begin{pmatrix}J&K\\0&I_k\end{pmatrix},
\qquad
D\Phi(a)^{-1}=
\begin{pmatrix}J^{-1}&-J^{-1}K\\0&I_k\end{pmatrix}.
\] The proof constructs the continuous linear equivalence from these
explicit matrices and verifies both inverse identities. The strict
inverse theorem gives a local inverse \(g\) of \(\Phi\), and the
analytic inverse result shows that this same inverse is analytic near
\(\Phi(a)=(0,t_0)\). Set \[
\psi(t)=\pi_Y g(0,t).
\] The two inverse identities and a sufficiently small product
neighborhood give the exact local graph statement. Composing the
derivative of \(g\) with parameter insertion and solved-coordinate
projection yields \(-J^{-1}K\). The chosen parameter neighborhood is
also restricted so that the inverse identities, analyticity and
inclusion \(\psi(U)\subseteq V\) hold there simultaneously.

\textbf{Transfer to a supplied continuous solution.} Suppose
\(g:\mathbb K^k\to\mathbb K^e\) is continuous at \(t_0\),
\(g(t_0)=y_0\), and \[
F(g(t),t)=0
\] for all \(t\) in some neighborhood of \(t_0\). Under the same
invertibility hypothesis at \(a\), continuity places \(g(t)\) in the
uniqueness neighborhood \(V\) for all sufficiently close parameters.
Hence \(g\) and \(\psi\) agree on a neighborhood of \(t_0\). The theorem
\mcode{continuous_selection_regularity} transfers both analyticity
and the strict derivative formula to \(g\).

There is also an explicit open-set version. Let \(U\) be open, let \(g\)
be continuous on \(U\), and suppose for every \(t\in U\) that \[
F(g(t),t)=0,
\qquad
J_t:=D_YF(g(t),t)\in\operatorname{GL}_e(\mathbb K).
\] Writing \(K_t=D_TF(g(t),t)\), the theorem
\mcode{continuous_selection_analyticOnNhd} proves \[
g\text{ is analytic at every }t\in U,
\qquad Dg(t)=-J_t^{-1}K_t\quad(t\in U),
\] with a strict Fréchet derivative at each point. The per-point
invertibility hypothesis is part of this theorem. This is a transfer
statement for continuous local root selections satisfying its
hypotheses; it does not assert equality of arbitrary global branch
choices.

\textbf{Integral p-adic specialization.} For any prime \(p\), let \[
F\in\mathbb Z_p[Y,T]^e,
\qquad a\in\mathbb Z_p^{e+k},
\qquad F(a)=0.
\] Let \(\iota:\mathbb Z_p\hookrightarrow\mathbb Q_p\) be the
coefficient inclusion, applied entrywise to points and matrices. Assume
only that \[
\iota(D_YF(a))\in\operatorname{GL}_e(\mathbb Q_p).
\] The theorem
\mcode{PolynomialRegularityIntegral.exists_analytic_integral_implicit}
constructs the preceding analytic graph in the ambient spaces
\(\mathbb Q_p^k\) and \(\mathbb Q_p^e\). Its root equation evaluates the
original integral polynomial coefficients using \(\iota\), and its base
derivative is \[
-\iota(D_YF(a))^{-1}\,\iota(D_TF(a)).
\] The source proves that polynomial evaluation, formal partial
differentiation and both Jacobian blocks commute with this inclusion. No
invertibility modulo \(p\) or integral inverse is required. This permits
a solved Jacobian that is singular modulo \(p\) while remaining
invertible over \(\mathbb Q_p\). The new specialization is an
ambient-field theorem; it does not treat \(\mathbb Z_p\) as a vector
space over \(\mathbb Q_p\) or assert a new explicit integral radius.

\begin{longtable}[]{@{}
  >{\raggedright\arraybackslash}p{(\columnwidth - 2\tabcolsep) * \real{0.5000}}
  >{\raggedright\arraybackslash}p{(\columnwidth - 2\tabcolsep) * \real{0.5000}}@{}}
\toprule\noalign{}
\begin{minipage}[b]{\linewidth}\raggedright
Source module
\end{minipage} & \begin{minipage}[b]{\linewidth}\raggedright
Mathematical role
\end{minipage} \\
\midrule\noalign{}
\endhead
\bottomrule\noalign{}
\endlastfoot
\mcode{PolynomialAnalyticEvaluation.lean} & Actual polynomial
analyticity and strict derivative equals the evaluated formal
Jacobian. \\
\mcode{PolynomialAnalyticInverse.lean} & Analyticity of the same local
inverse supplied by an invertible strict derivative. \\
\mcode{PolynomialAnalyticBlocks.lean} & Explicit augmented derivative
equivalence and its solved-coordinate restriction. \\
\mcode{PolynomialAnalyticImplicit.lean} & Open local root graph,
analyticity, base strict derivative and differentiability. \\
\mcode{PolynomialAnalyticSelection.lean} & Transfer to continuous root
selections, including the per-point open-set theorem. \\
\mcode{PolynomialRegularityIntegral.lean} & Integral coefficients and
centers included in the ambient p-adic field, with actual Jacobian
compatibility. \\
\end{longtable}

This block closes the local analytic regularity obligation stated in
\mcode{TARGET_CONTRACT.json}. It supplies no new uniform numerical
radius. The previous weighted lifting and joint deformation results
retain their own explicit radii and error estimates. It does not
establish a global analytic branch or a result at a point whose solved
Jacobian is singular over the fraction field. The proof uses established
analytic inverse machinery and makes no novelty claim for the general
implicit-function theorem.

The remaining Chapter IX work includes general singular persistence with
higher obstructions, native presentation invariance (13.2), global or
completed compatibility and descent (13.3), the outstanding sharper
arithmetic statements, and manuscript consolidation. None of those is
closed by this local regularity theorem. Release-level checks and audit
counts are recorded in the separate release verification and audit
files.


\clearpage\section{Restricted native progress and complete prose proofs}\label{app:native}
This appendix reproduces the mathematical parts of the 8 September 2026
native progress note. The polynomial obstruction, finite reconstruction,
rectangularity and topological arguments are prose proofs. The intrinsic
commutative-algebra obstruction is the stated classical Stacks specialization,
with its original source record and finiteness argument. None of these prose
arguments is promoted to a new Lean run or to a completion of unrestricted
Native 13.2 or 13.3.
\hypertarget{native:a-complete-polynomial-obstruction-theorem}{%
\subsection{2. A complete polynomial obstruction
theorem}\label{native:a-complete-polynomial-obstruction-theorem}}

\textbf{Scope.} Fix a prime \(p\), an integer \(L\ge1\), finite natural
numbers \(d,r\), and polynomials
\(F_1,\ldots,F_r\in\mathbb Z[X_1,\ldots,X_d]\). Let
\(\bar x\in(\mathbb Z/p^L\mathbb Z)^d\) satisfy every equation
\(F_i(\bar x)=0\). Choose any integer representative
\(x\in\mathbb Z^d\). Set

\[e_i=F_i(x)/p^L\pmod p,\qquad J_{ij}=\frac{\partial F_i}{\partial X_j}(x)\pmod p.\]

The divisions are exact integer divisions. Regard \(J\) as a linear map
\(\mathbb F_p^d\to\mathbb F_p^r\). Define

\[O_{F,\bar x}=\mathbb F_p^r/\operatorname{im}J,\qquad o_{F,\bar x}=[e].\]

\textbf{Theorem.} For this fixed polynomial presentation, the group and
distinguished class are independent of the chosen integer
representative. The class vanishes if and only if \(\bar x\) lifts to a
solution modulo \(p^{L+1}\). The obstruction group has
\(p^{r-\operatorname{rank}J}\) elements. When lifts exist, their set is
a torsor for \(\ker J\) and contains exactly
\(p^{d-\operatorname{rank}J}\) elements.

\textbf{Proof.} For every integer correction vector \(t\), expansion of
each monomial gives

\[F(x+p^Lt)\equiv F(x)+p^LJ(x)t\pmod {p^{L+1}}.\tag{1}\]

Indeed, each term involving at least two correction factors is divisible
by \(p^{2L}\), and \(2L\ge L+1\). Constant and linear correction terms
give the two displayed summands. Adding the monomial expansions proves
(1) for every polynomial, with no denominator or analytic convergence
issue.

Every element reducing to \(\bar x\) is uniquely expressible modulo
\(p^{L+1}\) as \(x+p^Lt\) with \(t\in\mathbb F_p^d\). By (1), it solves
the equations exactly when

\[e+Jt=0\quad\text{in }\mathbb F_p^r.\tag{2}\]

Equation (2) is soluble exactly when \(e\in\operatorname{im}J\): the
sign makes no difference because the image is a subspace. This is
equivalent to \([e]=0\).

Another representative has the form \(x'=x+p^Ls\). Since \(L\ge1\), its
Jacobian is the same modulo \(p\). Dividing (1) by \(p^L\) shows that
its residual digit is \(e'=e+Js\), so it gives exactly the same quotient
class. If \(t_0\) solves (2), all its solutions are \(t_0+\ker J\), and
no others. The finite cardinalities follow from rank-nullity. This also
covers no variables, no equations, zero Jacobian and singular solutions.
\(\square\)

The proof is a universal mathematical argument. The accompanying Python
checks are independently bounded evidence about an implementation, not a
substitute for that argument. The standalone Lean specialization does
not formalize multivariable polynomial expansion (1), unless its own
checked source explicitly says so.

\hypertarget{native:necessary-boundaries-and-exact-examples}{%
\subsubsection{Necessary boundaries and exact
examples}\label{native:necessary-boundaries-and-exact-examples}}

\begin{itemize}
\tightlist
\item
  \textbf{A failed singular lift:} \(p=3\), \(L=1\), \(F(X)=X^2-3\),
  \(x=0\). Here \(e=-1\ne0\) in \(\mathbb F_3\) and \(J=0\). There are
  no lifts of zero modulo 3 to a root modulo 9.
\item
  \textbf{Representative independence is a quotient assertion:} for
  \(F(X)=X^2-1\), \(p=3\), \(L=1\), representatives \(1\) and \(4\) have
  residual digits \(0\) and \(2\). Their common Jacobian is \(2\), with
  full image, and both give zero obstruction. Their corrected
  representatives \(1\) and \(10\) are the same root modulo 9.
\item
  \textbf{A successful step need not give indefinite continuation:} for
  \(F(X)=X^2-5\), \(p=2\), the root \(1\) modulo 2 has zero obstruction
  and lifts to \(1,3\) modulo 4. Both have nonzero obstruction at the
  next step. There is no root modulo 8 and no compatible infinite lift.
\item
  \textbf{The raw Jacobian cokernel need not be presentation
  independent:} the presentations \(F=(X)\) and \(G=(X,0)\) define the
  same solution system. Their raw obstruction groups at zero are
  respectively zero and \(\mathbb F_p\). Both actual distinguished
  obstructions are zero. Thus one cannot claim a canonical isomorphism
  of these two raw groups merely because the equations define the same
  seed.
\end{itemize}

\hypertarget{native:an-intrinsic-finite-obstruction-under-an-algebraic-restriction}{%
\subsection{3. An intrinsic finite obstruction under an algebraic
restriction}\label{native:an-intrinsic-finite-obstruction-under-an-algebraic-restriction}}

Let \(B\) be a finitely presented commutative \(\mathbb Z\)-algebra,
\(p\) prime and \(L\ge1\). For a ring map
\(f:B\to\mathbb Z/p^L\mathbb Z\), use its reduction to make
\(\mathbb F_p\) a \(B\)-module. There is a canonical class

\[o(f)\in\operatorname{Ext}^1_B(\mathrm{NL}_{B/\mathbb Z},\mathbb F_p)\]

whose vanishing is equivalent to a lift
\(B\to\mathbb Z/p^{L+1}\mathbb Z\). This follows from the ring-map
lifting theorem in
\href{https://stacks.math.columbia.edu/tag/08S5}{Stacks Project, Lemma
91.2.1, tag 08S5}. Specialize the source extension to
\(0\to0\to B\to B\) and the target extension to the kernel
\(p^L/p^{L+1}\cong\mathbb F_p\). The base maps are identities of
\(\mathbb Z\). The kernel is square zero because \(L\ge1\).

The comparison in
\href{https://stacks.math.columbia.edu/tag/00S1}{Stacks Project, Lemma
10.134.2, tag 00S1} identifies the complex from a polynomial
presentation with the canonical naive cotangent complex up to homotopy.
Thus this construction is intrinsic, up to the stated canonical
comparison.

The group is finite. Choose a finite presentation
\(P=\mathbb Z[X_1,\ldots,X_d]\twoheadrightarrow B\) with finitely
generated kernel \(I\). The presentation description in
\href{https://stacks.math.columbia.edu/tag/0GPU}{Stacks Project, Remark
91.2.4, tag 0GPU} identifies the group as

\[\operatorname{coker}\left(\operatorname{Hom}_B(\Omega_{P/\mathbb Z}\otimes_P B,\mathbb F_p)\longrightarrow\operatorname{Hom}_B(I/I^2,\mathbb F_p)\right).\]

If \(I\) has \(r\) chosen generators, evaluation on their classes
injects the latter Hom module into \(\mathbb F_p^r\). Hence its quotient
is finite. The relation module \(I/I^2\) retains dependencies between
equations, including redundant zero equations, which the raw target
\(\mathbb F_p^r\) does not.

\textbf{Status:} externally proved after the exact stated
specialization, with the finiteness argument above. This theorem is not
independently formalized in the accompanying Lean file. Extending it to
all operator/recurrence/module systems admitted by manuscript Definition
1.1 requires a specified encoding and a proof that its notion of lift is
preserved. That is the missing implication for applying it to native
Conjecture 13.2 without restriction.

\textbf{Source record.} The Stacks Project authors, live primary text
inspected 2026-09-08, tags 08S5, 0GPU and 00S1. Full theorem/proof text
was inspected, not only search snippets. No immutable source revision
was pinned or downloaded. The use is a classical specialization, not a
novelty search or a claim of a new obstruction theory. No broader paper
or source cache was changed.

\hypertarget{native:a-topological-obstruction-to-the-proposed-spine-formula}{%
\subsection{4. A topological obstruction to the proposed spine
formula}\label{native:a-topological-obstruction-to-the-proposed-spine-formula}}

Consider the recurrence \(a_{n+1}=a_n\) with integer initial value. Its
global parameter space is \(X=\mathbb Z\). Its local parameter spaces
are \(X_p=\mathbb Z_p\), and every local realization has the same
prime-independent recurrence. Give the local factors their usual
\(p\)-adic topologies and the product its product topology. For these
integral local factors, the restricted product relative to all
\(\mathbb Z_p\) is the ordinary product. Write

\[A=\mathbb Z_2,\qquad B=\prod_{p>2}\mathbb Z_p,\qquad D=\{(n,(n)_{p>2}):n\in\mathbb Z\}\subset A\times B.\]

\textbf{Theorem.} There are no Hausdorff space \(C\) and continuous maps
\(f:A\to C\), \(g:B\to C\) for which the equality locus

\[\{(a,b):f(a)=g(b)\}\]

is exactly \(D\). Thus the canonical local realizations of global
integer constant sequences cannot satisfy the manuscript's fiber-product
formula in this Hausdorff topological interpretation.

\textbf{Proof.} The equality locus is closed: it is the inverse image of
the diagonal of \(C\times C\) under the continuous map \(f\times g\),
and a Hausdorff space has closed diagonal.

The set \(D\) is dense in \(A\times B\). A nonempty basic open set
prescribes congruences at only finitely many primes and finite precision
at each. The moduli are pairwise coprime. The finite Chinese remainder
theorem supplies an integer meeting all those congruences, so every such
open set meets \(D\).

However, \(D\) is not the whole product. The element with 2-adic
coordinate zero and every odd-adic coordinate one is not diagonal. If an
integer had zero 2-adic coordinate, it would be divisible by every power
of two, forcing it to be zero. Its 3-adic coordinate would then be zero,
not one. Consequently \(D\) is a proper dense subset and is not closed.
It cannot be the equality locus above. \(\square\)

This is not asserted to refute the original, underspecified sentence
under every possible reading. It refutes the explicit and natural
interpretation just stated. The assumptions of continuity, Hausdorffness
and identification through the canonical local maps are essential to
this claim.

\hypertarget{native:a-precise-reconstruction-that-does-hold}{%
\subsubsection{A precise reconstruction that does
hold}\label{native:a-precise-reconstruction-that-does-hold}}

For a finite set \(S\) of primes, positive precisions \(L_p\), and a
fixed integer polynomial system \(F\), put

\[N=\prod_{p\in S}p^{L_p},\qquad\operatorname{Sol}_F(R)=\{x\in R^d:F(x)=0\}.\]

Then coefficientwise CRT gives a canonical bijection

\[\operatorname{Sol}_F(\mathbb Z/N\mathbb Z)\cong\prod_{p\in S}\operatorname{Sol}_F(\mathbb Z/p^{L_p}\mathbb Z).\]

The reason is exact: the CRT ring isomorphism preserves polynomial
evaluation, and zero in a finite product is equivalent to zero in every
component. For \(F\) empty, passing to all compatible finite congruences
identifies the profinite completion \(\widehat{\mathbb Z}\) with
\(\prod_p\mathbb Z_p\). The same finite reconstruction is compatible
with coefficient projections and integer recurrence matrices because
these operations commute with reduction and CRT.

Here the polynomial equations or recurrence matrices are fixed common
data. The theorem reconstructs finite or profinite objects. It does not
identify a ``smallest'' spine, impose additional spectral descent, or
recover global \(\mathbb Z\) from arbitrary tuples of local points. A
completed version of 13.3 must say exactly which object is completed and
specify its universal property.

\hypertarget{native:a-finite-compatibility-criterion-with-a-precise-universal-property}{%
\subsubsection{A finite compatibility criterion with a precise universal
property}\label{native:a-finite-compatibility-criterion-with-a-precise-universal-property}}

Fix finite sets \(A,B\) of local packets and an explicitly given
relation \(R\subseteq A\times B\). Say that \(R\) is rectangular when

\[R(a,b),\ R(a',b),\ R(a',b')\quad\Longrightarrow\quad R(a,b')\]

for all \(a,a'\in A\) and \(b,b'\in B\).

\textbf{Theorem.} The relation \(R\) can be represented exactly as

\[R(a,b)\quad\Longleftrightarrow\quad f(a)=g(b)\]

for maps from \(A\) and \(B\) to a common set if and only if \(R\) is
rectangular. In that case a canonical finite choice of common set is the
connected-component set \(C_R\) of the bipartite graph with vertex set
\(A\sqcup B\) and edges given by \(R\).

The component maps \(f_R:A\to C_R\) and \(g_R:B\to C_R\) have the
following universal property. For every set \(Y\) and maps \(h:A\to Y\),
\(k:B\to Y\) satisfying \(h(a)=k(b)\) whenever \(R(a,b)\), there is a
unique map \(u:C_R\to Y\) with \(u\circ f_R=h\) and \(u\circ g_R=k\).

\textbf{Proof.} Any equality representation is rectangular, since the
three assumed equalities imply \(f(a)=g(b)=f(a')=g(b')\).

Conversely, suppose \(R\) is rectangular. An edge always gives equal
component labels. If \(a\in A\) and \(b\in B\) have equal component
labels, there is a finite path between them. Its length is odd. A path
of length one is already an \(R\)-edge. A longer path begins with three
edges \(a-b_1-a_1-b_2\). Rectangularity gives the edge \(a-b_2\).
Replacing these three edges by that edge shortens the path by two.
Induction therefore yields \(R(a,b)\). This proves the exact equality
representation. Finiteness follows because the component set is a
quotient of the finite set \(A\sqcup B\).

For the universal property, combine \(h\) and \(k\) into a function on
\(A\sqcup B\). It is constant across every edge by hypothesis, hence
along every finite path. It therefore descends to a well-defined
function on components. The two required factorization identities hold
by construction. They determine the value on every component, proving
uniqueness. \(\square\)

For example, the three pairs \((a_0,b_0),(a_1,b_0),(a_1,b_1)\) without
\((a_0,b_1)\) fail rectangularity. Every equality-label representation
agreeing on those three pairs also admits the fourth pair, so it cannot
represent that relation exactly.

This is a precise finite gluing obligation: future arithmetic
applications must define \(A\), \(B\) and \(R\) and prove
rectangularity. The universal property holds for the component quotient
even when rectangularity fails, but then the equality locus strictly
enlarges \(R\). The theorem establishes neither a minimum number of
labels nor a smallest spine in the manuscript's unspecified category. It
carries no automatic topology, spectral information or Frobenius
compatibility. \textbf{Evidence: prose proof in this note, not a Lean
formalization.}

\clearpage\section{Recovered historical local polyphase manuscript}\label{app:recovered}
This is the complete substantive body of the reconstructed Annals textual
copy, preserved to retain its derivations and original parameter ranges.
The recovered text is not certified as the exact original TeX bytes. Its
historical theorem headings distinguish these statements from the current
numbered core. Preservation is not a blanket proof audit or Lean certificate:
in particular the broader Amice analytic pullback and depleted branches,
finite-field modal minimality and unrestricted infinite polyphase claims
retain their historical evidence scope. The current formal statement map
records only the individually identified interfaces in the main text.

The earlier reconstruction lacked all 14 original figure files. This
reading edition supplies new reconstructions at their retained labels and
captions. These are explanatory schematics and exact finite examples, not
recovered original artwork. The reproducible program
\nolinkurl{tools/reconstructed_figures.py} checks the displayed finite data.
The modulo-seven inverse plot uses the corrected positive tail
$6,5,2,6,5,2,\ldots$. A reconstructed illustration does not upgrade the
historical proof or formalization status. The original historical source
package, including all 14 vector figures and the bibliography, has now been
recovered and is preserved unchanged in the accompanying release. The current
reading edition retains its explicitly identified reconstructions and corrected
finite data. The 22 historical references now have recovered bibliographic
entries and checked identities. Superseded claims about the divisor-branch
finite-difference problem have local notices directing the reader to the
current theorem. The higher kernel-valuation and unrestricted
local-to-global questions remain open in their stated scopes.

\subsection{Introduction}\label{hist:sec:introduction}

Euler numbers carry several kinds of local information at once. Their indices have residue addresses. Their coefficient transforms have triangular depth. Their local expansions have Mahler degree. Their congruences depend on the valuation of the transfer kernel and on the valuation of the source value in the row being reconstructed. These pieces should not be merged into one undifferentiated congruence calculation.

This article begins with the Hurwitz product
\[
(a\starh b)_n
=
\sum_{r=0}^{n}
\binom nr a_rb_{n-r}.
\]
Over a field of characteristic zero it is multiplication of exponential generating functions. Over a finite field it realizes the additive group algebra through the exponential modes \(\epsilon_u=(u^n)_{n\geq0}\). This finite algebra gives an exact proof of the standard window period. Hurwitz series rings are treated in \cite{Keigher1997,BarberoCerrutiMurru2018}. The finite difference and umbral background may be found in \cite{Roman1984,Comtet1974,GrahamKnuthPatashnik1994}.

Fix an odd prime \(p\), a divisor \(m\mid p-1\), and put
\[
h=\frac{p-1}{m}.
\]
The reduced standard window inverse is defined by
\[
\sum_{q=0}^{n}
\binom{mn}{mq}e_q^{(m,p)}
=
\delta_{n0}
\quad\text{in }\FF_p.
\]
We prove that
\[
e_{n+h}^{(m,p)}=e_n^{(m,p)}
\qquad(n\geq1).
\]
The result first follows from subgroup symmetry in the finite Hurwitz algebra. A second proof shows that the standard window operator preserves the finite space of sequences that are \(h\) periodic after their constant term. The two arguments expose different parts of the same finite character structure.

The odd and dyadic Euler formulas share a common reconstruction framework. For each prime \(p\) we define an even source series \(H_p\) and an inverse kernel \(K_p\) satisfying
\[
K_p(t)H_p(t)=\operatorname{sech}t.
\]
At an odd prime,
\[
H_p(t)=\cosh^2(pt)\operatorname{sech}t,
\qquad
K_p(t)=\operatorname{sech}^2(pt).
\]
At the dyadic prime,
\[
H_2(t)=\frac{\sinh(2t)\sinh t}{2t^2},
\qquad
K_2(t)=\frac{4t^2}{\sinh^2(2t)}.
\]
The coefficient convolution is identical in form, but the kernels have different integral weights. This is the correct sense in which the local theories unify.

The dyadic kernel is substantially sharper than the separate coefficient bound alone suggests. If
\[
\frac{z^2}{\sinh^2z}
=
\sum_{r\geq0}\beta_rz^{2r}
\]
and
\[
\kappa_{2,r}=\beta_r2^{2r}(2r)!,
\]
then
\[
\vTwo(\kappa_{2,r})=4r-1
\qquad(r\geq1).
\]
Consequently, truncation after depth \(M\) is uniformly valid modulo
\[
2^{4M+3},
\]
and this exponent is optimal. The rowwise estimate records the additional binary carries in the binomial coefficient. Classical congruence theory for Euler families is extensive \cite{CarlitzLevine1960,Young1999,Young2001,Sun2008,Chen2004}. The formulas here are stated as exact coordinate identities. Historical priority requires a separate comparison of the individual theorems.

The second expansion concerns local interpolation. The exponent of the finite unit group moves every unit mode into a principal unit neighborhood. Binomial expansion then gives a Mahler series whose depth \(r\) coefficient gains at least \(Lr\) digits when the chosen jump is the identity modulo \(p^L\). The same argument applies to polynomial exponential modes and to invertible matrix state systems. Mahler interpolation and the Amice transform are classical \cite{Mahler1958,Amice1964,Robert2000}. The local unit structure and valuation tools used here are standard \cite{Koblitz1984,Serre1979,Washington1997,Gouvea1997}.

The odd Euler source contains one nonunit base. Removing that base gives an analytic unit source and an interpolation of a depleted Euler sequence on every even character branch. The dyadic source is already analytic on \(\Zp\) after a removable singularity is filled. It yields a direct analytic interpolation of \(E_{2n}\) on \(\mathbb Z_2\). These branches are then refined by the substitution
\[
x=Mq+s.
\]
The induced Mahler coefficient transform is finite and explicit. Its kernels compose exactly. At a prime power radix, they become sparse modulo the residue prime. Every monotone valuation ladder survives the refinement.

The remaining sections treat arithmetic lift obstructions, recurrence packets, and Chinese remainder reconstruction. The lifting obstruction is the residual defect modulo the image of the declared correction map. A recurrence packet is governed by a jump matrix \(C^N\). If one integral recurrence is fixed, all of its local realizations have one common source. These statements do not prove that arbitrary unrelated local objects have a global origin.

\begin{figure}[t]
\centering
\includegraphics[width=.95\linewidth,height=3.15in,keepaspectratio]{../../figures/reconstructed/architecture.pdf}
\caption{New reconstruction. The arithmetic coordinate chain. Residue address, local character, Mahler depth, transfer weight, and reconstruction are recorded separately.}
\label{hist:fig:architecture}
\end{figure}

Four boundaries remain explicit. A period modulo \(p\) does not by itself lift to all powers of \(p\). The full finite difference estimate for every congruential reciprocal packet still requires control of the triangular inverse. Exact higher valuations of the affine polyphase kernel are not yet known in general. Chinese remainder reconstruction and finite adelic synchronization do not prove a global source for arbitrary local systems.
\subsection{Hurwitz series and finite field subgroup inversion}\label{hist:sec:hurwitz}
\par\smallskip\noindent\textbf{Current status.} This dated statement is superseded for the stated divisor-branch family by Theorem~\ref{thm:arith-divisor-base}, with the $j!$ normalization and factorial valuation loss explicitly retained. It makes no claim for unrestricted families.\par

\subsubsection{The Hurwitz ring}

\begin{histdefinition}[Hurwitz series ring]\label{hist:def:hurwitz_ring}
Let \(R\) be a commutative ring with identity. The Hurwitz series ring \(\Hcal(R)\) is the set \(R^{\NN}\) with componentwise addition and product
\[
(a\starh b)_n
=
\sum_{r=0}^{n}
\binom nr a_r b_{n-r}.
\]
Its identity is
\[
\delta=(1,0,0,\ldots).
\]
\end{histdefinition}

When the factorials are invertible, the sequence \(a\) may be written as the exponential generating series
\[
A(t)=\sum_{n\geq0}a_n\frac{t^n}{n!},
\]
and the Hurwitz product becomes ordinary multiplication. The sequence definition remains valid in every characteristic.

\begin{histdefinition}[Exponential mode]\label{hist:def:exponential_mode}
For \(u\in R\), define
\[
\epsilon_u=(u^n)_{n\geq0},
\]
where \(0^0=1\).
\end{histdefinition}

\begin{histproposition}[Addition law for modes]\label{hist:prop:mode_addition}
For every \(u,v\in R\),
\[
\boxed{
\epsilon_u\starh\epsilon_v
=
\epsilon_{u+v}.
}
\]
\end{histproposition}

\begin{proof}
The \(n\)th coefficient is
\[
\sum_{r=0}^{n}
\binom nr u^rv^{n-r}
=
(u+v)^n
\]
by the binomial theorem.
\end{proof}

\subsubsection{The finite additive group algebra}

Let \(p\) be prime. Write \(\FF_p^+\) for the additive group of the finite field.

\begin{histtheorem}[Finite field mode algebra]\label{hist:thm:mode_group_algebra}
The map
\[
\Theta:
\FF_p[\FF_p^+]
\longrightarrow
\Hcal(\FF_p),
\qquad
[u]\longmapsto\epsilon_u,
\]
is an injective algebra homomorphism. Its image is the \(p\) dimensional subalgebra spanned by the modes \(\epsilon_u\).
\end{histtheorem}

\begin{proof}
The product law follows from \cref{hist:prop:mode_addition}. To prove injectivity, suppose
\[
\sum_{u\in\FF_p}c_u\epsilon_u=0.
\]
Looking at the coordinates \(n=0,1,\ldots,p-1\) gives a Vandermonde system with matrix
\[
(u^n)_{0\leq n\leq p-1,\ u\in\FF_p}.
\]
Its determinant is the product of all differences between distinct field elements, so it is nonzero. Hence every \(c_u\) vanishes.
\end{proof}

\begin{histproposition}[Augmentation criterion]\label{hist:prop:augmentation_criterion}
Let
\[
x=\sum_{u\in\FF_p}c_u[u]
\in
\FF_p[\FF_p^+].
\]
Then \(x\) is a unit if and only if
\[
\operatorname{aug}(x)
=
\sum_uc_u
\]
is nonzero. If \(J\) is the augmentation ideal, then
\[
J^p=0.
\]
\end{histproposition}

\begin{proof}
Choose a generator \(g\) of the cyclic additive group. Then
\[
\FF_p[\FF_p^+]
\cong
\FF_p[y]/(y^p-1)
=
\FF_p[y]/(y-1)^p.
\]
The augmentation ideal is generated by \(y-1\), so its \(p\)th power is zero. Every element has the form \(a+n\) with \(a\in\FF_p\) and \(n\in J\). It is invertible exactly when \(a\neq0\), which is the augmentation condition.
\end{proof}

\begin{histcorollary}[Finite inverse formula]\label{hist:cor:finite_inverse_formula}
If \(x\) has augmentation one and \(n=x-1\), then
\[
\boxed{
x^{-1}
=
\sum_{r=0}^{p-1}(-n)^r.
}
\]
\end{histcorollary}

\begin{proof}
By \cref{hist:prop:augmentation_criterion}, \(n^p=0\). The finite geometric sum gives the inverse.
\end{proof}

\subsubsection{Subgroup selectors}

Fix a divisor \(m\) of \(p-1\). Let
\[
H_m
=
\{\zeta\in\FF_p^\times:\zeta^m=1\}.
\]
This group has order \(m\).

\begin{histdefinition}[Multiple selector]\label{hist:def:multiple_selector}
Define \(a^{(m)}\in\Hcal(\FF_p)\) by
\[
a_n^{(m)}
=
\begin{cases}
1,&m\mid n,\\
0,&m\nmid n.
\end{cases}
\]
\end{histdefinition}

\begin{histlemma}[Subgroup average]\label{hist:lem:subgroup_average}
The multiple selector has the mode expansion
\[
\boxed{
a^{(m)}
=
\frac1m
\sum_{\zeta\in H_m}
\epsilon_\zeta.
}
\]
\end{histlemma}

\begin{proof}
The \(n\)th coefficient on the right is
\[
\frac1m
\sum_{\zeta\in H_m}\zeta^n.
\]
The character sum is one when \(m\mid n\) and zero otherwise.
\end{proof}

\begin{figure}[t]
\centering
\includegraphics[width=.95\linewidth,height=3.15in,keepaspectratio]{../../figures/reconstructed/subgroup_average.pdf}
\caption{New reconstruction. A subgroup average in the additive group algebra. The highlighted modes select coefficient degrees divisible by the subgroup order.}
\label{hist:fig:subgroup_average}
\end{figure}

Since the subgroup average has augmentation one, it has a unique Hurwitz inverse. Write
\[
b^{(m,p)}
=
(a^{(m)})^{-1}.
\]

\begin{histtheorem}[Support and positive degree period]\label{hist:thm:support_period}
The inverse \(b^{(m,p)}\) has the following properties.
\begin{enumerate}
\item \(b_n^{(m,p)}=0\) unless \(m\mid n\).
\item For every positive \(n\),
\[
\boxed{
b_{n+p-1}^{(m,p)}
=
b_n^{(m,p)}.
}
\]
\end{enumerate}
\end{histtheorem}

\begin{proof}
Multiplication by an element \(\eta\in H_m\) gives an automorphism of the additive group and therefore an automorphism
\[
\epsilon_u\longmapsto\epsilon_{\eta u}
\]
of the mode algebra. The selector \(a^{(m)}\) is fixed by every such automorphism, so its unique inverse is fixed as well.

Write
\[
b^{(m,p)}
=
\sum_{u\in\FF_p}c_u\epsilon_u.
\]
Invariance gives \(c_{\eta u}=c_u\). After the change of variable \(u\mapsto\eta u\),
\[
b_n^{(m,p)}
=
\eta^n b_n^{(m,p)}
\]
for every \(\eta\in H_m\). If \(m\nmid n\), some \(\eta\) satisfies \(\eta^n\neq1\), so the coefficient vanishes.

For positive \(n\), the zero mode contributes nothing. Every nonzero mode satisfies \(u^{p-1}=1\). Hence
\[
b_{n+p-1}^{(m,p)}
=
\sum_{u\neq0}c_uu^{n+p-1}
=
\sum_{u\neq0}c_uu^n
=
b_n^{(m,p)}.
\]
\end{proof}

\begin{histdefinition}[Compressed inverse]\label{hist:def:compressed_inverse}
Set
\[
e_n^{(m,p)}
=
b_{mn}^{(m,p)}.
\]
\end{histdefinition}

\begin{histproposition}[Compressed inverse recurrence]\label{hist:prop:compressed_inverse_recurrence}
The compressed coefficients satisfy
\[
\boxed{
\sum_{q=0}^{n}
\binom{mn}{mq}
e_q^{(m,p)}
=
\delta_{n0}.
}
\]
In particular, \(e_0^{(m,p)}=1\), and the positive coefficients are determined recursively because the diagonal is one.
\end{histproposition}

\begin{proof}
Evaluate the Hurwitz identity
\[
a^{(m)}\starh b^{(m,p)}=\delta
\]
at degree \(mn\). Both factors vanish outside degrees divisible by \(m\). Reindexing the surviving terms gives the displayed recurrence.
\end{proof}

\begin{histcorollary}[Uniform compressed period]\label{hist:cor:compressed_period}
Let
\[
h=\frac{p-1}{m}.
\]
Then
\[
\boxed{
e_{n+h}^{(m,p)}
=
e_n^{(m,p)}
\qquad(n\geq1).
}
\]
\end{histcorollary}

\begin{proof}
Since \(m(n+h)=mn+p-1\), the conclusion follows from \cref{hist:thm:support_period}.
\end{proof}

\begin{histcorollary}[Rational ordinary generating function]\label{hist:cor:periodic_ogf}
In \(\FF_p[[z]]\),
\[
\boxed{
\sum_{n\geq0}e_n^{(m,p)}z^n
=
1+
\frac{
\sum_{r=1}^{h}e_r^{(m,p)}z^r
}{1-z^h}.
}
\]
\end{histcorollary}

\begin{proof}
Separate the constant coefficient and group the positive indices into residue classes modulo \(h\).
\end{proof}

\begin{histtheorem}[Exponential form and minimal period]\label{hist:thm:minimal_period}
Let
\[
G_m=(\FF_p^\times)^m.
\]
This cyclic group has order \(h\). There are uniquely determined coefficients \(d_v\in\FF_p\), indexed by \(v\in G_m\), such that
\[
\boxed{
e_n^{(m,p)}
=
\sum_{v\in G_m}d_vv^n
\qquad(n\geq1).
}
\]
The minimal positive period is the least common multiple of the orders of those \(v\) for which \(d_v\neq0\). In particular, it divides \(h\).
\end{histtheorem}

\begin{proof}
The coefficients \(c_u\) in the proof of \cref{hist:thm:support_period} are constant on the cosets of \(H_m\). For \(u\neq0\), the value \(u^m\) depends only on the coset, and every element of \(G_m\) arises. Grouping the modes by cosets gives the displayed expansion.

Distinct elements of \(G_m\) define linearly independent functions on one full period, by the Vandermonde determinant. Therefore a positive integer \(r\) is a period exactly when \(v^r=1\) for every active mode. The least such \(r\) is the stated least common multiple.
\end{proof}

\begin{histexample}[A compressed inverse modulo seven]\label{hist:ex:compressed_mod_seven}
\par\noindent\textbf{Editorial numerical correction.} The positive-index standard-window tail printed below is erroneous. The defining recurrence gives $e_1=-1\equiv6$, $e_2=5$ and $e_3=-61\equiv2\pmod7$, so the correct positive tail is $6,5,2,6,5,2,\ldots$. The period-three conclusion remains correct. The original line is retained only as a corrected historical record.\par\medskip
Take \(p=7\) and \(m=2\), so \(h=3\). The recurrence gives
\[
e_n^{(2,7)}
=
1,6,4,1,6,4,\ldots
\qquad(n\geq1).
\]
The positive tail has exact period three.
\end{histexample}

\begin{figure}[t]
\centering
\includegraphics[width=.95\linewidth,height=3.15in,keepaspectratio]{../../figures/reconstructed/standard_periods.pdf}
\caption{New reconstruction. Several compressed inverses. Every row is periodic after the constant term, as proved in \cref{hist:cor:compressed_period}.}
\label{hist:fig:standard_periods}
\end{figure}


\subsection{Unit triangular coefficient transport}\label{hist:sec:triangular_transport}

The examples later in the chapter use different kernels, but their coefficient algebra has one common form. This section records that form once. Let \(R\) be a commutative algebra over \(\QQ\), and take all coefficients in \(R\).

Let
\[
H(t)=\sum_{r\geq0}h_r\frac{t^r}{r!},
\qquad
h_0=1.
\]
There is a unique inverse
\[
K(t)=H(t)^{-1}
=\sum_{r\geq0}k_r\frac{t^r}{r!}.
\]

\begin{histtheorem}[Exact coefficient transport]\label{hist:thm:exact_coefficient_transport}
Let
\[
A(t)=\sum_{n\geq0}A_n\frac{t^n}{n!},
\qquad
S(t)=H(t)A(t).
\]
Then
\[
\boxed{
S_n
=
\sum_{r=0}^{n}
\binom nr h_rA_{n-r}.
}
\]
The inverse coordinate formula is
\[
\boxed{
A_n
=
\sum_{r=0}^{n}
\binom nr k_rS_{n-r}.
}
\]
Both transforms are lower triangular with diagonal one.
\end{histtheorem}

\begin{proof}
The first identity is the coefficient rule for the product of exponential generating functions. Multiply by \(K(t)\) and use \(K(t)H(t)=1\) for the inverse identity. The diagonal coefficient in each transform is the constant coefficient of its kernel, which is one.
\end{proof}

\begin{histcorollary}[Composition of coefficient transports]\label{hist:cor:transport_composition}
If \(U_H\) and \(U_G\) are the triangular transforms associated with unit kernels \(H\) and \(G\), then
\[
\boxed{
U_HU_G=U_{HG}.
}
\]
In particular, the inverse of \(U_H\) is \(U_{H^{-1}}\).
\end{histcorollary}

\begin{proof}
Multiplication by \(G\), followed by multiplication by \(H\), is multiplication by \(HG\).
\end{proof}

\begin{histtheorem}[Filtered triangular reconstruction]\label{hist:thm:filtered_triangular_reconstruction}
Let \(p\) be prime. Suppose
\[
A_n
=
\sum_{r=0}^{n}
\binom nr
c_rp^{w(r)}S_{n-r},
\]
where \(c_r\in\Zp\), \(S_n\in\Zp\), and
\[
w(r)\in\NN.
\]
For every precision \(L\geq1\),
\[
\boxed{
A_n
\equiv
\sum_{\substack{0\leq r\leq n\\w(r)<L}}
\binom nr
c_rp^{w(r)}S_{n-r}
\pmod {p^L}.
}
\]
When the omitted set is nonempty, the tail has valuation at least
\[
\boxed{
\min_{\substack{0\leq r\leq n\\w(r)\geq L}}
\left(
w(r)+\vp\binom nr+\vp(S_{n-r})
\right).
}
\]
\end{histtheorem}

\begin{proof}
Every omitted summand belongs to \(p^L\Zp\). The sharper statement follows by taking the minimum of the exact lower bounds for the summands and applying the nonarchimedean triangle inequality.
\end{proof}

\begin{histremark}\label{hist:rem:triangular_freedom}
An arbitrary invertible triangular matrix can encode an arbitrary sequence. The transforms retained here have an explicit generating kernel and a visible filtration. Those two features make exact inversion and uniform truncation possible.
\end{histremark}

\subsection{Congruential reciprocal packets}\label{hist:sec:congruential}

\subsubsection{Residue sections of the exponential series}

For integers \(N\geq1\) and \(0\leq j<N\), define
\[
\Ecal_{N,j}(t)
=
\sum_{n\geq0}
\frac{t^{Nn+j}}{(Nn+j)!}.
\]
The forced factor \(t^j\) may be removed without changing the residue address.

\begin{histdefinition}[Congruential reciprocal packet]\label{hist:def:congruential_packet}
Define
\[
\boxed{
\Ccal_{N,j}(t)
=
\frac{t^j/j!}{\Ecal_{N,j}(t)}
=
\sum_{n\geq0}
C_{N,j}(n)
\frac{t^{Nn}}{(Nn)!}.
}
\]
\end{histdefinition}

\begin{histtheorem}[Triangular inverse system]\label{hist:thm:triangular_inverse}
The coefficients satisfy \(C_{N,j}(0)=1\) and
\[
\boxed{
\sum_{q=0}^{n}
\binom{Nn+j}{Nq}
C_{N,j}(q)
=
\delta_{n0}.
}
\]
For \(n\geq1\),
\[
\boxed{
C_{N,j}(n)
=
-
\frac{
\sum_{q=0}^{n-1}
\binom{Nn+j}{Nq}C_{N,j}(q)
}{
\binom{Nn+j}{Nn}
}.
}
\]
\end{histtheorem}

\begin{proof}
Multiply \(\Ccal_{N,j}(t)\) by \(\Ecal_{N,j}(t)\) and compare the coefficient of \(t^{Nn+j}/(Nn+j)!\). The last summand is the diagonal term, which is nonzero over \(\QQ\).
\end{proof}

\begin{histproposition}[A denominator bound]\label{hist:prop:denominator_bound}
Let
\[
D_{N,j}(n)
=
\prod_{r=1}^{n}
\binom{Nr+j}{Nr},
\qquad
D_{N,j}(0)=1.
\]
Then
\[
\boxed{
D_{N,j}(n)C_{N,j}(n)
\in
\ZZ.
}
\]
For every prime \(p\),
\[
\vp(C_{N,j}(n))
\geq
-
\sum_{r=1}^{n}
\vp\binom{Nr+j}{Nr}.
\]
\end{histproposition}

\begin{proof}
Use induction on \(n\) in the recurrence of \cref{hist:thm:triangular_inverse}. Every earlier denominator divides \(D_{N,j}(n-1)\), while the new diagonal denominator is the final factor of \(D_{N,j}(n)\). The valuation statement follows.
\end{proof}

\begin{histexample}[Euler and Bernoulli boundaries]\label{hist:ex:euler_bernoulli_boundaries}
For \((N,j)=(2,0)\),
\[
\Ccal_{2,0}(t)=\operatorname{sech}t,
\]
so
\[
C_{2,0}(n)=E_{2n}.
\]
For \((N,j)=(2,1)\),
\[
\Ccal_{2,1}(t)=\frac{t}{\sinh t},
\]
and
\[
C_{2,1}(n)
=
2\left(1-2^{2n-1}\right)B_{2n}.
\]
The usual convention \(B_1=-1/2\) is used. Euler and Bernoulli number conventions are reviewed in \cite{Comtet1974,Stanley2012}.
\end{histexample}

\subsubsection{Standard windows}

Fix a prime \(p\), an integer \(m\geq1\), and an address \(0\leq j<p\). Put
\[
N=mp.
\]
For arbitrary \(m\), let \(\widetilde e_n^{(m,p)}\in\FF_p\) be the unique sequence satisfying
\[
\boxed{
\sum_{q=0}^{n}
\binom{mn}{mq}
\widetilde e_q^{(m,p)}
=
\delta_{n0}.
}
\]
When \(m\mid p-1\), this sequence agrees with the compressed Hurwitz inverse in \cref{hist:prop:compressed_inverse_recurrence}.

\begin{histstandardfact}[Lucas congruence]\label{hist:fact:lucas}
If \(a=a_0+a_1p+\cdots\) and \(b=b_0+b_1p+\cdots\) are base \(p\) expansions, then
\[
\binom ab
\equiv
\prod_r\binom{a_r}{b_r}
\pmod p.
\]
This is the classical theorem of Lucas \cite{Lucas1878a}.
\end{histstandardfact}

\begin{histtheorem}[Standard window reduction]\label{hist:thm:standard_window_reduction}
For \(0\leq q\leq n\),
\[
\boxed{
\binom{mpn+j}{mpq}
\equiv
\binom{mn}{mq}
\pmod p.
}
\]
The diagonal coefficient is a unit modulo \(p\):
\[
\binom{mpn+j}{mpn}
\equiv1
\pmod p.
\]
\end{histtheorem}

\begin{proof}
Write
\[
mpn+j=p(mn)+j,
\qquad
mpq=p(mq).
\]
Lucas reduction separates the lowest base \(p\) digit and gives
\[
\binom{mpn+j}{mpq}
\equiv
\binom j0
\binom{mn}{mq}
\pmod p.
\]
The diagonal case is obtained by setting \(q=n\).
\end{proof}

\begin{histcorollary}[Standard window integrality]\label{hist:cor:standard_integrality}
Every coefficient \(C_{mp,j}(n)\) belongs to \(\Zp\).
\end{histcorollary}

\begin{proof}
The recurrence in \cref{hist:thm:triangular_inverse} starts with a \(p\) integral constant term. The binomial coefficients are integers, and every diagonal term is a \(p\) adic unit by \cref{hist:thm:standard_window_reduction}. Induction gives the result.
\end{proof}

\begin{histcorollary}[Address independence modulo the prime]\label{hist:cor:address_independence}
Let \(\widetilde e_n^{(m,p)}\) be the reduced standard inverse above. Then
\[
\boxed{
C_{mp,j}(n)
\equiv
\widetilde e_n^{(m,p)}
\pmod p
}
\]
for every \(0\leq j<p\).
\end{histcorollary}

\begin{proof}
Modulo \(p\), the recurrence in \cref{hist:thm:triangular_inverse} becomes
\[
\sum_{q=0}^{n}
\binom{mn}{mq}
C_{mp,j}(q)
=
\delta_{n0}.
\]
This is the defining recurrence for \(\widetilde e_n^{(m,p)}\). The diagonal is one, so the solution is unique.
\end{proof}

\begin{histtheorem}[Uniform standard window period]\label{hist:thm:standard_window_period}
Assume that \(p\) is odd and \(m\mid p-1\). Put
\[
h=\frac{p-1}{m}.
\]
Then for every standard window address \(0\leq j<p\),
\[
\boxed{
C_{mp,j}(n+h)
\equiv
C_{mp,j}(n)
\pmod p
\qquad(n\geq1).
}
\]
\end{histtheorem}

\begin{proof}
By \cref{hist:cor:address_independence}, every address reduces to the same reduced standard inverse. Under the divisibility assumption, that inverse agrees with the compressed Hurwitz inverse by \cref{hist:prop:compressed_inverse_recurrence}. Its positive tail has period \(h\) by \cref{hist:cor:compressed_period}.
\end{proof}

\begin{histcorollary}[Rational residual generating function]\label{hist:cor:standard_rational_ogf}
Under the hypotheses of \cref{hist:thm:standard_window_period}, for every standard address \(j\),
\[
\sum_{n\geq0}
\left(C_{mp,j}(n)\bmod p\right)z^n
=
1+
\frac{
\sum_{r=1}^{h}\widetilde e_r^{(m,p)}z^r
}{1-z^h}.
\]
\end{histcorollary}

\begin{proof}
Combine \cref{hist:cor:address_independence,hist:cor:periodic_ogf}.
\end{proof}


\begin{histtheorem}[Fourier invariance of the standard window operator]\label{hist:thm:standard_window_fourier}
Let \(p\) be odd, let \(m\mid p-1\), and let \(h=(p-1)/m\). Define
\[
(T_{m,p}f)_n
=
\sum_{q=0}^{n}\binom{mn}{mq}f_q
\quad\text{in }\FF_p.
\]
Let \(V_h\) be the vector space of sequences satisfying
\[
f_{n+h}=f_n
\qquad(n\geq1).
\]
Then
\[
\boxed{T_{m,p}(V_h)=V_h.}
\]
The restriction of \(T_{m,p}\) to \(V_h\) is an automorphism.
\end{histtheorem}

\begin{proof}
Choose \(d_0=f_h\) and \(d_n=f_n\) for \(n\geq1\). Then \(d\) is fully \(h\) periodic and \(f=d+c\delta_0\) for \(c=f_0-f_h\). Since \(T_{m,p}\delta_0\) is the constant sequence one, it is enough to treat \(d\). Expand
\[
d_q=\sum_{\chi\in\mu_h}c_\chi\chi^q.
\]
The \(m\)th power map on \(\FF_p^\times\) has image \(\mu_h\). Choose \(y\) with \(y^m=\chi\). The root filter gives
\[
T_{m,p}(\chi^\bullet)_n
=
\frac1m
\sum_{\zeta\in\mu_m}
(1+\zeta y)^{mn}.
\]
If \(1+\zeta y=0\), the corresponding term vanishes for positive \(n\). Otherwise
\[
\left((1+\zeta y)^m\right)^h
=
(1+\zeta y)^{p-1}
=1.
\]
Every summand is therefore \(h\) periodic for positive \(n\). Thus \(T_{m,p}\) preserves \(V_h\). It is unit lower triangular on every finite initial segment. Its restriction to the finite dimensional space determined by \(f_0,f_1,\ldots,f_h\) is injective and hence bijective.
\end{proof}

\begin{histcorollary}[A second proof of the standard window period]\label{hist:cor:standard_window_fourier_period}
The reduced inverse \(e^{(m,p)}\) and every standard address satisfy the period in \cref{hist:thm:standard_window_period}.
\end{histcorollary}

\begin{proof}
The sequence \(\delta=(1,0,0,\ldots)\) belongs to \(V_h\). Its inverse image under the automorphism in \cref{hist:thm:standard_window_fourier} also belongs to \(V_h\). Address independence modulo \(p\) gives the statement for every standard address.
\end{proof}

\begin{figure}[t]
\centering
\includegraphics[width=.95\linewidth,height=3.15in,keepaspectratio]{../../figures/reconstructed/standard_fourier.pdf}
\caption{New reconstruction. Finite field Fourier transfer for a standard window. Each character becomes a finite sum of positive index modes with period dividing \(h\).}
\label{hist:fig:standard_fourier}
\end{figure}
\begin{figure}[t]
\centering
\includegraphics[width=.95\linewidth,height=3.15in,keepaspectratio]{../../figures/reconstructed/standard_matrix.pdf}
\caption{New reconstruction. The standard window matrix for $p=7$, $m=2$, and $0\le j<7$. Reduction removes the address and retains diagonal one.}
\label{hist:fig:standard_matrix}
\end{figure}

\begin{histremark}\label{hist:rem:period_boundary}
The theorem gives a period modulo \(p\). A period modulo higher powers requires additional control of the triangular inverse and cannot be inferred from the finite field argument alone.

\end{histremark}

\subsection{One local Euler source and kernel calculus}\label{hist:sec:local_euler}

The odd and dyadic reconstructions have one convolutional form. Their sources, kernels, and valuation profiles are different. We first construct each chart and then state the common theorem.

\subsubsection{The odd source and its inverse}

Let \(q\) be a positive odd integer. Define
\[
S_n(q)
=
\sum_{a=0}^{q-1}
(-1)^a(2a+1)^n.
\]
Let the signed Euler numbers and Euler zigzag numbers be defined by
\[
\operatorname{sech}t
=
\sum_{n\geq0}E_n\frac{t^n}{n!}
\]
and
\[
\sec t+\tan t
=
\sum_{n\geq0}Z_n\frac{t^n}{n!}.
\]

\begin{histlemma}[Odd source generating function]\label{hist:lem:odd_source_egf}
For every positive odd \(q\),
\[
\boxed{
\sum_{n\geq0}S_n(q)\frac{t^n}{n!}
=
\e^{qt}\cosh(qt)\operatorname{sech}t.
}
\]
Its even part is
\[
\boxed{
\sum_{n\geq0}S_{2n}(q)\frac{t^{2n}}{(2n)!}
=
\cosh^2(qt)\operatorname{sech}t.
}
\]
\end{histlemma}

\begin{proof}
The finite geometric sum gives
\[
\sum_{a=0}^{q-1}(-1)^a\e^{(2a+1)t}
=
\e^t\frac{1+\e^{2qt}}{1+\e^{2t}}
=
\e^{qt}\cosh(qt)\operatorname{sech}t.
\]
Averaging at \(t\) and \(-t\) gives the even part.
\end{proof}

\begin{histlemma}[Zigzag inverse kernel]\label{hist:lem:local_zigzag_kernel}
\[
\boxed{
\operatorname{sech}^2t
=
\sum_{r\geq0}
(-1)^rZ_{2r+1}
\frac{t^{2r}}{(2r)!}.
}
\]
\end{histlemma}

\begin{proof}
The identity \(\tanh t=-i\tan(it)\) gives the odd zigzag expansion of \(\tanh t\). Differentiate.
\end{proof}

\begin{histtheorem}[Odd modulus Euler coordinate pair]\label{hist:thm:odd_coordinate_pair}
For every positive odd \(q\) and every even \(n\geq0\),
\[
\boxed{
E_n
=
\sum_{r=0}^{n/2}
(-1)^rZ_{2r+1}
\binom n{2r}
q^{2r}S_{n-2r}(q).
}
\]
The inverse transform is
\[
\boxed{
S_n(q)
=
E_n
+
\sum_{r=1}^{n/2}
2^{2r-1}
\binom n{2r}
q^{2r}E_{n-2r}.
}
\]
Both transforms are lower triangular with diagonal one.
\end{histtheorem}

\begin{proof}
The even source identity gives
\[
\operatorname{sech}t
=
\operatorname{sech}^2(qt)
\sum_{n\geq0}S_{2n}(q)\frac{t^{2n}}{(2n)!}.
\]
Insert \cref{hist:lem:local_zigzag_kernel} and compare coefficients. For the inverse, multiply \(\operatorname{sech}t\) by
\[
\cosh^2(qt)
=
1+
\sum_{r\geq1}2^{2r-1}q^{2r}\frac{t^{2r}}{(2r)!}.
\]
\end{proof}

\begin{histcorollary}[Odd modulus truncation]\label{hist:cor:odd_modulus_truncation}
For every \(M\geq0\), every positive odd \(q\), and every even \(n\),
\[
\boxed{
E_n
\equiv
\sum_{r=0}^{\min(M,n/2)}
(-1)^rZ_{2r+1}
\binom n{2r}q^{2r}S_{n-2r}(q)
\pmod {q^{2M+2}}.
}
\]
If \(p\mid q\) and \(a=\vp(q)\), the same truncation is valid modulo \(p^{2a(M+1)}\).
\end{histcorollary}

\begin{proof}
Every omitted term contains \(q^{2r}\) with \(r\geq M+1\). The local statement follows by taking the \(p\) adic valuation.
\end{proof}

\begin{histproposition}[Odd row weight]\label{hist:prop:odd_row_weight}
For an odd prime \(p\), the \(r\)th nonzero term in the reconstruction has valuation
\[
\boxed{
W_p(n,r)
=
2r+\vp(Z_{2r+1})
+\vp\binom n{2r}
+\vp\bigl(S_{n-2r}(p)\bigr).
}
\]
The binomial contribution is
\[
\boxed{
\vp\binom n{2r}
=
\frac{s_p(2r)+s_p(n-2r)-s_p(n)}{p-1},
}
\]
where \(s_p\) is the base \(p\) digit sum.
\end{histproposition}

\begin{proof}
The first identity separates the valuation of every factor. The second is Legendre's formula, equivalently Kummer's carry theorem \cite{IrelandRosen1990}.
\end{proof}

\subsubsection{The dyadic source and its inverse}

For even \(n\geq0\), define
\[
\boxed{
R_n
=
\frac{3^{n+2}-1}{4(n+1)(n+2)}.
}
\]
Its even exponential generating function is
\[
\boxed{
\sum_{n\text{ even}}R_n\frac{t^n}{n!}
=
\frac{\sinh(2t)\sinh t}{2t^2}.
}
\]
Define \(\beta_r\) by
\[
\boxed{
\frac{z^2}{\sinh^2z}
=
\sum_{r\geq0}\beta_rz^{2r}.
}
\]

\begin{histproposition}[Dyadic inverse coefficients]\label{hist:prop:dyadic_beta_unified}
For \(r\geq1\),
\[
\boxed{
\beta_r
=
-
\frac{2^{2r}(2r-1)B_{2r}}{(2r)!}.
}
\]
Consequently,
\[
\boxed{
\vTwo(\beta_r)=\sTwo(r)-1.
}
\]
Every \(R_n\) is a dyadic unit.
\end{histproposition}

\begin{proof}
Differentiate the Bernoulli expansion of \(\coth z\) and use \(-(\coth z)'=\sinh(z)^{-2}\). The von Staudt and Clausen theorem gives \(\vTwo(B_{2r})=-1\). Legendre's formula gives
\[
\vTwo((2r)!)=2r-\sTwo(r).
\]
For \(R_n\), put \(k=n+2\). Since \(k\) is even,
\[
\vTwo(3^k-1)=2+\vTwo(k).
\]
The denominator has the same valuation because \(n+1\) is odd.
\end{proof}

\begin{histtheorem}[Dyadic Euler coordinate pair]\label{hist:thm:dyadic_coordinate_pair}
For every even \(n\geq0\),
\[
\boxed{
E_n
=
\sum_{r=0}^{n/2}
\beta_r2^{2r}
\fall n{2r}
R_{n-2r}.
}
\]
The inverse transform is
\[
\boxed{
R_n
=
\sum_{r=0}^{n/2}
\frac{2^{4r+1}}{(2r+2)!}
\fall n{2r}
E_{n-2r}.
}
\]
\end{histtheorem}

\begin{proof}
Substitute \(z=2t\) in the inverse series and multiply
\[
\frac{4t^2}{\sinh^2(2t)}
\frac{\sinh(2t)\sinh t}{2t^2}
=
\operatorname{sech}t.
\]
Coefficient extraction gives the forward transform. The inverse follows from
\[
\left(\frac{\sinh(2t)}{2t}\right)^2
=
\sum_{r\geq0}
\frac{2^{4r+1}}{(2r+2)!}t^{2r}.
\]
\end{proof}

\subsubsection{The common source and kernel law}

For a prime \(p\), define
\[
A_{p,n}
=
\begin{cases}
S_n(p),&p\text{ odd},\\
R_n,&p=2\text{ and }n\text{ even},\\
0,&p=2\text{ and }n\text{ odd},
\end{cases}
\]
and
\[
\kappa_{p,r}
=
\begin{cases}
(-1)^rZ_{2r+1}p^{2r},&p\text{ odd},\\
\beta_r2^{2r}(2r)!,&p=2.
\end{cases}
\]

\begin{histtheorem}[Unified local Euler transfer]\label{hist:thm:unified_local_euler}
For every prime \(p\) and every even \(n\geq0\),
\[
\boxed{
E_n
=
\sum_{r=0}^{n/2}
\binom n{2r}
\kappa_{p,r}
A_{p,n-2r}.
}
\]
If
\[
H_p(t)
=
\sum_{m\geq0}A_{p,2m}\frac{t^{2m}}{(2m)!}
\]
and
\[
K_p(t)
=
\sum_{r\geq0}\kappa_{p,r}\frac{t^{2r}}{(2r)!},
\]
then
\[
\boxed{K_p(t)H_p(t)=\operatorname{sech}t.}
\]
For odd \(p\),
\[
H_p(t)=\cosh^2(pt)\operatorname{sech}t,
\qquad
K_p(t)=\operatorname{sech}^2(pt).
\]
For \(p=2\),
\[
H_2(t)=\frac{\sinh(2t)\sinh t}{2t^2},
\qquad
K_2(t)=\frac{4t^2}{\sinh^2(2t)}.
\]
\end{histtheorem}

\begin{proof}
The two source and kernel products were proved in \cref{hist:thm:odd_coordinate_pair,hist:thm:dyadic_coordinate_pair}. Their divided power coefficient laws have the displayed common form.
\end{proof}

\begin{histtheorem}[Unified inverse local transfer]\label{hist:thm:unified_inverse_local_euler}
Define
\[
\lambda_{p,r}
=
\begin{cases}
1,&p\text{ odd and }r=0,\\
2^{2r-1}p^{2r},&p\text{ odd and }r\geq1,\\
\dfrac{2^{4r+1}}{(2r+1)(2r+2)},&p=2.
\end{cases}
\]
Then for every prime \(p\) and every even \(n\geq0\),
\[
\boxed{
A_{p,n}
=
\sum_{r=0}^{n/2}
\binom n{2r}
\lambda_{p,r}
E_{n-2r}.
}
\]
Equivalently, if
\[
L_p(t)=K_p(t)^{-1},
\]
then
\[
\boxed{H_p(t)=L_p(t)\operatorname{sech}t,}
\]
where
\[
L_p(t)=\cosh^2(pt)
\]
for odd \(p\), and
\[
L_2(t)=\left(\frac{\sinh(2t)}{2t}\right)^2
\]
for \(p=2\).
\end{histtheorem}

\begin{proof}
For odd \(p\), this is the inverse identity in \cref{hist:thm:odd_coordinate_pair}. For \(p=2\), rewrite the inverse identity in \cref{hist:thm:dyadic_coordinate_pair} by using
\[
\fall n{2r}=(2r)!\binom n{2r}.
\]
The coefficient becomes
\[
\frac{2^{4r+1}(2r)!}{(2r+2)!}
=
\frac{2^{4r+1}}{(2r+1)(2r+2)}.
\]
The generating functions are the inverses of the kernels in \cref{hist:thm:unified_local_euler}.
\end{proof}

\begin{histcorollary}[Exact change between local source charts]\label{hist:cor:local_source_change}
\par\noindent\textbf{Editorial range correction.} This comparison holds for the even source subsequences, equivalently the evenized generating functions $H_p$. It does not identify their full odd-degree source values. For $p=3$, the original odd source has $A_{3,1}=S_1(3)=3$, whereas every odd dyadic source coefficient is zero; the even kernel $L_3K_2$ cannot change that odd coefficient. Read the coefficient consequence below only at even indices. The displayed identity $H_p=L_pK_\ell H_\ell$ and its proof remain valid.\par\medskip
For any two primes \(p\) and \(\ell\), the source sequences \(A_{p,n}\) and \(A_{\ell,n}\) are related by an exact unit lower triangular transformation obtained by composing the inverse kernel \(L_p\) with the reconstruction kernel \(K_\ell\).
\end{histcorollary}

\begin{proof}
Both source series reconstruct \(\operatorname{sech}t\). Hence
\[
H_p(t)=L_p(t)K_\ell(t)H_\ell(t).
\]
The multiplier has constant coefficient one, so \cref{hist:thm:exact_coefficient_transport} gives an exact unit lower triangular coordinate change.
\end{proof}

\begin{table}[t]
\centering
\caption{The common local Euler architecture and the parts that remain prime dependent.}
\label{hist:tab:local_chart_comparison}
\begin{tabular}{>{\raggedright\arraybackslash}p{.22\textwidth}>{\raggedright\arraybackslash}p{.31\textwidth}>{\raggedright\arraybackslash}p{.31\textwidth}}
\toprule
Feature & Odd prime chart & Dyadic chart\\
\midrule
Finite torsion address & Modulo \(p-1\) & Sign address modulo two\\
Principal units & \(1+p\mathbb Z_p\) & \(1+4\mathbb Z_2\)\\
Source & Alternating odd power sum \(S_n(p)\) & Weighted source \(R_n\)\\
Kernel & \(\operatorname{sech}^2(pt)\) & \(4t^2/\sinh^2(2t)\)\\
Positive depth weight & At least \(2r\), with zigzag and row factors & Exactly \(4r-1\), before the binomial carry\\
Analytic branch & Depleted unit source & Direct integral source\\
Natural refined step & \((p-1)p^a\) & \(2^{a+1}\)\\
\bottomrule
\end{tabular}
\end{table}

\begin{histexample}[The first odd reconstruction]\label{hist:ex:first_odd_reconstruction}
For \(p=3\) and \(n=2\),
\[
S_2(3)=1-9+25=17.
\]
Since \(Z_1=1\), \(Z_3=2\), and \(S_0(3)=1\), the odd coordinate formula gives
\[
E_2
=17-2\cdot9
=-1.
\]
The large source value and the depth one correction are both necessary.
\end{histexample}

\begin{histexample}[The first nontrivial dyadic row]\label{hist:ex:first_dyadic_reconstruction}
For \(n=4\),
\[
R_4=\frac{91}{15},
\qquad
R_2=\frac53,
\qquad
R_0=1.
\]
The exact kernel coefficients are
\[
\kappa_{2,1}=-\frac83,
\qquad
\kappa_{2,2}=\frac{128}{5}.
\]
Therefore
\[
E_4
=
\frac{91}{15}
-\frac83\binom42\frac53
+\frac{128}{5}
=5.
\]
The rational source values are dyadic units, while the transfer depth is carried by the kernel.
\end{histexample}

\begin{histremark}\label{hist:rem:no_dyadic_substitution}
The common convolution does not identify the integral charts. At an odd prime the finite character coordinate has order \(p-1\). At the dyadic prime it begins with the sign coordinate, and the principal unit subgroup begins at \(1+4\mathbb Z_2\).
\end{histremark}

\subsubsection{Exact dyadic transfer depth}

For \(r\geq1\), write
\[
\kappa_{2,r}
=
\beta_r2^{2r}(2r)!
=
-2^{4r}(2r-1)B_{2r}.
\]

\begin{histtheorem}[Exact dyadic kernel weight]\label{hist:thm:exact_dyadic_kernel_weight}
\par\noindent\textbf{Editorial index conventions.} The standalone kernel valuation holds for all $r\ge1$. Its reconstruction-term formula requires even $n\ge0$ and $1\le r\le n/2$. In the following truncation formulas only terms with $r\le n/2$ occur; terms that would use $R_{n-2r}$ at a negative index are omitted. An empty tail is exactly zero, and the minimum over an empty set of term valuations is $+\infty$. These conventions include $M\ge n/2$.\par\medskip
For every \(r\geq1\),
\[
\boxed{\vTwo(\kappa_{2,r})=4r-1.}
\]
The \(r\)th reconstruction term has exact valuation
\[
\boxed{
\vTwo\left(
\kappa_{2,r}\binom n{2r}R_{n-2r}
\right)
=
4r-1+\vTwo\binom n{2r}.
}
\]
Equivalently,
\[
\boxed{
W_2(n,r)
=
4r-1+\sTwo(r)+\sTwo(n-2r)-\sTwo(n).
}
\]
\end{histtheorem}

\begin{proof}
The Bernoulli formula gives \(\vTwo(B_{2r})=-1\). The source value is a unit. Kummer's binary carry formula gives the last expression.
\end{proof}

\begin{histtheorem}[Optimal dyadic truncation]\label{hist:thm:optimal_dyadic_truncation}
Let
\[
T_{2,M}(n)
=
\sum_{r=0}^{\min(M,n/2)}
\kappa_{2,r}\binom n{2r}R_{n-2r}.
\]
Then
\[
\boxed{
E_n\equiv T_{2,M}(n)\pmod {2^{4M+3}}.
}
\]
More precisely,
\[
\boxed{
\vTwo\bigl(E_n-T_{2,M}(n)\bigr)
\geq
\min_{M+1\leq r\leq n/2}
\left(
4r-1+\vTwo\binom n{2r}
\right).
}
\]
The uniform exponent \(4M+3\) is best possible.
\end{histtheorem}

\begin{proof}
Every omitted term has valuation at least \(4(M+1)-1\). For optimality, take \(n=2M+2\). There is one omitted term. Its binomial coefficient and source value are one, so its valuation is exactly \(4M+3\).
\end{proof}

The first three sharp truncations are
\[
\begin{aligned}
E_n&\equiv R_n&&\pmod 8,\\
E_n&\equiv R_n-\frac83\binom n2R_{n-2}&&\pmod {128},\\
E_n&\equiv R_n-\frac83\binom n2R_{n-2}
+\frac{128}{5}\binom n4R_{n-4}&&\pmod {2048}.
\end{aligned}
\]

\begin{figure}[t]
\centering
\includegraphics[width=.95\linewidth,height=3.15in,keepaspectratio]{../../figures/reconstructed/local_kernel_weights.pdf}
\caption{New reconstruction. Odd and dyadic kernels share one convolution law but have different valuation profiles. The positive depth dyadic kernel has exact slope four and offset minus one.}
\label{hist:fig:local_kernel_weights}
\end{figure}

\begin{figure}[t]
\centering
\includegraphics[width=.95\linewidth,height=3.15in,keepaspectratio]{../../figures/reconstructed/odd_dyadic_geometry.pdf}
\caption{New reconstruction. The local Euler charts share source and kernel reconstruction. Their torsion addresses, principal unit depths, source corrections, and precision growth remain different.}
\label{hist:fig:odd_dyadic_geometry}
\end{figure}
\subsection{Mahler transfer and local analytic branches}\label{hist:sec:mahler}

The finite character coordinate resolves the torsion part of a local unit. A sufficiently long arithmetic jump moves the remaining unit into a principal unit neighborhood. The coarse variable then has a convergent Mahler expansion with an explicit valuation slope.

\subsubsection{Finite Mahler coordinates}

For a function \(f\) on the nonnegative integers, define
\[
\Delta f(n)=f(n+1)-f(n).
\]

\begin{histtheorem}[Finite Mahler reconstruction]\label{hist:thm:finite_mahler_unified}
Given \(f(0),\ldots,f(N)\), put \(c_r=\Delta^rf(0)\). Then
\[
\boxed{
f(n)=\sum_{r=0}^{n}c_r\binom nr
\qquad(0\leq n\leq N).
}
\]
For every \(s\) with \(0\leq s\leq N\),
\[
\boxed{
\Delta^sf(n)
=
\sum_{r=s}^{N}c_r\binom n{r-s}
\qquad(0\leq n\leq N-s).
}
\]
\end{histtheorem}

\begin{proof}
Let
\[
P(x)=\sum_{r=0}^{N}c_r\binom xr.
\]
Newton interpolation gives \(P(n)=f(n)\) for \(0\leq n\leq N\). Iterating
\[
\Delta\binom xr=\binom x{r-1}
\]
gives
\[
\Delta^sP(n)=\sum_{r=s}^{N}c_r\binom n{r-s}.
\]
When \(0\leq n\leq N-s\), every value used in \(\Delta^sf(n)\) lies among the prescribed data, so \(\Delta^sf(n)=\Delta^sP(n)\).
\end{proof}

\begin{histcorollary}[Mahler valuation ledger]\label{hist:cor:mahler_valuation_ledger}
If \(c_r\in\Zp\) and \(\vp(c_r)\geq w(r)\), then
\[
\boxed{
\vp(\Delta^sf(n))
\geq
\min_{r\geq s}
\left(
w(r)+\vp\binom n{r-s}
\right).
}
\]
\end{histcorollary}

\begin{proof}
Apply the nonarchimedean triangle inequality to \cref{hist:thm:finite_mahler_unified}.
\end{proof}

The infinite interpolation theorem states that every continuous function on \(\Zp\) has a unique convergent Mahler expansion \cite{Mahler1958,Robert2000}. We use that theorem as a standard input when the local variable is allowed to range over all of \(\Zp\).

\subsubsection{Carmichael and Mahler transfer for unit modes}

\begin{histstandardfact}[Local unit decomposition]\label{hist:fact:local_unit_decomposition}
For odd \(p\), every unit has a unique factorization
\[
u=\omega(u)\langle u\rangle,
\qquad
\omega(u)^{p-1}=1,
\qquad
\langle u\rangle\in1+p\mathbb Z_p.
\]
Every dyadic unit has a unique factorization
\[
u=(-1)^{\epsilon(u)}\langle u\rangle,
\qquad
\epsilon(u)\in\{0,1\},
\qquad
\langle u\rangle\in1+4\mathbb Z_2.
\]
The first factor gives the finite character address. The second factor admits logarithmic and Mahler coordinates.
\end{histstandardfact}

For odd \(p\) and \(L\geq1\), put
\[
\Lambda_{p,L}=(p-1)p^{L-1}.
\]
For \(p=2\) and \(L\geq3\), put
\[
\Lambda_{2,L}=2^{L-2}.
\]
These are exponents of the unit groups modulo the declared prime powers.

\begin{histtheorem}[One unit mode]\label{hist:thm:carmichael_mahler_scalar}
Let \(u\in\Zp^\times\), let \(H=\Lambda_{p,L}\), and put
\[
q_{u,L}=u^H-1\in p^L\Zp.
\]
For every integer address \(j\) and every \(x\in\Zp\),
\[
\boxed{
u^{Hx+j}
=
u^j\sum_{r\geq0}q_{u,L}^{\,r}\binom xr.
}
\]
The \(r\)th Mahler coefficient has valuation at least \(Lr\).
\end{histtheorem}

\begin{proof}
Write \(u^{Hx+j}=u^j(1+q_{u,L})^x\) and use the convergent binomial series on \(\Zp\).
\end{proof}

\begin{histexample}[Paired principal unit expansions]\label{hist:ex:paired_unit_expansions}
At the odd prime three,
\[
2^{2x+j}
=
2^j4^x
=
2^j\sum_{r\geq0}3^r\binom xr.
\]
At the dyadic prime,
\[
3^{2x+j}
=
3^j9^x
=
3^j\sum_{r\geq0}8^r\binom xr.
\]
The formal expansions are identical. Their first local slopes are one and three because \(4-1=3\) while \(9-1=8\).
\end{histexample}

\begin{histcorollary}[Basic odd and dyadic difference gains]\label{hist:cor:basic_local_difference_gains}
For odd \(p\),
\[
\boxed{
\Delta_{(p-1)p^a}^{\,r}u^n
=
u^n\left(u^{(p-1)p^a}-1\right)^r,
\qquad
\vp\left(\Delta_{(p-1)p^a}^{\,r}u^n\right)\geq(a+1)r.
}
\]
For odd \(u\in\mathbb Z_2^\times\),
\[
\boxed{
\Delta_{2^{a+1}}^{\,r}u^n
=
u^n\left(u^{2^{a+1}}-1\right)^r,
\qquad
\vTwo\left(\Delta_{2^{a+1}}^{\,r}u^n\right)\geq(a+3)r.
}
\]
\end{histcorollary}

\begin{proof}
One step difference multiplies a unit mode by \(u^H-1\). Fermat reduction and lifting the exponent give the odd estimate. For odd \(u\),
\[
\vTwo\left(u^{2^{a+1}}-1\right)
=
\vTwo(u^2-1)+a
\geq a+3.
\]
\end{proof}

\subsubsection{Polynomial exponential modes}

\begin{histtheorem}[Polynomial exponential Mahler transfer]\label{hist:thm:polynomial_exponential_mahler}
Let
\[
f(n)=P(n)u^n
\]
with \(P\in\Zp[n]\) of degree \(d\). Let \(H=\Lambda_{p,L}\), fix an address \(j\), and write
\[
Q_j(x)=P(Hx+j),
\qquad
y=u^H,
\qquad q=y-1.
\]
Then
\[
\boxed{
f(Hx+j)=\sum_{r\geq0}c_{j,r}\binom xr,
}
\]
where
\[
\boxed{
c_{j,r}
=
u^j
\sum_{k=0}^{\min(d,r)}
\binom rk
\Delta^kQ_j(0)y^kq^{r-k}.
}
\]
Consequently,
\[
\boxed{
\vp(c_{j,r})\geq L\max(r-d,0).
}
\]
For \(R\geq d\), truncation at depth \(R\) is valid modulo \(p^{L(R+1-d)}\) uniformly on \(\Zp\).
\end{histtheorem}

\begin{proof}
On the polynomial factor the difference operator acts as \(q+y\Delta\). The two terms commute. Expand \((q+y\Delta)^r\) and evaluate at zero. Since \(\Delta^kQ_j=0\) for \(k>d\), every surviving term contains at least \(q^{r-d}\).
\end{proof}

\subsubsection{Matrix state systems and finite adelic synchronization}

Let
\[
v_{n+1}=Cv_n,
\qquad
a_n=w^TC^nv_0,
\]
with \(C\in M_d(\Zp)\) invertible modulo \(p\).

\begin{histtheorem}[Matrix Carmichael and Mahler transfer]\label{hist:thm:matrix_carmichael_mahler}
\par\noindent\textbf{Editorial hypothesis correction.} The original displayed statement leaves the vectors untyped. Its integrality conclusion requires $v_0,w\in\Z_p^d$, $L\ge1$, and $H,j,n\in\N$ with $H>0$; the value at $Hx+j$ for $x\in\Z_p$ means the matrix-binomial interpolation. With $p=3$, $C=4$, $L=1$, $w=1$ and $v_0=1/3$, one has $\Delta_1a_0=1$, which is not in $3\Z_3$. The printed claim without these restrictions is therefore not an as-written certificate. The full original proof follows under the corrected hypotheses.\par\medskip
Assume
\[
C^H=I+p^LB
\]
with \(B\in M_d(\Zp)\). For every address \(j\) and every \(x\in\Zp\),
\[
\boxed{
a_{Hx+j}
=
\sum_{r\geq0}
p^{Lr}\binom xr
w^TC^jB^rv_0.
}
\]
The series converges on \(\Zp\), and
\[
\boxed{
\Delta_H^ra_n
\equiv0
\pmod {p^{Lr}}.
}
\]
\end{histtheorem}

\begin{proof}
Use
\[
C^{Hx+j}=C^j(I+p^LB)^x
\]
and the matrix binomial series. One difference in the coarse variable multiplies the series by \(p^LB\).
\end{proof}

\begin{histtheorem}[Finite adelic transfer]\label{hist:thm:finite_adelic_transfer}
\par\noindent\textbf{Editorial hypothesis correction.} Read this statement with $v_0,w\in\Z^d$, $H>0$ and $j,n,r\in\N$, and positive integer $L_p$. These restrictions make the integral congruences and powers well typed. Allowing $v_0=1/3$ in the one-prime example above loses the advertised digit; invertibility of $C$ alone does not restore it. The original proof is retained under these hypotheses.\par\medskip
Let
\[
Q=\prod_{p\in S}p^{L_p}
\]
for a finite set of primes. Let \(C\in M_d(\ZZ)\) be invertible modulo \(Q\), and choose \(H\) such that
\[
C^H=I+QB
\]
with \(B\in M_d(\ZZ)\). Then for every nonnegative integer \(x\),
\[
\boxed{
a_{Hx+j}
=
\sum_{r=0}^{x}Q^r\binom xr
w^TC^jB^rv_0.
}
\]
For every \(r\),
\[
\boxed{
\Delta_H^ra_n\equiv0\pmod {Q^r}.
}
\]
\end{histtheorem}

\begin{proof}
Apply the integer matrix binomial theorem and reduce at every selected prime.
\end{proof}

\begin{histexample}[A simultaneous Fibonacci jump]\label{hist:ex:fibonacci_finite_adelic}
Let
\[
C=
\begin{pmatrix}
1&1\\
1&0
\end{pmatrix}.
\]
The matrix \(C\) is invertible modulo twenty four, and
\[
C^{24}-I
=
24
\begin{pmatrix}
3126&1932\\
1932&1194
\end{pmatrix}.
\]
For every scalar coordinate of the state sequence \(C^nv_0\),
\[
\boxed{
\Delta_{24}^{\,r}a_n
\equiv0
\pmod {24^r}.
}
\]
One integral jump therefore gives simultaneous congruence ladders at the primes two and three.
\end{histexample}

\subsubsection{Amice pullback}

\begin{histdefinition}[Amice pullback]\label{hist:def:amice_pullback_unified}
For a formal exponential generating series
\[
F(z)=\sum_{n\geq0}a_n\frac{z^n}{n!},
\]
define
\[
\widehat F(T)=F(\log(1+T)).
\]
\end{histdefinition}

\begin{histtheorem}[Derivative conjugacy]\label{hist:thm:amice_derivative_unified}
\[
\boxed{
\widehat{F'}(T)
=
(1+T)\frac{d}{dT}\widehat F(T).
}
\]
\end{histtheorem}

\begin{proof}
Differentiate the composition \(F(\log(1+T))\).
\end{proof}

\begin{histproposition}[Weighted convergence]\label{hist:prop:amice_weighted_convergence_unified}
If
\[
\widehat F(T)=\sum_{r\geq0}c_rT^r
\]
and
\[
\vp(c_r)+r\alpha\longrightarrow+\infty,
\]
then the series converges whenever \(\vp(T)\geq\alpha\).
\end{histproposition}

\begin{proof}
The valuation of the \(r\)th term is \(\vp(c_r)+r\vp(T)\), which tends to positive infinity.
\end{proof}

\subsubsection{Odd prime depletion}

For an odd prime \(p\), the source \(S_n(p)\) contains one nonunit base. Put
\[
\sigma_p=(-1)^{(p-1)/2}
\]
and
\[
S_n^\times(p)=S_n(p)-\sigma_pp^n.
\]

\begin{histtheorem}[Depleted odd prime reconstruction]\label{hist:thm:depleted_odd_reconstruction}
For every even \(n\),
\[
\boxed{
\left(1-\sigma_pp^n\right)E_n
=
\sum_{r=0}^{n/2}
(-1)^rZ_{2r+1}
\binom n{2r}p^{2r}S_{n-2r}^\times(p).
}
\]
\end{histtheorem}

\begin{proof}
The removed source mode contributes \(\sigma_p\cosh(pt)\). Multiplication by \(\operatorname{sech}^2(pt)\) gives \(\sigma_p\operatorname{sech}(pt)\). Subtract it from the full source and kernel product.
\end{proof}

Fix an even address \(j\in\{0,2,\ldots,p-3\}\). For \(x\in\Zp\), define
\[
\begin{aligned}
\mathscr E_{p,j}^\times(x)
=
\sum_{r\geq0}
&(-1)^rZ_{2r+1}p^{2r}
\binom{(p-1)x+j}{2r}\\
&\times
\sum_{\substack{0\leq a<p\\2a+1\neq p}}
(-1)^a
(2a+1)^{j-2r}
\left((2a+1)^{p-1}\right)^x.
\end{aligned}
\]

\begin{histtheorem}[Odd prime analytic Euler branch]\label{hist:thm:odd_analytic_branch}
The series defining \(\mathscr E_{p,j}^\times\) converges in the Tate algebra on \(\Zp\). For every nonnegative integer \(x\),
\[
\boxed{
\mathscr E_{p,j}^\times(x)
=
\left(1-\sigma_pp^{(p-1)x+j}\right)
E_{(p-1)x+j}.
}
\]
\end{histtheorem}

\begin{proof}
Every retained base is a unit, and its \((p-1)\)st power lies in \(1+p\Zp\). The polynomial
\[
p^{2r}\binom{(p-1)x+j}{2r}
\]
has Gauss norm at most
\[
p^{-2r+\vp((2r)!)}.
\]
This tends to zero because \(p\) is odd. The remaining factors have bounded Gauss norm. Thus the series converges in the Tate algebra. At a nonnegative integer \(x\), the binomial factor vanishes for large \(r\), and \cref{hist:thm:depleted_odd_reconstruction} gives the value.
\end{proof}

\begin{histremark}\label{hist:rem:depleted_factor_boundary}
The analytic function is defined by its unit source series. The expression \(p^{(p-1)x+j}\) in the interpolation identity is evaluated only at nonnegative integers. No analytic function \(x\mapsto p^x\) on all of \(\Zp\) is used.
\end{histremark}

\subsubsection{Direct dyadic interpolation}

For \(x\in\mathbb Z_2\), define
\[
\boxed{
\mathscr R(x)
=
\frac{9^{x+1}-1}{8(x+1)(2x+1)}.
}
\]
The apparent singularity at \(x=-1\) is removable, with value \(-\log(9)/8\).

\begin{histlemma}[Analytic dyadic source]\label{hist:lem:analytic_dyadic_source}
The function \(\mathscr R\) belongs to the dyadic Tate algebra and takes values in \(\mathbb Z_2\). For every nonnegative integer \(n\),
\[
\boxed{\mathscr R(n)=R_{2n}.}
\]
\end{histlemma}

\begin{proof}
Put \(y=x+1\) and let \(\ell=\log 9\). Since \(\vTwo(\ell)=3\),
\[
\frac{9^y-1}{8y}
=
\frac{\ell}{8}
\sum_{m\geq0}
\frac{\ell^m}{(m+1)!}y^m.
\]
The coefficient of \(y^m\) has valuation
\[
3m-\vTwo((m+1)!),
\]
which is nonnegative and tends to positive infinity. The quotient therefore belongs to \(\mathbb Z_2\langle y\rangle\). The factor \(2x+1=1+2x\) is a unit in \(\mathbb Z_2\langle x\rangle\), with inverse \(\sum_{m\geq0}(-2x)^m\). Hence \(\mathscr R\) is integral and analytic. The integer values follow from the definition of \(R_{2n}\).
\end{proof}

Define
\[
\boxed{
\mathscr E_2(x)
=
\sum_{r\geq0}
\kappa_{2,r}
\binom{2x}{2r}
\mathscr R(x-r).
}
\]

\begin{histtheorem}[Direct dyadic Euler interpolation]\label{hist:thm:direct_dyadic_interpolation}
The series defining \(\mathscr E_2\) converges in the dyadic Tate algebra. For every nonnegative integer \(n\),
\[
\boxed{\mathscr E_2(n)=E_{2n}.}
\]
Truncation after depth \(M\) is uniformly valid modulo \(2^{4M+3}\) on \(\mathbb Z_2\).
\end{histtheorem}

\begin{proof}
The monomial Gauss norm of \(\binom{2x}{2r}\) is at most \(2^{\vTwo((2r)!)}\). Since
\[
\vTwo(\kappa_{2,r})=4r-1,
\]
the product tends to zero in Gauss norm. Translation by the integer \(r\) preserves the unit disc, so \(\mathscr R(x-r)\) remains an integral analytic function with uniformly bounded supremum norm. The series therefore converges in the Tate algebra.

At a nonnegative integer \(n\), every term with \(r>n\) vanishes and \cref{hist:thm:dyadic_coordinate_pair} remains. For arbitrary \(x\in\mathbb Z_2\), the integer valued polynomial \(\binom{2x}{2r}\) and the source value \(\mathscr R(x-r)\) lie in \(\mathbb Z_2\). Every omitted term after depth \(M\) is therefore divisible by \(2^{4M+3}\). This proves the uniform truncation statement.
\end{proof}

\begin{histtheorem}[Unified local interpolation statement]\label{hist:thm:unified_local_interpolation}
The odd branch \(\mathscr E_{p,j}^\times\) and the dyadic branch \(\mathscr E_2\) are analytic local realizations of the source and kernel law in \cref{hist:thm:unified_local_euler}. At integer arguments they recover, respectively,
\[
\left(1-\sigma_pp^n\right)E_n
\]
on each even odd prime character branch and
\[
E_{2n}
\]
on the dyadic even branch. The two realizations have distinct source corrections and distinct kernel valuation slopes.
\end{histtheorem}

\begin{proof}
Combine \cref{hist:thm:odd_analytic_branch,hist:thm:direct_dyadic_interpolation} with the source and kernel definitions in \cref{hist:thm:unified_local_euler}.
\end{proof}

\subsection{Polyphase refinement of local Mahler coordinates}\label{hist:sec:polyphase}

Polyphase refinement replaces one local branch by finitely many digit branches while preserving exact synthesis. For sequences it is residue analysis. For local functions it is the digit tree of \(\mathbb Z_p\). In Mahler coordinates it becomes a finite triangular kernel.

\subsubsection{Discrete analysis and exact synthesis}

For a sequence \(a=(a_n)_{n\geq0}\) and \(M\geq2\), define
\[
a_r^{[M]}(q)=a_{Mq+r}
\qquad(0\leq r<M).
\]
The analysis map is
\[
\mathcal P_Ma
=
\left(a_0^{[M]},\ldots,a_{M-1}^{[M]}\right).
\]
Synthesis interleaves the branches by the rule
\[
a_{Mq+r}=a_r^{[M]}(q).
\]

\begin{histtheorem}[Exact discrete polyphase reconstruction]\label{hist:thm:discrete_polyphase_reconstruction}
The analysis map \(\mathcal P_M\) is a bijection between sequences and ordered families of \(M\) sequences. If
\[
A(z)=\sum_{n\geq0}a_nz^n
\]
and
\[
A_r^{[M]}(w)=\sum_{q\geq0}a_{Mq+r}w^q,
\]
then
\[
\boxed{
A(z)=\sum_{r=0}^{M-1}z^rA_r^{[M]}(z^M).
}
\]
\end{histtheorem}

\begin{proof}
Every nonnegative integer has a unique expression \(Mq+r\) with \(0\leq r<M\). Group the terms of the ordinary generating function by this residue.
\end{proof}

\begin{histtheorem}[Nested discrete refinement]\label{hist:thm:nested_discrete_polyphase}
For positive integers \(M,N\),
\[
\boxed{
\left(a_r^{[M]}\right)_s^{[N]}
=
a_{r+Ms}^{[MN]}.
}
\]
For a mixed radix path \((M_1,\ldots,M_\ell)\) with digits \(0\leq s_i<M_i\), the final address is
\[
\boxed{
r(\boldsymbol s)
=
s_1+M_1s_2+\cdots+
\left(\prod_{i=1}^{\ell-1}M_i\right)s_\ell.
}
\]
Only the product radix and accumulated address determine the final branch.
\end{histtheorem}

\begin{proof}
Direct substitution gives
\[
\left(a_r^{[M]}\right)_s^{[N]}(q)
=
a_{M(Nq+s)+r}.
\]
Iterate this identity along the mixed radix path.
\end{proof}

\subsubsection{Carry convolution and the shift matrix}

Let \(c=a*b\) be the ordinary Cauchy convolution. Write \(A_r\), \(B_r\), and \(C_r\) for the phase ordinary generating functions.

\begin{histproposition}[One carry convolution law]\label{hist:prop:polyphase_carry_convolution}
For every residue \(r\),
\[
\boxed{
C_r(w)
=
\sum_{u+v=r}A_u(w)B_v(w)
+
w\sum_{u+v=r+M}A_u(w)B_v(w).
}
\]
The first sum has no phase carry. The second has one coarse index carry.
\end{histproposition}

\begin{proof}
Write the input indices as \(Mq+u\) and \(Mq'+v\). Since \(0\leq u,v<M\), their sum has residue \(u+v\) when \(u+v<M\), and residue \(u+v-M\) with one additional coarse index when \(u+v\geq M\).
\end{proof}

\begin{histtheorem}[Polyphase shift matrix]\label{hist:thm:polyphase_shift_matrix}
Let \(E\) be the scalar shift \((Ea)_n=a_{n+1}\), and let \(E_q\) shift the coarse index. Then
\[
\boxed{
\mathcal P_M(Ea)
=
\mathsf S_M(E_q)\mathcal P_Ma,
}
\]
where
\[
\mathsf S_M(E_q)
=
\begin{pmatrix}
0&1&0&\cdots&0\\
0&0&1&\cdots&0\\
\vdots&&&\ddots&\vdots\\
0&0&0&\cdots&1\\
E_q&0&0&\cdots&0
\end{pmatrix}.
\]
It also satisfies
\[
\boxed{\mathsf S_M(E_q)^M=E_qI_M.}
\]
\end{histtheorem}

\begin{proof}
For \(r<M-1\),
\[
(Ea)_{Mq+r}=a_{Mq+r+1}=a_{r+1}^{[M]}(q).
\]
For \(r=M-1\),
\[
(Ea)_{Mq+M-1}=a_{M(q+1)}=E_qa_0^{[M]}(q).
\]
After \(M\) shifts every phase has advanced by one coarse index.
\end{proof}

\subsubsection{Local branches and digit reconstruction}

For \(a\geq1\) and \(0\leq s<p^a\), define
\[
f_{a,s}(x)=f(p^ax+s).
\]

\begin{histtheorem}[Exact local branch reconstruction]\label{hist:thm:local_branch_reconstruction}
The family \(\{f_{a,s}\}_{0\leq s<p^a}\) determines \(f\) uniquely on \(\mathbb Z_p\). More precisely, if \(y\in s+p^a\mathbb Z_p\), then
\[
\boxed{
f(y)=f_{a,s}\left(\frac{y-s}{p^a}\right).
}
\]
For \(a,b\geq1\),
\[
\boxed{
(f_{a,s})_{b,t}
=
f_{a+b,s+p^at}.
}
\]
Thus local refinement is exactly the \(p\) adic digit tree.
\end{histtheorem}

\begin{proof}
The cosets \(s+p^a\mathbb Z_p\) form a disjoint partition of \(\mathbb Z_p\). The first formula is the inverse affine change of variables on one coset. The second follows from
\[
p^a(p^bx+t)+s
=
p^{a+b}x+p^at+s.
\]
\end{proof}

\subsubsection{The affine Mahler kernel}

Let
\[
\Phi_M(z)=(1+z)^M-1.
\]
For \(M\geq1\), \(s\geq0\), and \(r,k\geq0\), define
\[
\boxed{
\Theta_{M,s}(r,k)
=
[z^r](1+z)^s\Phi_M(z)^k.
}
\]
Equivalently,
\[
\boxed{
\Theta_{M,s}(r,k)
=
\sum_{j=0}^{k}
(-1)^{k-j}\binom kj\binom{Mj+s}{r}.
}
\]
It vanishes unless
\[
k\leq r\leq Mk+s.
\]

\begin{histproposition}[Edge coefficients of the affine kernel]\label{hist:prop:affine_kernel_edges}
For every \(M\geq1\), \(s\geq0\), and \(k\geq0\),
\[
\boxed{
\Theta_{M,s}(k,k)=M^k,
\qquad
\Theta_{M,s}(Mk+s,k)=1.
}
\]
For binary radix and zero address,
\[
\boxed{
\Theta_{2,0}(k+j,k)
=
2^{k-j}\binom kj
\qquad(0\leq j\leq k).
}
\]
Hence
\[
\boxed{
\vTwo\left(\Theta_{2,0}(k+j,k)\right)
=
k-j+\vTwo\binom kj.
}
\]
\end{histproposition}

\begin{proof}
The lowest degree term of \(\Phi_M(z)\) is \(Mz\), and its highest degree term is \(z^M\). Multiplication by \((1+z)^s\) gives the two edge coefficients. When \(M=2\) and \(s=0\),
\[
\Phi_2(z)^k
=
(2z+z^2)^k
=
z^k(2+z)^k.
\]
Taking the coefficient of \(z^{k+j}\) gives the binary formula.
\end{proof}

\begin{histtheorem}[Mahler polyphase transfer]\label{hist:thm:mahler_polyphase_transfer}
Let
\[
f(x)=\sum_{r\geq0}c_r\binom xr
\]
be a convergent Mahler expansion. For
\[
g_s(q)=f(Mq+s),
\]
one has
\[
\boxed{
g_s(q)=\sum_{k\geq0}d_{s,k}\binom qk,}
\]
where
\[
\boxed{
d_{s,k}
=
\sum_{r=k}^{Mk+s}
c_r\Theta_{M,s}(r,k).
}
\]
Every refined coefficient uses a finite window of the original Mahler coefficients.
\end{histtheorem}

\begin{proof}
Use
\[
(1+z)^{Mq+s}
=
(1+z)^s\left(1+\Phi_M(z)\right)^q
=
\sum_{k\geq0}\binom qk(1+z)^s\Phi_M(z)^k.
\]
Taking the coefficient of \(z^r\) gives the affine expansion of \(\binom{Mq+s}{r}\). Since \(c_r\to0\) and every \(d_{s,k}\) uses indices \(r\geq k\), the transformed coefficients also tend to zero.
\end{proof}

\begin{histtheorem}[Kernel composition]\label{hist:thm:polyphase_kernel_composition}
For positive integers \(M,N\),
\[
\boxed{
\sum_{k\geq0}
\Theta_{M,s}(r,k)
\Theta_{N,t}(k,\ell)
=
\Theta_{MN,Mt+s}(r,\ell).
}
\]
The sum is finite.
\end{histtheorem}

\begin{proof}
Apply \cref{hist:thm:mahler_polyphase_transfer} first to \(x=Mq+s\) and then to \(q=Ny+t\). The composite substitution is \(x=MNy+Mt+s\). Uniqueness of Mahler coefficients gives the identity.
\end{proof}

\subsubsection{Valuation transport}

\begin{histtheorem}[Polyphase valuation ledger]\label{hist:thm:polyphase_valuation_ledger}
For every refined coefficient,
\[
\boxed{
\vp(d_{s,k})
\geq
\min_{k\leq r\leq Mk+s}
\left(
\vp(c_r)+\vp(\Theta_{M,s}(r,k))
\right).
}
\]
If \(\vp(c_r)\geq w(r)\) for a nondecreasing function \(w\), then
\[
\boxed{\vp(d_{s,k})\geq w(k).}
\]
Thus every monotone Mahler precision ladder survives arbitrary polyphase refinement.
\end{histtheorem}

\begin{proof}
Apply the nonarchimedean triangle inequality to the finite coefficient formula. Since \(r\geq k\), monotonicity gives \(w(r)\geq w(k)\), and the affine kernel is integral.
\end{proof}

\begin{histtheorem}[Prime power sparsification]\label{hist:thm:prime_power_sparsification}
\par\noindent\textbf{Editorial hypothesis correction.} The affine-kernel congruence below needs no change, but its consequence for transformed Mahler coefficients requires $c_r\in\Z_p$ for every $r$. Convergence of a $\Q_p$-valued Mahler expansion alone is insufficient. For $p=2$, $M=2$, $s=0$ and $f(x)=x/2$, one has $c_1=1/2$, $c_2=0$, while $g(q)=q$ has $d_1=1$; thus $d_1\not\equiv c_2\pmod2$. The coefficient consequence is read with integral original coefficients.\par\medskip
Let \(M=p^a\). Then
\[
\boxed{
\Theta_{p^a,s}(r,k)
\equiv
\binom{s}{r-p^ak}
\pmod p,
}
\]
where the binomial coefficient is zero outside its natural range. Consequently,
\[
\boxed{
d_{s,k}
\equiv
\sum_{u=0}^{s}
\binom su
c_{p^ak+u}
\pmod p.
}
\]
In particular,
\[
\boxed{d_{0,k}\equiv c_{p^ak}\pmod p.}
\]
\end{histtheorem}

\begin{proof}
In characteristic \(p\),
\[
(1+z)^{p^a}-1\equiv z^{p^a}.
\]
Substitute this congruence in the definition of \(\Theta_{p^a,s}\).
\end{proof}

\begin{figure}[t]
\centering
\includegraphics[width=.95\linewidth,height=3.15in,keepaspectratio]{../../figures/reconstructed/mahler_polyphase_kernel.pdf}
\caption{New reconstruction. The affine Mahler kernel has finite triangular support. At a prime power radix, its reduction is confined to a short Frobenius window.}
\label{hist:fig:mahler_polyphase_kernel}
\end{figure}

\subsubsection{Precision amplification}

\begin{histtheorem}[Scalar polyphase precision amplification]\label{hist:thm:scalar_polyphase_precision}
Let \(u\in\mathbb Z_p^\times\), let \(a,r\geq0\), and put
\[
H_a=
\begin{cases}
(p-1)p^a,&p\text{ odd},\\
2^{a+1},&p=2.
\end{cases}
\]
Then
\[
\boxed{
\Delta_{H_a}^{\,r}u^n
=
u^n\left(u^{H_a}-1\right)^r.
}
\]
For odd \(p\),
\[
\boxed{
\vp\left(\Delta_{H_a}^{\,r}u^n\right)
\geq(a+1)r.
}
\]
For \(p=2\) and odd \(u\),
\[
\boxed{
\vTwo\left(\Delta_{H_a}^{\,r}u^n\right)
\geq(a+3)r.
}
\]
Equivalently, every additional prime power digit increases the certified Mahler slope by one.
\end{histtheorem}

\begin{proof}
One difference with step \(H_a\) multiplies the unit mode by \(u^{H_a}-1\). Iterate. The valuation estimates are those of \cref{hist:cor:basic_local_difference_gains}.
\end{proof}

\begin{histlemma}[Matrix power amplification]\label{hist:lem:matrix_power_amplification}
\par\noindent\textbf{Editorial hypothesis correction.} Require $L\ge1$. At $L=0$ the standalone claim is false: for $p=3$, $d=1$ and $X=2=1+B$, one has $X^3-1=7\notin3\Z_3$. With $L\ge1$, every higher binomial term in the retained proof has the required valuation, including $p=2$.\par\medskip
If
\[
X=I+p^LB
\]
with \(B\in M_d(\Zp)\), then for every \(a\geq0\) there is \(B_a\in M_d(\Zp)\) such that
\[
\boxed{X^{p^a}=I+p^{L+a}B_a.}
\]
\end{histlemma}

\begin{proof}
Induct on \(a\). Expand the \(p\)th power of \(I+p^{L+a}B_a\). The linear term has valuation at least \(L+a+1\), and every higher term has at least that valuation.
\end{proof}

\begin{histtheorem}[Matrix polyphase precision amplification]\label{hist:thm:matrix_polyphase_precision}
\par\noindent\textbf{Editorial hypothesis correction.} Require $L\ge1$, $C,B\in M_d(\Z_p)$, $v_0,w\in\Z_p^d$, $H>0$, and $a,j,n,r\in\N$. These are the hypotheses of the two corrected matrix statements used by the retained proof. Without integral observation vectors or positive $L$, the preceding counterexamples invalidate the claimed precision.\par\medskip
Assume
\[
C^H=I+p^LB.
\]
Put \(H_a=Hp^a\). Then
\[
C^{H_a}=I+p^{L+a}B_a
\]
for an integral matrix \(B_a\). Every refined packet
\[
b_{a,j}(x)=a_{H_ax+j}
\]
has Mahler coefficients satisfying
\[
\boxed{\vp(c_{a,j,r})\geq(L+a)r.}
\]
Equivalently,
\[
\boxed{
\Delta_{H_a}^{\,r}a_n
\equiv0
\pmod {p^{(L+a)r}}.
}
\]
\end{histtheorem}

\begin{proof}
Combine \cref{hist:lem:matrix_power_amplification,hist:thm:matrix_carmichael_mahler}.
\end{proof}

\begin{histexample}[Refined dyadic Fibonacci ladder]\label{hist:ex:refined_fibonacci_ladder}
For the Fibonacci companion matrix
\[
C=
\begin{pmatrix}
1&1\\
1&0
\end{pmatrix},
\]
one has
\[
C^{12}
=I+8
\begin{pmatrix}
29&18\\
18&11
\end{pmatrix}.
\]
Therefore
\[
\boxed{
\Delta_{12\cdot2^a}^{\,r}F_n
\equiv0
\pmod {2^{(a+3)r}}.
}
\]
Every dyadic refinement doubles the step and adds one guaranteed binary digit at each Mahler depth.
\end{histexample}

\begin{figure}[t]
\centering
\includegraphics[width=.95\linewidth,height=3.15in,keepaspectratio]{../../figures/reconstructed/polyphase_precision_tree.pdf}
\caption{New reconstruction. Every additional prime power refinement digit increases the certified Mahler slope of a principal unit matrix packet by one.}
\label{hist:fig:polyphase_precision_tree}
\end{figure}

\subsubsection{Character and polyphase coordinates}

For odd \(p\), write
\[
\boxed{
n=j+(p-1)(s+p^aq),
}
\]
with \(0\leq j<p-1\) and \(0\leq s<p^a\). The address \(j\) records the finite character. The address \(s\) records the principal unit digit. The variable \(q\) carries the remaining Mahler expansion.

At the dyadic prime, write
\[
\boxed{
n=j+2(s+2^aq),
}
\]
with \(j\in\{0,1\}\) and \(0\leq s<2^a\). The sign coordinate replaces the odd prime character coordinate.

\begin{histproposition}[Local phase lattice]\label{hist:prop:local_phase_lattice}
For a unit mode \(u^n\), the odd prime coordinate gives
\[
\boxed{
u^n
=
u^{j+(p-1)s}
\left(u^{(p-1)p^a}\right)^q.
}
\]
The dyadic coordinate gives
\[
\boxed{
u^n
=
u^{j+2s}
\left(u^{2^{a+1}}\right)^q.
}
\]
Thus finite character phase, local digit address, and Mahler depth are separate coordinates.
\end{histproposition}

\begin{proof}
Substitute the two displayed decompositions of \(n\).
\end{proof}

\begin{histcorollary}[Polyphase refinement of the Euler branches]\label{hist:cor:euler_polyphase_refinement}
The analytic branches in \cref{hist:thm:odd_analytic_branch,hist:thm:direct_dyadic_interpolation} admit exact refinements under \(x=p^aq+s\). Their refined Mahler coefficients are obtained from \(\Theta_{p^a,s}\). The original monotone valuation ladders are preserved, and unit mode or invertible matrix representations gain one certified digit per refinement level.
\end{histcorollary}

\begin{proof}
Apply \cref{hist:thm:mahler_polyphase_transfer,hist:thm:polyphase_valuation_ledger,hist:cor:basic_local_difference_gains,hist:thm:matrix_polyphase_precision}.
\end{proof}

\begin{histremark}\label{hist:rem:polyphase_boundary}
Exact refinement does not imply that every branch is simpler. A simplification theorem requires an additional invariant, such as Frobenius sparsity, valuation gain, spectral collision, or a vanishing transfer block.
\end{histremark}
\subsection{Arithmetic lift obstructions}\label{hist:sec:lifting}

A congruence class may solve an equation to one precision and fail at the next. The correct obstruction is not an arbitrary residual vector. It is the residual vector modulo the image of the declared correction map.

\subsubsection{Linear lifting}

Let \(p\) be a prime. Let
\[
M\in M_{r,d}(\Zp),
\qquad
b\in\Zp^r.
\]
Suppose that \(x_L\in\Zp^d\) satisfies
\[
Mx_L\equiv b\pmod {p^L}
\]
for some \(L\geq1\). Define the reduced residual
\[
\rho_L(x_L)
=
\frac{Mx_L-b}{p^L}
\pmod p
\in\FF_p^r.
\]
Write \(\overline M\) for the reduction of \(M\) modulo \(p\).

\begin{histdefinition}[Linear lift obstruction]\label{hist:def:linear_lift_obstruction}
The linear obstruction at depth \(L\) is
\[
\boxed{
\operatorname{Ob}^{\mathrm{lin}}_L(x_L)
=
[\rho_L(x_L)]
\in
\coker(\overline M).
}
\]
\end{histdefinition}

\begin{histtheorem}[One level linear lifting]\label{hist:thm:linear_lifting}
A class represented by \(x_L\) lifts to a solution modulo \(p^{L+1}\) if and only if
\[
\boxed{
\operatorname{Ob}^{\mathrm{lin}}_L(x_L)=0.
}
\]
More precisely, a correction \(y\in\FF_p^d\) gives a lift
\[
x_{L+1}=x_L+p^Ly
\pmod {p^{L+1}}
\]
if and only if
\[
\boxed{
\overline My=-\rho_L(x_L).
}
\]
When one correction exists, the set of all correction classes is a coset of \(\ker\overline M\).
\end{histtheorem}

\begin{proof}
Modulo \(p^{L+1}\),
\[
M(x_L+p^Ly)-b
=
Mx_L-b+p^LMy.
\]
After division by \(p^L\) and reduction modulo \(p\), the lift condition becomes
\[
\rho_L(x_L)+\overline My=0.
\]
This equation is solvable exactly when the residual lies in the image of \(\overline M\). If \(y_0\) is one solution, every solution is \(y_0+z\) with \(z\in\ker\overline M\).
\end{proof}

\begin{histproposition}[Representative independence]\label{hist:prop:linear_obstruction_independence}
The class \(\operatorname{Ob}^{\mathrm{lin}}_L(x_L)\) depends only on the solution class of \(x_L\) modulo \(p^L\).
\end{histproposition}

\begin{proof}
Replace \(x_L\) by \(x_L+p^Lz\). The reduced residual changes by \(\overline Mz\). Its class in the cokernel is unchanged.
\end{proof}

\begin{histcorollary}[Exact number of linear children]\label{hist:cor:linear_child_count}
If \(\rho=\rank_{\FF_p}\overline M\), then a solvable depth \(L\) class has exactly
\[
\boxed{
p^{d-\rho}
}
\]
children modulo \(p^{L+1}\).
\end{histcorollary}

\begin{proof}
The correction set is a coset of a vector space of dimension \(d-\rho\).
\end{proof}

\begin{histexample}[A visible linear obstruction]\label{hist:ex:linear_lift_obstruction}
Take \(p=3\), \(L=1\), and
\[
M=
\begin{pmatrix}
1&1\\
0&0
\end{pmatrix}.
\]
The class \(x_1=(0,0)\) solves
\[
Mx_1\equiv b\pmod3
\]
for both
\[
b_1=(-3,0)
\qquad\text{and}\qquad
b_2=(0,-3).
\]
For \(b_1\), the reduced residual is \((1,0)\), which belongs to the image of \(\overline M\). There are three correction classes, given by
\[
y_1+y_2=2
\pmod3.
\]
For \(b_2\), the residual is \((0,1)\), whose cokernel class is nonzero. No lift modulo nine exists because the second row of \(M\) remains zero.
\end{histexample}

\begin{figure}[t]
\centering
\includegraphics[width=.95\linewidth,height=3.15in,keepaspectratio]{../../figures/reconstructed/lift_obstruction.pdf}
\caption{New reconstruction. A residual vector lifts precisely when its class vanishes in the cokernel of the reduced correction map.}
\label{hist:fig:lift_obstruction}
\end{figure}

\subsubsection{Polynomial lifting}

Let
\[
F=(F_1,\ldots,F_r):\Zp^d\longrightarrow\Zp^r
\]
be a polynomial map.\par\noindent\textbf{Editorial coefficient clarification.} Here polynomial means $F_i\in\Z_p[X_1,\ldots,X_d]$, not merely an integer-valued polynomial with coefficients in $\Q_p$. This makes the Jacobian and Taylor remainder integral and their reduction well typed.\par Suppose
\[
F(x_L)\in p^L\Zp^r.
\]
Define
\[
\rho_L^F(x_L)
=
\frac{F(x_L)}{p^L}
\pmod p.
\]
Let \(J_F(x_L)\) denote the Jacobian matrix and write \(\overline J_F(x_L)\) for its reduction modulo \(p\).

\begin{histdefinition}[Polynomial lift obstruction]\label{hist:def:polynomial_lift_obstruction}
The polynomial obstruction at depth \(L\) is
\[
\boxed{
\operatorname{Ob}^{F}_L(x_L)
=
[\rho_L^F(x_L)]
\in
\coker\bigl(\overline J_F(x_L)\bigr).
}
\]
\end{histdefinition}

\begin{histtheorem}[One level polynomial lifting]\label{hist:thm:polynomial_lifting}
Let \(L\geq1\). A class represented by \(x_L\) lifts to a zero of \(F\) modulo \(p^{L+1}\) if and only if
\[
\boxed{
\operatorname{Ob}^{F}_L(x_L)=0.
}
\]
A correction \(y\in\FF_p^d\) gives a lift if and only if
\[
\boxed{
\overline J_F(x_L)y
=
-\rho_L^F(x_L).
}
\]
When a correction exists, the correction classes form a coset of
\[
\ker\overline J_F(x_L).
\]
\end{histtheorem}

\begin{proof}
The polynomial Taylor expansion gives
\[
F(x_L+p^Ly)
\equiv
F(x_L)+p^LJ_F(x_L)y
\pmod {p^{L+1}}.
\]
Every term of degree at least two in the increment contains \(p^{2L}\), which is divisible by \(p^{L+1}\). Divide by \(p^L\), reduce modulo \(p\), and apply the same cokernel criterion as in \cref{hist:thm:linear_lifting}.
\end{proof}

\begin{histproposition}[Polynomial representative independence]\label{hist:prop:polynomial_obstruction_independence}
The class \(\operatorname{Ob}^{F}_L(x_L)\) depends only on the class of \(x_L\) modulo \(p^L\).
\end{histproposition}

\begin{proof}
If \(x_L\) is replaced by \(x_L+p^Lz\), Taylor expansion shows that the reduced residual changes by an element of the image of \(\overline J_F(x_L)\). The reduced Jacobian itself is unchanged because \(L\geq1\).
\end{proof}

\begin{histcorollary}[Simple root lifting]\label{hist:cor:simple_root_lifting}
Assume \(r=d\) and
\[
\det J_F(x_L)\not\equiv0\pmod p.
\]
Then the depth \(L\) solution has exactly one child modulo \(p^{L+1}\). Iterating gives a unique \(p\) adic zero above the initial class.
\end{histcorollary}

\begin{proof}
The reduced Jacobian is invertible, so every residual has one correction. The resulting compatible sequence is Cauchy in \(\Zp^d\), and continuity of \(F\) gives a zero in the limit. Uniqueness holds at every depth.
\end{proof}

\begin{histexample}[A singular polynomial obstruction]\label{hist:ex:singular_polynomial_obstruction}
Let \(p=3\) and consider the class \(x_1=0\) modulo three. For
\[
F_0(x)=x^2,
\]
the residual and the reduced derivative both vanish. Every correction \(y\in\FF_3\) gives
\[
F_0(3y)\equiv0\pmod9,
\]
so the class has three children. For
\[
F_1(x)=x^2+3,
\]
the reduced derivative is still zero, but the reduced residual is one. Its class in the cokernel is nonzero, so no lift modulo nine exists.
\end{histexample}

This is the usual simple root mechanism behind Hensel lifting. The obstruction formulation also covers singular Jacobians, where some residual classes lift and others do not \cite{Serre1979,Gouvea1997,Robert2000}.

\subsubsection{Packetwise obstruction splitting}

\begin{histtheorem}[Packet decomposition of a lift obstruction]\label{hist:thm:packet_lift_obstruction}
Suppose
\[
\FF_p^d
=
\bigoplus_{a\in A}U_a,
\qquad
\FF_p^r
=
\bigoplus_{a\in A}V_a,
\]
and suppose a linear correction map \(D:\FF_p^d\to\FF_p^r\) satisfies
\[
D(U_a)\subseteq V_a
\]
for every \(a\). Write \(D_a=D|_{U_a}\). Then
\[
\boxed{
\coker D
\cong
\bigoplus_{a\in A}\coker D_a.
}
\]
A residual lifts if and only if every packet obstruction vanishes.
\end{histtheorem}

\begin{proof}
The image of \(D\) is the direct sum of the images of the maps \(D_a\). Taking the quotient of the codomain packet decomposition gives the stated cokernel decomposition.
\end{proof}

\begin{histremark}\label{hist:rem:lifting_boundary}
The obstruction theorems solve one declared local lifting step. They do not prove that a compatible branch exists across unrelated primes, that a global Diophantine point exists, or that an arithmetic representation is automorphic.
\end{histremark}

\subsection{Recurrence packets and a common integral spine}\label{hist:sec:recurrences}

Linear recurrences provide a setting where residue subsequences, local completions, and a common global source can all be described exactly. Background on recurrence sequences may be found in \cite{EverestVanDerPoortenShparlinskiWard2003,Stanley2012}.

\subsubsection{Companion jumps}

Let \(R\) be a commutative ring. Let
\[
v_{n+1}=Cv_n,
\qquad
C\in M_d(R).
\]
For a positive integer \(N\) and an address \(j\in\{0,\ldots,N-1\}\), define the residue state packet
\[
u_k^{(j)}=v_{Nk+j}.
\]

\begin{histtheorem}[Companion jump recurrence]\label{hist:thm:companion_jump}
Every residue state packet satisfies
\[
\boxed{
u_{k+1}^{(j)}=C^Nu_k^{(j)}.}
\]
More explicitly,
\[
\boxed{
u_k^{(j)}=C^j(C^N)^kv_0.}
\]
\end{histtheorem}

\begin{proof}
The state relation gives \(v_n=C^nv_0\). Substitute \(n=Nk+j\).
\end{proof}

\begin{figure}[t]
\centering
\includegraphics[width=.95\linewidth,height=3.15in,keepaspectratio]{../../figures/reconstructed/recurrence_packets.pdf}
\caption{New reconstruction. Fibonacci residues $F_{3k+j}$, with $F_0=0$ and $F_1=1$. Each packet evolves by the same jump matrix \(C^3\) and has a different initial address.}
\label{hist:fig:recurrence_packets}
\end{figure}

We use the resultant convention
\[
\Res_y(f,g)
=
\prod_{f(\alpha)=0}g(\alpha)
\]
when \(f\) is monic, with roots counted with multiplicity.

\begin{histtheorem}[Characteristic polynomial of the jump]\label{hist:thm:jump_characteristic}
Let
\[
\chi_C(y)=\det(yI-C).
\]
Then
\[
\boxed{
\chi_{C^N}(x)
=
\Res_y\bigl(\chi_C(y),x-y^N\bigr).
}
\]
The identity holds over the coefficient ring after interpreting the resultant universally.
\end{histtheorem}

\begin{proof}
After passage to an algebraic closure, let \(\alpha_1,\ldots,\alpha_d\) be the eigenvalues of \(C\), counted with algebraic multiplicity. The resultant is
\[
\prod_{i=1}^{d}(x-\alpha_i^N),
\]
which is the characteristic polynomial of \(C^N\). Both sides are polynomial expressions with integer coefficients in the entries of \(C\), so the identity descends to every commutative coefficient ring.
\end{proof}

\begin{figure}[t]
\centering
\includegraphics[width=.95\linewidth,height=3.15in,keepaspectratio]{../../figures/reconstructed/recurrence_root_power.pdf}
\caption{New reconstruction. Passing to every \(N\)th term raises each characteristic root to its \(N\)th power.}
\label{hist:fig:recurrence_root_power}
\end{figure}

\subsubsection{Second order recurrences}

\begin{histtheorem}[Residue subsequences of a second order recurrence]\label{hist:thm:second_order_jump}
Suppose
\[
x_{n+2}=Px_{n+1}-Qx_n
\]
over a commutative ring. Define \(V_0=2\), \(V_1=P\), and
\[
V_{n+2}=PV_{n+1}-QV_n.
\]
For every positive integer \(N\) and every address \(j\), the subsequence
\[
u_k=x_{Nk+j}
\]
satisfies
\[
\boxed{
u_{k+2}=V_Nu_{k+1}-Q^Nu_k.}
\]
\end{histtheorem}

\begin{proof}
The companion matrix
\[
C=
\begin{pmatrix}
P&-Q\\
1&0
\end{pmatrix}
\]
has trace \(P\) and determinant \(Q\). The trace of \(C^N\) is \(V_N\), while its determinant is \(Q^N\). Cayley Hamilton applied to \(C^N\) gives
\[
(C^N)^2-V_NC^N+Q^NI=0.
\]
Apply this identity to the residue state packet and take its scalar coordinate.
\end{proof}

\begin{histcorollary}[Fibonacci residue packets]\label{hist:cor:fibonacci_packets}
Let \(F_0=0\), \(F_1=1\), and \(F_{n+2}=F_{n+1}+F_n\). Let \(L_N\) denote the Lucas number. For every \(N\geq1\) and every address \(j\),
\[
\boxed{
F_{N(k+2)+j}
=
L_NF_{N(k+1)+j}
-(-1)^NF_{Nk+j}.
}
\]
\end{histcorollary}

\begin{proof}
Apply \cref{hist:thm:second_order_jump} with \(P=1\) and \(Q=-1\). Then \(V_N=L_N\).
\end{proof}

\subsubsection{Finite and local periodicity}

\begin{histtheorem}[Periodicity of an invertible recurrence modulo a prime power]\label{hist:thm:recurrence_periodicity}
Let
\[
C\in M_d(\ZZ),
\qquad
p\nmid\det C.
\]
For every \(L\geq1\), every state orbit modulo \(p^L\) is purely periodic. Its period divides the order of \(C\) in
\[
\operatorname{GL}_d(\ZZ/p^L\ZZ).
\]
If \(p\mid\det C\), every orbit is still eventually periodic.
\end{histtheorem}

\begin{proof}
When \(p\nmid\det C\), multiplication by \(C\) is a permutation of the finite state set \((\ZZ/p^L\ZZ)^d\). Every orbit of a permutation is periodic, and applying the group order relation for \(C\) gives the divisibility statement. Without invertibility, repeated application of a map on a finite set eventually repeats a state, after which the orbit is periodic.
\end{proof}

\begin{histcorollary}[Residue packet periods]\label{hist:cor:residue_packet_periods}
Under the invertibility hypothesis of \cref{hist:thm:recurrence_periodicity}, every residue packet with jump \(N\) has period dividing the order of \(C^N\) modulo \(p^L\).
\end{histcorollary}

\begin{proof}
Use the jump matrix from \cref{hist:thm:companion_jump}.
\end{proof}

\subsubsection{The common integral spine}

\begin{histtheorem}[Base change of an integral recurrence]\label{hist:thm:integral_recurrence_spine}
Let
\[
C\in M_d(\ZZ),
\qquad
v_0\in\ZZ^d.
\]
The integral orbit
\[
v_n=C^nv_0
\]
determines, by base change, compatible orbits in every ring \(R\) equipped with a homomorphism \(\ZZ\to R\). In particular, the reductions modulo every positive integer and the orbits in every \(\Zp^d\) share the same companion matrix, characteristic polynomial, and integer recurrence relations.
\end{histtheorem}

\begin{proof}
Every coordinate of \(C^nv_0\) is an integer polynomial in the entries of \(C\) and \(v_0\). Applying a ring homomorphism commutes with matrix multiplication and preserves every polynomial relation with integer coefficients.
\end{proof}

\begin{histcorollary}[Prime independent recurrence spine]\label{hist:cor:prime_independent_spine}
An integral companion system gives a canonical common source for all of its local realizations. No choice of coefficient prime is required to define the recurrence or its residue packets.
\end{histcorollary}

\begin{proof}
Apply \cref{hist:thm:integral_recurrence_spine} to the companion matrix and its initial state.
\end{proof}

\begin{histremark}\label{hist:rem:spine_boundary}
This theorem supplies a prime independent spine for integral linear recurrences. It does not prove that every unrelated family of local arithmetic packets arises from one global object.
\end{histremark}

\subsection{Chinese remainder packet refinement}\label{hist:sec:crt}

The Chinese remainder theorem gives an exact factorization of residue addresses. We state it at the level of sequence projectors so that it can be combined directly with the earlier coordinate systems.

\begin{histdefinition}[Residue projector]\label{hist:def:residue_projector}
For a sequence \(a=(a_n)_{n\geq0}\), a positive integer \(N\), and an address \(r\) modulo \(N\), define
\[
(P_{N,r}a)_n
=
\begin{cases}
a_n,&n\equiv r\pmod N,\\
0,&n\not\equiv r\pmod N.
\end{cases}
\]
\end{histdefinition}

\begin{histproposition}[Complete residue packet system]\label{hist:prop:residue_packet_system}
For fixed \(N\), the projectors \(P_{N,r}\) satisfy
\[
P_{N,r}P_{N,s}
=
\delta_{rs}P_{N,r}
\]
and
\[
\boxed{
\sum_{r\bmod N}P_{N,r}=I.
}
\]
\end{histproposition}

\begin{proof}
Every nonnegative integer belongs to exactly one residue class modulo \(N\).
\end{proof}

\begin{histtheorem}[Coprime packet factorization]\label{hist:thm:crt_packet_factorization}
Let \(A\) and \(B\) be coprime positive integers. Let \(r_A\) be an address modulo \(A\) and \(r_B\) an address modulo \(B\). Let \(r\) be the unique address modulo \(AB\) satisfying
\[
r\equiv r_A\pmod A,
\qquad
r\equiv r_B\pmod B.
\]
Then
\[
\boxed{
P_{AB,r}
=
P_{A,r_A}P_{B,r_B}
=
P_{B,r_B}P_{A,r_A}.
}
\]
\end{histtheorem}

\begin{proof}
The product retains exactly the indices satisfying both congruences. By the Chinese remainder theorem, those indices form one residue class modulo \(AB\) \cite{IrelandRosen1990}.
\end{proof}

\begin{histcorollary}[Several coprime moduli]\label{hist:cor:many_crt_packets}
Let \(N_1,\ldots,N_s\) be pairwise coprime and put
\[
N=\prod_{j=1}^{s}N_j.
\]
For compatible addresses \(r_j\),
\[
\boxed{
P_{N,r}
=
\prod_{j=1}^{s}P_{N_j,r_j},
}
\]
where \(r\) is the unique combined address modulo \(N\).
\end{histcorollary}

\begin{proof}
Iterate \cref{hist:thm:crt_packet_factorization}.
\end{proof}

\begin{figure}[t]
\centering
\includegraphics[width=.95\linewidth,height=3.15in,keepaspectratio]{../../figures/reconstructed/crt_packets.pdf}
\caption{New reconstruction. A residue modulo a product of coprime moduli is the exact intersection of its local residue packets.}
\label{hist:fig:crt_packets}
\end{figure}

\begin{histstandardfact}[Finite local reconstruction]\label{hist:fact:finite_local_reconstruction}
Let
\[
N=\prod_{j=1}^{s}p_j^{a_j}
\]
be the prime power factorization of a positive integer. The map
\[
\ZZ/N\ZZ
\longrightarrow
\prod_{j=1}^{s}\ZZ/p_j^{a_j}\ZZ
\]
is a ring isomorphism. Thus a compatible collection of prime power residues reconstructs one residue modulo \(N\) \cite{IrelandRosen1990}.
\end{histstandardfact}

\begin{histproposition}[Dyadic and odd residue refinement]\label{hist:prop:dyadic_odd_refinement}
Let \(N=2^am\) with \(m\) odd. Every residue packet modulo \(N\) factors uniquely as the product of one dyadic packet modulo \(2^a\) and one odd packet modulo \(m\).
\end{histproposition}

\begin{proof}
The moduli \(2^a\) and \(m\) are coprime. Apply \cref{hist:thm:crt_packet_factorization}.
\end{proof}

\begin{histremark}\label{hist:rem:crt_weight_boundary}
The residue address factors exactly, but valuation weights do not become identical across primes. The odd coordinate system of \cref{hist:sec:local_euler} and the dyadic coordinate system of \cref{hist:sec:local_euler} remain genuinely different integral charts.
\end{histremark}

\subsection{Verification and exact boundaries}\label{hist:sec:boundaries}

The proofs in the article are symbolic. Exact computation tests formulas, finite ranges, and the source implementation. It does not replace a uniform proof.

\begin{table}[t]
\centering
\caption{Status of the principal local arithmetic blocks.}
\label{hist:tab:status}
\begin{tabular}{>{\raggedright\arraybackslash}p{.33\textwidth}>{\raggedright\arraybackslash}p{.24\textwidth}>{\raggedright\arraybackslash}p{.31\textwidth}}
\toprule
Block & Status & Main boundary\\
\midrule
Standard window period & Proved modulo the residue prime & Higher prime powers require inverse estimates\\
Unified local Euler transfer & Proved for every prime & Odd and dyadic integral charts remain different\\
Dyadic kernel and truncation & Exact weight and optimal uniform modulus proved & Row cancellation may give stronger values\\
Carmichael and Mahler transfer & Proved for scalar, polynomial exponential, matrix, and finite adelic systems & A jump congruence must first be available\\
Local Euler interpolation & Proved for depleted odd branches and the direct dyadic branch & The odd nonunit factor is only evaluated at integer interpolation points\\
Affine polyphase transfer & Exact analysis, synthesis, carry law, kernel composition, residual sparsity, and valuation preservation proved & General exact kernel valuations remain open\\
Lift obstruction and recurrence spine & Proved for the declared local problems and integral recurrences & Unrestricted global gluing is not proved\\
\bottomrule
\end{tabular}
\end{table}

\subsubsection{The remaining congruential difference problem}

Let \(C_{mp,j}(n)\) be a congruential reciprocal coefficient and let
\[
h=\frac{p-1}{m}.
\]
The unit mode calculations suggest bounds for
\[
\vp\left(\Delta_h^sC_{mp,j}(n)\right).
\]
The theorems of \cref{hist:sec:mahler,hist:sec:polyphase} determine the gain before the congruential triangular inverse is applied. They do not yet determine every valuation loss introduced by that inverse. The full family bound therefore remains open.

\subsubsection{Higher valuations of the affine kernel}

The residual formula
\[
\Theta_{p^a,s}(r,k)
\equiv
\binom{s}{r-p^ak}
\pmod p
\]
is exact. A closed formula for
\[
\vp\left(\Theta_{p^a,s}(r,k)\right)
\]
in terms of base \(p\) carries and the three digit addresses is not proved here. Such a formula would sharpen the min plus valuation ledger and could distinguish optimal refinement trees.

\subsubsection{The remaining global gluing problem}

The finite adelic theorem synchronizes a finite set of primes when one integral state system is already fixed. The recurrence theorem gives a common local source under the same assumption. Chinese remainder projectors reconstruct finite residue addresses. None of these results proves that an arbitrary collection of local analytic branches has a common global arithmetic origin.

\begin{figure}[t]
\centering
\includegraphics[width=.95\linewidth,height=3.15in,keepaspectratio]{../../figures/reconstructed/claim_boundary.pdf}
\caption{New reconstruction. The proved local coordinate theory and the remaining inverse valuation, affine kernel, and global gluing problems.}
\label{hist:fig:claim_boundary}
\end{figure}

\subsubsection{Historical scope}

The article uses classical work on Euler congruences, finite differences, local fields, Hurwitz series, Mahler interpolation, Lucas congruences, and recurrence sequences \cite{CarlitzLevine1960,Young1999,Young2001,Sun2008,Chen2004,Lucas1878a,LidlNiederreiter1997,EverestVanDerPoortenShparlinskiWard2003}. The exact formulas and local transfer statements should be compared individually with that literature before any historical priority claim is made.

\subsection{Conclusion}\label{hist:sec:conclusion}

The local arithmetic theory now has one common architecture. The finite character coordinate identifies the torsion phase. The source and kernel product reconstructs the Euler coefficient. The kernel valuation records certified precision. Carmichael and Mahler coordinates describe the remaining principal unit motion. Polyphase analysis resolves additional local digits, records carries and shifts exactly, and transports Mahler coefficients through a finite affine kernel.

The odd and dyadic cases agree at the level of convolution and differ at the level of integral geometry. At odd primes the source must be depleted before it becomes a unit mode interpolation. At the dyadic prime the source extends directly to an analytic integral function. The dyadic kernel has exact weight \(4r-1\), giving the optimal uniform modulus \(2^{4M+3}\) after depth \(M\).

Three tasks remain outside the proved region. The first is a complete finite difference estimate after the congruential triangular inverse. The second is an exact higher valuation theory for the affine Mahler kernel. The third is a global existence theorem for local systems that do not already arise from one integral source.

\par\smallskip\noindent\textbf{Current status.} This dated statement is superseded for the stated divisor-branch family by Theorem~\ref{thm:arith-divisor-base}, with the $j!$ normalization and factorial valuation loss explicitly retained. It makes no claim for unrestricted families.\par



\clearpage\section{Complete mathematical checkpoint supplements}\label{app:supplements}
These eleven supplements reproduce the complete mathematical bodies of their
successive checkpoints. The numerical inventories, coverage fractions and
then-pending statements belong to each original checkpoint. They are not
current total counts. In particular the later all-residue and Nishibuchi
completion supersede earlier pending labels. The final arithmetic and
polynomial scopes are the ones stated in the main text. All derivations,
examples, formal-interface explanations and original qualification paragraphs
remain available below.
\clearpage\subsection{First Lean Block Supplement}\label{app:sup1}
\noindent\textbf{Checkpoint scope.} This complete source body retains the status of its dated original supplement. Earlier pending or partial-coverage labels are historical; the final core records their resolution.\par\medskip
\subsubsection*{Scope and result}
This block uses the 28-page \emph{status-corrected} Paper IX manuscript. That PDF is preserved unchanged. The release supplies a first formalization companion, not a completed or rewritten arithmetic paper.

\textbf{Seven of 24 named manuscript-proved items now have Lean coverage in the stated scopes. Seventeen remain pending.} Three classical inputs and two conjectures are registered separately. The development contains \textbf{89 explicit theorem declarations and 22 definitions across eight modules}. Supporting lemmas and concrete instances are included; these are not 89 independent discoveries.

\begin{center}
\begin{tabular}{@{}p{.15\textwidth}p{.73\textwidth}@{}}
\toprule
Item & Covered mathematical content\\\midrule
1.2 & Residue reconstruction and divisor refinement\\
2.1 & Actual unit-EGF coordinate equivalence and transported projectors\\
2.2 & Filtered finite coefficient reconstruction, including $\Z_p$\\
3.2 & Constructed unique triangular inverse and original section identity\\
3.3 & Universal denominator clearing\\
8.2 & Lucas standard-window reduction and the diagonal unit\\
13.1 & Exact normalization counterexample and sharp factorial loss\\
\bottomrule
\end{tabular}
\end{center}

\subsubsection*{The correction is preserved}
The normalized packet is not the same object as the packet multiplied by $j!$. The checked recurrence gives
\[
 C_{6,3}(1)=-\frac1{84},\qquad C_{6,3}(2)=\frac{703}{5460},
\]
\[
 C_{6,3}(2)-C_{6,3}(1)=\frac{64}{455},\qquad
 v_3\!\left(\frac{64}{455}\right)=0.
\]
Thus the unrestricted normalized no-loss claim fails already at a valid one-step parameter instance. Multiplication by $3!$ restores one factor of three in this instance. The general corrected finite-difference theorem is \textbf{not} inferred from this calculation and is not yet formalized here.

\subsubsection*{Verification standard}
All eight modules passed a fresh cumulative build with compiler warnings treated as errors. The audit resolved all 89 explicit theorems and all 22 explicit definitions. Their transitive axiom dependencies are confined to \mcode{propext}, \mcode{Classical.choice} and \mcode{Quot.sound}, or subsets. Three invalid-proof controls were rejected. The proof, scope, and numerical-regression records remain distinct.
\par\medskip
\pageheading{1. Residue packets and actual EGF coordinates}{Source items 1.2 and 2.1}
For a sequence $a:\N\to R$, define the actual selector
\[
 (P_{N,j}a)(n)=\begin{cases}a(n),&n\bmod N=j,\\0,&\text{otherwise.}\end{cases}
\]
The source proves idempotence for every selector, orthogonality for distinct standard residue labels, and, when $N>0$,
\[
 \sum_{j=0}^{N-1}P_{N,j}a=a.
\]
If $N\mid M$ and $M>0$, it also proves
\[
 P_{N,j}a=\sum_{\substack{0\le r<M\\r\bmod N=j}}P_{M,r}a.
\]
The formal power-series versions are proved coefficientwise for actual \mcode{PowerSeries} objects. There is no finite truncation of the sequence domain in these theorems.

\subsection*{A constructed inverse, not an inverse hypothesis}
Let $K$ be a characteristic-zero field. Define
\[
 \mathcal E(a)=\sum_{n\ge0}\frac{a_n}{n!}t^n,\qquad
 \mathcal D(A)_n=n![t^n]A.
\]
Both inverse laws $\mathcal D\mathcal E=I$ and $\mathcal E\mathcal D=I$ are proved. The ordinary power-series product then yields the genuine binomial convolution
\[
 \mathcal D\bigl(\mathcal E(h)\mathcal E(a)\bigr)_n
 =\sum_{r=0}^{n}\binom nr h_r a_{n-r}.
\]
For $h_0=1$, the formal inverse of $H=\mathcal E(h)$ is used to construct a sequence equivalence
\[
 U_H a=\mathcal D(H\mathcal E(a)),\qquad
 U_H^{-1}a=\mathcal D(H^{-1}\mathcal E(a)).
\]
The equivalence has proved left and right inverse laws. The inverse's constant coefficient is proved to be one. The associated lower-triangular entries have diagonal one and vanish above the diagonal.

For the actual conjugated selectors
\[
 \widehat P_{N,j}=U_HP_{N,j}U_H^{-1},
\]
idempotence, orthogonality, and finite reconstruction are proved from the constructed maps. No projector identity for the conjugates is inserted as an assumption.

\textbf{Formal boundary.} These are formal series and coefficient identities. They do not establish convergence for an analytic EGF. The characteristic-zero field assumption is explicit; unrestricted factorial division in an arbitrary ring is not asserted.

\textbf{Principal targets.}\par
\target{ResiduePackets.packet_reconstruction}\par
\target{ResiduePackets.packet_refinement}\par
\target{EGF.coordinateEquiv, EGF.coordinate_forward, EGF.coordinate_inverse}\par
\target{EGF.transported_reconstruction}
\par\medskip
\pageheading{2. Triangular inversion and denominator control}{Source items 3.2 and 3.3}
For a general lower-triangular coefficient array $T$ over a field with nonzero diagonal, the code defines forward substitution by well-founded recursion:
\[
 a_n=\frac{b_n-\sum_{q<n}T_{nq}a_q}{T_{nn}}.
\]
It proves that this constructed sequence solves every row and that every sequence solving those rows equals it. No existence or uniqueness premise is used.

\subsection*{The congruential coefficients}
Over $\Q$, specialize to
\[
 T_{nq}=\binom{Nn+j}{Nq},\qquad b_n=\delta_{n0}.
\]
The diagonal is positive as an integer, hence nonzero. The resulting sequence $C_{N,j}(n)$ satisfies
\[
 C_{N,j}(0)=1,\qquad
 \sum_{q=0}^{n}\binom{Nn+j}{Nq}C_{N,j}(q)=\delta_{n0}.
\]
The displayed recurrence in the manuscript and the unique-solution theorem are consequences of this definition.

The formalization does not stop at a recurrence. In the variable $z=t^N$, construct
\[
 K_{N,j}(z)=\sum_{n\ge0}\frac{j!}{(Nn+j)!}z^n,\qquad
 F_{N,j}(z)=\sum_{n\ge0}\frac{C_{N,j}(n)}{(Nn)!}z^n.
\]
It proves the factorial coefficient identity, then $F_{N,j}K_{N,j}=1$ and therefore $F_{N,j}=K_{N,j}^{-1}$.

For $N>0$ and $0\le j<N$, the independently defined original exponential section is
\[
 E_{N,j}(t)=\sum_{d\equiv j\ ({\rm mod}\ N)}\frac{t^d}{d!}.
\]
Using the actual power-series expansion homomorphism $z\mapsto t^N$, the code proves
\[
 E_{N,j}(t)=\frac{t^j}{j!}K_{N,j}(t^N),\qquad
 \boxed{F_{N,j}(t^N)E_{N,j}(t)=\frac{t^j}{j!}.}
\]
Thus the recurrence is connected back to the original normalized EGF definition. The section vanishes at zero when $j>0$; no inverse of that nonunit section is assumed.

\subsection*{Universal clearing}
Define
\[
 D_{N,j}(n)=\prod_{r=1}^{n}\binom{Nr+j}{Nr},\qquad D_{N,j}(0)=1.
\]
For every $n$, the code proves $D_{N,j}(n)C_{N,j}(n)\in\Z$. The general proof promotes each earlier clearing denominator to the common product, then cancels the new nonzero diagonal. This is a denominator bound, not a claim that every normalized coefficient is integral.
\par\medskip
\pageheading{3. Filtered tails and the standard window}{Source items 2.2 and 8.2}
\subsection*{Filtration at a finite coefficient row}
Let $R$ be any commutative ring. For $w\in\N\cup\{\infty\}$ define
\[
 p^{[w]}=\begin{cases}p^w,&w<\infty,\\0,&w=\infty.\end{cases}
\]
The infinity convention expresses a vanishing coefficient, not an infinite exponent operation. For any finite index set, every term with weight at least $L$ is divisible by $p^L$. A finite-sum proof yields
\[
 \boxed{p^L\ \bigm|\ 
 \sum_{r=0}^{n}\binom nr c_rp^{[w(r)]}S_{n-r}
 -\sum_{\substack{0\le r\le n\\w(r)<L}}
     \binom nr c_rp^{[w(r)]}S_{n-r}.}
\]
The result is specialized to actual $p$-adic integers $R=\Z_p$ for prime $p$. Divisibility there is precisely the asserted membership in $p^L\Z_p$.

For an individual row the sum is finite, so the proof does not require $w(r)\to\infty$. This is a stronger coefficient-row scope, not a theorem about an infinite limiting process. A supporting lemma also proves a $p^{2M+2}$ tail bound for a sum with visible factors $p^{2m}$. It is not counted as either Euler supercongruence theorem: their generating-function reconstruction identities remain pending.

\subsection*{Actual Lucas reduction}
For prime $p$, all $m,n,q\in\N$, and $0\le j<p$, the checked binomial congruence is
\[
 \boxed{\binom{mpn+j}{mpq}\equiv\binom{mn}{mq}\pmod p.}
\]
The proof applies Mathlib's Lucas low-digit reduction to the actual natural binomial coefficients. It does not assume a pre-reduced row. The statement also permits $q>n$. If $m>0$, both binomial coefficients then vanish.

Taking $q=n$ gives
\[
 \binom{mpn+j}{mpn}\equiv1\pmod p.
\]
Nondivisibility by $p$ and invertibility of that diagonal in $\Z/p\Z$ are proved. The separate supporting lemma $p\nmid j!$ for $j<p$ establishes the standard-window factorial-unit condition.

\textbf{Dependency boundary.} Lucas' theorem is a named classical input in the manuscript and is reused from the pinned Mathlib library. This block does not independently reconstruct Lucas' full digit theorem. The residue-independence result (8.3), uniform standard-window period (8.5), and the global finite-difference theorem (7.5) remain pending.

\textbf{Principal targets.}\par
\target{Filtered.padic_filtered_reconstruction}\par
\target{StandardWindow.lucas_row_reduction}\par
\target{StandardWindow.standard_diagonal_unit}\par
\target{StandardWindow.factorial_valuation_zero}
\par\medskip
\pageheading{4. Sharp normalization boundary}{Source item 13.1}
The construction fixes the two coefficient conventions:
\[
 C_{N,j}(n),\qquad \widetilde C_{N,j}(n)=j!\,C_{N,j}(n).
\]
The witness uses $p=3$, $m=2$, $N=mp=6$, $j=3$, $n=1$ and difference step $h=(p-1)/m=1$. The parameter conditions are checked explicitly, including odd primality, $m\mid p-1$ and $j<N$.

The first row gives $84C_{6,3}(1)+1=0$. The second row, evaluated by the same recursive definition, gives $C_{6,3}(2)=703/5460$. Therefore
\[
 \Delta_1 C_{6,3}(1)=\frac{64}{455}\ne0.
\]
Neither numerator nor denominator is divisible by three, so its actual rational 3-adic valuation is zero. A further theorem negates the universal one-step no-loss statement by applying it to this valid tuple and deriving $1\le0$.

For the unnormalized convention,
\[
 \Delta_1\widetilde C_{6,3}(1)=3!\,\frac{64}{455}=\frac{384}{455},
 \qquad v_3\!\left(\frac{384}{455}\right)=1.
\]
The factorial's valuation is also proved:
\[
 v_3(3!)=1,\qquad
 v_3\bigl(\Delta_1C_{6,3}(1)\bigr)
 =1-v_3(3!)=0.
\]
This proves sharpness at the displayed normalization boundary. It does not establish the corrected global inequality for all parameters, and does not assert that every nonunit-factorial parameter tuple must violate every no-loss specialization. All valuation calculations here concern nonzero rationals.

\subsubsection*{Work remaining in the named proved register}
\begin{center}\small
\begin{tabular}{@{}p{.19\textwidth}p{.69\textwidth}@{}}
\toprule
Items & Pending mathematical content\\\midrule
4.1--4.3 & Odd-power generating function, Euler reconstruction, odd-prime ladder\\
5.1--5.4 & Dyadic inverse and source integrality, reconstruction, weighted ladder\\
6.2 & Finite unit-decomposition certificates\\
7.2--7.5 & Finite Mahler reconstruction, weighted differences, unit modes, global odd-prime differences\\
8.3, 8.5 & Standard-window residue independence and uniform period\\
9.1, 9.2 & Companion jumps and Fibonacci packet recurrence\\
10.1 & Arithmetic CRT packet factorization\\\bottomrule
\end{tabular}\end{center}
The two conjectures, 13.2 and 13.3, remain unproved and outside the 24-item register. The 17 pending entries are manuscript proof claims awaiting this project's Lean verification, not conjectures newly created by the formalization process.
\par\medskip
\pageheading{5. Verification and reproducibility}{Exact source pins, observed checks and retained failures}
\begin{center}\small
\begin{tabular}{@{}lrr@{}}\toprule
Module & Theorems & Definitions\\\midrule
ResiduePackets & 8 & 2\\
Triangular & 14 & 3\\
EGFCoordinates & 21 & 6\\
FilteredReconstruction & 7 & 1\\
CongruentialEuler & 16 & 7\\
ExponentialSection & 5 & 3\\
StandardWindow & 6 & 0\\
NormalizationBoundary & 12 & 0\\\midrule
Total & 89 & 22\\\bottomrule
\end{tabular}\end{center}

\textbf{Lean run.} All eight modules were compiled into a new empty output directory under Lean 4.34.0-rc2. All nine dependency revisions matched the pinned manifest. Mathlib's revision is\par
{\footnotesize\mcode{aeb72f1e9981b6d35b426b16dca7759be3e7ba06}.}

The 111 explicit-declaration axiom reports matched the inventory. A false natural equality, an admitted placeholder, and the false normalized bound were all rejected. The placeholder control exercises the warnings-as-errors policy; \mcode{sorryAx} is also excluded by the separate axiom whitelist. The 21 infrastructure and release tests passed.

\textbf{Independent exact regressions.} Python integer and rational mirrors passed 120 triangular systems, 45 congruential parameter pairs (495 rows), 120 EGF-coordinate cases, 45 residue-refinement cases, 240 filtered tails, 10,824 Lucas rows, the normalization witness and eight invalid-input controls. These are not executions extracted from Lean or reruns of the original historical numerical package.

\textbf{Failures retained.} Of 29 captured development checks, 15 passed and 14 failed before repair. A malformed generic-field draft triggered a release-candidate compiler panic after proof errors. The repaired sources passed the final clean run. Development products were not reused in that run. An earlier cumulative invocation was interrupted before completion and is retained separately. One initial Python scanner-test assertion had the wrong expected whitespace count and was repaired; its failure log is retained.

\textbf{Trust and source boundary.} Mathlib's precompiled dependency cache was reused, not rebuilt from source. No earlier paper was rebuilt. Only the status-corrected primary PDF was recovered for this manuscript baseline, not its original TeX, experiment package, or complete global finite-difference supplement. Formal-to-prose correspondence is a reviewed crosswalk, not an automated semantic-equivalence theorem. Native Windows was not tested.

\textbf{Reproduce.} With Python 3.10+ and the separate extracted Linux portable kit:
\begin{verbatim}
bash Run-Portable.sh /path/to/Komorebi_Portable_Lean_4.34.0-rc2
\end{verbatim}
The ZIP includes sources, inventories, the unchanged manuscript, this supplement, exact-regression scripts, tests, and successful and failed logs. The large runtime and local build products are not included.

\clearpage\subsection{Second Lean Block Supplement}\label{app:sup2}
\noindent\textbf{Checkpoint scope.} This complete source body retains the status of its dated original supplement. Earlier pending or partial-coverage labels are historical; the final core records their resolution.\par\medskip
\subsubsection*{Scope and result}
This release continues the 28-page \emph{status-corrected} Paper IX manuscript. It completes formal coverage of Lemma 4.1, Theorem 4.2 and Corollary 4.3. The original manuscript and all eight first-block proof modules remain unchanged.

\textbf{Ten of 24 named manuscript-proved items are now covered in their recorded scopes. Fourteen remain pending.} Three classical inputs and two conjectures are tracked separately. The cumulative corpus has \textbf{163 named theorem declarations and 37 explicit definitions in 13 modules}. The 74 new theorems include auxiliary identities and concrete examples; they are not 74 independent discoveries.

\begin{center}
\begin{tabular}{@{}p{.13\textwidth}p{.63\textwidth}r@{}}
\toprule
Item & New formal coverage & Theorems\\\midrule
4.1 & Literal odd-power EGF and its even part & 18\\
4.2 & Integral Euler and signed zigzag coefficient bridges & 29\\
4.2, 4.3 & Exact reconstruction, matrix rows and truncation & 16\\
Example & Exact scalar anchors and sharpness controls & 11\\\midrule
 & Total added across five modules & 74\\\bottomrule
\end{tabular}
\end{center}
The five added modules are:\par
{\small\mcode{OddPowerSeries}, \mcode{EulerCoefficients}, \mcode{ZigzagCoefficients},\newline
\mcode{OddEulerReconstruction} and \mcode{OddEulerExamples}.}

\subsubsection*{A useful scope strengthening}
The algebraic proof needs \emph{oddness}, not primality. The exact reconstruction and the integer divisibility statement are therefore checked for every positive odd length $p$, including $p=1$ and odd composites. The manuscript's odd-prime statement follows as a specialization. This is a sharpening of assumptions in this development, not a historical-priority claim. Even lengths are not included; a checked row-zero counterexample shows why.

\subsubsection*{What is not completed}
The weighted dyadic theory and the general corrected odd-prime finite-difference theorem remain pending. No continuum convergence or analytic summability theorem is inferred from the formal power-series equalities in this block. The historical numerical study is not being relabelled as a fresh experiment.
\par\medskip
\pageheading{1. The actual odd-power generating function}{Lemma 4.1; manuscript page 10}
Work in the rational formal power-series ring $\Q[[t]]$. Define
\[
 S_n(p)=\sum_{a=0}^{p-1}(-1)^a(2a+1)^n\in\Z,
 \qquad F_p(t)=\sum_{n\ge0}S_n(p)\frac{t^n}{n!}.
\]
The formal exponential is Mathlib's \mcode{PowerSeries.exp}; scalar substitution is its actual \mcode{rescale} homomorphism. We define
\[
 \exp_a(t)=e^{at},\qquad
 \cosh_a(t)=\frac{e^{at}+e^{-at}}2,\qquad
 \sech(t)=\cosh_1(t)^{-1}.
\]
The constant term of $\cosh_a$ is proved to be one, so the reciprocal used here is a genuine unit inverse.

\subsection*{From finite sums to a series identity}
Coefficient extraction proves
\[
 F_p(t)=\sum_{a=0}^{p-1}(-1)^a e^{(2a+1)t}.
\]
An induction on the actual finite summation length establishes the telescoping identity
\[
 F_p(t)\cosh t=\frac{1-(-1)^p e^{2pt}}2.
\]
For odd $p$ this gives, without a geometric-series convergence assumption,
\[
 \boxed{F_p(t)=\frac{1+e^{2pt}}2\,\sech t.}
\]
This is an equality of complete formal series, not a fixed-degree polynomial check.

\subsection*{Even projection is identified coefficientwise}
For any formal series $f$, let $\operatorname{EvenPart}(f)=(f(t)+f(-t))/2$. Its $n$th coefficient is proved to be the coefficient of $f$ when $n$ is even and zero otherwise. Rescaling the inverse law proves $\sech(-t)=\sech(t)$. The exponential multiplication identity then gives
\[
 \cosh^2(pt)=\frac{2+e^{2pt}+e^{-2pt}}4,
\]
and hence
\[
 \boxed{\sum_{r\ge0}S_{2r}(p)\frac{t^{2r}}{(2r)!}
       =\cosh^2(pt)\,\sech t.}
\]
Both the full and even EGF formulas are derived from the original integer sum. Neither is supplied as a hypothesis.

\textbf{Principal targets.}\par
\target{OddPower.source_as_exponentials}\par
\target{OddPower.source_times_cosh}\par
\target{OddPower.odd_source_egf, OddPower.odd_source_even_egf}
\par\medskip
\pageheading{2. Integral Euler coefficients and the zigzag bridge}{Coefficient foundations for Theorem 4.2}
The first block already constructed $C_{N,j}(r)$ by the manuscript's triangular recurrence. At $(N,j)=(2,0)$, every diagonal binomial coefficient equals one. Its denominator-clearing product is therefore one, and $C_{2,0}(r)$ is an integer.

The source defines $E_{2r}$ as the integer numerator of that rational coefficient and proves that its rational cast equals $C_{2,0}(r)$. Set $E_{2r+1}=0$. The actual exponential-section identity from the first block proves
\[
 \boxed{\sum_{n\ge0}E_n\frac{t^n}{n!}=\sech t.}
\]
Thus the integer sequence is connected to the intended signed Euler family, rather than merely assigned the name Euler.

\subsection*{The inverse coefficients are integral by construction}
Define the integer binomial convolution
\[
 B_n=\sum_{r=0}^{n}\binom nr E_rE_{n-r},\qquad b_m=B_{2m}.
\]
The product formula for actual EGFs gives
\[
 \sum_{n\ge0}B_n\frac{t^n}{n!}=\sech^2 t.
\]
Odd coefficients vanish by parity. In particular, every $b_m$ is an integer; this is the integrality required by the later congruence.

\subsection*{Connecting to the chapter's $Z_n$ convention}
The chapter defines $Z_n$ through $\sec t+\tan t$. The Lean construction uses Mathlib's actual formal sine and cosine series and defines
\[
 \sec t=(\cos t)^{-1},\quad \tan t=(\sin t)\sec t,
 \quad Z_n=n![t^n](\sec t+\tan t).
\]
It proves the sine and cosine derivative identities, $\sin^2t+\cos^2t=1$, and
\[
 (\tan t)'=\sec^2t.
\]
Sign substitution in the squared variable carries the hyperbolic cosine kernel to the actual cosine kernel. The same homomorphism carries their inverses and their squares. Extracting the even coefficient in the derivative of tangent yields
\[
 \boxed{b_m=(-1)^m Z_{2m+1}.}
\]
This formula is \emph{proved}, not assumed as an input to reconstruction. The formalization follows the generating-function definition of zigzag numbers; it does not additionally formalize their interpretation as counts of alternating permutations.

\textbf{Principal targets.}\par
\target{EulerCoefficients.euler_egf, EulerCoefficients.inverse_coeff_cast}\par
\target{Zigzag.derivative_tan, Zigzag.inverse_is_signed_zigzag}
\par\medskip
\pageheading{3. Exact reconstruction and the triangular matrix}{Theorem 4.2; manuscript page 11}
For odd $p$, rescaling the already proved inverse law gives
\[
 \sech(pt)\cosh(pt)=1.
\]
Multiplying the even-source identity by $\sech^2(pt)$ therefore proves
\[
 \sech^2(pt)\left(\sum_{r\ge0}S_{2r}(p)\frac{t^{2r}}{(2r)!}\right)=\sech t.
\]
Actual EGF multiplication now gives the integer coefficient formula
\[
 \boxed{E_{2r}=\sum_{m=0}^{r}\binom{2r}{2m}b_m p^{2m}S_{2(r-m)}(p).}
\]
The formal proof first extracts the sum over all indices and then proves that the odd summands vanish. A checked even-index summation identity supplies the displayed change of variables.

Substituting the independently proved signed-zigzag identity recovers the manuscript's notation. For every even degree $n$,
\[
 \boxed{E_n=\sum_{m=0}^{n/2}(-1)^m Z_{2m+1}\binom n{2m}p^{2m}S_{n-2m}(p).}
\]
The Lean version uses $n=2r$ to avoid ambiguous natural-number division. Every even natural degree has that form. Primality is not used in this equality.

\subsection*{The row action is connected to the matrix entries}
Reversing the finite sum and using binomial symmetry gives the actual source-index row action
\[
 E_{2r}=\sum_{q=0}^{r}U_{r,q}(p)S_{2q}(p),
\]
where
\[
 U_{r,q}(p)=
 \begin{cases}
 \displaystyle\binom{2r}{2q}b_{r-q}p^{2(r-q)},&q\le r,\\
 0,&q>r.
 \end{cases}
\]
The row-action equality, the diagonal identity $U_{r,r}=1$, and vanishing above the diagonal are all checked. The arithmetic weight is also explicit:
\[
 p^{2(r-q)}\mid U_{r,q}(p)\quad\text{in }\Z,\qquad 0\le q\le r.
\]
No inverse-coordinate identity, unit diagonal or weight bound is assumed for these rows.

\textbf{Principal targets.}\par
\target{OddEuler.odd_reconstruction_series}\par
\target{OddEuler.odd_euler_reconstruction}\par
\target{OddEuler.odd_euler_reconstruction_zigzag}\par
\target{OddEuler.row_action, OddEuler.row_diagonal, OddEuler.row_weight}
\par\medskip
\pageheading{4. Uniform truncation and a sharp exact example}{Corollary 4.3 and Example 4.4}
For row $2r$ and depth $M$, the actual truncation is
\[
 T_M(p,2r)=\sum_{m=0}^{\min(M,r)}\binom{2r}{2m}b_m p^{2m}S_{2(r-m)}(p).
\]
Every omitted term has $m\ge M+1$. Its other factors are integers, so the first-block omitted-tail theorem, now applied to the \emph{proved Euler identity}, gives
\[
 \boxed{p^{2M+2}\mid E_{2r}-T_M(p,2r)\quad\text{in }\Z.}
\]
This is not divisibility in $\Q$, where a nonzero divisor would make the assertion vacuous. The statement holds for every row and depth. For an odd prime it is precisely the claimed supercongruence. If $M\ge r$, complete retention is proved exact.

\subsection*{The displayed manuscript row is kernel-checked}
The scalar anchors are derived from the same general coefficient recurrence:
\[
 (E_0,E_2,E_4,E_6,E_8)=(1,-1,5,-61,1385),
\]
\[
 (b_0,b_1,b_2,b_3,b_4)=(1,-2,16,-272,7936).
\]
The literal source sums at $p=5$ give
\[
 (S_0,S_2,S_4,S_6,S_8)=(1,49,4705,428689,37665985).
\]
Consequently,
\[
 T_2(5,8)=2\,731\,001\,385,\qquad
 E_8-T_2(5,8)=-2\,731\,000\,000.
\]
The source proves both
\[
 5^6\mid E_8-T_2(5,8),\qquad 5^7\nmid E_8-T_2(5,8).
\]
Thus the guaranteed power is attained in this example. No stronger uniform exponent is inferred.

\subsection*{Oddness cannot be silently dropped}
For length two, $S_0(2)=1-1=0$, whereas $E_0=1$. The row-zero reconstruction already fails. The source proves this failure, and the release verifier separately attempts and rejects an invalid proof of the even-length identity.

\textbf{Principal targets.}\par
\target{OddEuler.odd_euler_supercongruence}\par
\target{OddEuler.odd_prime_supercongruence, OddEuler.full_depth_exact}\par
\target{OddEulerExamples.manuscript_error_sharp}\par
\target{OddEulerExamples.even_length_reconstruction_fails}
\par\medskip
\pageheading{5. Verification, reproducibility and remaining work}{Cumulative release evidence}
\begin{center}
\begin{tabular}{@{}p{.70\textwidth}r@{}}
\toprule
Check & Result\\\midrule
Fresh cumulative proof-module compilation & 13/13\\
Named theorem declarations audited & 163/163\\
Explicit definitions audited & 37/37\\
Exact dependency-revision checks & 9/9\\
Invalid proof controls rejected as intended & 5/5\\
Infrastructure and release-contract unit tests & 33 passed\\
New captured development attempts & 15\\
Development attempts failing before correction & 8\\\bottomrule
\end{tabular}
\end{center}
The pinned runtime is Lean 4.34.0-rc2 and Mathlib revision
\begin{center}\footnotesize\mcode{aeb72f1e9981b6d35b426b16dca7759be3e7ba06}.\end{center}
Every final module compiled with warnings treated as errors. All 200 expected axiom reports were resolved and used only \mcode{propext}, \mcode{Classical.choice} and \mcode{Quot.sound}, or subsets. The audit is of explicit project declarations. Mathlib's compiled cache was reused, not independently rebuilt.

The negative controls cover a false equality, a proof placeholder, the old false normalized finite-difference bound, the false even-length extension, and a false extra digit of precision. Missing imports and timeouts do not count as successful rejection. The placeholder is rejected by warnings-as-errors, and \mcode{sorryAx} is separately excluded from the axiom audit.

\subsection*{Independent finite arithmetic checks}
Exact Python mirrors passed 234 reconstruction rows and 2,340 truncation-depth cases for odd lengths 1 through 35 and even degrees 0 through 24. They also checked 468 coefficients of each generating-function identity, 13 signed-zigzag bridges and seven invalid-input controls. These are rational/integer calculations, not executions extracted from Lean. The first-block mirror suite was rerun separately. The original historical numerical suite was not recovered or rerun.

\subsection*{Provenance and limits}
All eight upstream proof files, the original 28-page corrected manuscript and the first supplement are preserved. The initial clean baseline and a preliminary successful second-block build are separate from the final build that includes the explicit matrix row-action theorem. Eight failed development drafts and their diagnostics are retained, not counted as successful proofs. An unsupported interactive container invocation was replaced by successful noninteractive runtime recovery; every archive-part hash and the combined hash matched.

\textbf{Fourteen manuscript-proved items remain pending.} The next coherent target is the separately weighted dyadic reconstruction. The global finite-difference theorem, three classical inputs and two conjectures have not been promoted by this block.

\textbf{Reproduce.} Run \mcode{bash Run-Portable.sh /path/to/extracted/kit}. This performs unit tests, both independent regression suites and a fresh Lean build. The kit is separate; native Windows is not tested. Precise commands, source hashes, raw compiler logs, scope entries and failures accompany the release. Manuscript-to-Lean correspondence remains a mathematical review, not an automated semantic-equivalence theorem.

\clearpage\subsection{Third Lean Block Supplement}\label{app:sup3}
\noindent\textbf{Checkpoint scope.} This complete source body retains the status of its dated original supplement. Earlier pending or partial-coverage labels are historical; the final core records their resolution.\par\medskip
\subsubsection*{Scope and result}
This release continues the unchanged 28-page \emph{status-corrected} Paper IX manuscript and the checked second block. It completes the recorded scopes of Lemmas 5.1 and 5.2, Theorem 5.3 and Corollary 5.4. All 13 upstream proof modules and both earlier mathematical supplements remain byte-identical.

\textbf{Fourteen of 24 named manuscript-proved items now have formal coverage; ten remain pending.} Three classical inputs and two conjectures are tracked separately. The cumulative corpus has \textbf{219 named theorem declarations and 53 explicit definitions in 17 modules}. The 56 added theorems include supporting identities and exact examples, not 56 independent discoveries.

\begin{center}
\begin{tabular}{@{}p{.48\textwidth}>{\raggedright\arraybackslash}p{.18\textwidth}rr@{}}
\toprule
New module & Scope & Theorems & Defs.\\\midrule
\mcode{DyadicIntegrality} & 5.1, 5.2 & 21 & 8\\
\mcode{DyadicSeries} & 5.1, 5.3 & 11 & 4\\
\mcode{DyadicReconstruction} & 5.3, 5.4 & 10 & 4\\
\mcode{DyadicExamples} & Exact controls & 14 & 0\\\midrule
Total added & Four named items & 56 & 16\\\bottomrule
\end{tabular}
\end{center}

\subsubsection*{The completed mathematical chain}
The inverse coefficients are defined from a genuine unit power series and connected to $z^2/\sinh^2z$. The source is the chapter's explicit rational formula, not an abstract integral input. Its valuation calculation proves the stronger statement that \emph{every even source coefficient is a $2$-adic unit}. Formal-series multiplication and coefficient extraction then yield the Euler identity with the correct falling factorial. The omitted-tail theorem gives divisibility in the actual completed integer ring $\Z_2$.

The clean cumulative run compiled all 17 modules, matched all 272 expected transitive axiom reports and rejected eight invalid-proof controls. All 49 infrastructure and release tests passed. Exact-rational regression checks are recorded separately from the universal proofs.

\subsubsection*{Completion boundary}
This is a formalization companion, not a revised primary manuscript. The unnumbered Bernoulli closed form for the inverse coefficients is not universally formalized here. Formal-series identities do not establish analytic convergence. The corrected global odd-prime finite-difference theorem remains pending. Historical numerical evidence is not presented as a fresh reproduction.
\par\medskip
\pageheading{1. Integral coefficients and an exact unit source}{Lemmas 5.1 and 5.2; manuscript page 14}
Define the rational subring
\[
 \mathcal O_2=\{q\in\Q:|q|_2\leq1\}.
\]
Its ring laws follow from the actual nonarchimedean norm. For every $q\in\mathcal O_2$, the code constructs its element in Mathlib's completed ring $\Z_2$, and proves that the underlying value in $\Q_2$ is the rational cast of $q$. No integrality claim is hidden in a coercion.

\subsection*{The unit kernel and its reciprocal}
Use a squared formal variable $u$ and define
\[
 k_m=\frac{2^{2m+1}}{(2m+2)!},\qquad
 K(u)=\sum_{m\geq0}k_mu^m,\qquad
 \beta_m=[u^m]K(u)^{-1}.
\]
The constant term is proved to be one. The factorial valuation bound gives
\[
 v_2(k_m)=2m+1-v_2((2m+2)!)\geq0.
\]
Thus every $k_m$ belongs to $\mathcal O_2$. A general coefficientwise subring theorem is proved by strong induction using the actual unit-inverse coefficient formula. Applied to $K$, it establishes
\[
 \boxed{\beta_m\in\mathcal O_2\quad\text{for every }m\geq0.}
\]
The reciprocal's coefficients are therefore proved integral; the proof does not assume them integral in order to obtain the final congruence. Their connection to the intended hyperbolic function is proved on the next page.

\subsection*{The source has valuation exactly zero}
At even degree $2r$, define the literal rational source
\[
 R_{2r}=\frac{3^{2r+2}-1}{4(2r+1)(2r+2)}.
\]
Both numerator and denominator are nonzero. The pinned Mathlib lifting-the-exponent theorem gives
\[
 v_2(3^{2r+2}-1)=2+v_2(2r+2).
\]
Because $2r+1$ is odd, the denominator has the same valuation. Consequently,
\[
 \boxed{v_2(R_{2r})=0,\qquad |R_{2r}|_2=1.}
\]
The code proves actual \mcode{IsUnit} status for the constructed $\Z_2$ element. This strengthens the manuscript's integrality conclusion without changing its source. The argument applies at $r=0$, where $R_0=1$, as well as at every later row.

\textbf{Principal targets.}\par
\target{Dyadic.kernel_coefficient_integral, Dyadic.beta_integral}\par
\target{Dyadic.source_valuation, Dyadic.sourcePadic_isUnit}
\par\medskip
\pageheading{2. Connecting the definitions to actual formal series}{Generating-function foundations for Theorem 5.3}
Work in $\Q[[t]]$. The formal exponential and rescaling homomorphism are the same checked objects used in the second block. Define
\[
 \sinh_a(t)=\frac{e^{at}-e^{-at}}2,\qquad
 \cosh_a(t)=\frac{e^{at}+e^{-at}}2.
\]
Let $B(t)=K(t^2)^{-1}=\sum_{m\geq0}\beta_mt^{2m}$ and let
\[
 F(t)=\sum_{r\geq0}R_{2r}\frac{t^{2r}}{(2r)!}.
\]
These series are constructed independently from the definitions on the preceding page.

\subsection*{No inverse of a zero-constant series is assumed}
Coefficient extraction proves the cleared identity
\[
 t^2K(t^2)=\frac{\cosh(2t)-1}{2}.
\]
A separate identity from the formal exponential gives
\[
 \sinh^2t=\frac{\cosh(2t)-1}{2}.
\]
Since $K(t^2)$ has constant term one, its reciprocal exists as a unit. Multiplication then proves
\[
 \boxed{B(t)\sinh^2t=t^2.}
\]
This is the precise formal meaning of $B(t)=t^2/\sinh^2t$. The forced quadratic zero is removed first; the proof does not claim that $\sinh t$ itself is an invertible formal power series.

\subsection*{The explicit rational source supplies its own identity}
From the formula for $R_{2r}$, coefficient arithmetic proves
\[
 \boxed{t^2 F(t)=\frac{\cosh(3t)-\cosh t}{4}.}
\]
The product relation for the formal hyperbolic cosines yields
\[
 F(t)\cosh t=K(4t^2).
\]
This equality is proved after multiplication by $t^2$ and cancellation of that monomial in the formal-series ring. It is not an assumed source-transform relation.

Rescale the reciprocal by $t\mapsto2t$. The unit inverse law and the preceding identity now give
\[
 \boxed{B(2t)F(t)=\sech t.}
\]
The right side is connected to the same integral signed Euler sequence constructed in the previous release. These are equalities of full formal series, not fixed-degree polynomial experiments.

\textbf{Principal targets.}\par
\target{DyadicSeries.beta_original_generating_function}\par
\target{DyadicSeries.source_shift_identity, DyadicSeries.source_times_cosh}\par
\target{DyadicSeries.dyadic_reconstruction_series}
\par\medskip
\pageheading{3. Exact reconstruction and genuine dyadic congruence}{Theorem 5.3 and Corollary 5.4; manuscript pages 14--15}
The reciprocal coefficients $\beta_m$ use an \emph{ordinary} power expansion. In contrast, the source and Euler series use exponential coefficients. Extracting degree $2r$ from $B(2t)F(t)=\sech t$ therefore gives
\[
 \boxed{E_{2r}=\sum_{m=0}^{r}\beta_m2^{2m}\fall{2r}{2m}R_{2(r-m)},}
 \qquad
 \fall{2r}{2m}=\frac{(2r)!}{(2r-2m)!}.
\]
The formal proof uses the actual natural-number descending factorial and proves its factorial-quotient identity for $2m\leq2r$. The odd coefficient terms vanish by a proved parity argument. The binomial EGF product then supplies the falling factorial, rather than treating it as an extra assumed factor.

Both row indexing by $r$ and the literal even-degree version with hypothesis $\operatorname{Even}(n)$ are checked. All rows include degree zero. The exact identity is a rational equality with the Euler coefficient's integer-to-rational cast made explicit.

\subsection*{The finite truncation}
For $M\geq0$, define
\[
 T_M(2r)=\sum_{m=0}^{\min(M,r)}\beta_m2^{2m}\fall{2r}{2m}R_{2(r-m)}.
\]
The integral rational factors are embedded into $\Z_2$ using their proved witnesses. The cast of this $\Z_2$ truncation into $\Q_2$ is proved to equal the cast of the literal rational truncation.

Each omitted term has $m\geq M+1$. Its remaining factors lie in $\Z_2$, so the checked weighted-tail theorem proves
\[
 \boxed{2^{2M+2}\mid E_{2r}-T_M(2r)\quad\text{in }\Z_2.}
\]
This is not divisibility in $\Q$ or $\Q_2$, where division by any nonzero scalar would make the claim vacuous. A second theorem supplies the equivalent rational statement
\[
 \boxed{\frac{E_{2r}-T_M(2r)}{2^{2M+2}}\in\mathcal O_2.}
\]
It follows from actual $\Z_2$ divisibility and the proved rational/completed-field bridge. No precision assertion is passed in as a hypothesis.

\subsection*{Endpoints and scope}
When $M\geq r$, the sum retains every term and $T_M(2r)=E_{2r}$ exactly, in both rational and dyadic coordinates. The zero-error case is included; no finite valuation of zero is used. This is a separate dyadic construction, not substitution of $p=2$ into the odd-length theorem.

\textbf{Principal targets.}\par
\target{DyadicReconstruction.dyadic_reconstruction}\par
\target{DyadicReconstruction.dyadic_supercongruence}\par
\target{DyadicReconstruction.normalized_error_integral}
\par\medskip
\pageheading{4. Exact examples and adversarial checks}{Example 5.5 and controls against incorrect normalizations}
The actual reciprocal recurrence gives
\[
 (\beta_0,\beta_1,\beta_2,\beta_3,\beta_4)
 =\left(1,-\frac13,\frac1{15},-\frac2{189},\frac1{675}\right).
\]
These values are obtained from the general formal-series definition. They are not substituted as independent input data for the reconstruction theorem. The source formula gives
\[
 (R_0,R_2,R_4,R_6,R_8)
 =\left(1,\frac53,\frac{91}{15},\frac{205}{7},\frac{7381}{45}\right).
\]
Using the retained Euler anchor $E_8=1385$, Lean checks the literal decomposition
\[
 E_8=\frac{7381}{45}-\frac{6560}{3}+\frac{163072}{15}
       -\frac{204800}{9}+\frac{229376}{15}=1385.
\]

\subsection*{Truncations and actual valuations}
\begin{center}
\begin{tabular}{@{}crr@{}}
\toprule
Depth $M$ & $T_M(8)$ & Guaranteed exponent $2M+2$\\\midrule
0 & $7381/45$ & 2\\[3pt]
1 & $-91019/45$ & 4\\[3pt]
2 & $398197/45$ & 6\\[3pt]
3 & $-208601/15$ & 8\\[3pt]
4 & $1385$ & 10\\\bottomrule
\end{tabular}
\end{center}
The checked nonzero error valuations are
\[
 E_8-T_0(8)=\frac{54944}{45},\qquad v_2(E_8-T_0(8))=5,
\]
\[
 E_8-T_3(8)=\frac{229376}{15},\qquad v_2(E_8-T_3(8))=15.
\]
They exceed the uniform guarantees; these two examples do not establish a universally stronger exponent.

\subsection*{What the negative controls rule out}
Omitting the rescaling $2^{2m}$ already breaks the degree-two identity: the altered sum is $1$, whereas $E_2=-1$. Confusing ordinary beta coefficients with EGF coefficients would replace $\beta_1=-1/3$ by $-2/3$; a checked inequality excludes that normalization. A deliberately false claim that the depth-zero degree-eight error has valuation at least six is rejected.

These three negative proof files are separate from the active mathematical modules. The verifier also retains the five earlier invalid controls, including the normalized all-residue counterexample. Rejection of an attempted proof is an engineering control; the actual counterexample theorems establish the mathematical failures.

\textbf{Principal targets.}\par
\target{DyadicExamples.beta_one, DyadicExamples.source_anchor_values}\par
\target{DyadicExamples.manuscript_error_zero_valuation}\par
\target{DyadicExamples.manuscript_error_three_valuation}
\par\medskip
\pageheading{5. Verification, provenance and remaining work}{Actual results for the v0.3.0 source snapshot}
\begin{center}
\begin{tabular}{@{}p{.67\textwidth}r@{}}
\toprule
Check & Recorded result\\\midrule
Fresh local Lean modules & 17 / 17\\
Explicit theorem declarations & 219 / 219\\
Explicit definitions & 53 / 53\\
Matched transitive axiom reports & 272 / 272\\
Exact dependency revisions & 9 / 9\\
Invalid-proof controls rejected & 8 / 8\\
Infrastructure and release tests & 49 / 49\\
Independent exact reconstruction rows & 49\\
Independent truncation-depth cases & 2,548\\\bottomrule
\end{tabular}
\end{center}
The clean cumulative run used \textbf{Lean 4.34.0-rc2}, with compiler warnings treated as errors. Mathlib is pinned to
\begin{center}
\mcode{aeb72f1e9981b6d35b426b16dca7759be3e7ba06}.
\end{center}
Every expected axiom report uses only \mcode{propext}, \mcode{Classical.choice} and \mcode{Quot.sound}, or a subset. No additional project axiom or admitted proof is permitted. The verifier audits the explicit declaration inventory, not every generated declaration or all imported Mathlib proofs.

\textbf{Independent regressions.} Exact Python integers and rational fractions check even degrees through 96. The 2,548 truncation cases include 1,176 nonzero errors and 1,372 exact zeros. Zero is handled separately. Kernel and reciprocal integrality and source unit status are checked in 49 cases each. Ninety-seven series coefficients agree. Forty-eight finite comparisons with the Bernoulli closed form pass, but do not prove that unnumbered formula universally. Both older independent Python suites were also rerun. None of these mirrors is Lean-extracted code or a reproduction of the original historical experiments.

\textbf{Development failures.} Eight of 15 captured development checks failed before repair. Their exact inactive source snapshots and diagnostics are retained separately. The seven successful development checks do not substitute for the clean cumulative release run. All 13 upstream proof files, the original corrected PDF and both earlier supplements remain unchanged. Previous metadata and original test copies are preserved under \mcode{provenance/second_block/}.

\textbf{Ten pending named items.} These are 6.2 (finite units), 7.2--7.5 (Mahler and finite differences), 8.3 and 8.5 (residue independence and period), 9.1 and 9.2 (recurrence jumps), and 10.1 (CRT factorization). Three classical inputs and both conjectures remain separate. The corrected global finite-difference claim is not promoted by the dyadic proof.

\textbf{Trust and reproduction.} The pinned compiled Mathlib cache was reused, not independently rebuilt. The integral-source proof invokes the pinned library's LTE theorem; this project does not reprove LTE. Papers I--VIII and the historical numerical suite were not rerun. The source hashes, commands, axiom output and raw logs accompany the release. The external release audit separately records archive integrity and post-packaging replay. Manuscript correspondence remains a mathematical review, not an automated semantic-equivalence theorem.

Run \mcode{bash Run-Portable.sh /path/to/extracted/kit}. The runtime is separate from this ZIP. The full source-level scope and target names are in \mcode{coverage.json} and the third-block ledger.

\clearpage\subsection{Fourth Lean Block Supplement}\label{app:sup4}
\noindent\textbf{Checkpoint scope.} This complete source body retains the status of its dated original supplement. Earlier pending or partial-coverage labels are historical; the final core records their resolution.\par\medskip
\subsubsection*{Scope and result}
This release continues the unchanged 28-page status-corrected manuscript and the checked third block. It completes the recorded scopes of Theorem 7.2, Corollary 7.3 and Theorem 7.4. All 17 upstream proof files and the three earlier mathematical supplements remain unchanged.

\textbf{Seventeen of 24 named manuscript-proved items now have formal coverage; seven remain pending.} Three classical inputs and two conjectures are tracked separately. The cumulative project contains \textbf{272 named theorem declarations and 60 explicit definitions across 21 modules}.

\begin{center}
\begin{tabular}{@{}p{.48\textwidth}>{\raggedright\arraybackslash}p{.18\textwidth}rr@{}}
\toprule
New module & Scope & Theorems & Defs.\\\midrule
\mcode{FiniteMahler} & 7.2 & 19 & 3\\
\mcode{PadicDifferenceBounds} & 7.3 & 12 & 2\\
\mcode{UnitModeDifferences} & 7.4 & 9 & 0\\
\mcode{MahlerExamples} & Exact controls & 13 & 2\\\midrule
Total added & Three items & 53 & 7\\\bottomrule
\end{tabular}
\end{center}
The 53 additions include supporting lemmas and examples, not 53 independent discoveries. The classical Gregory--Newton and binomial-difference results are reused from pinned Mathlib. The new modules connect those results to the actual Paper IX coefficient definitions, interpolation scopes, valuation bounds and unit modes.

\subsubsection*{Two explicit specification boundaries}
A finite sample through $N$ determines $\Delta^s f(n)$ only when $n+s\leq N$. For arbitrary orders, the finite-coefficient identity describes the constructed interpolant, not every unknown continuation of the sample. All three versions are separately checked.

The additive valuation is extended to $\ord_p(0)=\infty$. This avoids incorrectly applying Mathlib's finite-valued default at zero to an inequality that must also cover cancellation and identically zero differences.

\subsubsection*{Verification status}
The fresh cumulative Lean run compiled all 21 modules with warnings as errors, matched all 332 explicit declaration axiom reports, and rejected 12 invalid-proof controls. All 74 infrastructure and release tests passed. Independent exact-rational mirrors are recorded separately from the universal proofs. The archive is also subjected to a fresh-extraction rebuild and integrity audit; its separate report records that outcome.
\par\medskip
\pageheading{1. Reconstruction and finite-data scope}{Theorem 7.2; manuscript page 17}
Let $R$ be any commutative ring. The implementation uses the actual forward-difference operator
\[
 \Delta_h f(n)=f(n+h)-f(n),\qquad c_r=\Delta_1^r f(0).
\]
For a full sequence $f:\N\to R$, pinned Mathlib's Gregory--Newton formula gives
\[
 \boxed{f(n)=\sum_{r=0}^{n}c_r\binom nr.}
\]
The source applies the theorem to these exact definitions. A separate explicit formula relates an $s$-th difference to its original sampled values, so the data dependence is inspectable:
\[
 \Delta_h^s f(n)=\sum_{j=0}^{s}(-1)^{s-j}\binom sj f(n+jh).
\]
Applying Newton reconstruction to $\Delta_1^s f$ and composing the iterated operators proves
\[
 \boxed{\Delta_1^s f(n)=\sum_{j=0}^{n}c_{s+j}\binom nj.}
\]
The apparently infinite coefficient sum is finite at each natural argument. Neither continuity nor convergence is needed for these identities.

\subsection*{An independently defined interpolant}
For a finite coefficient list, define
\[
 I_N(c)(n)=\sum_{r=0}^{N}c_r\binom nr.
\]
The source proves exact coefficient extraction:
\[
 \Delta_1^s I_N(c)(0)=
 \begin{cases}c_s,&s\leq N,\\0,&s>N.\end{cases}
\]
It follows that $I_N(c)$ has the stated shifted-coefficient formula at every order, with zero coefficients beyond $N$. Every difference of order greater than $N$ vanishes. Its coefficients are uniquely determined by its values at $0,\ldots,N$.

\subsection*{What the finite sample does and does not determine}
For $c_r=\Delta_1^r f(0)$ with $r\leq N$, the source proves
\[
 I_N(c)(n)=f(n)\quad(n\leq N),\qquad
 \Delta_1^s I_N(c)(n)=\Delta_1^s f(n)\quad(n+s\leq N).
\]
The condition is essential. The sequences $f(n)=n$ and $g(n)=0$ agree on the sample $\{0\}$, but $\Delta f(0)=1$ and $\Delta g(0)=0$. This is a checked counterexample to unrestricted inference outside the observed stencil.

\textbf{Principal targets.}\par
\target{Mahler.finite_mahler_reconstruction, Mahler.coefficient_interpolant}\par
\target{Mahler.difference_from_coefficients, Mahler.finite_sample_difference}
\par\medskip
\pageheading{2. Weighted bounds in the actual ring $\Z_p$}{Corollary 7.3; manuscript page 18}
Let $p$ be prime. For $x\in\Z_p$, define
\[
 \ord_p(x)=\begin{cases}\infty,&x=0,\\v_p(x),&x\neq0.\end{cases}
\]
The nonzero branch uses Mathlib's actual $p$-adic valuation. The implementation proves its connection to genuine integral divisibility:
\[
 \boxed{L\leq\ord_p(x)\quad\Longleftrightarrow\quad p^L\mid x
 \ \text{in }\Z_p,\qquad L\in\N.}
\]
This is not divisibility in $\Q_p$, where a nonzero scalar can divide any field element. The zero case is proved separately, so zero does not spuriously violate a positive lower bound.

The source also proves multiplicativity of this extended order, the nonarchimedean lower bound for addition, and its finite-sum consequence. The natural-number cast is connected to the actual integer binomial valuation.

\subsection*{The binomial valuation is retained}
Let $f:\N\to\Z_p$, let $c_r=\Delta_1^r f(0)$, and suppose
\[
 w(s+j)\leq\ord_p(c_{s+j})\quad(0\leq j\leq n),
 \qquad w:\N\to\N\cup\{\infty\}.
\]
Set
\[
 B_p(w;n,s)=\min_{0\leq j\leq n}
 \left\{w(s+j)+v_p\!\binom nj\right\}.
\]
Every binomial in this range is nonzero. The code proves
\[
 \boxed{B_p(w;n,s)\leq\ord_p(\Delta_1^s f(n)).}
\]
This is the manuscript minimum after the substitution $r=s+j$. Coefficients outside this range have zero binomial multiplier. Weights equal to infinity are allowed, and a completely vanishing difference has order infinity.

For a uniform coefficient hypothesis, the simpler conclusion follows:
\[
 \boxed{\inf_{j\geq0} w(s+j)\leq\ord_p(\Delta_1^s f(n)).}
\]
For each finite $L\leq B_p(w;n,s)$, a separate theorem yields $p^L\mid\Delta_1^s f(n)$ inside $\Z_p$. Another specialization handles arbitrary finite interpolants with infinite weights beyond their degree.

\textbf{Principal targets.}\par
\target{PadicBounds.nat_le_order_iff, PadicBounds.weighted_difference_bound}\par
\target{PadicBounds.weighted_difference_divisibility}\par
\target{PadicBounds.coefficient_tail_bound, PadicBounds.interpolant_weighted_bound}
\par\medskip
\pageheading{3. Unit modes and residue-field reduction}{Theorem 7.4; manuscript page 18}
For $u$ in any commutative ring, and every natural step $h$, iteration proves
\[
 \boxed{\Delta_h^s(u^n)=u^n(u^h-1)^s.}
\]
The proof unfolds the actual forward-difference operator and uses the power multiplication law. No valuation premise or geometric-difference identity is assumed. More generally, a proved divisibility $p^k\mid u^h-1$ implies
\[
 p^{ks}\mid\Delta_h^s(u^n).
\]

\subsection*{Why a $p$-adic unit supplies one factor of $p$}
Now let $u\in\Z_p^\times$. The actual ring homomorphism
\[
 \pi:\Z_p\longrightarrow\mathbb F_p
\]
 sends $u$ to a nonzero unit. Fermat's finite-field theorem therefore gives
\[
 \pi(u)^{p-1}=1,\qquad \pi(u^{p-1}-1)=0.
\]
Using the proved kernel identity $\ker\pi=p\Z_p$, the source derives
\[
 \boxed{p\mid u^{p-1}-1\quad\text{inside }\Z_p.}
\]
The conclusion is not supplied as an abstract certificate. It follows from the unit hypothesis and the standard residue-field construction.

Combining it with the geometric identity establishes
\[
 \boxed{p^s\mid\Delta_{p-1}^s(u^n),\qquad
 s\leq\ord_p(\Delta_{p-1}^s(u^n)).}
\]
These statements hold for every prime, including $p=2$ with step one. Zero differences and order zero are included. Every multiple of the Fermat step satisfies the same lower bound.

\subsection*{Finite integral linear combinations}
For any finite index set $S$, coefficients $a_i\in\Z_p$ and units $u_i\in\Z_p^\times$, linearity of the actual iterated operator gives
\[
 \Delta_h^s\!\left(\sum_{i\in S}a_i u_i^n\right)
 =\sum_{i\in S}a_i u_i^n(u_i^h-1)^s.
\]
At $h=p-1$, the result is divisible by $p^s$. The empty sum is permitted. This does not prove that the general congruential Euler packets of Theorem 7.5 admit such a representation.

\textbf{Principal targets.}\par
\target{UnitModes.geometric_difference, UnitModes.unit_fermat_divisibility}\par
\target{UnitModes.unit_mode_valuation_gain}\par
\target{UnitModes.finite_unit_modes_valuation_gain}
\par\medskip
\pageheading{4. Exact examples and adversarial checks}{The examples use the same definitions as the general results}
\subsection*{A polynomial anchor}
The actual difference table for $f(n)=n^2$ yields
\[
 (c_0,c_1,c_2,c_3)=(0,1,2,0).
\]
The finite interpolant therefore reconstructs $I_2(c)(7)=49$. This checks the sign and indexing conventions, not just the names of the defined objects.

\subsection*{A sharp binomial-weight improvement}
Let
\[
 F(n)=81\binom n1+3\binom n2+3\binom n3+81\binom n4.
\]
At $n=3$ and $s=1$, the four coefficient valuations are $(4,1,1,4)$, while the binomial valuations are $(0,1,1,0)$. Thus the refined bound is
\[
 B_3(w;3,1)=\min(4,2,2,4)=2,
\]
stronger than the coefficient-only lower bound one. The source checks both the minimum and the actual result:
\[
 \boxed{\Delta F(3)=180,\qquad\ord_3(180)=2.}
\]
The valuation equality is proved in the actual completed ring $\Z_3$, with an independent exact integer divisibility check as well.

\subsection*{A sharp unit-mode gain}
For $p=5$ and $u=2$,
\[
 \Delta_4^2(u^n)\big|_{n=0}=(2^4-1)^2=225.
\]
The source proves $5^2\mid225$ and $5^3\nmid225$. This example attains the guaranteed gain; no additional uniform power of $p$ is claimed.

\subsection*{What fails when hypotheses or conventions are dropped}
The nonunit $u=3$ at $p=3$ gives $\Delta_2(u^n)|_{n=0}=8$, with no factor of three. A constant unit has every positive-order difference zero, hence infinite extended valuation. The finite-sample counterexample on page 2 distinguishes a true interpolant identity from an unsupported claim about unknown continuation.

The release verifier adds four negative Lean controls: an unrestricted finite-sample inference, zero assigned finite extended valuation, nonunit valuation gain, and an extra unit-mode precision claim. All are required to fail with substantive proof diagnostics. Import errors or timeouts do not qualify as successful rejection tests.

\textbf{Source.} The module \mcode{MahlerExamples} contains these exact certificates. Principal target names are listed in the fourth-block ledger.
\par\medskip
\pageheading{5. Verification record and remaining work}{Evidence is retained separately from conclusions}
\subsection*{Cumulative proof checks}
Lean 4.34.0-rc2 rebuilt all 21 local modules into a fresh directory with warnings treated as errors. All 272 theorem declarations and 60 explicit definitions matched their expected transitive axiom reports. Only \mcode{propext}, \mcode{Classical.choice} and \mcode{Quot.sound}, or subsets thereof, occurred. All nine dependency pins matched. Twelve invalid-proof controls were rejected.

All 74 infrastructure and release tests passed. Historical version-specific tests use immutable metadata snapshots. Current source counts, preserved upstream hashes, exact newly covered manuscript items, remaining items and substantive mathematical boundary cases are checked by the active suite.

Ten targeted development runs preceded the cumulative build. Six failed before repair, due to elaboration, coercions, finite-sum normalization, an inferred index type and tactic or linter diagnostics. Those drafts and compiler outputs remain in a separate inactive evidence directory. No failed draft is counted as a passing theorem.

\subsection*{Independent exact-rational mirrors}
The new standard-library Python program uses repeated difference tables and independently evaluated binomial sums. It passed 2,079 reconstructed values, 15,939 in-range sample differences, 2,600 arbitrary-interpolant differences, 900 weighted cases, 966 unit-mode cases and 350 finite-mode-sum cases. Among the weighted cases, 522 differences were zero. It also rejected 189 out-of-range stencils and 15 nonunit inputs.

These are finite tests on rational examples, not executable extraction of Lean functions, a substitute for a universal proof, or a new run of the manuscript's historical experiments.

\subsection*{Remaining named manuscript-proved items}
\begin{center}
\begin{tabular}{@{}ll@{}}\toprule
Item & Remaining formalization\\\midrule
6.2 & Finite unit decomposition certificates\\
7.5 & Global odd-prime Euler finite-difference theorem\\
8.3 & Standard-window residue independence\\
8.5 & Uniform standard-window period\\
9.1 & Companion-matrix jump theorem\\
9.2 & Fibonacci packet recurrence\\
10.1 & Arithmetic CRT packet factorization\\\bottomrule
\end{tabular}
\end{center}
The three classical inputs and two conjectures remain separate. Mathlib's pinned compiled cache was reused, not independently rebuilt. No infinite Mahler convergence theorem or prime-independent arithmetic spine is proved here. The original manuscript and three earlier supplements remain unchanged; the finite-data and zero-valuation clarifications accompany them in this release.

\textbf{Reproduce.} Run \mcode{bash Run-Portable.sh /path/to/kit}. The report in \path{verification/fourth_release/} records the cumulative build. The separately distributed release audit records archive integrity and fresh-extraction replay.

\clearpage\subsection{Fifth Lean Block Supplement}\label{app:sup5}
\noindent\textbf{Checkpoint scope.} This complete source body retains the status of its dated original supplement. Earlier pending or partial-coverage labels are historical; the final core records their resolution.\par\medskip
\subsubsection*{Scope and result}
This release continues the unchanged 28-page status-corrected manuscript and the checked fourth block. It completes Proposition 6.2, Corollary 8.3 and Theorem 10.1 in the scopes stated below. All 21 upstream proof files and the four earlier mathematical supplements remain unchanged.

\textbf{Twenty of 24 named manuscript-proved items now have formal coverage; four remain pending.} Three classical inputs and two conjectures are tracked separately. The cumulative project contains \textbf{333 named theorem declarations and 77 explicit definitions across 28 modules}.

\begin{center}
\begin{tabular}{@{}p{.43\textwidth}p{.22\textwidth}rr@{}}
\toprule
New module & Scope & Theorems & Defs.\\\midrule
\mcode{UnitTriangular} & Integral solver & 5 & 1\\
\mcode{WindowIndependence} & 8.3 & 11 & 3\\
\mcode{FiniteUnits} & 6.2, prime powers & 13 & 4\\
\mcode{DyadicUnits} & 6.2, sign branch & 10 & 4\\
\mcode{UnitRefinement} & Finite compatibility & 4 & 2\\
\mcode{CRTPackets} & 10.1 & 5 & 1\\
\mcode{FiniteArithmeticExamples} & Exact controls & 13 & 2\\\midrule
Total added & Three items & 61 & 17\\\bottomrule
\end{tabular}
\end{center}
The 61 new theorems include supporting lemmas and examples, not 61 independent discoveries. Pinned Mathlib supplies finite-field, unit-group, ideal-power, $p$-adic residue and Chinese-remainder results. The new modules connect these to the actual packet coefficients, explicit finite lifts and existing selectors.

\subsubsection*{The specification matters}
Rational coefficients are not reduced through an imaginary map $\Q\to\F_p$. The proof first constructs their integral lifts in $\Z_p$ and establishes equality in $\Q_p$. Root representatives are actual units at the requested precision, not unchanged residues that need not satisfy the root equation modulo $p^\ell$.

\subsubsection*{Verification status}
The fresh cumulative Lean build compiled all 28 modules with warnings treated as errors, matched 410 declaration axiom reports, and rejected all 17 invalid-proof controls. All 103 infrastructure and release tests passed. Independent finite regressions are recorded separately from the proofs. A separate release audit records the archive's fresh-extraction rebuild and integrity checks.
\par\medskip
\pageheading{1. Integral reduction of the original coefficients}{Corollary 8.3; manuscript pages 19--20}
Fix a prime $p$, an integer $m\geq1$ and a window address $0\leq j<p$. The earlier project defines the rational packet coefficients by the actual recurrence
\[
 \sum_{q=0}^{n}\binom{mpn+j}{mpq}C_{mp,j}(q)=\delta_{n0}.
\]
The previous Lucas theorem proves
\[
 \binom{mpn+j}{mpq}\equiv\binom{mn}{mq}\pmod p,\qquad
 \binom{mpn+j}{mpn}\equiv1\pmod p.
\]
The last congruence makes each diagonal entry a unit in the actual completed ring $\Z_p$.

\subsection*{First construct the integral solution}
For a ring $R$, a triangular row system $T$, right-hand side $b$, and proved unit pivots $d_n=T_{nn}$, define
\[
 s_n=d_n^{-1}\left(b_n-\sum_{q<n}T_{nq}s_q\right).
\]
The new ring-level solver uses well-founded recursion. Its row identities and unique-solution property are proved. The source also proves compatibility of this solver with ring homomorphisms into fields.

Specializing to $R=\Z_p$ constructs $c^{(p)}_{m,j}(n)$. Mapping to $\Q_p$ and comparing with the original rational forward substitution gives
\[
 \boxed{c^{(p)}_{m,j}(n)=C_{mp,j}(n)\quad\text{in }\Q_p.}
\]
Thus every original standard-window coefficient is $p$-integral, for every row. Integrality is a conclusion, not an additional hypothesis on the coefficients.

\subsection*{Then apply the residue map}
Independently define $e^{(m,p)}$ over $\F_p$ by the unit-diagonal system
\[
 \sum_{q=0}^{n}\binom{mn}{mq}e^{(m,p)}_q=\delta_{n0}.
\]
Its recurrence and uniqueness are checked. Applying the actual homomorphism $\pi:\Z_p\to\F_p$ to the integral row, and then Lucas reduction, proves
\[
 \boxed{\pi(c^{(p)}_{m,j}(n))=e^{(m,p)}_n.}
\]
For $j,k<p$, a separate ideal statement gives
\[
 p\mid c^{(p)}_{m,j}(n)-c^{(p)}_{m,k}(n)\quad\text{inside }\Z_p.
\]
The formal solver also covers $m=0$; the manuscript uses $m\geq1$. No global finite-difference theorem or period assertion enters this proof.

\textbf{Principal targets.}\par
\target{WindowIndependence.coeffPadic_coe, .standard_window_integral}\par
\target{WindowIndependence.residue_independence, .reduced_unique}\par
\target{WindowIndependence.two_window_addresses_congruent}
\par\medskip
\pageheading{2. Explicit finite prime-power decomposition}{Proposition 6.2; odd-prime branch}
Let $p$ be prime, $\ell\geq1$, and
\[
 a\in\bigl(\Z/p^\ell\Z\bigr)^\times.
\]
Define actual units in the same ring by
\[
 \boxed{\omega_\ell(a)=a^{p^{\ell-1}},\qquad
 \langle a\rangle_\ell=a\,\omega_\ell(a)^{-1}.}
\]
The Lean source writes $\ell=L+1$, so zero precision is excluded without truncated natural-number subtraction.

\subsection*{The root equation and residue are derived}
Euler's unit-group exponent gives
\[
 \omega_\ell(a)^{p-1}
 =a^{p^{\ell-1}(p-1)}=a^{\varphi(p^\ell)}=1.
\]
Reduction to the actual field $\F_p$ commutes with powers. Its Frobenius identity gives
\[
 \overline{\omega_\ell(a)}
 =\bar a^{p^{\ell-1}}=\bar a.
\]
Consequently,
\[
 \boxed{a=\omega_\ell(a)\langle a\rangle_\ell,\quad
 \omega_\ell(a)^{p-1}=1,\quad
 \overline{\langle a\rangle_\ell}=1.}
\]
The manuscript's finite existence conclusion is obtained directly. The proof does not assume a Hensel lift or an already supplied root representative.

\subsection*{Uniqueness requires a lifting argument}
The source proves the actual kernel statement
\[
 \bar x=0\quad\Longleftrightarrow\quad p\mid x
 \quad\text{in }\Z/p^\ell\Z.
\]
For $x\equiv y\pmod p$, the checked ideal-power lifting theorem gives
\[
 x^{p^{\ell-1}}-y^{p^{\ell-1}}\in(p)^\ell=(0).
\]
If $w^{p-1}=1$, then $w^p=w$ and hence $w^{p^k}=w$ for every $k$. Therefore two roots with the same residue are equal. The decomposition is unique whenever its second factor reduces to one. Multiplicativity of both factors is also proved.

For $p=3,\ell=2,a=2$, the checked values are
\[
 \omega_2(2)=8,\qquad \langle2\rangle_2=7,\qquad
 8^2=1,\quad8\cdot7=2\quad\text{in }\Z/9\Z.
\]
The unlifted representative $2$ is not a square root of one modulo nine.

\textbf{Principal targets.}\par
\target{FiniteUnits.finite_unit_decomposition, .power_lift_congruence}\par
\target{FiniteUnits.omega_unique, .decomposition_unique}\par
\target{FiniteUnits.omega_mul, .principal_mul}
\par\medskip
\pageheading{3. Precision compatibility and the dyadic sign}{Finite refinements and Proposition 6.2's dyadic branch}
For $1\leq k\leq\ell$, let
\[
 \rho_{\ell,k}:\Z/p^\ell\Z\longrightarrow\Z/p^k\Z
\]
be the actual reduction homomorphism. It induces a map on units. The explicit construction satisfies
\[
 \boxed{\rho_{\ell,k}(\omega_\ell(a))
 =\omega_k(\rho_{\ell,k}(a)),}
\]
\[
 \boxed{\rho_{\ell,k}(\langle a\rangle_\ell)
 =\langle\rho_{\ell,k}(a)\rangle_k.}
\]
The proof uses the exponent factorization
$p^{\ell-1}=p^{k-1}p^{\ell-k}$ and the already proved Frobenius invariance of a lifted root. Principal factors also satisfy
\[
 \langle a\rangle_\ell^{\,p^{\ell-1}}=1.
\]
At precision one the principal factor is exactly one. These are coherent finite certificates; no infinite inverse limit or topological unit-group equivalence is constructed in this release.

\subsection*{Why the mod-four branch is separate}
Now let $\ell\geq2$ and $a\in(\Z/2^\ell\Z)^\times$. The actual units of $\Z/4\Z$ are exhausted by $1$ and $-1$. Define
\[
 \varepsilon_\ell(a)=
 \begin{cases}1,&a\equiv1\pmod4,\\-1,&a\equiv3\pmod4,\end{cases}
 \qquad
 \langle a\rangle_{2,\ell}=\varepsilon_\ell(a)a.
\]
The source uses $\ell=L+2$. It proves
\[
 \boxed{a=\varepsilon_\ell(a)\langle a\rangle_{2,\ell},\quad
 \varepsilon_\ell(a)\in\{1,-1\},\quad
 \langle a\rangle_{2,\ell}\equiv1\pmod4.}
\]
The sign is uniquely determined by the residue modulo four, and cancellation gives uniqueness of the principal factor.

The checked example $a=3$ modulo eight has sign $7=-1$ and principal factor $5$:
\[
 7\cdot5=3\pmod8,\qquad5=1\pmod4.
\]
The code separately proves $1=-1$ in $\Z/2\Z$ but $1\neq-1$ in $\Z/4\Z$. Modulo two alone cannot identify the sign. The generic prime-power construction at $p=2$ only asks for a principal residue modulo two and is not silently substituted for this stronger mod-four factorization.

\textbf{Principal targets.}\par
\target{UnitRefinement.omega_compatible, .principal_compatible, .principal_power}\par
\target{DyadicUnits.dyadic_unit_decomposition, .decomposition_unique}
\par\medskip
\pageheading{4. CRT addresses and the actual packet selectors}{Theorem 10.1; manuscript page 22}
Let $a\geq0$, let $M$ be odd, and put $N=2^aM$. Coprimality is derived from oddness:
\[
 \gcd(2^a,M)=1.
\]
The construction uses the actual Chinese-remainder ring equivalence
\[
 \boxed{\Z/N\Z\ \simeq\ \Z/2^a\Z\times\Z/M\Z.}
\]
Its bijectivity and component formula are checked. On a natural representative $r$, the components are precisely its reductions modulo the two factors, not arbitrary coordinate labels.

\subsection*{The algebraic bridge to projectors}
For natural indices $n,r$, the source proves
\[
 n\equiv r\pmod{2^aM}
 \quad\Longleftrightarrow\quad
 \bigl(n\equiv r\pmod{2^a}\text{ and }n\equiv r\pmod M\bigr).
\]
For the original coefficient selector
\[
 (P_{N,r}f)(n)=
 \begin{cases}f(n),&n\bmod N=r,\\0,&\text{otherwise},\end{cases}
\]
take the canonical label $0\leq r<N$. The logical equivalence yields the actual equality of sequence operators
\[
 \boxed{P_{N,r}f=
 P_{2^a,\,r\bmod2^a}\bigl(P_{M,\,r\bmod M}f\bigr).}
\]
This statement holds for coefficients in any type with zero. A separate theorem lifts the equality coefficientwise to the existing formal power-series selectors over any semiring.

The cases $a=0$ and $M=1$ are included. The corresponding one-element residue ring is not excluded. Oddness of $M$ already implies positivity.

\subsection*{What is not asserted}
CRT factors addresses, not the visibility weights or valuation filtrations associated with those addresses. Nothing in this proof gives an equal-weight decomposition at every prime or an additional analytic convergence theorem.

Coprimality cannot be discarded. The distinct classes $0$ and $2$ modulo four both map to $(0,0)$ when both components are reduced modulo two. The source checks this counterexample. It also checks a composite-modulus failure of the Fermat exponent and a coefficient outside the standard window that is not $3$-adically integral.

\textbf{Principal targets.}\par
\target{CRTPackets.address_bijective, .address_components, .residue_equivalence}\par
\target{CRTPackets.packet_factorization, .series_packet_factorization}\par
\target{FiniteArithmeticExamples.outside_window_not_integral}
\par\medskip
\pageheading{5. Verification, regressions and remaining work}{Cumulative release evidence}
The fresh build started with an empty output directory under Lean 4.34.0-rc2 and Mathlib revision
\[
 \mcode{aeb72f1e9981b6d35b426b16dca7759be3e7ba06}.
\]
All nine dependency revisions matched the runtime manifest and repository heads.

\begin{center}
\begin{tabular}{@{}p{.68\textwidth}r@{}}
\toprule
Check & Result\\\midrule
Modules compiled, warnings treated as errors & 28/28\\
Named theorem declarations & 333\\
Explicit definitions audited separately & 77\\
Expected axiom reports matched & 410/410\\
Invalid-proof controls rejected & 17/17\\
Infrastructure and release tests & 103 passed\\\bottomrule
\end{tabular}
\end{center}
Dependencies are among \mcode{propext}, \mcode{Classical.choice} and \mcode{Quot.sound}, or subsets thereof. The source scan excludes admitted proofs, new axioms and selected trust bypasses. The five new rejection controls cover root lifting, primality, dyadic signs, integrality and coprimality.

\subsection*{Independent finite checks}
The exact Python mirrors passed 2,132 standard-window coefficient rows, 1,476 prime-power unit cases, 555 exhaustive root-uniqueness checks, 4,594 precision comparisons and 576 multiplicativity comparisons. The dyadic suite checked 454 decompositions and their uniqueness. CRT covered 1,240 addresses and 6,200 selector cases. Thirteen invalid inputs were rejected.

These are independently implemented finite arithmetic checks, not code extracted from Lean or replacements for its universal proofs. Their deterministic mathematical payload is replayed separately from elapsed time in the fresh-extraction release audit.

\subsection*{Failures and preservation}
Eleven of 18 development checks failed before repair, on lemma interfaces, coercions, finite cases and simplification. Their inactive snapshots and diagnostics are retained. All 21 upstream proof files and four earlier supplements are unchanged. Edition-specific test updates retain archived originals.

\subsection*{Four manuscript-proved items remain pending}
\begin{center}
\begin{tabular}{@{}lp{.79\textwidth}@{}}
7.5 & Global odd-prime Euler finite-difference theorem\\
8.5 & Uniform standard-window period\\
9.1 & Companion jump theorem\\
9.2 & Fibonacci packet recurrence\\
\end{tabular}
\end{center}
Three classical inputs and two conjectures remain separate. Mathlib's cache was reused, not independently rebuilt. Earlier papers and historical experiments were not rerun. The preserved manuscript must be read with the cumulative supplements and amendments.

\clearpage\subsection{Sixth Lean Block Supplement}\label{app:sup6}
\noindent\textbf{Checkpoint scope.} This complete source body retains the status of its dated original supplement. Earlier pending or partial-coverage labels are historical; the final core records their resolution.\par\medskip
\subsubsection*{Scope and result}
This release completes Theorem 9.1 and Corollary 9.2 of the unchanged 28-page status-corrected manuscript. The preceding 28 proof modules and five supplements remain unchanged.

\textbf{Twenty-two of 24 named manuscript-proved items now have formal coverage; two remain pending.} Three classical inputs and two conjectures are tracked separately. The cumulative corpus has \textbf{382 named theorem declarations and 84 explicit definitions across 31 modules}.

\begin{center}
\begin{tabular}{@{}p{.42\textwidth}p{.24\textwidth}rr@{}}
\toprule
New module & Scope & Theorems & Defs.\\\midrule
\mcode{CompanionPackets} & General jump and recurrence & 17 & 2\\
\mcode{FibonacciPackets} & Actual Fibonacci specialization & 18 & 4\\
\mcode{RecurrenceExamples} & Examples and boundary cases & 14 & 1\\\midrule
Total added & Two named items & 49 & 7\\\bottomrule
\end{tabular}
\end{center}
These counts include supporting lemmas and examples, not 49 independent mathematical discoveries. The principal statements are universal over their explicit inputs, not finite numerical observations.

\subsubsection*{The main mathematical connection}
A sampled state $u_k=v_{Nk+j}$ evolves by the actual jump matrix $C^N$. Its recurrence is derived from the characteristic polynomial of that matrix. No scalar recurrence or spectral decomposition is supplied as a substitute for the desired conclusion.

For Fibonacci numbers, the state is independently defined from Mathlib's actual sequence. The Lucas sequence is independently defined by its initial values and recurrence. The trace and determinant of $C^N$ are then proved to supply the sampled recurrence coefficients.

\subsubsection*{Important qualification}
A $d$-dimensional state system supplies a recurrence of \textbf{order at most $d$}. Its minimal recurrence need not have order $d$. The original statement is read as the characteristic-polynomial recurrence, not a minimal-order claim. The formal equations also handle singular matrices, zero-dimensional spaces and zero step.

\par\medskip
\pageheading{1. Actual matrix orbits and arithmetic jumps}{Theorem 9.1; manuscript page 22}
Let $R$ be a commutative ring, $I$ a finite coordinate type and $d=|I|$. For $C\in M_I(R)$ and $v_0\in R^I$, define
\[
 \operatorname{orbit}(C,v_0,n)=C^nv_0,
 \qquad u_k=\operatorname{packet}(C,v_0,N,j,k)=C^{Nk+j}v_0.
\]
The implementation uses actual matrix powers and matrix--vector multiplication. A separately supplied sequence $v$ with
\[
 v_{n+1}=Cv_n
\]
is proved by induction to equal this power orbit. Thus the source-sequence formulation is connected to the construction.

\subsection*{The jump formula}
Associativity and the laws for powers of one matrix give
\[
 \boxed{u_k=C^j(C^N)^kv_0=(C^N)^k(C^jv_0).}
\]
Only powers of the same matrix are interchanged. The argument does not assume that arbitrary matrices commute. In particular,
\[
 u_{k+1}=C^Nu_k,\qquad u_{k+a}=(C^N)^a u_k.
\]
The step matrix is $C^N$, not $C$. A checked example with the Fibonacci matrix and $N=2$ rejects replacing the jump with a single original step.

\subsection*{Refinement of an arithmetic address}
Taking another subsequence with step $M$ and offset $s$ yields
\[
 u_{Mk+s}=v_{(NM)k+(Ns+j)}.
\]
The formal proof establishes this identity for the actual packet function by checking the index arithmetic. It is not inferred from a diagram or a naming convention.

\subsection*{Endpoints and interpretation}
The algebraic identities allow $N=0$, in which case $u_k=v_j$ is constant. Positive-modulus residue packets are the special case $N>0$ and $0\le j<N$. The formula also permits offsets beyond that standard range. No inverse matrix, diagonalization, eigenvalue, or characteristic-zero assumption is needed.

\textbf{Principal targets.}\par
\target{Companion.orbit_of_recurrence}\par
\target{Companion.companion_jump, Companion.source_companion_jump}\par
\target{Companion.packet_step, Companion.packet_advance}\par
\target{Companion.packet_refinement, Companion.packet_zero_step}
\par\medskip
\pageheading{2. From the characteristic polynomial to a recurrence}{Theorem 9.1, continued}
Put $D=C^N$ and let
\[
 \chi_D(X)=\det(XI-D)=X^d+\sum_{a=0}^{d-1}c_aX^a
\]
be its actual characteristic polynomial. The implementation reuses the pinned Mathlib Cayley--Hamilton theorem, expands polynomial evaluation into a finite sum and proves
\[
 \sum_{a=0}^{d}\chi_D[a]D^a=0,
 \qquad \chi_D[d]=1.
\]
Here $\chi_D[a]$ denotes the coefficient at index $a$. The identities hold over any commutative coefficient ring; the zero-ring case is handled explicitly by the proof's subsingleton branch.

\subsection*{The complete vector recurrence}
Apply the matrix identity to the actual state $u_k$. The already proved advance identity identifies each $D^a u_k$ with $u_{k+a}$. Consequently,
\[
 \sum_{a=0}^{d}\chi_D[a]u_{k+a}=0,
 \qquad
 \boxed{u_{k+d}=-\sum_{a=0}^{d-1}c_a u_{k+a}.}
\]
This is the desired recurrence, with its monic leading coefficient and all earlier coefficients supplied by $\chi_D$. It is not an assumption on the packet sequence.

Every coordinate and every $R$-linear scalar observation $\lambda:R^I\to R$ satisfies the corresponding equation:
\[
 \lambda(u_{k+d})=-\sum_{a=0}^{d-1}c_a\lambda(u_{k+a}).
\]
A separate source-sequence theorem transports the recurrence back to the original $v_{Nk+j}$ notation.

\subsection*{Order versus minimal order}
The characteristic polynomial gives a recurrence of order at most $d$. For example, $C=I_2$ gives a degree-two characteristic polynomial, but every orbit is constant and obeys $u_{k+1}=u_k$. Both statements are checked. No minimal-polynomial classification is claimed.

For $d=0$, the unique empty vector is zero and the formula remains valid. For singular matrices, Cayley--Hamilton still applies: invertibility is irrelevant. The nilpotent example $C=\left(\begin{smallmatrix}0&1\\0&0\end{smallmatrix}\right)$ has $C^2=0$ and $\chi_C=X^2$.

\textbf{Principal targets.}\par
\target{Companion.charpoly_sum_powers, Companion.charpoly_coeff_dimension}\par
\target{Companion.packet_charpoly_annihilation}\par
\target{Companion.packet_charpoly_recurrence}\par
\target{Companion.packet_coordinate_recurrence}\par
\target{Companion.packet_observable_recurrence}
\par\medskip
\pageheading{3. Fibonacci and Lucas coefficients are derived}{Corollary 9.2; manuscript page 22}
Define the actual integer matrix and state
\[
 C=\begin{pmatrix}1&1\\1&0\end{pmatrix},
 \qquad v_n=\begin{pmatrix}F_{n+1}\\F_n\end{pmatrix},
 \qquad v_0=\begin{pmatrix}1\\0\end{pmatrix}.
\]
The values $F_n$ are Mathlib's \mcode{Nat.fib}, cast into $\Z$. Its basic recurrence proves $v_{n+1}=Cv_n$. The second coordinate of the packet state is therefore proved to be the intended $F_{Nk+j}$, not a newly assumed sequence.

\subsection*{An independent Lucas sequence}
Define
\[
 L_0=2,\quad L_1=1,\quad L_{n+2}=L_n+L_{n+1}.
\]
Direct finite matrix arithmetic proves $C^2=C+I$. Multiplication by $C^n$ gives $C^{n+2}=C^n+C^{n+1}$. Taking traces shows that $\operatorname{tr}(C^n)$ satisfies the same recurrence and initial values as $L_n$. Induction yields
\[
 \boxed{\operatorname{tr}(C^N)=L_N.}
\]
The implementation separately proves
\[
 \det C=-1,\qquad \det(C^N)=(-1)^N,
 \qquad L_N=2F_{N+1}-F_N.
\]
No Binet roots or diagonalizability argument is needed.

\subsection*{The sampled Fibonacci recurrence}
The actual two-by-two characteristic-polynomial formula gives
\[
 \chi_{C^N}(X)=X^2-L_NX+(-1)^N.
\]
Insert these coefficients into the general packet recurrence and take the second coordinate:
\[
 \boxed{F_{N(k+2)+j}=L_NF_{N(k+1)+j}-(-1)^NF_{Nk+j}.}
\]
The full two-coordinate state satisfies the same equation. All subtraction is in $\Z$, not truncated natural-number subtraction. The result holds for all natural $N,j,k$, including zero step and zero initial values.

\textbf{Principal targets.}\par
\target{FibonacciPackets.state_is_orbit, FibonacciPackets.packet_is_fibonacci}\par
\target{FibonacciPackets.trace_power_is_lucas, FibonacciPackets.det_power}\par
\target{FibonacciPackets.jump_charpoly}\par
\target{FibonacciPackets.fibonacci_packet_recurrence}
\par\medskip
\pageheading{4. Exact examples and meaningful boundary checks}{Example 9.3 and supporting results}
The independently defined Lucas recurrence gives
\[
 (L_0,L_1,L_2,L_3,L_4,L_5)=(2,1,3,4,7,11).
\]
For $N=3$, the general theorem therefore specializes to
\[
 u_{k+2}=4u_{k+1}+u_k.
\]
The implementation proves this for every offset $j$ and every $k$, not only the displayed sample. At $j=2$, direct evaluation of the actual Fibonacci sequence gives
\[
 (F_2,F_5,F_8,F_{11},F_{14},F_{17})=(1,5,21,89,377,1597).
\]
This is the exact manuscript example.

\subsection*{The parity sign cannot be removed}
For $N=2$, the recurrence is
\[
 u_{k+2}=3u_{k+1}-u_k.
\]
At $j=1,k=0$, the correct equation is $5=3\cdot2-1$. Replacing the minus sign with a plus would assert $5=7$. The source proves the inequality, and a separate invalid-proof control attempts and rejects the false equation.

At $N=1$ the ordinary Fibonacci recurrence is recovered. At $N=0$, $L_0=2$ and the packet is constant, so the same formula becomes $u=2u-u$.

\subsection*{Uniqueness from the first two packet values}
For a fixed $N,j$, any integer sequence satisfying the proved sampled recurrence and starting at $F_j,F_{N+j}$ agrees with $F_{Nk+j}$ for all $k$. This is proved by two-step induction. It is a consequence of the established coefficients and recurrence, not a premise used to identify Fibonacci values initially.

\subsection*{Independent finite mirrors}
A pure-Python exact-arithmetic implementation computes characteristic polynomials by the permutation determinant and states by repeated matrix application. It checks 280 matrix systems over integers, rationals and several modular rings, including composite moduli and the zero ring. Separately it computes sampled Fibonacci sequences recursively and compares them with direct Fibonacci evaluation. These are useful regression checks, not Lean extraction or universal proofs.

\textbf{Principal targets.}\par
\target{RecurrenceExamples.manuscript_sequence}\par
\target{RecurrenceExamples.manuscript_step_three, RecurrenceExamples.step_two}\par
\target{FibonacciPackets.fibonacci_packet_unique}\par
\target{RecurrenceExamples.even_step_wrong_sign}\par
\target{RecurrenceExamples.wrong_jump_matrix}
\par\medskip
\pageheading{5. Verification, provenance and remaining work}{Cumulative checks and exact scope}
\textbf{Lean.} The fresh cumulative run uses Lean 4.34.0-rc2 and Mathlib revision
\begin{center}\small\mcode{aeb72f1e9981b6d35b426b16dca7759be3e7ba06}.\end{center}
All 31 modules compile with warnings treated as errors. The inventory contains 382 theorem declarations and 84 explicit definitions; all 466 expected axiom reports must match. Allowed dependencies are \mcode{propext}, \mcode{Classical.choice} and \mcode{Quot.sound}, or subsets. The source-policy check excludes admitted proofs, added axioms and selected trust bypasses. All nine package revision pins are checked.

\textbf{Controls and tests.} Twenty invalid-proof controls are expected to fail for mathematical reasons, not missing imports or timeouts. The three new controls test the Fibonacci sign, the jump matrix and the trace coefficient. All 131 infrastructure and release tests pass. These include 28 new sixth-block tests and preserved earlier arithmetic tests.

\begin{center}
\begin{tabular}{@{}lr@{}}
\toprule
Independent exact finite check & Cases\\\midrule
Matrix systems across seven coefficient domains & 280\\
Jump identities and characteristic recurrences & 1,680 each\\
One-step packet transitions & 1,400\\
Nested packet addresses & 280\\
Fibonacci packet sequences & 357\\
Fibonacci scalar recurrence steps & 11,067\\
Lucas trace, determinant and Fibonacci identities & 81\\
Invalid numerical inputs rejected & 10\\\bottomrule
\end{tabular}
\end{center}
\textbf{Development failures.} Four Lean development checks failed before correction. The errors concerned a matrix-associativity step, recursive simplification, coefficient simplification and an ineffective coordinate for a counterexample. Their source snapshots and compiler logs are retained. An initial unit-test run failed two stale baseline-count assertions; those expected counts were updated, not the mathematical tests weakened.

\textbf{Trust boundary.} Pinned Mathlib supplies Cayley--Hamilton, characteristic-polynomial facts, matrix algebra and the basic Fibonacci recurrence. Its compiled cache is reused, not independently rebuilt. Python mirrors do not execute extracted Lean definitions. The historical numerical suite and earlier papers are not rerun. The original manuscript and all 28 upstream proof files remain unchanged.

\textbf{Remaining pair.} Theorem 7.5, the global odd-prime Euler finite-difference estimate, and Corollary 8.5, the uniform standard-window period, remain pending. The two conjectures are still unproved. This block does not use those claims as assumptions.

Detailed compiler logs, inventories, amendments, the completed item crosswalk, and fresh-extraction integrity evidence accompany the release. Manuscript-to-Lean correspondence is a mathematical review, not an automatically certified semantic equivalence.

\clearpage\subsection{Seventh Lean Block Supplement}\label{app:sup7}
\noindent\textbf{Checkpoint scope.} This complete source body retains the status of its dated original supplement. Earlier pending or partial-coverage labels are historical; the final core records their resolution.\par\medskip
\subsubsection*{Scope and result}
This block completes Corollary 8.5 of the unchanged status-corrected manuscript by an independent finite-field argument. It does \textbf{not} assume the still-pending global finite-difference theorem, item 7.5.

\textbf{Twenty-three of 24 named manuscript-proved items now have formal coverage.} The cumulative corpus contains \textbf{439 named theorem declarations and 94 explicit definitions across 38 modules}. Three classical inputs and two conjectures remain separately tracked. All 31 upstream proof modules and the six earlier supplements remain unchanged.

\begin{center}
\begin{tabular}{@{}p{.54\textwidth}rr@{}}
\toprule
New module & Theorems & Definitions\\\midrule
\mcode{HurwitzConvolution} & 13 & 3\\
\mcode{HurwitzInverse} & 6 & 1\\
\mcode{StandardWindowPeriod} & 8 & 1\\
\mcode{DifferenceFiltration} & 4 & 0\\
\mcode{GlobalEulerBoundary} & 7 & 1\\
\mcode{DworkCriterion} & 7 & 4\\
\mcode{PeriodExamples} & 12 & 0\\\midrule
Total added & 57 & 10\\\bottomrule
\end{tabular}
\end{center}
The counts include supporting lemmas, conditional implications and exact examples. They are not counts of independent discoveries. The conditional implications do not promote their antecedents to proved facts.

\subsubsection*{The independent route}
In characteristic $p$, both the first and $p$th coefficient shifts obey a Leibniz rule for Hurwitz convolution. A section supported at multiples of $m$, with $m\mid p-1$, has equal first and $p$th shifts. Its constructed reciprocal inherits this equality. Identifying that reciprocal with the existing reduced Euler sequence gives its period on the positive-index tail.

\subsubsection*{What remains}
Theorem 7.5 asks for powers as high as $p^{\min(k,mn)}$ and covers nonstandard residue classes. A period modulo $p$ does not establish that stronger result. This release proves a genuine standard-window mod-$p$ fragment and precise conditional interfaces for the remaining argument, while keeping item 7.5 pending.

\par\medskip
\pageheading{1. Coefficient products and reciprocal differentiation}{A characteristic-$p$ argument without factorial division}
For sequences $f,g:\N\to K$, where $K$ is a field, define
\[
 (f\star g)(d)=\sum_{i=0}^{d}\binom di f(i)g(d-i),
 \qquad (U^a f)(d)=f(d+a).
\]
These are Hurwitz coefficients. The notation does not divide by $d!$ in $\F_p$. In particular, the implementation is not ordinary power-series multiplication with undefined inverse factorials in positive characteristic.

The identity sequence is $\mathbf1(d)=\delta_{d0}$. The new source proves the actual higher product rule
\[
 U^a(f\star g)=\sum_{i=0}^{a}\binom ai
       (U^i f)\star(U^{a-i}g).
\]
The argument starts with Pascal's identity for the first shift, then inducts on $a$. All sums are finite at each coefficient.

\subsection*{The prime-shift cancellation}
For prime $p$ and characteristic $p$, every interior binomial coefficient in row $p$ vanishes. Hence
\[
 U^p(f\star g)=(U^pf)\star g+f\star(U^pg).
\]
The first-shift rule has exactly the same form. Both are checked for the literal coefficient product.

Suppose $g(0)=1$. The source constructs $f$ by the full triangular recurrence
\[
 f(d)=\delta_{d0}-\sum_{i<d}\binom di f(i)g(d-i)
\]
and proves $f\star g=\mathbf1$. It also proves cancellation: if $a\star g=0$, then $a=0$, by induction on the first nonzero coefficient.

If $U^pg=Ug$, differentiate the inverse equation with both operators and subtract. The terms involving $f\star(Ug)$ cancel, leaving
\[
 (U^pf-Uf)\star g=0.
\]
The proved cancellation law yields
\[
 \boxed{U^pf=Uf.}
\]
The section identity is an explicit premise of this general reciprocal lemma; the next page proves that premise for the specific section used by the Euler sequence.

\textbf{Principal targets.}\par
\target{Hurwitz.shift_convolution, Hurwitz.prime_shift_convolution}\par
\target{Hurwitz.convolution_left_eq_zero}\par
\target{HurwitzInverse.inverse_convolution}\par
\target{HurwitzInverse.inverse_prime_difference}

\par\medskip
\pageheading{2. Identification with the reduced Euler sequence}{Corollary 8.5; manuscript page 21}
Work over $\F_p=\Z/p\Z$. For positive $m\mid p-1$, take the literal section
\[
 g_m(d)=\begin{cases}1,&m\mid d,\\0,&m\nmid d.\end{cases}
\]
It satisfies $g_m(0)=1$. Since $p\equiv1\pmod m$, the conditions $m\mid d+p$ and $m\mid d+1$ are equivalent. The source therefore proves $U^pg_m=Ug_m$.

Let $f_m$ be the reciprocal constructed on the previous page. A separate strong-induction proof establishes its actual support:
\[
 m\nmid d\quad\Longrightarrow\quad f_m(d)=0.
\]
This support is not included as an assumption of the construction.

\subsection*{The same residual sequence, not a renamed substitute}
Evaluate $f_m\star g_m=\mathbf1$ at $d=mn$. Reindex the surviving summation indices by $i=mq$ using the proved support. The result is
\[
 \sum_{q=0}^{n}\binom{mn}{mq}f_m(mq)=\delta_{n0}.
\]
The fifth block independently constructed the unique sequence $e_n^{(m,p)}$ satisfying this triangular system. Applying that uniqueness theorem gives
\[
 \boxed{f_m(mn)=e_n^{(m,p)}.}
\]
Thus the reciprocal really is connected to the manuscript's Lucas-reduced inverse.

\subsection*{The positive-tail period}
Write $h=(p-1)/m$. For $n\ge1$, evaluate $U^pf_m=Uf_m$ at the nonnegative degree $mn-1$:
\[
 f_m(mn+p-1)=f_m(mn).
\]
Since $mn+p-1=m(n+h)$, the identified sequence satisfies
\[
 \boxed{e_{n+h}^{(m,p)}=e_n^{(m,p)}\quad(n\ge1).}
\]
Induction gives $e_{n+uh}^{(m,p)}=e_n^{(m,p)}$ for every $u\ge0$. This proves a period, not necessarily the least positive period. The proof also includes $p=2$, where $m=1$ and $h=1$.

\textbf{Principal targets.}\par
\target{WindowPeriod.section_prime_shift}\par
\target{WindowPeriod.inverse_off_support, WindowPeriod.inverse_on_support}\par
\target{WindowPeriod.uniform_standard_window_period}\par
\target{WindowPeriod.multiple_period}

\par\medskip
\pageheading{3. Actual packet congruences and sharp boundaries}{Consequences for the standard window; not the full Theorem 7.5}
The fifth block constructed integral lifts $\widehat C_{mp,j}(n)\in\Z_p$ of the original rational coefficients when $0\le j<p$, and proved that their residues equal $e_n^{(m,p)}$.

Applying the actual residue map and its kernel ideal to the new period gives
\[
 \boxed{p\mid\widehat C_{mp,j}(n+uh)-\widehat C_{mp,j}(n)
       \quad\hbox{in }\Z_p}
\]
for $n\ge1$ and $u\ge0$. Every higher positive difference is a difference of terms already divisible by $p$. Consequently,
\[
 p\mid\Delta_{uh}^{k}\widehat C_{mp,j}(n)\qquad(k\ge1,\ n\ge1).
\]
A ring-homomorphism lemma identifies the image of these differences in $\Q_p$ with the differences of the original rational coefficient function. This is not vacuous field divisibility. The valuation version uses order infinity at zero.

\subsection*{Three checked boundaries}
\textbf{The initial row.} For $(p,m)=(3,2)$, the same reduced recurrence gives $e_0=1$ and $e_1=-1=2$ in $\F_3$. Its proposed period is one, so the positive-tail result cannot in general include row zero.

\textbf{The divisor hypothesis.} At $(p,m)=(5,3)$, where $3\nmid4$, the independently computed reduced row is
\[
 (e_0,e_1,e_2,e_3)=(1,4,4,2).
\]
Replacing $h$ by the integer quotient $\lfloor4/3\rfloor=1$ does not give a period: $e_3\ne e_2$.

\textbf{No automatic extra power.} The original rational triangular coefficients satisfy
\[
 C_{6,0}(1)=-1,\quad C_{6,0}(2)=923,\quad
 v_3\bigl(C_{6,0}(2)-C_{6,0}(1)\bigr)=v_3(924)=1.
\]
A period modulo $3$ therefore cannot simply be reported as a period modulo $9$. This is consistent with the global theorem's proposed first-difference bound; it is not a counterexample to Theorem 7.5.

\textbf{Principal targets.}\par
\target{WindowPeriod.standard_window_period_divisibility}\par
\target{GlobalEulerBoundary.standard_window_all_orders_mod_p}\par
\target{GlobalEulerBoundary.lifted_difference_coe}\par
\target{PeriodExamples.period_can_fail_at_zero}\par
\target{PeriodExamples.first_difference_precision_is_one}

\par\medskip
\pageheading{4. Exact interfaces for the remaining global argument}{Conditional results are not counted as completion of item 7.5}
Let $p$ be prime and $m$ a positive divisor of $p-1$. For a sequence $e$ in a commutative coefficient ring, define its actual packed Hurwitz sequence
\[
 X(d)=\begin{cases}e(d/N),&N\mid d,\\0,&N\nmid d,\end{cases}
 \qquad N=mp.
\]
Let $L=p(p-1)$ and $\delta=U^{p+L}-U^p=U^{p^2}-U^p$. The source proves the operator dictionary
\[
 (\delta^sX)(d)=\Delta_L^s X(d+ps).
\]
At the valid nonnegative coefficient degree $d=mpn-ps$, with $s\le mn$, this becomes
\[
 (\delta^sX)(mpn-ps)=\Delta_h^s e(n),\qquad h=(p-1)/m.
\]
The packed sequence and difference operations are definitions. Their coefficient correspondence is proved before any divisibility assumption is introduced.

\subsection*{A sufficient condition, explicitly conditional}
Define flatness by the actual coefficientwise divisibility assertion
\[
 \delta^sX(d)\in p^sR\qquad(s,d\ge0).
\]
The new theorem proves that this hypothesis implies
\[
 \boxed{p^{\min(k,mn)}\mid\Delta_h^k e(n).}
\]
First, coefficient extraction proves the triangular region $k\le mn$. A separate induction on difference order proves that this region is equivalent to the complete saturated estimate. This implication is useful over $\Z_p$; interpreting it over a field would not supply the intended integral content.

\subsection*{Normalization is also an explicit implication}
Using the existing definition $E^{(N,j)}_{Nn}=j!C_{N,j}(n)$, the source proves that multiplication by $j!$ commutes with every iterated difference. If the unnormalized difference has precision $A$, the normalized difference has precision $A-v_p(j!)$. Precision is defined as
\[
 x=0\quad\hbox{or}\quad v_p(x)\ge A,
\]
so zero and negative factorial-loss thresholds are handled correctly. For $j<p$, the loss is zero.

\textbf{Not established here.} The general unnormalized Euler family has not been proved to possess the required flat integral lift, especially across singular residue intervals. The retrieved 18-page supporting prose proof uses weighted Amice/Mahler bounds, singular ratios and nonsingular propagation. Those obligations remain, and the full arbitrary-multiple global estimate is not promoted by these interfaces.

\par\medskip
\pageheading{5. Verification record and reproducibility boundary}{Evidence accompanies the source, not only the summary}
The release verifier compiles the cumulative 38-module project into an initially empty output directory, with warnings treated as errors. It checks all \textbf{439 theorem declarations and 94 explicit definitions}, requiring all \textbf{533 expected axiom reports} and allowing only \mcode{propext}, \mcode{Classical.choice} and \mcode{Quot.sound}, or subsets thereof.

The pinned runtime is Lean 4.34.0-rc2. The Mathlib revision is
\begin{center}\footnotesize\mcode{aeb72f1e9981b6d35b426b16dca7759be3e7ba06}.\end{center}
Nine dependency heads are compared with their recorded manifest pins. The compiled Mathlib cache is reused, not independently rebuilt. The source checker rejects admitted proofs and selected trust-bypass constructs; the transitive axiom audit is a separate check.

\subsection*{Negative controls and infrastructure}
The verifier retains its earlier 20 invalid-proof controls and adds three: a false row-zero period, a false nondivisor period, and an unjustified extra power of the modulus. Missing imports, timeouts and compiler panics do not count as successful rejections.

The infrastructure suite contains \textbf{162 tests}. New tests verify the one-item coverage transition, preservation of the 31 upstream sources, the explicit conditional flatness premise, resolution of partial targets, and continued pending status of 7.5. Historical census tests use their frozen release snapshots. Current census assertions use the new inventory.

\subsection*{Independent exact finite checks}
The new rational and finite-field mirrors passed 4,848 reciprocal coefficient rows, 2,592 higher-Leibniz cases, 4,520 reciprocal prime-shift cases, 3,720 period comparisons, and 1,476 standard-window rational rows. They also checked 1,104 rational period differences, 620 higher differences modulo $p$, and 586 Dwork extraction interfaces. Eleven invalid input configurations were rejected.

A further 4,096 finite tests of the remaining global bound passed. \textbf{These do not prove that bound and do not change the coverage count.} All mirrors are independent Python implementations, not execution extracted from Lean. The original historical numerical suite was not rerun.

\subsection*{Scope, preservation and failures}
The original manuscript and all six earlier supplements remain unchanged. The new amendment replaces the period's dependence on the still-pending global theorem with the direct proof. Development failures, source snapshots and logs are kept separately from the release run. A second clean build from a freshly extracted archive is recorded in the companion release audit.

Run \mcode{bash Run-Portable.sh /path/to/kit} to repeat the current unit tests, cumulative finite mirrors and Lean verification. Run \mcode{python scripts/check_manifest.py} to check the delivered file snapshot. The manifest does not claim that every prose sentence is equivalent to its Lean type; that correspondence remains a mathematical review.

\clearpage\subsection{Eighth Lean Block Supplement}\label{app:sup8}
\noindent\textbf{Checkpoint scope.} This complete source body retains the status of its dated original supplement. Earlier pending or partial-coverage labels are historical; the final core records their resolution.\par\medskip
\subsubsection*{Scope and result}
The all-residue Theorem 7.5 remains pending. This release proves its \textbf{full saturated higher-power estimate throughout the standard window}, rather than only a period modulo $p$. For every prime $p$, positive $m\mid p-1$, $0\le j<p$, and natural $n,k,u$, put $h=(p-1)/m$. Then
\[
\boxed{\quad
\Delta_{uh}^{k}C_{mp,j}(n)=0\quad\text{or}\quad
v_p\!\left(\Delta_{uh}^{k}C_{mp,j}(n)\right)\ge\min\{k,mn\}.
\quad}
\]
The unnormalized family $j!C_{mp,j}(n)$ satisfies the same bound. Both statements concern the original rational coefficient function, identified with its actual integral $p$-adic lift. They include $p=2$, $k=0$, $n=0$, and the additional harmless case $u=0$.

\textbf{Coverage remains 23 of 24 named manuscript-proved items.} Three classical inputs and two conjectures remain separate. The new result covers every permissible residue when $m=1$, but not the range $p\le j<mp$ when $m>1$.

\begin{center}\begin{tabular}{@{}p{.58\textwidth}rr@{}}\toprule
New module & Theorems & Definitions\\\midrule
\mcode{DworkAlgebra} & 7 & 3\\
\mcode{DworkClosure} & 20 & 2\\
\mcode{HurwitzDworkClosure} & 15 & 2\\
\mcode{ResidueZeroDwork} & 9 & 2\\
\mcode{DifferenceStepMultiples} & 5 & 1\\
\mcode{EulerDefect} & 4 & 2\\
\mcode{StandardWindowDwork} & 14 & 1\\
\mcode{StandardWindowPrecision} & 5 & 0\\
\mcode{DworkWindowExamples} & 8 & 0\\\midrule
Total added & 87 & 13\\\bottomrule
\end{tabular}\end{center}
The cumulative corpus has \textbf{526 explicit theorem declarations and 107 explicit definitions in 47 modules}. Counts include support lemmas and examples, not independent discoveries. All 38 upstream proof modules and seven earlier supplements are preserved.

\par\medskip
\pageheading{1. The characteristic-zero product rule}{The normalized defect is not an ordinary derivation}
Let $A$ be a commutative algebra over a characteristic-zero field $K$, and let $U$ be an actual $K$-derivation. The source defines
\[
 B=\frac{U^{p^2}-U^p}{p},\qquad
 M_q(f,g)=\sum_{a=1}^{q-1}\frac{1}{p}\binom qa
          (U^a f)(U^{q-a}g),\quad q=p,p^2.
\]
The coefficients of each $M_q$ are proved integral: the middle entries of a prime-power binomial row are divisible by $p$. In Lean they are implemented as natural quotients, with exact division proved before the product identity is used. Induction on derivative order establishes the full Leibniz formula and hence
\[
\boxed{B(fg)=B(f)g+fB(g)+M_{p^2}(f,g)-M_p(f,g).}
\]
The middle terms are essential. At $p=3$, $f=t$ and $g=t^2$, the constant coefficient of $B(fg)$ is $-2$, whereas that of $B(f)g+fB(g)$ is zero. This inequality is itself Lean-checked.

\subsection*{All orders, without a circular closure assumption}
Let $I\subset A$ be a subring stable under $U$. Define
\[
 \operatorname{Bounded}_r(f)\iff B^s f\in I\ (0\le s\le r),
 \qquad \operatorname{Flat}(f)\iff B^s f\in I\ (s\ge0).
\]
The definitions contain no product or inverse closure requirement. Since $B$ commutes with $U$, derivatives of a bounded element remain bounded. Induction on $r$ in the displayed product rule then proves closure under multiplication. The finite middle sums use the lower-order induction hypothesis, not the desired next-order conclusion.

For an actual inverse $fg=1$, applying the product rule and multiplying by $g$ gives
\[
 B(g)=-g^2B(f)-gM_{p^2}(f,g)+gM_p(f,g).
\]
If $g\in I$, induction on $r$ proves boundedness of $g$ from boundedness of $f$. At each successor stage the preceding stage supplies boundedness of $g$ needed on the right. Passing to all finite orders proves reciprocal closure for an integral inverse.

\textbf{Principal targets.}\par
\target{DworkAlgebra.B_mul}\par
\target{DworkClosure.bounded_mul, DworkClosure.bounded_inverse}\par
\target{DworkClosure.flat_mul, DworkClosure.flat_inverse}

\par\medskip
\pageheading{2. The actual integral Hurwitz lattice}{Connecting differential expressions to coefficient divisibility}
The specialization works inside $\Q_p[[t]]$ and uses factorial-scaled coefficients
\[
 a_f(d)=d![t^d]f,\qquad
 I_p=\{f\in\Q_p[[t]]:a_f(d)\in\Z_p\text{ for every }d\}.
\]
The source constructs $I_p$ as a subring. Actual exponential-series multiplication gives
\[
 a_{fg}(d)=\sum_{r=0}^{d}\binom dr a_f(r)a_g(d-r),
 \qquad a_{Uf}(d)=a_f(d+1).
\]
These identities prove multiplication and derivative stability of the lattice. Membership does not require ordinary power-series coefficients to lie in $\Z_p$.

\subsection*{The reciprocal is genuinely integral}
For $f\in I_p$ with constant coefficient one, the usual power-series reciprocal is used. Its inverse equation yields the exact coefficient recurrence
\[
 a_{f^{-1}}(d)=\delta_{d0}-\sum_{r=0}^{d-1}
     \binom dr a_{f^{-1}}(r)a_f(d-r).
\]
Strong induction proves that all of these coefficients are integral. Combining this construction with the preceding all-order theorem proves flatness of the \emph{actual reciprocal}. No flatness or integrality of that reciprocal is assumed in the normalized-series specialization.

\subsection*{The earlier coefficient criterion is recovered exactly}
Define the literal coefficient defect
\[
 (\delta a)(d)=a(d+p^2)-a(d+p).
\]
The new source proves, for every $s,d$,
\[
 a_{B^s f}(d)=p^{-s}(\delta^s a_f)(d).
\]
For an integral coefficient $x$, it proves the bridge
\[
\boxed{p^{-s}x\in\Z_p\quad\Longleftrightarrow\quad p^s\mid x\text{ in }\Z_p.}
\]
Thus the differential definition is equivalent to the existing coefficient-level Dwork criterion. Division occurs in $\Q_p$, while the assertion after division is membership in $\Z_p$. It is not divisibility in the field. Zero coefficients satisfy every finite precision condition.

\textbf{Principal targets.}\par
\target{HurwitzDworkClosure.integral_inverse}\par
\target{HurwitzDworkClosure.flat_product, HurwitzDworkClosure.flat_reciprocal}\par
\target{HurwitzDworkClosure.decode_B_iterate}\par
\target{HurwitzDworkClosure.normalized_mem_iff, HurwitzDworkClosure.flat_encode_iff}

\par\medskip
\pageheading{3. A nonsingular Euler-operator argument}{The standard interval without an analytic power-series-in-the-index premise}
On Hurwitz coefficients, the actual Euler operator plus $c$ is
\[
 (\Lambda_c a)(d)=(d+c)a(d).
\]
Direct induction proves its commutator with every defect iterate:
\[
 \delta^s\Lambda_c a
 =\Lambda_c\delta^s a
  +sp\bigl(pU^{p^2}-U^p\bigr)\delta^{s-1}a,
 \qquad s\ge1.
\]
Suppose $a$ is integral and supported at degrees divisible by $p$, and $p\nmid c$. Every supported pivot $d+c$ is proved to be a unit of $\Z_p$. The defect preserves this support. Induction in the commutator identity therefore proves
\[
 \operatorname{Flat}(\Lambda_c a)\Longrightarrow\operatorname{Flat}(a).
\]
This is a reflection theorem on the integral supported domain. It does not assume flatness of the proposed preimage, and does not assert invertibility at a singular pivot $p\mid c$.

\subsection*{Apply it to the explicit normalized sections}
Set $N=mp$ and, for $0\le j<p$, define integral Hurwitz coefficients
\[
 g_j(d)=
 \begin{cases}
 \displaystyle\binom{d+j}{d}^{-1},&N\mid d,\\[1mm]
 0,&N\nmid d.
 \end{cases}
\]
The inverse binomial is constructed using the actual unit pivots already proved by Lucas reduction. Its image in $\Q_p$ is proved equal to the displayed reciprocal.

For $j=0$, $g_0$ is the literal $N$-divisibility indicator. Because $N\mid p(p-1)$, the two shifts in $\delta g_0$ agree and $\delta g_0=0$. It is flat at every order. The actual binomial identity gives
\[
 \Lambda_{j+1}g_{j+1}=(j+1)g_j.
\]
As $1\le j+1<p$, the nonsingular reflection theorem inductively proves flatness of every $g_j$ in the standard window. There is no uniform-convergence assumption in this argument.

\textbf{Principal targets.}\par
\target{EulerDefect.iterated_commutator, EulerDefect.flat_of_euler_flat}\par
\target{StandardWindowDwork.kernel_coe, StandardWindowDwork.kernel_euler_step}\par
\target{StandardWindowDwork.kernel_flat}

\par\medskip
\pageheading{4. From these sections to the original Euler packets}{Actual reciprocal, actual rational sequence, every order and step multiple}
The factorial-ratio identity is proved coefficient by coefficient:
\[
 \sum_{d\ge0}g_j(d)\frac{t^d}{d!}
 =j!\sum_{r\ge0}\frac{t^{Nr}}{(Nr+j)!}.
\]
The right side is precisely the normalized source kernel from the first block, not a replacement family. The earlier rational identity for its reciprocal is transported to $\Q_p[[t]]$. A separate encoding theorem identifies the reciprocal's coefficients with the already constructed $\Z_p$ lifts of the original $C_{N,j}(r)$.

All-order reciprocal closure therefore gives flatness of the actual packed Euler coefficient sequence. The earlier coefficient-extraction theorem now has a \emph{proved antecedent} for every $j<p$, and yields
\[
 p^{\min(k,mn)}\mid\Delta_h^k C_{mp,j}(n)\quad\text{in }\Z_p,
 \qquad h=(p-1)/m.
\]
A new finite telescoping argument proves the extension to arbitrary natural multiples of $h$. Its first-step identity is
\[
 f(n+uh)-f(n)=\sum_{a=0}^{u-1}\Delta_h f(n+ah).
\]
Commutation of differences and induction on their order transfer any fixed ideal bound valid on a forward tail. The required precision $\min(k,mn)$ is nondecreasing as the starting row advances. Hence
\[
 \boxed{p^{\min(k,mn)}\mid\Delta_{uh}^{k}C_{mp,j}(n)\quad(0\le j<p).}
\]

The scalar-extension bridge identifies these differences with the original rational coefficient function. The rational theorem explicitly allows zero differences, avoiding the library's finite valuation convention at zero. Since $j!$ is a unit for $j<p$, both normalized and unnormalized versions have the same precision.

\textbf{Scope.} The proof includes all primes, all rows and orders, and all natural step multipliers. For $m=1$, the standard window is the complete allowed residue range. For larger $m$, residues $p\le j<mp$ are not covered by this theorem. The earlier one-power period is not the source of the new higher-power estimate.

\textbf{Principal targets.}\par
\target{StandardWindowDwork.window_is_reciprocal, StandardWindowDwork.standard_window_flat}\par
\target{StandardWindowDwork.standard_window_global_bound}\par
\target{StandardWindowPrecision.normalized_window_precision}\par
\target{StandardWindowPrecision.unnormalized_window_precision}\par
\target{StandardWindowPrecision.modulus_p_complete_family}

\par\medskip
\pageheading{5. Examples, verification and the remaining gap}{The coverage count is unchanged; the unresolved theorem now has a stronger proved part}
The original rational recurrence gives
\[
 C_{6,2}(1)=-\frac1{28},\quad C_{6,2}(2)=\frac{425}{364},\quad
 C_{6,2}(3)=-\frac{2039759}{2660}.
\]
Consequently, the new source proves
\[
 \Delta_1^2 C_{6,2}(1)=-\frac{949959}{1235},\qquad
 v_3\!\left(\Delta_1^2 C_{6,2}(1)\right)=2.
\]
This attains the bound $\min(2,2\cdot1)=2$. A claim of a third power is rejected. The separate cubic example proves that the normalized differential defect is not an ordinary derivation. Controls also reject an asserted unit pivot at $c=p$ and the prime-power binomial-divisibility formula at a composite base.

\subsection*{Recorded checks}
The cumulative verifier rebuilds all 47 modules and audits all 633 explicit declarations: 526 theorems and 107 definitions. Warnings are errors, unexpected axiom dependencies are rejected, and 27 deliberately invalid proof controls are tested. The allowed transitive axioms remain \mcode{propext}, \mcode{Classical.choice} and \mcode{Quot.sound}. Exact results, compiler commands and source hashes are in the accompanying machine-readable reports.

All 189 infrastructure and release tests passed. Independent exact-arithmetic mirrors checked 3,040 normalized and 3,040 unnormalized bounds in 95 standard-window families, 6,596 source defects and 6,596 reciprocal defects, 75 product identities, 360 inverse-closure cases, 720 Euler commutators and 150 telescoping identities. They are finite Python mirrors, not code extracted from Lean and not substitutes for the universal proofs.

Failed development snapshots and diagnostics are preserved separately. Historical source files and earlier supplements remain unchanged. Mathlib's pinned compiled cache is reused, not independently rebuilt; the original historical numerical suite is not rerun. The manuscript-to-formal correspondence remains a mathematical review.

\subsection*{What is still unproved}
The all-residue Theorem 7.5 requires the additional range $p\le j<mp$. Its supporting prose proof treats singular ratios at $c=p,2p,\ldots,(m-1)p$ through weighted Amice estimates and a dual Mahler pairing, then propagates across the intervening indices. Those actual-family estimates and the full remaining propagation are not supplied by this release. A unit argument at $p\nmid c$ cannot be used at $p\mid c$.

\textbf{Release boundary: 23/24 named manuscript-proved items, with Theorem 7.5 still pending.} The two conjectures remain unproved. The proved standard-window result is an infinite family, but no numerical proportion of the remaining work is inferred from the width of that window.

\clearpage\subsection{Singular Amice Supplement}\label{app:sup9}
\noindent\textbf{Checkpoint scope.} This complete source body retains the status of its dated original supplement. Earlier pending or partial-coverage labels are historical; the final core records their resolution.\par\medskip
\subsubsection*{Scope and result}
Let $p$ be prime, let $m\mid p-1$, and put $N=mp$. For $1\le a<m$, put $c=ap$ and define the original exponential section and ratio by
\[
 H_{N,c}(t)=\sum_{n\ge0}\frac{t^{Nn+c}}{(Nn+c)!},\qquad
 R_c(t)=\frac{tH_{N,c-1}(t)}{H_{N,c}(t)}.
\]
The quotient is interpreted after removing its forced zero. This block proves, for the actual rational coefficients $\alpha_d$,
\[
 R_{ap}(\log(1+T))=\sum_{d\ge0}\alpha_dT^d,
 \qquad
 \boxed{\alpha_d=0\ \text{or}\ v_p(\alpha_d)\ge-\left\lfloor d/p\right\rfloor.}
\]
Equivalently, every $p^{\lfloor d/p\rfloor}\alpha_d$ belongs to the actual completed ring $\Z_p$. The zero-safe alternative is explicit in the formal rational statement.

This is the weighted estimate of Proposition 4.7 in the supporting argument. It is \textbf{not yet the singular higher finite-difference theorem}. The full all-residue Paper IX Theorem 7.5 remains pending, and full manuscript coverage stays at \textbf{23 of 24}.

\begin{center}\begin{tabular}{lrr}\toprule
Inventory & v0.8.0 & This release\\\midrule
Proof modules &47&57\\
Explicit theorem declarations &526&630\\
Explicit definitions audited separately &107&128\\
Fully covered manuscript items &23/24&23/24\\\bottomrule
\end{tabular}\end{center}
The 104 added declarations include 30 declarations in four previously compiled research modules, now included in the cumulative audit, and 74 declarations in six new modules. They include supporting lemmas and examples, not independent discoveries. All 47 earlier checked proof modules are preserved byte for byte.

\textbf{No hidden estimate.} The final weighted-ratio theorem assumes primality, the modulus divisibility and the specified residue bounds. It does not assume weighted section membership, reciprocal closure, the ratio identity, a Mahler pairing or the desired higher-difference bound. Pages 2--5 describe the constructed objects and proved bridges.
\par\medskip
\subsubsection*{1. Integral cyclic coefficients}
Work first in an arbitrary commutative ring $R$. If $z^p=1$, set
\[
 I=(p,z-1),\qquad P=(p).
\]
Expanding $((z-1)+1)^p=1$ and using divisibility of every interior binomial coefficient in prime row $p$ gives
\[
 (z-1)^p\in PI.
\]
Every other degree-$p$ product of the two generators of $I$ also belongs to $PI$. The formal ideal proof therefore establishes
\[
 \boxed{I^p\subseteq PI.}
\]
For $q=p\lfloor q/p\rfloor+(q\bmod p)$, strong induction and the factorial valuation identity give
\[
 \boxed{I^q\subseteq P^{v_p(q!)}}.
\]
In detail, write $b=\lfloor q/p\rfloor$. Then
\[
 I^q\subseteq P^b I^b\subseteq P^{b+v_p(b!)}=P^{v_p(q!)}.
\]
No integral-domain hypothesis or division in a field is used here.

Now take the integral cyclic group ring
\[
 R_N=\Z[\Z/N\Z],\qquad x=[1],\qquad x^N=1,
\]
and assume $N\mid p(p-1)$. Thus $(x^{p-1})^p=1$. Each genuine consecutive block
\[
 \prod_{i=0}^{p-1}\bigl(x-(pb+i)\bigr)
\]
belongs to $(p,x^{p-1}-1)$. This follows from the actual polynomial identity
\[
 \prod_{i\in\F_p}(X-i)=X^p-X
\]
modulo $p$, not from a hypothesized valuation of the product.

For the falling factorial $(x)_d=\prod_{i=0}^{d-1}(x-i)$, grouping into consecutive blocks gives
\[
 (x)_d\in I^{\lfloor d/p\rfloor}
 \subseteq\bigl(p^{v_p(\lfloor d/p\rfloor!)}\bigr).
\]
Every integral residue coordinate therefore satisfies
\[
 \boxed{p^{v_p(\lfloor d/p\rfloor!)}\mid [(r)](x)_d\quad\text{in }\Z.}
\]
Dividing a coordinate by $d!$ is done only afterward, using
$v_p(d!)=\lfloor d/p\rfloor+v_p(\lfloor d/p\rfloor!)$. This establishes the weighted denominator estimate before any field-valued divisibility could make it vacuous.

\textbf{Modules.} \mcode{CyclotomicFiltration}, \mcode{PrimeBlockPolynomial}, and \mcode{CyclicAmiceCoefficients}.
\par\medskip
\subsubsection*{2. Connection to the actual section}
An integral coordinate estimate becomes relevant only after connecting it to the manuscript's generating functions. For a residue $r\in\Z/N\Z$, define
\[
 H_r(t)=\sum_{d\ge0}\mathbf1_{d\equiv r\ (N)}\frac{t^d}{d!},
 \qquad F_r(T)=H_r(\log(1+T)).
\]
The implementation uses the actual formal logarithm and substitution in $\Q[[T]]$. It proves
\[
 H_r'=H_{r-1},\qquad (1+T)F_r'=F_{r-1},
\]
and hence the coefficient recurrence
\[
 (d+1)\coef{T}{d+1}F_r
 =\coef{T}{d}F_{r-1}-d\coef{T}{d}F_r.
\]
The initial values are $F_r(0)=\mathbf1_{r=0}$.

Independently, let $b_{d,r}$ be the residue-$r$ coordinate of $(x)_d$ in $R_N$. Actual group-ring multiplication gives
\[
 b_{0,r}=\mathbf1_{r=0},\qquad
 b_{d+1,r}=b_{d,r-1}-d\,b_{d,r}.
\]
Induction comparing these two recurrences proves
\[
 \boxed{d!\coef{T}{d}H_r(\log(1+T))=b_{d,r}.}
\]
Thus the cyclic coefficient is not an unrelated stand-in for a section coefficient. For the canonical label $0\le c<N$, the formal section is also proved equal to the original \mcode{ExponentialSection.}\allowbreak\mcode{originalSection} used in the existing Euler project.

\subsection*{Weighted membership}
The class used in the formalization is
\[
 \Wp=\left\{F\in\Q_p[[T]]:
 p^{\lfloor d/p\rfloor}\coef{T}{d}F\in\Z_p\ \text{for every }d\right\}.
\]
The integral-coordinate estimate from page 2 now proves
\[
 \boxed{H_r(\log(1+T))\in\Wp}
\]
for every residue $r$, including the singular ones. Rational coefficients are mapped into $\Q_p$ by the actual coefficient ring homomorphism.

The same source package retains the earlier independent formal identity
\[
 d!\coef{T}{d}\exp(x\log(1+T))=(x)_d
\]
over a commutative $\Q$-algebra. The section bridge above does not rely on selecting or averaging roots in a cyclotomic extension.

\textbf{Principal targets.}\par \mcode{factorial_coeff_section}; \mcode{section_weighted_integral};\par \mcode{weighted_section}.
\par\medskip
\subsubsection*{3. Weighted closure and the normalized denominator}
Let $w(d)=\lfloor d/p\rfloor$. Since $w(i)+w(d-i)\le w(d)$, the coefficient formula for a product proves closure of $\Wp$ under multiplication. Additive closure is immediate. For $U\in\Wp$ with $U(0)=1$, strong induction in the actual reciprocal coefficient recurrence proves
\[
 \boxed{U^{-1}\in\Wp.}
\]
This inverse property is a theorem, not an additional field in a structure encoding weighted series.

\subsection*{Differentiation and the Euler field}
The derivative coefficient is $(d+1)a_{d+1}$. If $p\nmid d+1$, the adjacent weights agree. If $p\mid d+1$, the integer factor $d+1$ compensates for the one-unit weight increase. Therefore ordinary differentiation preserves $\Wp$.

For the actual logarithm, its coefficient denominator is $d$, and the proved inequality
\[
 v_p(d)\le\lfloor d/p\rfloor\qquad(d>0)
\]
implies $\log(1+T)\in\Wp$. Consequently
\[
 L_TF=\log(1+T)(1+T)F'
\]
also preserves $\Wp$.

\subsection*{Removing a singular forced zero}
Put $c=ap<N$ and define the actual normalized kernel
\[
 K_c(t)=c!\sum_{n\ge0}\frac{t^{Nn}}{(Nn+c)!},\qquad K_c(0)=1.
\]
Let $\ell(T)=\log(1+T)/T$, defined by shifting logarithm coefficients, not by inverting $T$. Set
\[
 U(T)=\ell(T)^cK_c(\log(1+T)).
\]
The formal proof establishes $U(0)=1$ and the exact factorization
\[
 H_{N,c}(\log(1+T))=\frac{T^c}{c!}U(T).
\]
It follows coefficientwise that
\[
 \coef{T}{d}U=c!\coef{T}{d+c}H_{N,c}(\log(1+T)).
\]
Because $p^a\mid(ap)!$ and $\lfloor(d+ap)/p\rfloor=a+\lfloor d/p\rfloor$, the already proved section estimate gives $U\in\Wp$. Thus $U^{-1}\in\Wp$ as well. The factorial divisibility is derived from the pinned Legendre valuation theorem.

\textbf{Boundary.} The original positive-address section has zero constant coefficient. It is not the series to which unit inverse closure is applied.
\par\medskip
\subsubsection*{4. The singular ratio estimate}
For a positive address $c$, the rational ratio is independently constructed from the original kernels:
\[
 R_c(t)=c\,K_{c-1}(t)K_c(t)^{-1}.
\]
The scalar factor $c$ is necessary. The source proves the cleared identity
\[
 \boxed{R_c(t)H_{N,c}(t)=tH_{N,c-1}(t).}
\]
This connects the constructed object to the original unnormalized ratio $A_{c-1}/A_c$. It does not assign the desired transformed formula as its definition.

After actual logarithmic substitution and the chain rule,
\[
 R_c(\log(1+T))H_{N,c}(\log(1+T))
 =L_TH_{N,c}(\log(1+T)).
\]
Insert the proved factorization $H_{N,c}(\log(1+T))=T^cU/c!$. Leibniz differentiation and cancellation of that nonzero section give
\[
 \boxed{R_c(\log(1+T))
 =c(1+T)\ell(T)+(L_TU)U^{-1}.}
\]
All differentiation, substitution, nonzeroness and cancellation obligations are proved in the source.

For $c=ap$, the second summand belongs to $\Wp$ by the closure results on page 4. The first does too: shifting the logarithm coefficient loses at most one unit of the floor weight, so $p\ell(T)\in\Wp$, and the integer $a$ supplies the remaining factor. Therefore
\[
 \boxed{R_{ap}(\log(1+T))\in\Wp.}
\]
The final rational coefficient theorem states the bound with an explicit zero alternative. No finite-valued default valuation at zero is used as the meaning of unlimited precision.

\subsection*{Checks that prevent a normalization mistake}
The source proves that the actual ratio has constant coefficient $c$, not one. In the concrete case $N=6,c=3$, the universal theorem gives weighted membership and the constant coefficient is proved to be $3$.

It also derives $\coef{T}{3}H_{3,0}(\log(1+T))=1/6$ from the actual differential recurrence and checks its valuation $v_3(1/6)=-1$. Thus weighted membership does not mean ordinary coefficient integrality. This value attains the allowed floor weight at degree three.

\textbf{Principal target.} \mcode{SingularRatioAmice.paper_singular_amice_estimate}. The theorem's assumptions contain only the declared arithmetic parameter conditions.
\par\medskip
\subsubsection*{5. Verification and the remaining theorem}
The cumulative inventory consists of 57 modules, 630 explicit theorem declarations and 128 explicit definitions. The release verifier compiles every module in dependency order into a new directory, treats warnings as errors, and requests a transitive axiom report for all 758 listed declarations. The allowed axioms remain \mcode{propext}, \mcode{Classical.choice} and \mcode{Quot.sound}.

The verifier retains the previous 27 invalid-proof controls and adds two: incorrectly promoting a weighted coefficient to an integral one, and incorrectly normalizing the singular ratio constant to one. Missing imports, panics and timeouts do not qualify as successful negative controls. The preserved source checker rejects proof placeholders, added axioms and selected kernel-bypass constructs.

The Python infrastructure suite now contains 209 tests. Four old cumulative-count assertions initially failed against the extended inventory. Their original files and failed test output are preserved; the historical tests now inspect the frozen eighth-block inventory where appropriate, while current-corpus tests enforce the new counts.

\subsection*{Independent finite regressions}
The new exact-rational regression program is a separate implementation, not Lean extraction. Its checked scope is:
\begin{center}\begin{tabular}{lr}\toprule
Finite check & Count\\\midrule
Integral cyclic-coordinate numerator tests &17,775\\
Corresponding section-weight tests &17,775\\
Direct logarithmic composition comparisons &24\\
Singular ratio families &13\\
Singular ratio coefficient comparisons &597\\
Weighted product, inverse and derivative cases &120\\\bottomrule
\end{tabular}\end{center}
For $p=3,N=6,c=3$, this independent computation obtains the degree-six ratio coefficient $1/10080$, with valuation $-2$, attaining the allowed weight. This particular coefficient calculation is recorded as a finite regression, not as a new Lean theorem. A violating modulus $p=3,N=4$ yields a section coefficient of valuation $-4$ at degree nine, showing that the modulus divisibility hypothesis cannot simply be discarded.

\subsection*{What has not been proved by this block}
To obtain the singular higher-difference theorem, the weighted coefficient estimate must still be connected to the dual Mahler polynomial filtration and its integral pairing. The argument must then establish the required ratio flatness, propagate between singular indices, and assemble the original all-residue Euler family. None of these obligations is replaced by the finite regression tests.

The full manuscript register remains 23/24. Neither conjecture is resolved. The original manuscript is unchanged. Mathlib's pinned compiled cache is reused, not independently rebuilt. Frozen compiler logs, source hashes, failed drafts and fresh-extraction results accompany the release; its machine-readable audit is authoritative about the runs actually performed.

\clearpage\subsection{All Residue Completion Supplement}\label{app:sup10}
\noindent\textbf{Checkpoint scope.} This complete source body retains the status of its dated original supplement. Earlier pending or partial-coverage labels are historical; the final core records their resolution.\par\medskip
\subsubsection*{Scope and final statement}
Let $p$ be a prime, let $m>0$ divide $p-1$, put $N=mp$ and $h=(p-1)/m$, and let $0\le j<N$. The original rational packet is defined by
\[
 \sum_{q=0}^{n}\binom{Nn+j}{Nq}C_{N,j}(q)=\delta_{n0},
 \qquad E_j(n)=j!\,C_{N,j}(n).
\]
For every $n,k,u\in\N$, the new formal result establishes
\[
 \boxed{E_j(n)\in\Z_p,\qquad
 \Delta_{uh}^{k}E_j(n)\in p^{\min(k,mn)}\Z_p.}
\]
Here $\Delta_h f(n)=f(n+h)-f(n)$. The corresponding rational valuation statement is zero-safe:
\[
 \Delta_{uh}^{k}E_j(n)=0\quad\text{or}\quad
 v_p\bigl(\Delta_{uh}^{k}E_j(n)\bigr)\ge\min(k,mn).
\]
For the normalized coefficients, the exact correction is retained:
\[
 \boxed{\Delta_{uh}^{k}C_{N,j}(n)=0\quad\text{or}\quad
 v_p\bigl(\Delta_{uh}^{k}C_{N,j}(n)\bigr)\ge\min(k,mn)-v_p(j!).}
\]
The manuscript's odd-prime, positive-$u$ theorem is a specialization. The formal statement also includes zero row, zero order, zero multiplier, and the $p=2,m=1$ subcase. This does not assert a different or broader external dyadic conjecture.

\begin{center}\begin{tabular}{lrr}\toprule
Inventory & Prior v0.9.0 & Current source corpus\\\midrule
Named manuscript-proved items covered &23/24&24/24\\
Active mathematical modules &57&70\\
Explicit theorem declarations &630&753\\
Explicit definitions audited separately &128&148\\\bottomrule
\end{tabular}\end{center}
The 123 added theorem declarations include supporting lemmas and examples. They are not 123 independent discoveries. Three classical inputs and two conjectures remain separately tracked. The original manuscript and all 57 upstream proof modules are preserved unchanged.

\textbf{Principal declaration.} \path{KomorebiPaperIX.Theorem75.global_finite_difference} states the integrality and both precision bounds together, for the original coefficient functions. Its hypotheses are primality, $m>0$, $m\mid p-1$ and $j<mp$, not an assumed flatness or weightedness certificate.
\par\medskip
\subsubsection*{1. Stability of the weighted differential ring}
The recovered checkpoint used the actual weighted subring
\[
 \Wp=\left\{F(T)=\sum_{d\ge0}a_dT^d:\
 p^{\lfloor d/p\rfloor}a_d\in\Z_p\ \text{for every }d\right\}.
\]
Its coefficient condition is identified with the one in the previously checked singular-Amice development. The failed \mcode{ThetaWeighted} draft is retained separately; the repaired active source checks the complete formal-series statement below.

Put
\[
 \Theta=(1+T)\frac{d}{dT},\qquad
 B_\Theta=\frac{\Theta^{p^2}-\Theta^p}{p}.
\]
The coefficient formula
\[
 [T^d]\Theta F=(d+1)a_{d+1}+da_d
\]
proves $\Theta\Wp\subseteq\Wp$. If the floor weight increases between $d$ and $d+1$, then $p\mid d+1$, which compensates for that increase.

\textbf{The normalized defect is not assumed to be a derivation.}
Its proved product rule contains the two middle Leibniz sums:
\[
 B_\Theta(FG)=(B_\Theta F)G+F(B_\Theta G)+\mathcal B_p(F,G).
\]
The coefficients of $\mathcal B_p$ are the actual integral quotients
$\binom{p^2}{a}/p$ and $\binom pa/p$. Derivative stability and this correction formula show that the elements $F\in\Wp$ satisfying $B_\Theta F\in\Wp$ are closed under the finite ring operations.

The scalar input is proved for every natural $b$, including nonunits:
\[
 \boxed{p^2\mid b^{p^2}-b^p\quad\text{in }\Z.}
\]
Because $\Theta(1+T)^b=b(1+T)^b$, the binomial expansion of $T^p=((1+T)-1)^p$ then proves stability of the divided generator $T^p/p$. For $d=pq+r$, $0\le r<p$, every weighted monomial has the exact factorization
\[
 a_dT^d=(p^qa_d)\left(\frac{T^p}{p}\right)^qT^r.
\]
The scalar $p^qa_d$ is integral by the actual weight condition. Thus all weighted monomials, and all finite weighted polynomials, satisfy the defect condition.

Finally, $[T^d]B_\Theta F$ depends only on coefficients through $d+p^2$. Replacing $F$ by this finite prefix proves the result for the entire formal series, without an analytic convergence assumption:
\[
 \boxed{B_\Theta\Wp\subseteq\Wp,\qquad B_\Theta^s\Wp\subseteq\Wp\ (s\ge0).}
\]
\textbf{Modules.} \mcode{WeightedAmiceRing}, \mcode{FrobeniusSquare}, \mcode{ThetaWeighted}.
\par\medskip
\subsubsection*{2. Actual substitution and singular higher differences}
Let
\[
 \mathcal E(F)(t)=F(e^t-1),\qquad D=\frac{d}{dt},\qquad
 B_D=\frac{D^{p^2}-D^p}{p}.
\]
All substitutions use Mathlib's actual formal exponential, logarithm and substitution maps. The chain rule gives
\[
 D\mathcal E(F)=\mathcal E(\Theta F),\qquad
 B_D^s\mathcal E(F)=\mathcal E(B_\Theta^sF).
\]
The implementation also proves the inverse-substitution identity
\[
 \mathcal E\bigl(f(\log(1+T))\bigr)=f(t).
\]

The relevant integral subring is the \emph{Hurwitz} subring
\[
 \Hp=\{f\in\Q_p[[t]]:\ n![t^n]f\in\Z_p\ \text{for every }n\}.
\]
It is not coefficientwise integrality of ordinary power-series coefficients. The precise connection is
\[
 n![t^n]\mathcal E(F)=[T^0]\Theta^nF.
\]
Since every iterate of $\Theta$ preserves $\Wp$, the right side is integral. Therefore $\mathcal E(\Wp)\subseteq\Hp$. Applying the same statement to every $B_\Theta^sF$ proves
\[
 \boxed{f(\log(1+T))\in\Wp\quad\Longrightarrow\quad
 B_D^sf\in\Hp\text{ for every }s.}
\]
This is the proved transfer that was missing in the earlier attempt. No weighted Mahler pairing is assumed; the formal differential route supplies an alternative bridge.

\subsection*{Application to the actual singular ratios}
The preceding v0.9.0 source independently constructs
\[
 H_{N,c}(t)=\sum_{n\ge0}\frac{t^{Nn+c}}{(Nn+c)!},\quad
 A_c(t)=t^{-c}H_{N,c}(t),\quad R_c=\frac{A_{c-1}}{A_c}.
\]
The notation $t^{-c}H_{N,c}$ means removing its proved initial zero coefficients. For each $c=ap$, $1\le a<m$, that source proves
$R_{ap}(\log(1+T))\in\Wp$. The new transfer consequently proves flatness of the \emph{same} rational ratio after scalar extension to $\Q_p$.

The actual ratio is also proved to be an expansion in $t^N$, so its coefficients vanish outside multiples of $N$. Constructed $\Z_p$ lifts and the established coefficient criterion now give, for its literal Hurwitz coefficients $r_a(n)$,
\[
 \boxed{\Delta_{uh}^{k}r_a(n)\in p^{\min(k,mn)}\Z_p.}
\]
The separate rational precision theorem is obtained through the coefficient-lift equality. Divisibility in the field $\Q_p$ is never substituted for this assertion.

\textbf{Modules.} \mcode{WeightedDworkTransfer}, \mcode{SeriesDifferenceBridge}, \mcode{RatioSeriesSupport}.
\par\medskip
\subsubsection*{3. Nonsingular propagation without an assumed limit}
For consecutive original sections, the already defined ratios satisfy
\[
 R_cH_{N,c}=tH_{N,c-1},\qquad H_{N,c}'=H_{N,c-1}.
\]
Differentiating the first relation, using the relation for $R_{c-1}$, and cancelling a section whose leading coefficient is proved nonzero gives the actual Riccati equation
\[
 \boxed{tR_c'=R_c(1+R_{c-1}-R_c),\qquad 2\le c<N.}
\]
It is not an equation imposed on a replacement ratio. Write
\[
 q=R_{c-1}-(c-1),\qquad y=R_c-c.
\]
Both have zero constant coefficient. With $\Lambda_c=tD+c$, the equation becomes
\[
 \boxed{\Lambda_c y=cq+yq-y^2.}
\]

When $p\nmid c$, define the actual inverse coefficientwise:
\[
 [t^d]\Lambda_c^{-1}f=\frac{[t^d]f}{d+c}.
\]
On $p$-supported degrees, $d+c$ is a genuine $p$-adic unit. The proof establishes Hurwitz integrality of this solution, its support, both inverse identities, and preservation of full flatness using the earlier all-order Euler commutator theorem. The singular pivots $c=ap$ are never declared to be units.

Now define explicit iterations
\[
 y_0=0,\qquad y_{r+1}=\Lambda_c^{-1}(cq+y_rq-y_r^2).
\]
If $q$ is flat and $p$-supported, every iterate has those properties. The crucial link to the actual ratio is an exact finite-prefix theorem:
\[
 \boxed{[t^d]y_r=[t^d]y\qquad(d\le r).}
\]
Indeed,
\[
 (cq+zq-z^2)-(cq+yq-y^2)=(z-y)(q-z-y).
\]
The second factor has zero constant term. Multiplication therefore raises the established agreement prefix by one; coefficientwise division by $d+c$ preserves that prefix.

For fixed $s,d$, the coefficient $d![t^d]B_D^sy$ depends only on the coefficients of $y$ through $d+p^2s$. Choosing the matching finite iterate proves its integrality. This establishes flatness of the actual $y$, and hence of $R_c$, without postulating flatness of the solution or invoking an unproved analytic fixed-point theorem.

\textbf{Modules.} \mcode{EulerSeriesInverse}, \mcode{RiccatiFlatPropagation}, \mcode{RatioRiccati}.
\par\medskip
\subsubsection*{4. All-residue assembly and original coefficients}
The standard-window source and inverse results supply $R_1$ and $F_{N,0}$. Induction over $1\le c<N$ now has two proved cases. At $p\mid c$, the singular weighted estimate and the new substitution theorem give flatness. At $p\nmid c$, the preceding Riccati propagation applies. Thus every actual ratio $R_c$ is flat.

Let $K_j(t)$ be the normalized section kernel with constant coefficient one. The constructed unnormalized Euler series is
\[
 F_{N,j}(t)=j!\,K_j(t)^{-1}.
\]
The source/inverse identities and the factorial successor identity prove
\[
 F_{N,j+1}=F_{N,j}R_{j+1},\qquad
 F_{N,j}=F_{N,0}\prod_{c=1}^{j}R_c.
\]
The product formula is not an assembly hypothesis. Proved flatness of all factors and the original product-closure theorem establish flatness of the full family.

Coefficient extraction identifies, literally,
\[
 (Nn)![t^{Nn}]F_{N,j}=j!\,C_{N,j}(n).
\]
The support outside multiples of $N$ is proved separately. The constructed integral lift therefore represents the original rational coefficient function. The existing saturation criterion and the proved step-multiple argument give the ideal bound at every order, row and multiplier. Multiplication by $j!$ commutes with finite differences, giving precisely the displayed normalization loss.

\subsection*{Exact examples outside the standard window}
For $(p,m,j)=(3,2,3)$, the source derives
\[
 C_{6,3}(1)=-\frac1{84},\quad C_{6,3}(2)=\frac{703}{5460},\quad
 C_{6,3}(3)=-\frac{223199}{7980}.
\]
Consequently,
\[
 \Delta_1^2\bigl(3!C_{6,3}\bigr)(1)=-\frac{1464768}{8645},\qquad
 v_3\!\left(-\frac{1464768}{8645}\right)=2.
\]
This attains $\min(2,2\cdot1)=2$. At the regular address $j=4$ beyond the singular boundary, the source likewise proves
\[
 \Delta_1^2\bigl(4!C_{6,4}\bigr)(1)=-\frac{102745512}{2377375},
\]
and applies the universal exponent-two bound to this actual coefficient function.

The earlier no-loss counterexample is retained: $v_3(C_{6,3}(2)-C_{6,3}(1))=0$. Proving the corrected theorem does not make that false normalized claim true.

\textbf{Modules.} \mcode{AllRatiosDwork}, \mcode{AllResidueEuler}, \mcode{Theorem75}, \mcode{AllResidueExamples}.
\par\medskip
\subsubsection*{5. Verification record and claim boundary}
Each local module is rebuilt in a fresh output directory, with warnings treated as errors. A live \mcode{\#print axioms} audit checks every registered declaration, permitting only \mcode{propext}, \mcode{Classical.choice} and \mcode{Quot.sound}, or subsets.

\begin{center}\begin{tabular}{lr}\toprule
Required release check & Inventory\\\midrule
Local proof modules &70\\
Named theorem declarations &753\\
Explicit definitions &148\\
Expected axiom reports &901\\
Pinned dependency revisions &9\\
Deliberate invalid-proof controls &33\\
Infrastructure and release tests &235\\\bottomrule
\end{tabular}\end{center}
The accompanying JSON verification record, not the existence of this table, determines whether the actual run passed. The final package also includes a fresh-extraction replay record, checksums, the declaration inventory and the updated manuscript coverage ledger.

The negative controls preserve the earlier rejected false equalities, placeholders, invalid normalizations and unit pivots. Four new controls demand an extra power in the sharp singular example, remove the factorial loss, misuse a composite base, or center a ratio at the wrong constant. An infrastructure timeout or missing import is not an accepted negative result.

\subsection*{Independent finite checks}
The standard-library Python mirror checks 328 complete residue families for primes $2,3,5,7,11$, including 233 higher-residue families. It checks 4,264 integral unnormalized coefficients, 22,137 normalized and 22,137 unnormalized difference bounds, and 1,640 comparisons with a separately computed ordinary reciprocal. Further checks cover 108 weighted differential iterations, 205 integer Frobenius-square cases and 180 exact Picard agreement prefixes.

These are finite exact-arithmetic checks, not extracted Lean code and not substitutes for the universal formal proofs. The old normalization counterexample and the new sharp exponent-two example remain visible.

\subsection*{Preservation and limitations}
All 57 earlier proof modules and the original 28-page status-corrected manuscript are retained unchanged. The new supplement and amendments give the current proof route. Historical failed drafts and interrupted recovery attempts are retained separately from accepted builds. Historical tests that assert an older unfinished status now read that frozen historical status; current tests require the actual unconditional capstone theorem before accepting completion.

The portable Lean and Mathlib versions are pinned. Mathlib's compiled dependency cache is reused, not independently rebuilt. The historical numerical experiment suite and Papers I--VIII are not rerun. The two native conjectures remain unproved. No independent historical-priority assessment is made. Manuscript-to-Lean correspondence remains mathematical review, not a theorem equating natural-language prose with formal types.

\clearpage\subsection{Nishibuchi conjecture completion}\label{app:sup11}
\noindent\textbf{Checkpoint scope.} This complete source body retains the status of its dated original supplement. Earlier pending or partial-coverage labels are historical; the final core records their resolution.\par\medskip
\subsubsection*{1. The complete external statement}
For $N>0$ and $0\le j<N$, use the original unnormalized convention
\[
 \sum_{d\ge0}\EE_d^{(N,j)}\frac{t^d}{d!}
 =\left(\sum_{n\ge0}\frac{t^{Nn}}{(Nn+j)!}\right)^{-1}.
\]
The existing packet convention satisfies
$\EE_{Nn}^{(N,j)}=j!C_{N,j}(n)$. The formalization defines the degree-indexed coefficients through the actual factorial-scaled series coefficients, proves their packet identification, and checks the cleared generating identity with the original exponential section.

\begin{theorem}
Let $p$ be prime, $m>0$, and assume $m\mid p-1$ or $m=2$. Put $N=mp$ and $q=\operatorname{lcm}(2,p-1)$. Every $\EE_{Nn}^{(N,j)}$ is $p$-adically integral. For every $r\ge1$,
\[
\boxed{\EE_{Nn+qp^r}^{(N,j)}\equiv\EE_{Nn}^{(N,j)}\pmod{p^r}
       \qquad(0\le j<N,\ n\ge r).}
\]
In particular, the positive sufficient threshold $n_0=r$ works uniformly in all allowed parameters.
\end{theorem}

This is the parameter range and displacement of Nishibuchi's Conjecture 3.9 in arXiv:2605.10677v2. The explicit threshold also establishes the suggested bound $n_0\le r$. The original paper already proves the $p=2,m=2,j=0$ period as a special case, with all starting rows; this extension does not claim that special case as a new discovery.

\textbf{Principal declaration}\par
\code{KomorebiPaperIX.NishibuchiCompletion.conjecture_3_9}

The theorem type includes only the stated arithmetic hypotheses. There is no assumed congruence, weighted estimate, or flatness certificate. Precision is the zero-safe predicate $x=0$ or $v_p(x)\ge r$. Its connection to genuine ideal divisibility in $\Z_p$ is used in the proof.

\textbf{Claim boundary.} This establishes a sufficient threshold and a period, not the earliest starting index, the least period, or a pointwise optimal valuation formula. The correspondence with the externally written statement is a mathematical review of definitions and quantifiers, not an automated semantic-equivalence theorem. The two native conjectures in Paper IX remain unchanged and unproved.

\par\medskip
\subsubsection*{2. The missing dyadic branch}
For odd primes, the extra option $m=2$ already divides $p-1$. Therefore the only allowed family outside the preceding all-residue theorem is $p=2,m=2$, with $N=4$.

The new proof establishes a higher-difference estimate for every one of its four addresses:
\[
\boxed{v_2\!\left(\Delta_u^k C_{4,j}(n)\right)\ge\min(k,n)
       \qquad(0\le j<4).}
\]
All natural $u,k,n$ are allowed, including zero step, zero order and zero differences. The same bound holds for $j!C_{4,j}(n)$. Both a genuine $\Z_2$ divisibility theorem and a valuation theorem for the original rational sequence are included.

\subsection*{An explicit source instead of an assumed regularity condition}
Let $g_j(n)=\binom{4n+j}{4n}^{-1}$. Its four formulas are
\[
\begin{aligned}
 g_0(n)&=1,\\
 g_1(n)&=(4n+1)^{-1},\\
 g_2(n)&=((4n+1)(2n+1))^{-1},\\
 g_3(n)&=3((4n+1)(2n+1)(4n+3))^{-1}.
\end{aligned}
\]
Every denominator is odd. The code constructs the actual unit inverses in $\Z_2$ and proves these identities against the original binomial coefficients, including $n=0$.

Define a Kummer sequence over a commutative ring by
\[
\PP_a(f)\quad\Longleftrightarrow\quad
 a^k\mid\Delta_1^k f(n)\ \text{for all }k,n.
\]
The exact discrete product rule
\[
\Delta_1^s(fg)(n)=\sum_{i=0}^s\binom si
     \Delta_1^i f(n)\,\Delta_1^{s-i}g(n+i)
\]
is proved by induction. It proves closure under products. Applied to $fg=1$, the same formula isolates
$\Delta_1^s g(n)f(n+s)$; induction and multiplication by $g(n+s)$ prove closure under actual pointwise inverses.

An affine sequence whose slope is divisible by $a$ satisfies $\PP_a$. Consequently the four displayed source sequences satisfy $\PP_2$. No numerical interpolation or finite set of samples is used in that deduction.

\textbf{Source modules:} \code{KummerSequences}, \code{DyadicFourKernel}.

\par\medskip
\subsubsection*{3. The dyadic differential operator and reciprocity}
The operator required here is
\[
\boxed{B_4=\frac{D^8-D^4}{2},\qquad D=\frac{d}{dt}.}
\]
It is not the preceding $p=2$ specialization $(D^4-D^2)/2$. It is also not an ordinary derivation. The new algebra module proves the actual product formula, including the middle derivative sums whose coefficients are $\binom 8i/2$ and $\binom 4i/2$.

Write $\mathcal H_{\mathrm{int}}$ for the subring of $\Q_2[[t]]$ whose factorial-scaled coefficients belong to $\Z_2$. Flatness means $B_4^sF\in\mathcal H_{\mathrm{int}}$ for every $s$. The proof of multiplicative and reciprocal closure is induction on $s$, using the full product correction. It separately proves that the actual reciprocal has integral Hurwitz coefficients when $F(0)=1$.

Pack $g_j$ into degrees divisible by four and form its actual EGF. The formal coefficient identity is
\[
G_j(t)=\sum_{n\ge0}g_j(n)\frac{t^{4n}}{(4n)!}
       =j!\sum_{n\ge0}\frac{t^{4n}}{(4n+j)!}.
\]
Its inverse is therefore the original normalized packet EGF. The source Kummer property gives flatness of $G_j$. The proved reciprocal closure gives flatness of $G_j^{-1}$.

For a packed coefficient sequence $e$, the exact extraction formula is
\[
\bigl(U^8-U^4\bigr)^k\operatorname{pack}_4(e)(4n-4k)
 =\Delta_1^k e(n),\qquad k\le n.
\]
Here $U$ is the coefficient shift. The nonnegative extraction index is why the first estimate covers $k\le n$. The existing saturation lemma extends it to all $k$ with exponent $\min(k,n)$, and the step-multiple lemma extends step one to every natural $u$.

\subsection*{Identity with the original rational coefficients}
The integral lift is constructed from the decoded coefficients of the actual reciprocal. The code proves
\[
 \operatorname{coe}(e_j(n))=C_{4,j}(n)\quad\text{in }\Q_2,
\]
and proves that the coefficient map commutes with every difference stencil. Multiplying by $j!$ yields the unnormalized family. The valuation result is not divisibility in a field.

\textbf{Source modules:} \code{CubicDworkAlgebra}, \code{CubicDworkClosure};\par
\code{DyadicFourDwork} and \code{DyadicFourCongruences}.

\par\medskip
\subsubsection*{4. From differences to the exact external period}
Let $f$ take values in a commutative ring and suppose
\[
 p^{\min(k,r)}\mid\Delta_h^k f(n)\qquad\text{for every }k.
\]
The proved finite Newton expansion is
\[
 f(n+p^{r-1}h)-f(n)=
 \sum_{k=1}^{p^{r-1}}\binom{p^{r-1}}k\Delta_h^k f(n).
\]
For $k\ge r$, the difference already supplies $p^r$. For $1\le k<r$, the code proves the integer divisibility
\[
 p^r\mid\binom{p^{r-1}}k p^k.
\]
It follows from $k\binom{p^{r-1}}k=p^{r-1}\binom{p^{r-1}-1}{k-1}$ and $v_p(k)<k$. Thus every summand is divisible by $p^r$, including the endpoint $r=1$.

\subsection*{Divisor branch}
For $m\mid p-1$, the preceding all-residue theorem provides the needed differences as soon as $r\le mn$. With $q=\operatorname{lcm}(2,p-1)$, the required packet step $q/m$ is proved to be an integer multiple of $(p-1)/m$. The degree conversion is checked exactly:
\[
 mp\left(n+p^{r-1}\frac qm\right)=mpn+qp^r.
\]
This retains the stronger sufficient condition $r\le mn$ on the divisor branch. The uniform choice $n\ge r$ implies it for every positive $m$.

\subsection*{Exceptional branch}
For $p=2,m=2$, $q=2$ and the packet displacement is $2^{r-1}$. The new dyadic bound supplies the hypothesis when $n\ge r$, and
\[
 4\bigl(n+2^{r-1}\bigr)=4n+2^{r+1}.
\]
The $p=2,m=1$ branch is also checked using $q=2$, not an incorrectly substituted $q=p-1=1$.

\subsection*{An exact arithmetic check, not an optimality claim}
The independent rational recurrence gives
$C_{4,1}(1)=-1/5$ and $C_{4,1}(2)=121/45$. Their difference is $26/9$, of $2$-adic valuation one. This agrees with the new bound at $n=k=u=1$. The source also contains universal instantiations of the degree-indexed theorem at all four dyadic addresses and at an odd singular address.

\textbf{Source modules:} \code{PeriodFromDifferences}, \code{NishibuchiCompletion}, \code{ConjectureExamples}.

\par\medskip
\subsubsection*{5. Verification, provenance and boundary}
The release has 79 active modules: 70 unchanged baseline modules and nine new modules. Its inventory contains 836 named theorem declarations, 162 explicit definitions and two abbreviations. The transitive audit includes all 1,000 of these declarations. The increase of 83 named theorems includes supporting lemmas and examples, not 83 independent discoveries.

A clean cumulative run compiled all \textbf{79 modules}, matched all \textbf{1,000 axiom reports}, verified all \textbf{nine dependency revisions}, and correctly rejected all \textbf{37 invalid-proof controls}. The statement-contract file separately checked the external parameter disjunction, integrality, positive threshold, and exact degree displacement. Warnings were treated as errors.

The capstone depends only on \code{propext}, \code{Classical.choice}, and \code{Quot.sound}. The 25 extension-infrastructure tests passed. Exact-rational checks passed 6,500 higher-difference cases in each dyadic normalization, 588 dyadic periods and 1,690 divisor-branch periods. These finite checks are not the universal proof. A fresh-extraction replay, when recorded in the separate release audit, is additional evidence beyond the source build.


\subsection*{Reproduction}
The portable runtime is the pinned Linux x86-64 Lean 4.34.0-rc2 kit. It is not duplicated inside this source archive. From the extracted release, run
\begin{verbatim}
bash Run-Portable.sh /path/to/Komorebi_Portable_Lean_4.34.0-rc2
\end{verbatim}
The launcher checks file hashes, runs 25 extension-infrastructure tests and the new exact-arithmetic regressions, then builds every module in a new output directory and audits every registered declaration. It also compiles an explicit statement contract and negative controls. The final check verifies the released snapshot again.

The source-policy scan excludes proof placeholders and additional project axioms. Compiler errors, timeouts, missing imports and invalid negative-control preludes are not accepted as successful theorem verification. Failed development snapshots are retained separately from active proof sources.

\subsection*{Limits of the completion claim}
The pinned Mathlib compiled cache is reused, not rebuilt from source. No independent third-party checker or peer review was run. The original numerical experiment suite and Papers I--VIII were not rerun. Prior verification evidence is retained as historical evidence and is not relabeled as fresh testing. The original Paper IX manuscript still requires cumulative supplements and amendments; this extension is not a rewritten primary manuscript.

The external Conjecture 3.9 is distinct from Paper IX's native Conjectures 13.2 and 13.3. Those two remain unproved. The present result does not assert least periods, earliest thresholds, or an optimal exponent for every individual parameter tuple. It makes no historical-priority claim relative to other unpublished work.

\subsection*{Reference}
Yuta Nishibuchi. \emph{$p$-adic Congruences of Generalized Euler Numbers and Relations to Even Zeta Values}. arXiv:2605.10677v2, revised 1 June 2026. Definition 2.1, Conjecture 3.9 and Proposition 3.10. The version identifier is fixed for the statement comparison.


\clearpage\def\IXMainAtlas{1}\def\AtlasFigDir{../../figures/worked}% Shared worked examples. Exact data and figure generation: tools/worked_examples.py.
\begingroup
\ifdefined\IXMainAtlas
\section{Worked arithmetic and lifting atlas}\label{app:worked-atlas}
These examples connect the coefficient calculus with the correction machinery. They instantiate the statements in the main text rather than add new universal claims. Each displayed rational value, residue, and finite reconstruction is checked by the accompanying exact script.
\fi
\newcommand{\AtlasHead}[2]{\ifdefined\AtlasStarted\clearpage\fi\ifdefined\IXMainAtlas\subsection{#1}\else\section{#1}\fi\label{#2}\def\AtlasStarted{1}}
\newcommand{\AtlasImage}[1]{\begin{center}\includegraphics[width=.91\linewidth]{\AtlasFigDir/#1.pdf}\end{center}}
\AtlasHead{Compute the reciprocal before estimating it}{atlas:reciprocal}
For width $N$ and address $j$, the normalized source and its reciprocal are
\[
 G_{N,j}(t)=j!\sum_{n\ge0}\frac{t^{Nn}}{(Nn+j)!},\qquad
 G_{N,j}(t)^{-1}=\sum_{n\ge0}C_{N,j}(n)\frac{t^{Nn}}{(Nn)!}.
\]
Multiplying the two series and clearing the factorials gives an exact triangular recurrence:
\[
 C_{N,j}(0)=1,\qquad
 C_{N,j}(n)=-\frac{\displaystyle\sum_{r=0}^{n-1}
 \binom{Nn+j}{Nr}C_{N,j}(r)}{\binom{Nn+j}{Nn}}.
\]
\AtlasImage{reciprocal_rows}
The table is for $N=4$. The first two address-zero rows illustrate the computation without any valuation estimate:
\[
 C_{4,0}(1)=-1,\qquad
 C_{4,0}(2)=-\bigl(1+\binom84(-1)\bigr)=69.
\]
At address one, the same recurrence yields $C_{4,1}(1)=-1/5$ and $C_{4,1}(2)=121/45$. These are exact rational coefficients, not rounded approximations. The second implementation in the supplied script inverts the ordinary power series first, then multiplies by $(Nn)!$. Agreement between the two constructions checks the factorial placement.

\AtlasHead{The normalization loss is a real boundary}{atlas:normalization}
The original Euler coefficient is $E_{Nn}^{(N,j)}=j!C_{N,j}(n)$. Multiplication by $j!$ changes the available $p$-adic precision. That change cannot be discarded outside the unit-factorial window.

Take $p=3$, $N=6$, $j=3$, and compare rows one and two. Exact inversion gives
\[
 C_{6,3}(1)=-\frac1{84},\quad
 C_{6,3}(2)=\frac{703}{5460},\quad
 C_{6,3}(2)-C_{6,3}(1)=\frac{64}{455}.
\]
\AtlasImage{normalization_loss}
Neither $64$ nor $455$ is divisible by three, so the normalized difference has valuation zero. The original coefficient difference is $3!\cdot64/455=384/455$, of valuation one.

This is the counterexample behind the corrected factorial-loss statement. A bound for the original Euler normalization must lose $v_p(j!)$ when translated to $C_{N,j}$. For $0\le j<p$, that loss is zero. The distinction concerns which sequence the theorem measures, not a numerical weakness of the calculation.

\AtlasHead{Odd primes require the sign twist}{atlas:sign}
The prime-power transfer controls a sign-twisted reciprocal sequence at odd primes. It must not be read as an ordinary-difference theorem for the untwisted coefficients.

At $N=p=3$ and $j=0$, put $A(n)=(-1)^nC_{3,0}(n)$. The first three rows are
\AtlasImage{odd_prime_sign}
The ordinary first difference at row zero is $-1-1=-2$. It is a $3$-adic unit, so an asserted lower bound of one for that difference would be false. The twisted first difference is zero, whose mathematical valuation is $+\infty$. At the next row, the twisted difference is $19-1=18$, of valuation two.

Equivalently, the untwisted sequence is acted on by $T_H=S^H-(-1)^H$, where $S$ advances one row. At odd step one this is $S+1$, not $S-1$. The zero in the example is handled by a zero-safe divisibility statement rather than a finite valuation equality. The explicit sign and the zero convention are both part of the theorem's interface.

\AtlasHead{Newton coordinates turn precision into a finite computation}{atlas:newton}
Let $b_k=\Delta_1^kC_{4,0}(0)$. The first coefficients obtained from the difference triangle are
\input{generated/worked_newton_values.tex}
The exact Newton identity recovers each nonnegative integer row:
\[
 C_{4,0}(n)=\sum_{k=0}^{n}b_k\binom nk.
\]
\AtlasImage{newton_cutoff}
For this address the proved law is $v_2(b_k)=2k-s_2(k)$ for $k\ge1$. Every $k\ge5$ therefore has $v_2(b_k)\ge8$. Since $\binom nk$ is integral, those terms vanish modulo $256$. Thus, for every $n\ge0$,
\[
 C_{4,0}(n)\equiv1+254\binom n1+72\binom n2+
 176\binom n3+128\binom n4\pmod{256}.
\]
The uniform truncation follows from the coefficient theorem. The script separately checks the exact Newton reconstruction and this reduction for rows $0$ through $32$. Those finite checks validate the worked arithmetic without replacing the theorem.

\AtlasHead{A compatible root is built one digit at a time}{atlas:smooth}
Take $F(x)=x^2-7$ over the $3$-adic integers and start with $x_1=1$ modulo three. If $F(x_L)=3^Lr_L$, the next representative has the form $x_{L+1}=x_L+3^Lt_L$. Its correcting digit satisfies
\[
 r_L+2x_Lt_L\equiv0\pmod3.
\]
\AtlasImage{smooth_tower}
For example, $F(1)/3=-2$ and $2x_1=2$, so the next digit is one and $x_2=4$. Then $F(4)/9=1$, again giving digit one and $x_3=13$. At the next stage the residual is divisible by three, so digit zero is correct. A zero digit is not a stalled construction.

Every center remains congruent to one modulo three. Hence $2x_L$ stays a unit, and the correction equation has a unique digit at every stage. The representatives are compatible and their increments are divisible by $3^L$. Completeness gives a limit, and continuity of the polynomial gives a root of $F$. The finite table exhibits the construction. The unit derivative is what supports indefinite continuation.

\AtlasHead{A singular residual can stop a lift}{atlas:blocked}
Now take $F(x)=x^2-3$ at the residue $x=0$ modulo three. The residual is $r=F(0)/3=-1$ and the Jacobian is zero modulo three. The next-digit equation would require
\[
 -1+0\cdot t\equiv0\pmod3,
\]
which has no solution.
\AtlasImage{singular_block}
The three representatives of the old root class modulo nine all leave residual six modulo nine. Thus the root class exists at precision three but does not lift to precision nine. This is an exact obstruction, not failure of an iterative search.

A singular Jacobian alone does not imply failure. For $F(x)=x^2$ at the same center, the residual is zero and all three candidate digits lift to roots modulo nine. The obstruction is the residual class in the cokernel of the actual Jacobian. It is this class, rather than the rank by itself, that decides whether a next lift exists.

\AtlasHead{Separate corrections need not be one shared correction}{atlas:shared}
Consider the two-component system $F(x)=(x,x-3)$ at $x=0$, with $q=3$ and $L=1$. Both components vanish modulo three. After writing a candidate as $y=3t$, its residual vector modulo three is
\[
 \overline r+J\overline t=(0,2)+(t,t)=(t,t-1).
\]
\AtlasImage{shared_obstruction}
The first coordinate can be corrected by $t=0$. The second can be corrected by $t=1$. No digit corrects both. The available residuals are $(0,2)$, $(1,0)$, and $(2,1)$, none of which is $(0,0)$.

The combined Jacobian maps $t$ to the diagonal vector $(t,t)$. Its image is smaller than the product of its coordinate images. Checking the two projected obstructions independently therefore loses the compatibility condition. The common-correction theorem retains the full residual vector and the shared digit space, which is exactly the information needed to detect this failure.

\AtlasHead{A derivative weight can cost one digit of root precision}{atlas:weighted}
Compare the linear equations $F(x)=3x=0$ and $G_A(x)=3x+3^A=0$, where $A\ge2$. Their coefficient difference has $3$-adic valuation $A$, but their roots differ by
\[
 x_A-0=-3^{A-1},\qquad v_3(x_A)=A-1.
\]
\AtlasImage{weighted_motion}
The derivative has valuation one. Solving the correction equation divides by three, so one digit of coefficient precision is lost in root displacement. The calculation is exact for every displayed $A$ and also follows immediately from the two linear equations.

This example explains why weighted lifting and deformation estimates track a Jacobian loss rather than silently treating every invertible derivative over $\mathbb Q_p$ as a unit over $\mathbb Z_p$. A valuation-zero derivative would not incur this particular loss. Nonlinear systems additionally require their stated neighborhood, residual, and right-inverse hypotheses.

\medskip\noindent\textbf{Reproducibility.} The eight examples use exact rational and integer arithmetic in \nolinkurl{tools/worked_examples.py}. The release record states the executed finite ranges. This atlas adds no Lean endpoint and changes no formal claim in the companion.
\endgroup
\clearpage
\begin{thebibliography}{99}
\bibitem{Nishibuchi}
Yuta Nishibuchi,
\emph{$p$-adic Congruences of Generalized Euler Numbers and Relations to Even
Zeta Values}, arXiv:2605.10677v2, 1 June 2026.
\url{https://arxiv.org/abs/2605.10677v2}.
The comparison to Conjecture 3.9 uses the source review and fixed-version
record preserved with this edition.

\bibitem{gessel}
Ira M. Gessel, \emph{Applications of the classical umbral calculus},
arXiv:math/0108121v3, 22 November 2001, Section 9, in particular the Frobenius Euler congruence and the derivation following Theorem 9.2.
\url{https://arxiv.org/abs/math/0108121}.
The normalization and quantifiers of the ordinary Euler input were checked against the primary PDF in the strengthening checkpoint.

\bibitem{DLMF24}
NIST Digital Library of Mathematical Functions, \S24.10(i),
Eq.\ 24.10.1 (Bernoulli denominator theorem).
\url{https://dlmf.nist.gov/24.10.E1}. The official primary entry was inspected
for the positive even-index denominator specialization in this edition.

\bibitem{Mathlib}
The Mathlib Community,
\emph{Mathlib}, source revision
\nolinkurl{aeb72f1e9981b6d35b426b16dca7759be3e7ba06},
used with Lean 4.34.0-rc2.
In particular, the polynomial, strict inverse-function and analytic
inverse-function developments cited in Section~\ref{sec:regularity}.
The exact local source paths and dependency records accompany the release.

\bibitem{PaperIXArchive}
Komorebi P.,
\emph{Projector-Resolved Seed Calculus IX: original manuscript, formal
supplements and cumulative continuation sources}, preserved source collection,
consolidated 8 September 2026.
The original 28-page status-corrected manuscript and all subsequent release
records remain in the accompanying archive.

\medskip\noindent\textbf{Recovered historical references.}
The bibliography underlying Appendix~\ref{app:recovered} has been recovered
from the original source package. Its 22 citation keys are retained below,
with bibliographic identities checked against publisher and author records.
This repairs the source attributions; it does not enlarge the historical
chapter's stated proof or formalization scope.
% Generated from references_historical.bib using amsplain, with full repeated authors and DOI links.
\bibitem{Amice1964}
Yvette Amice, \emph{{Interpolation {$p$}-adique}}, Bulletin de la Soci\'et\'e
  Math\'ematique de France \textbf{92} (1964), 117--180.
\url{https://doi.org/10.24033/bsmf.1606}.

\bibitem{BarberoCerrutiMurru2018}
Stefano Barbero, Umberto Cerruti, and Nadir Murru, \emph{{Some Combinatorial
  Properties of the Hurwitz Series Ring}}, Ricerche di Matematica \textbf{67}
  (2018), no.~2, 491--507.
\url{https://doi.org/10.1007/s11587-017-0336-x}.

\bibitem{CarlitzLevine1960}
Leonard Carlitz and Jack Levine, \emph{{Some Problems Concerning Kummer's
  Congruences for the Euler Numbers and Polynomials}}, Transactions of the
  American Mathematical Society \textbf{96} (1960), no.~1, 23--37.
\url{https://doi.org/10.1090/s0002-9947-1960-0115971-x}.

\bibitem{Chen2004}
Kwang-Wu Chen, \emph{{Congruences for Euler Numbers}}, The Fibonacci Quarterly
  \textbf{42} (2004), no.~2, 128--140.
\url{https://doi.org/10.1080/00150517.2004.12428429}.

\bibitem{Comtet1974}
Louis Comtet, \emph{{Advanced Combinatorics: The Art of Finite and Infinite
  Expansions}}, D. Reidel Publishing Company, Dordrecht, 1974.
\url{https://doi.org/10.1007/978-94-010-2196-8}.

\bibitem{EverestVanDerPoortenShparlinskiWard2003}
Graham Everest, Alf van~der Poorten, Igor Shparlinski, and Thomas Ward,
  \emph{{Recurrence Sequences}}, Mathematical Surveys and Monographs, vol. 104,
  American Mathematical Society, Providence, 2003.
\url{https://doi.org/10.1090/surv/104}.

\bibitem{Gouvea1997}
Fernando~Q. Gouv\^ea, \emph{{{$p$}-adic Numbers: An Introduction}}, second ed.,
  Universitext, Springer, Berlin, 1997.
\url{https://doi.org/10.1007/978-3-642-59058-0}.

\bibitem{GrahamKnuthPatashnik1994}
Ronald~L. Graham, Donald~E. Knuth, and Oren Patashnik, \emph{{Concrete
  Mathematics: A Foundation for Computer Science}}, second ed., Addison-Wesley,
  Reading, Massachusetts, 1994.

\bibitem{IrelandRosen1990}
Kenneth Ireland and Michael Rosen, \emph{{A Classical Introduction to Modern
  Number Theory}}, second ed., Graduate Texts in Mathematics, vol.~84,
  Springer, New York, 1990.
\url{https://doi.org/10.1007/978-1-4757-2103-4}.

\bibitem{Keigher1997}
William~F. Keigher, \emph{{On the Ring of Hurwitz Series}}, Communications in
  Algebra \textbf{25} (1997), no.~6, 1845--1859.
\url{https://doi.org/10.1080/00927879708825957}.

\bibitem{Koblitz1984}
Neal Koblitz, \emph{{{$p$}-adic Numbers, {$p$}-adic Analysis, and
  Zeta-Functions}}, second ed., Graduate Texts in Mathematics, vol.~58,
  Springer, New York, 1984.
\url{https://doi.org/10.1007/978-1-4612-1112-9}.

\bibitem{LidlNiederreiter1997}
Rudolf Lidl and Harald Niederreiter, \emph{{Finite Fields}}, second ed.,
  Encyclopedia of Mathematics and its Applications, vol.~20, Cambridge
  University Press, Cambridge, 1997.
\url{https://doi.org/10.1017/CBO9780511525926}.

\bibitem{Lucas1878a}
\'Edouard Lucas, \emph{{Th\'eorie des Fonctions Num\'eriques Simplement
  P\'eriodiques}}, American Journal of Mathematics \textbf{1} (1878), no.~2,
  184--196.
\url{https://doi.org/10.2307/2369308}.

\bibitem{Mahler1958}
Kurt Mahler, \emph{{An Interpolation Series for Continuous Functions of a
  {$p$}-adic Variable}}, Journal f\"ur die Reine und Angewandte Mathematik
  \textbf{199} (1958), 23--34.
\url{https://doi.org/10.1515/crll.1958.199.23}.

\bibitem{Robert2000}
Alain~M. Robert, \emph{{A Course in {$p$}-adic Analysis}}, Graduate Texts in
  Mathematics, vol. 198, Springer, New York, 2000.
\url{https://doi.org/10.1007/978-1-4757-3254-2}.

\bibitem{Roman1984}
Steven Roman, \emph{{The Umbral Calculus}}, Pure and Applied Mathematics, vol.
  111, Academic Press, New York, 1984.

\bibitem{Serre1979}
Jean-Pierre Serre, \emph{{Local Fields}}, Graduate Texts in Mathematics,
  vol.~67, Springer, New York, 1979.
\url{https://doi.org/10.1007/978-1-4757-5673-9}.

\bibitem{Stanley2012}
Richard~P. Stanley, \emph{{Enumerative Combinatorics. Volume 1}}, second ed.,
  Cambridge Studies in Advanced Mathematics, vol.~49, Cambridge University
  Press, Cambridge, 2012.
\url{https://doi.org/10.1017/CBO9781139058520}.

\bibitem{Sun2008}
Zhi-Hong Sun, \emph{{Congruences Involving Bernoulli and Euler Numbers}},
  Journal of Number Theory \textbf{128} (2008), no.~2, 280--312.
\url{https://doi.org/10.1016/j.jnt.2007.03.003}.

\bibitem{Washington1997}
Lawrence~C. Washington, \emph{{Introduction to Cyclotomic Fields}}, second ed.,
  Graduate Texts in Mathematics, vol.~83, Springer, New York, 1997.
\url{https://doi.org/10.1007/978-1-4612-1934-7}.

\bibitem{Young1999}
Paul~Thomas Young, \emph{{Congruences for Bernoulli, Euler, and Stirling
  Numbers}}, Journal of Number Theory \textbf{78} (1999), no.~2, 204--227.
\url{https://doi.org/10.1006/jnth.1999.2401}.

\bibitem{Young2001}
Paul~Thomas Young, \emph{{Kummer Congruences for Values of Bernoulli and Euler
  Polynomials}}, Acta Arithmetica \textbf{99} (2001), no.~3, 277--288.
\url{https://doi.org/10.4064/aa99-3-5}.

\end{thebibliography}
\end{document}
